\documentclass[a4paper,12pt]{book}

\usepackage[utf8]{inputenc}
\usepackage[T1]{fontenc}
\usepackage[french]{babel}
\usepackage{amsmath,amssymb,amsthm,mathtools}
\usepackage{tcolorbox}
\tcbuselibrary{theorems,skins,breakable}
\usepackage{tikz,tikz-cd,pgfplots}
\pgfplotsset{compat=1.18}
\usetikzlibrary{patterns,arrows.meta,calc,decorations.markings}
\usepackage{graphicx,float,caption}
\usepackage{enumitem,booktabs,array,longtable}
\usepackage{hyperref}
\usepackage{makeidx}
\usepackage{minitoc}
\usepackage{esint}
\makeindex

% --- Math Operators ---
\DeclareMathOperator{\Hom}{Hom}
\DeclareMathOperator{\End}{End}
\DeclareMathOperator{\Aut}{Aut}
\DeclareMathOperator{\im}{im}
\DeclareMathOperator{\Id}{Id}
\DeclareMathOperator{\Ker}{Ker}
\DeclareMathOperator{\coker}{coker}
\DeclareMathOperator{\Spec}{Spec}
\DeclareMathOperator{\Tr}{Tr}
\DeclareMathOperator{\supp}{supp}
\DeclareMathOperator{\diam}{diam}
\DeclareMathOperator{\sgn}{sgn}
\DeclareMathOperator{\conv}{conv}
\DeclareMathOperator{\codim}{codim}
\DeclareMathOperator{\spn}{span}
\DeclareMathOperator{\dom}{dom}
\DeclareMathOperator{\ran}{ran}

\newcommand{\N}{\mathbb{N}}
\newcommand{\Z}{\mathbb{Z}}
\newcommand{\Q}{\mathbb{Q}}
\newcommand{\R}{\mathbb{R}}
\newcommand{\C}{\mathbb{C}}
\newcommand{\K}{\mathbb{K}}
\newcommand{\F}{\mathbb{F}}
\newcommand{\dd}{\mathrm{d}}
\newcommand{\norm}[1]{\left\| #1 \right\|}
\newcommand{\abs}[1]{\left| #1 \right|}
\newcommand{\inner}[2]{\left\langle #1, #2 \right\rangle}
\newcommand{\weakto}{\rightharpoonup}
\newcommand{\weakstarto}{\overset{*}{\rightharpoonup}}

% --- Theorem Environments ---
\theoremstyle{definition}
\newtheorem{definition}{Définition}[chapter]
\newtheorem{exemple}[definition]{Exemple}
\newtheorem{exercice}{Exercice}[chapter]

\theoremstyle{plain}
\newtheorem{theoreme}[definition]{Théorème}
\newtheorem{proposition}[definition]{Proposition}
\newtheorem{lemme}[definition]{Lemme}
\newtheorem{corollaire}[definition]{Corollaire}

\theoremstyle{remark}
\newtheorem{remarque}[definition]{Remarque}
\newtheorem{notation}[definition]{Notation}

\tcolorboxenvironment{definition}{enhanced,breakable,colback=blue!5!white,colframe=blue!75!black,fonttitle=\bfseries}
\tcolorboxenvironment{theoreme}{enhanced,breakable,colback=red!5!white,colframe=red!75!black,fonttitle=\bfseries}
\tcolorboxenvironment{proposition}{enhanced,breakable,colback=red!5!white,colframe=red!60!black,fonttitle=\bfseries}
\tcolorboxenvironment{lemme}{enhanced,breakable,colback=yellow!5!white,colframe=yellow!70!black,fonttitle=\bfseries}
\tcolorboxenvironment{corollaire}{enhanced,breakable,colback=orange!5!white,colframe=orange!70!black,fonttitle=\bfseries}
\tcolorboxenvironment{exemple}{enhanced,breakable,colback=green!5!white,colframe=green!60!black,fonttitle=\bfseries}
\tcolorboxenvironment{exercice}{enhanced,breakable,colback=purple!5!white,colframe=purple!60!black,fonttitle=\bfseries}
\tcolorboxenvironment{remarque}{enhanced,breakable,colback=gray!5!white,colframe=gray!50!black,fonttitle=\bfseries}

% --- Document Info ---
\begin{document}

\begin{titlepage}
\begin{tikzpicture}[remember picture, overlay]
  \fill[left color=blue!15!white, right color=blue!5!white]
    (current page.south west) rectangle (current page.north east);
  \fill[color=orange!70!yellow]
    ([yshift=-2cm]current page.north west) rectangle ([yshift=-2.3cm]current page.north east);
  \fill[color=blue!80!black, rounded corners=8pt]
    ([xshift=3cm, yshift=-4cm]current page.north west) rectangle
    ([xshift=-3cm, yshift=-10cm]current page.north east);
  \node[anchor=center, text=white, font=\Huge\bfseries]
    at ([yshift=-6cm]current page.north) {Analyse Fonctionnelle};
  \node[anchor=center, text=orange!80!yellow, font=\Large]
    at ([yshift=-7.5cm]current page.north) {Notes de Cours / Lecture Notes};
  \node[anchor=center, text=white!80, font=\large]
    at ([yshift=-8.8cm]current page.north) {Master --- 2025--2026};
  \node[anchor=center, text=blue!80!black, font=\Large\itshape]
    at ([yshift=-12cm]current page.north) {Ya\'e Ulrich Gaba};
  \draw[orange!70!yellow, line width=2pt]
    ([xshift=5cm, yshift=-13.5cm]current page.north west) --
    ([xshift=-5cm, yshift=-13.5cm]current page.north east);
  \node[anchor=center, text width=12cm, align=center, text=blue!60!black, font=\itshape]
    at ([yshift=-16cm]current page.north) {Des espaces de Banach aux alg\`ebres d'op\'erateurs~: un voyage \`a travers l'analyse fonctionnelle moderne.};
  \node[anchor=south, text=gray, font=\small]
    at ([yshift=1.5cm]current page.south) {\today};
\end{tikzpicture}
\end{titlepage}
\dominitoc
\tableofcontents

% ============================================================
\chapter*{Pr\'eface}
\addcontentsline{toc}{chapter}{Pr\'eface}
\markboth{Pr\'eface}{Pr\'eface}

L'analyse fonctionnelle est n\'ee d'une contradiction : les m\'ethodes de l'analyse classique, con\c{c}ues pour la droite r\'eelle et $\R^n$, devaient s'\'etendre \`a des espaces de fonctions, c'est-\`a-dire \`a des espaces de dimension infinie. Ce qui s'y joue n'est pas seulement technique. Dans un espace de dimension infinie, la boule unit\'e n'est jamais compacte, deux normes \'equivalentes peuvent ne pas exister, et la notion m\^eme de convergence se d\'edouble entre topologie forte et topologie faible.

\medskip

Ce cours construit progressivement les outils qui ont permis de surmonter ces obstacles : espaces de Banach et de Hilbert, th\'eor\`emes fondamentaux (Hahn--Banach, Banach--Steinhaus, application ouverte, graphe ferm\'e), th\'eorie spectrale des op\'erateurs compacts et auto-adjoints. Les applications vis\'ees sont l'analyse harmonique, la m\'ecanique quantique, et la th\'eorie moderne des \'equations aux d\'eriv\'ees partielles.

\medskip

\noindent\textbf{Pr\'erequis.} Topologie g\'en\'erale (compacit\'e, convergence), th\'eorie de la mesure (espaces $L^p$), alg\`ebre lin\'eaire.

\vfill
\hfill\textit{L'auteur}

% ============================================================

%%%%%%%%%%%%%%%%%%%%%%%%%%%%%%%%%%%%%%%%%%%%%%%%%%%%%%%%%%%%%%%%%%%%%
%%%%%%%%%%%%%%%%%%%%%%%%%%%%%%%%%%%%%%%%%%%%%%%%%%%%%%%%%%%%%%%%%%%%%
%%  CHAPITRE 1 : ESPACES DE BANACH --- FONDEMENTS
%%%%%%%%%%%%%%%%%%%%%%%%%%%%%%%%%%%%%%%%%%%%%%%%%%%%%%%%%%%%%%%%%%%%%
%%%%%%%%%%%%%%%%%%%%%%%%%%%%%%%%%%%%%%%%%%%%%%%%%%%%%%%%%%%%%%%%%%%%%

\chapter{Espaces de Banach --- Fondements}
\minitoc
\vspace{1cm}

L'analyse fonctionnelle\index{analyse fonctionnelle} est l'étude des espaces vectoriels de dimension infinie munis de structures topologiques compatibles avec les opérations algébriques. Ce premier chapitre établit les fondations de la théorie en introduisant les espaces vectoriels normés\index{espace vectoriel normé} et les espaces de Banach\index{espace de Banach}, qui constituent le cadre naturel pour l'étude des équations fonctionnelles, des équations aux dérivées partielles et de la théorie spectrale.

Historiquement, la notion d'espace de Banach a été introduite par Stefan Banach\index{Banach, Stefan} dans sa thèse de 1920, publiée en 1922. Son ouvrage fondateur \emph{Théorie des opérations linéaires}\index{Théorie des opérations linéaires} (1932) a posé les bases de l'analyse fonctionnelle moderne. Les contributions de Banach, Hahn\index{Hahn, Hans}, Schauder\index{Schauder, Juliusz} et leurs contemporains de l'école polonaise de Lwów\index{école de Lwów} ont profondément transformé l'analyse mathématique du XXe siècle.

Dans ce chapitre, le corps de base $\K$ désigne $\R$ ou $\C$ sauf mention explicite du contraire.

%====================================================================
\section{Espaces vectoriels normés}
%====================================================================

\begin{definition}[Norme]\label{def:norme}
\index{norme}
Soit $E$ un espace vectoriel sur le corps $\K$ (où $\K = \R$ ou $\K = \C$). Une \textbf{norme} sur $E$ est une application $\norm{\cdot} : E \to [0, +\infty[$ vérifiant :
\begin{enumerate}[label=(\roman*)]
    \item \textbf{Séparation}\index{norme!séparation} : $\norm{x} = 0 \iff x = 0$ ;
    \item \textbf{Homogénéité}\index{norme!homogénéité} : $\norm{\lambda x} = \abs{\lambda} \norm{x}$ pour tout $\lambda \in \K$ et tout $x \in E$ ;
    \item \textbf{Inégalité triangulaire}\index{inégalité triangulaire} : $\norm{x + y} \leq \norm{x} + \norm{y}$ pour tous $x, y \in E$.
\end{enumerate}
Le couple $(E, \norm{\cdot})$ est appelé \textbf{espace vectoriel normé} (e.v.n.).
\end{definition}

\begin{remarque}\label{rem:distance-induite}
Toute norme $\norm{\cdot}$ sur $E$ induit une distance\index{distance} $d(x,y) = \norm{x - y}$, qui est \emph{invariante par translation} : $d(x+a, y+a) = d(x,y)$ pour tout $a \in E$. Cette distance munit $E$ d'une structure d'espace métrique\index{espace métrique}.
\end{remarque}

\begin{notation}\label{not:boules}
On note $B(x,r) = \{y \in E : \norm{y - x} < r\}$ la boule ouverte\index{boule ouverte} de centre $x$ et de rayon $r > 0$, et $\overline{B}(x,r) = \{y \in E : \norm{y - x} \leq r\}$ la boule fermée\index{boule fermée}. La sphère unité\index{sphère unité} est $S_E = \{x \in E : \norm{x} = 1\}$. On écrit $B_E = \overline{B}(0,1)$ pour la boule unité fermée\index{boule unité fermée}.
\end{notation}

\begin{definition}[Semi-norme]\label{def:semi-norme}
\index{semi-norme}
Une \textbf{semi-norme} sur $E$ est une application $p : E \to [0,+\infty[$ vérifiant les conditions (ii) et (iii) de la 
ef{def:norme}, mais pas nécessairement (i). Autrement dit, on peut avoir $p(x) = 0$ pour $x \neq 0$.
\end{definition}

\subsection{Exemples fondamentaux}

\begin{exemple}[Les espaces $\R^n$ et $\C^n$]\label{ex:Rn}
\index{espace!Rn@$\R^n$}
Sur $\K^n$, on dispose de la famille de normes\index{norme!$\ell^p$ sur $\K^n$} suivante. Pour $1 \leq p < \infty$ et $x = (x_1, \ldots, x_n) \in \K^n$ :
\[
    \norm{x}_p = \left( \sum_{k=1}^{n} \abs{x_k}^p \right)^{1/p},
\]
et pour $p = \infty$ :
\[
    \norm{x}_\infty = \max_{1 \leq k \leq n} \abs{x_k}.
\]
Les trois cas les plus courants sont :
\begin{itemize}
    \item $\norm{x}_1 = \abs{x_1} + \cdots + \abs{x_n}$ \quad (norme $\ell^1$)\index{norme!$\ell^1$} ;
    \item $\norm{x}_2 = \sqrt{\abs{x_1}^2 + \cdots + \abs{x_n}^2}$ \quad (norme euclidienne)\index{norme!euclidienne} ;
    \item $\norm{x}_\infty = \max_k \abs{x_k}$ \quad (norme sup)\index{norme!sup}.
\end{itemize}

\begin{center}
\begin{tikzpicture}[scale=1.8]
    \draw[->] (-1.5,0) -- (1.5,0) node[right] {$x_1$};
    \draw[->] (0,-1.5) -- (0,1.5) node[above] {$x_2$};
    \draw[blue,thick] (1,0) -- (0,1) -- (-1,0) -- (0,-1) -- cycle;
    \draw[red,thick] (0,0) circle (1);
    \draw[green!60!black,thick] (-1,-1) rectangle (1,1);
    \node[blue] at (0.7,0.85) {$\ell^1$};
    \node[red] at (1.15,0.5) {$\ell^2$};
    \node[green!60!black] at (1.2,1.15) {$\ell^\infty$};
\end{tikzpicture}
\captionof{figure}{Boules unités dans $\R^2$ pour les normes $\ell^1$, $\ell^2$ et $\ell^\infty$.}\label{fig:boules-Rn}
\end{center}
\end{exemple}

\begin{exemple}[Les espaces $\ell^p$]\label{ex:ell-p}
\index{espace!$\ell^p$}\index{ell-p@$\ell^p$}
Pour $1 \leq p < \infty$, l'espace $\ell^p$ est défini par :
\[
    \ell^p = \left\{ x = (x_n)_{n \geq 1} \subset \K : \sum_{n=1}^{\infty} \abs{x_n}^p < \infty \right\},
\]
muni de la norme $\norm{x}_{\ell^p} = \bigl( \sum_{n=1}^{\infty} \abs{x_n}^p \bigr)^{1/p}$. Pour $p = \infty$ :
\[
    \ell^\infty = \left\{ x = (x_n)_{n \geq 1} \subset \K : \sup_{n \geq 1} \abs{x_n} < \infty \right\}, \quad \norm{x}_{\ell^\infty} = \sup_{n \geq 1} \abs{x_n}.
\]
L'espace $c_0$\index{espace!$c_0$} des suites convergeant vers $0$ est un sous-espace fermé de $\ell^\infty$.
\end{exemple}

\begin{exemple}[L'espace $C(K)$]\label{ex:CK}
\index{espace!$C(K)$}\index{C(K)@$C(K)$}
Soit $K$ un espace topologique compact\index{compact}. L'espace $C(K)$ des fonctions continues $f : K \to \K$ est un espace vectoriel normé pour la norme uniforme\index{norme!uniforme} :
\[
    \norm{f}_\infty = \max_{t \in K} \abs{f(t)}.
\]
\end{exemple}

\begin{exemple}[Aperçu des espaces $L^p$]\label{ex:Lp-preview}
\index{espace!$L^p$}\index{Lp@$L^p$}
Soit $(\Omega, \mathcal{A}, \mu)$ un espace mesuré\index{espace mesuré}. Pour $1 \leq p < \infty$ :
\[
    L^p(\Omega, \mu) = \left\{ f \text{ mesurable} : \int_\Omega \abs{f}^p \, \dd\mu < \infty \right\} \Big/ \sim,
\]
où $f \sim g$ si $f = g$ $\mu$-presque partout, muni de $\norm{f}_{L^p} = \bigl( \int_\Omega \abs{f}^p \, \dd\mu \bigr)^{1/p}$.
\end{exemple}

\subsection{Inégalités fondamentales}

\begin{proposition}[Inégalité de Hölder discrète]\label{prop:holder-discret}
\index{inégalité de Hölder!discrète}
Soient $p, q \in [1, +\infty]$ des exposants conjugués\index{exposants conjugués}, c'est-à-dire $\frac{1}{p} + \frac{1}{q} = 1$. Pour tous $x \in \ell^p$ et $y \in \ell^q$ :
\[
    \sum_{n=1}^{\infty} \abs{x_n y_n} \leq \norm{x}_{\ell^p} \norm{y}_{\ell^q}.
\]
\end{proposition}

\begin{proof}
On utilise l'\textbf{inégalité de Young}\index{inégalité de Young} : pour $a, b \geq 0$, $ab \leq \frac{a^p}{p} + \frac{b^q}{q}$. Si $\norm{x}_{\ell^p} = 0$ ou $\norm{y}_{\ell^q} = 0$, le résultat est trivial. Sinon, posons $\tilde{x}_n = \abs{x_n}/\norm{x}_{\ell^p}$ et $\tilde{y}_n = \abs{y_n}/\norm{y}_{\ell^q}$. L'inégalité de Young donne $\tilde{x}_n \tilde{y}_n \leq \frac{\tilde{x}_n^p}{p} + \frac{\tilde{y}_n^q}{q}$. En sommant :
\[
    \sum_{n=1}^{\infty} \tilde{x}_n \tilde{y}_n \leq \frac{1}{p} \sum_{n=1}^{\infty} \tilde{x}_n^p + \frac{1}{q} \sum_{n=1}^{\infty} \tilde{y}_n^q = \frac{1}{p} + \frac{1}{q} = 1.
\]
En multipliant par $\norm{x}_{\ell^p}\norm{y}_{\ell^q}$, on obtient le résultat.
\end{proof}

\begin{proposition}[Inégalité de Minkowski discrète]\label{prop:minkowski-discret}
\index{inégalité de Minkowski!discrète}
Pour $1 \leq p \leq \infty$ et $x, y \in \ell^p$ : $\norm{x + y}_{\ell^p} \leq \norm{x}_{\ell^p} + \norm{y}_{\ell^p}$.
\end{proposition}

\begin{proof}
Les cas $p = 1$ et $p = \infty$ sont immédiats. Pour $1 < p < \infty$, on écrit $\abs{x_n + y_n}^p \leq \abs{x_n + y_n}^{p-1}(\abs{x_n} + \abs{y_n})$. En sommant et en appliquant Hölder avec $q = p/(p-1)$ :
\begin{align*}
    \sum_n \abs{x_n + y_n}^p &\leq \left(\sum_n \abs{x_n + y_n}^{(p-1)q}\right)^{1/q} \norm{x}_p + \left(\sum_n \abs{x_n + y_n}^{(p-1)q}\right)^{1/q} \norm{y}_p \\
    &= \left(\sum_n \abs{x_n + y_n}^p\right)^{1/q} \bigl(\norm{x}_p + \norm{y}_p\bigr).
\end{align*}
On divise par $\bigl(\sum_n \abs{x_n + y_n}^p\bigr)^{1/q}$ pour conclure, puisque $1 - 1/q = 1/p$.
\end{proof}

%====================================================================
\section{Convergence et séries dans les espaces normés}
%====================================================================

\begin{definition}[Convergence]\label{def:convergence-evn}
\index{convergence!dans un e.v.n.}
Soit $(E, \norm{\cdot})$ un e.v.n. Une suite $(x_n)_{n \geq 0}$ dans $E$ \textbf{converge} vers $x \in E$ si $\norm{x_n - x} \to 0$ quand $n \to \infty$. On écrit $x_n \to x$ ou $\lim_{n\to\infty} x_n = x$.
\end{definition}

\begin{definition}[Suite de Cauchy]\label{def:cauchy}
\index{suite de Cauchy}
Une suite $(x_n)$ dans $(E, \norm{\cdot})$ est une \textbf{suite de Cauchy} si :
\[
    \forall \varepsilon > 0, \quad \exists N \in \N, \quad \forall m, n \geq N, \quad \norm{x_m - x_n} < \varepsilon.
\]
\end{definition}

\begin{proposition}\label{prop:cv-implique-cauchy}
Toute suite convergente est de Cauchy. La réciproque est fausse en général.
\end{proposition}

\begin{proof}
Si $x_n \to x$, alors $\norm{x_m - x_n} \leq \norm{x_m - x} + \norm{x - x_n} \to 0$. Pour la réciproque, considérer $E = \Q$ muni de la valeur absolue : la suite des approximations décimales de $\sqrt{2}$ est de Cauchy mais ne converge pas dans $\Q$.
\end{proof}

\begin{definition}[Séries dans un e.v.n.]\label{def:series-evn}
\index{série!dans un e.v.n.}
Soit $(x_n)_{n \geq 0}$ une suite dans $E$. La \textbf{série} $\sum x_n$ est dite \textbf{convergente}\index{série!convergente} si la suite des sommes partielles $S_N = \sum_{n=0}^{N} x_n$ converge dans $E$. La série est \textbf{absolument convergente}\index{série!absolument convergente} si $\sum_{n=0}^{\infty} \norm{x_n} < \infty$.
\end{definition}

\begin{remarque}\label{rem:abs-cv-dim-finie}
En dimension finie, toute série absolument convergente est convergente. En dimension infinie, c'est une propriété caractéristique de la complétude (voir le 
ef{thm:abs-cv-completude}).
\end{remarque}

\begin{proposition}[Continuité des opérations]\label{prop:continuite-operations}
\index{continuité!des opérations vectorielles}
Soit $(E, \norm{\cdot})$ un e.v.n. Les applications suivantes sont continues :
\begin{enumerate}[label=(\alph*)]
    \item L'addition $(x, y) \mapsto x + y$ de $E \times E$ dans $E$ ;
    \item La multiplication par un scalaire $(\lambda, x) \mapsto \lambda x$ de $\K \times E$ dans $E$ ;
    \item La norme $x \mapsto \norm{x}$ de $E$ dans $\R$.
\end{enumerate}
Autrement dit, un e.v.n. est un \textbf{espace vectoriel topologique}\index{espace vectoriel topologique}.
\end{proposition}

\begin{proof}
(a) $\norm{(x+y) - (x_0 + y_0)} \leq \norm{x - x_0} + \norm{y - y_0}$.

(b) $\norm{\lambda x - \lambda_0 x_0} = \norm{(\lambda - \lambda_0)x + \lambda_0(x - x_0)} \leq \abs{\lambda - \lambda_0}\norm{x} + \abs{\lambda_0}\norm{x - x_0}$.

(c) $\abs{\norm{x} - \norm{y}} \leq \norm{x - y}$ par l'inégalité triangulaire inversée\index{inégalité triangulaire inversée}.
\end{proof}

\begin{definition}[Normes équivalentes]\label{def:normes-equivalentes}
\index{normes équivalentes}
Deux normes $\norm{\cdot}_a$ et $\norm{\cdot}_b$ sur un espace vectoriel $E$ sont \textbf{équivalentes} s'il existe des constantes $c, C > 0$ telles que
\[
    c \norm{x}_a \leq \norm{x}_b \leq C \norm{x}_a \quad \text{pour tout } x \in E.
\]
Des normes équivalentes définissent les mêmes notions de convergence, de suites de Cauchy, d'ouverts, de fermés et de complets.
\end{definition}

\begin{exemple}[Comparaison des normes $\ell^p$ sur $\K^n$]\label{ex:comparaison-normes-Kn}
\index{normes équivalentes!sur $\K^n$}
Sur $\K^n$, on a les inégalités :
\[
    \norm{x}_\infty \leq \norm{x}_2 \leq \norm{x}_1 \leq n \norm{x}_\infty, \qquad
    \norm{x}_\infty \leq \norm{x}_2 \leq \sqrt{n}\norm{x}_\infty.
\]
Ces inégalités sont optimales : l'égalité dans $\norm{x}_1 \leq n\norm{x}_\infty$ est atteinte pour $x = (1,1,\ldots,1)$ et l'égalité dans $\norm{x}_2 \leq \sqrt{n}\norm{x}_\infty$ est aussi atteinte pour ce vecteur.
\end{exemple}

\begin{definition}[Parties bornées, fermées, denses]\label{def:parties-evn}
\index{partie bornée}\index{partie dense}\index{sous-espace!dense}
Soit $(E, \norm{\cdot})$ un e.v.n.
\begin{itemize}
    \item Une partie $A \subset E$ est \textbf{bornée}\index{borné} si $\sup_{x \in A}\norm{x} < \infty$.
    \item Une partie $D \subset E$ est \textbf{dense} dans $E$ si $\overline{D} = E$, c'est-à-dire si pour tout $x \in E$ et tout $\varepsilon > 0$, il existe $d \in D$ tel que $\norm{x - d} < \varepsilon$.
    \item Un e.v.n. est \textbf{séparable}\index{séparable} s'il contient un sous-ensemble dénombrable dense.
\end{itemize}
\end{definition}

\begin{exemple}[Séparabilité de $\ell^p$]\label{ex:ell-p-separable}
\index{espace!$\ell^p$!séparabilité}
Pour $1 \leq p < \infty$, l'espace $\ell^p$ est séparable. En effet, l'ensemble des suites à support fini et à coefficients dans $\Q$ (ou $\Q + i\Q$ si $\K = \C$) est dénombrable et dense. En revanche, $\ell^\infty$ n'est \emph{pas} séparable : l'ensemble des suites à valeurs dans $\{0,1\}$ est non dénombrable et les éléments sont à distance $1$ les uns des autres.
\end{exemple}

%====================================================================
\section{Espaces de Banach : définition et complétude}
%====================================================================

\begin{definition}[Espace de Banach]\label{def:banach}
\index{espace de Banach}\index{Banach, Stefan}
Un espace vectoriel normé $(E, \norm{\cdot})$ est un \textbf{espace de Banach} s'il est \textbf{complet}\index{complet}, c'est-à-dire si toute suite de Cauchy dans $E$ converge dans $E$.
\end{definition}

\begin{theoreme}[Complétude de $\ell^p$]\label{thm:ell-p-complet}
\index{espace!$\ell^p$!complétude}\index{complétude!de $\ell^p$}
Pour tout $1 \leq p \leq \infty$, l'espace $\ell^p$ est un espace de Banach.
\end{theoreme}

\begin{proof}
\textbf{Cas $1 \leq p < \infty$.} Soit $(x^{(k)})_{k \geq 1}$ une suite de Cauchy dans $\ell^p$, où $x^{(k)} = (x_n^{(k)})_{n \geq 1}$. Pour chaque $n$ fixé :
\[
    \abs{x_n^{(k)} - x_n^{(j)}} \leq \norm{x^{(k)} - x^{(j)}}_{\ell^p} \to 0 \quad \text{quand } k, j \to \infty.
\]
Donc $(x_n^{(k)})_{k \geq 1}$ est de Cauchy dans $\K$, qui est complet. Posons $x_n = \lim_{k \to \infty} x_n^{(k)}$.

\textbf{Étape 1 : $x = (x_n) \in \ell^p$.} Soit $\varepsilon > 0$. Il existe $K$ tel que pour $k, j \geq K$ : $\norm{x^{(k)} - x^{(j)}}_p < \varepsilon$. Pour tout $N \in \N$ :
\[
    \sum_{n=1}^{N} \abs{x_n^{(k)} - x_n^{(j)}}^p < \varepsilon^p.
\]
En faisant $j \to \infty$ (à $N$ et $k$ fixés) :
\[
    \sum_{n=1}^{N} \abs{x_n^{(k)} - x_n}^p \leq \varepsilon^p.
\]
Puis en faisant $N \to \infty$ : $\norm{x^{(k)} - x}_p \leq \varepsilon$. Par l'inégalité triangulaire :
\[
    \norm{x}_p \leq \norm{x - x^{(k)}}_p + \norm{x^{(k)}}_p \leq \varepsilon + \norm{x^{(K)}}_p < \infty.
\]
Donc $x \in \ell^p$.

\textbf{Étape 2 : $x^{(k)} \to x$ dans $\ell^p$.} On a montré que $\norm{x^{(k)} - x}_p \leq \varepsilon$ pour tout $k \geq K$, ce qui donne $x^{(k)} \to x$.

\textbf{Cas $p = \infty$.} Soit $(x^{(k)})$ de Cauchy dans $\ell^\infty$. Pour chaque $n$, $(x_n^{(k)})_k$ est de Cauchy dans $\K$, donc converge vers $x_n$. Soit $\varepsilon > 0$ et $K$ tel que $\norm{x^{(k)} - x^{(j)}}_\infty < \varepsilon$ pour $k, j \geq K$. Pour tout $n$ : $\abs{x_n^{(k)} - x_n^{(j)}} < \varepsilon$. En faisant $j \to \infty$ : $\abs{x_n^{(k)} - x_n} \leq \varepsilon$ pour tout $n$, donc $\norm{x^{(k)} - x}_\infty \leq \varepsilon$. Comme $\norm{x}_\infty \leq \norm{x - x^{(K)}}_\infty + \norm{x^{(K)}}_\infty < \infty$, on a $x \in \ell^\infty$ et $x^{(k)} \to x$.
\end{proof}

\begin{theoreme}[Complétude de $C(K)$]\label{thm:CK-complet}
\index{espace!$C(K)$!complétude}\index{complétude!de $C(K)$}
Soit $K$ un espace topologique compact. L'espace $(C(K), \norm{\cdot}_\infty)$ est un espace de Banach.
\end{theoreme}

\begin{proof}
Soit $(f_k)_{k \geq 1}$ une suite de Cauchy dans $C(K)$. Pour chaque $t \in K$ :
\[
    \abs{f_k(t) - f_j(t)} \leq \norm{f_k - f_j}_\infty \to 0,
\]
donc $(f_k(t))_k$ converge dans $\K$. Posons $f(t) = \lim_{k\to\infty} f_k(t)$.

\textbf{Convergence uniforme.} Soit $\varepsilon > 0$ et $K_0$ tel que $\norm{f_k - f_j}_\infty < \varepsilon$ pour $k, j \geq K_0$. Pour tout $t \in K$ et $k \geq K_0$, en faisant $j \to \infty$ dans $\abs{f_k(t) - f_j(t)} < \varepsilon$, on obtient $\abs{f_k(t) - f(t)} \leq \varepsilon$. Ceci étant vrai pour tout $t$, on a $\norm{f_k - f}_\infty \leq \varepsilon$.

\textbf{Continuité de $f$.} Comme limite uniforme de fonctions continues\index{convergence uniforme!et continuité}, $f$ est continue. En effet, pour $t_0 \in K$ et $\varepsilon > 0$, on choisit $k$ tel que $\norm{f_k - f}_\infty < \varepsilon/3$, puis un voisinage $V$ de $t_0$ tel que $\abs{f_k(t) - f_k(t_0)} < \varepsilon/3$ pour $t \in V$. Alors pour $t \in V$ :
\[
    \abs{f(t) - f(t_0)} \leq \abs{f(t) - f_k(t)} + \abs{f_k(t) - f_k(t_0)} + \abs{f_k(t_0) - f(t_0)} < \varepsilon.
\]
Donc $f \in C(K)$ et $f_k \to f$ dans $C(K)$.
\end{proof}

\begin{remarque}
Les espaces $c_0$\index{espace!$c_0$!complétude} et $c$\index{espace!$c$} (suites convergentes) sont des sous-espaces fermés de $\ell^\infty$, donc sont des espaces de Banach.
\end{remarque}

\begin{proposition}[Sous-espaces fermés et complétude]\label{prop:sous-espace-ferme-complet}
\index{sous-espace!fermé!et complétude}
Soit $E$ un espace de Banach et $F$ un sous-espace vectoriel de $E$. Alors $F$ est un espace de Banach (pour la norme induite) si et seulement si $F$ est fermé dans $E$.
\end{proposition}

\begin{proof}
$(\Rightarrow)$ Si $F$ est complet et $(x_n) \subset F$ converge vers $x \in E$, alors $(x_n)$ est de Cauchy dans $F$, donc converge dans $F$ : $x \in F$.

$(\Leftarrow)$ Si $F$ est fermé et $(x_n) \subset F$ est de Cauchy, alors $(x_n)$ converge vers un $x \in E$ (car $E$ est complet). Comme $F$ est fermé, $x \in F$.
\end{proof}

\begin{exemple}[{L'espace $C^k([a,b])$}]\label{ex:Ck}
\index{espace!$C^k$}\index{C-k@$C^k([a,b])$}
L'espace $C^k([a,b])$ des fonctions $k$ fois continûment dérivables sur $[a,b]$, muni de la norme
\[
    \norm{f}_{C^k} = \sum_{j=0}^{k} \norm{f^{(j)}}_\infty = \sum_{j=0}^{k} \max_{t \in [a,b]} \abs{f^{(j)}(t)},
\]
est un espace de Banach. En effet, si $(f_n)$ est de Cauchy dans $C^k$, alors pour chaque $0 \leq j \leq k$, $(f_n^{(j)})$ est de Cauchy dans $C([a,b])$, donc converge uniformément vers une fonction continue $g_j$. On vérifie que $g_0^{(j)} = g_j$ par le théorème de dérivation des limites uniformes\index{dérivation!des limites uniformes}.
\end{exemple}

\begin{exemple}[Espace des suites à décroissance rapide]\label{ex:schwartz-discret}
\index{décroissance rapide}\index{espace!$s$}
L'espace
\[
    s = \left\{ x = (x_n)_{n \geq 1} : \sup_{n \geq 1} n^k \abs{x_n} < \infty \text{ pour tout } k \in \N \right\}
\]
est un sous-espace de $\ell^p$ pour tout $p \geq 1$, mais il n'est pas fermé (et donc pas un espace de Banach pour la norme $\ell^p$). On peut le munir d'une topologie d'espace de Fréchet\index{espace de Fréchet} définie par la famille de semi-normes $p_k(x) = \sup_n n^k\abs{x_n}$.
\end{exemple}

%====================================================================
\section{Espaces de dimension finie}
%====================================================================

\begin{theoreme}[Équivalence des normes en dimension finie]\label{thm:equiv-normes-dim-finie}
\index{normes équivalentes!en dimension finie}\index{dimension finie!équivalence des normes}
Sur un espace vectoriel $E$ de dimension finie, toutes les normes sont équivalentes.
\end{theoreme}

\begin{proof}
Soit $\dim E = n$ et $(e_1, \ldots, e_n)$ une base de $E$. Munissons $E$ de la norme $\norm{\cdot}_\infty$ définie par $\norm{\sum_{i=1}^n \alpha_i e_i}_\infty = \max_{1\leq i\leq n} \abs{\alpha_i}$. Il suffit de montrer que toute norme $\norm{\cdot}$ sur $E$ est équivalente à $\norm{\cdot}_\infty$.

\textbf{Majoration.} Pour $x = \sum_{i=1}^n \alpha_i e_i$ :
\[
    \norm{x} \leq \sum_{i=1}^n \abs{\alpha_i}\norm{e_i} \leq \left(\sum_{i=1}^n \norm{e_i}\right) \norm{x}_\infty = C \norm{x}_\infty,
\]
où $C = \sum_{i=1}^n \norm{e_i}$. Ceci montre que $\norm{\cdot}$ est continue pour la topologie de $\norm{\cdot}_\infty$.

\textbf{Minoration.} Considérons la sphère unité $S = \{x \in E : \norm{x}_\infty = 1\}$. L'application $\varphi : (E, \norm{\cdot}_\infty) \to \R$ définie par $\varphi(x) = \norm{x}$ est continue (par la majoration ci-dessus) et $S$ est compact\index{compact!en dimension finie} (car $S$ est identifié à un sous-ensemble fermé borné de $\K^n$, qui est compact par le théorème de Heine-Borel\index{Heine-Borel, théorème de}). Donc $\varphi$ atteint son minimum $m$ sur $S$. Comme $\norm{x} > 0$ pour $x \in S$ (car $x \neq 0$), on a $m > 0$. Pour tout $x \neq 0$ :
\[
    \norm{x} = \norm{x}_\infty \cdot \norm*{\frac{x}{\norm{x}_\infty}} \geq m \norm{x}_\infty.
\]
Donc $m \norm{x}_\infty \leq \norm{x} \leq C \norm{x}_\infty$.
\end{proof}

\begin{corollaire}\label{cor:dim-finie-banach}
\index{dimension finie!complétude}
Tout espace vectoriel normé de dimension finie est un espace de Banach.
\end{corollaire}

\begin{proof}
Soit $E$ de dimension $n$. Fixons une base et identifions $E$ à $\K^n$. Toute norme sur $E$ est équivalente à $\norm{\cdot}_2$, et $(\K^n, \norm{\cdot}_2)$ est complet. Comme la complétude est préservée par équivalence de normes\index{normes équivalentes!et complétude}, $E$ est complet.
\end{proof}

\begin{theoreme}[Compacité de la boule unité fermée]\label{thm:compacite-boule}
\index{boule unité fermée!compacité}\index{dimension finie!compacité de la boule}
Soit $E$ un espace vectoriel normé. La boule unité fermée $B_E$ est compacte si et seulement si $\dim E < \infty$.
\end{theoreme}

\begin{proof}
$(\Leftarrow)$ Si $\dim E < \infty$, alors $B_E$ est un fermé borné de $E \simeq \K^n$, donc compact par Heine-Borel.

$(\Rightarrow)$ On montre la contraposée : si $\dim E = \infty$, alors $B_E$ n'est pas compacte. On utilise le lemme de Riesz (
ef{lem:riesz}) pour construire une suite sans sous-suite convergente.

On construit par récurrence une suite $(x_n)$ dans $S_E$ telle que $\norm{x_m - x_n} \geq 1/2$ pour $m \neq n$. Posons $F_1 = \spn\{x_1\}$ avec $x_1 \in S_E$ quelconque. Par le lemme de Riesz appliqué au sous-espace propre fermé $F_1$, il existe $x_2 \in S_E$ tel que $\norm{x_2 - y} \geq 1/2$ pour tout $y \in F_1$. On continue : $F_n = \spn\{x_1, \ldots, x_n\}$ est fermé (dimension finie), et le lemme de Riesz donne $x_{n+1} \in S_E$ tel que $d(x_{n+1}, F_n) \geq 1/2$. Comme $x_m \in F_n$ pour $m \leq n$, on a $\norm{x_{n+1} - x_m} \geq 1/2$.

Cette suite n'admet aucune sous-suite de Cauchy, donc aucune sous-suite convergente, et $B_E$ n'est pas séquentiellement compacte\index{compacité séquentielle}, donc pas compacte (dans un espace métrique, les deux notions coïncident).
\end{proof}

\begin{lemme}[Lemme de Riesz]\label{lem:riesz}
\index{Riesz, lemme de}\index{lemme de Riesz (géométrie)}
Soit $E$ un espace vectoriel normé et $F$ un sous-espace vectoriel fermé de $E$ avec $F \neq E$. Pour tout $\theta \in \,]0, 1[$, il existe $x_\theta \in E$ tel que $\norm{x_\theta} = 1$ et $d(x_\theta, F) \geq \theta$.
\end{lemme}

\begin{proof}
Soit $y \in E \setminus F$. Comme $F$ est fermé, $d = d(y, F) > 0$. Par définition de l'infimum, il existe $z \in F$ tel que $\norm{y - z} \leq d/\theta$. Posons :
\[
    x_\theta = \frac{y - z}{\norm{y - z}}.
\]
Alors $\norm{x_\theta} = 1$. Pour tout $f \in F$ :
\[
    \norm{x_\theta - f} = \frac{1}{\norm{y - z}} \norm{y - z - \norm{y-z} f} = \frac{1}{\norm{y-z}} \norm{y - \underbrace{(z + \norm{y-z}f)}_{\in F}} \geq \frac{d}{\norm{y-z}} \geq \theta.
\]
\end{proof}

\begin{corollaire}\label{cor:sous-espace-propre}
\index{sous-espace!propre fermé}
Si $F$ est un sous-espace propre fermé d'un espace de Banach $E$, alors $F$ est d'intérieur vide.
\end{corollaire}

\begin{proof}
Si $F$ contenait une boule ouverte $B(x_0, r)$, alors pour tout $y \in E$ avec $\norm{y} < r$, on aurait $x_0 + y \in F$, donc $y \in F$, et par homogénéité $F = E$ : contradiction.
\end{proof}

\begin{theoreme}[Sous-espace de dimension finie]\label{thm:dim-finie-ferme}
\index{dimension finie!sous-espace fermé}
Tout sous-espace vectoriel de dimension finie d'un e.v.n. est fermé.
\end{theoreme}

\begin{proof}
Soit $F$ un sous-espace de dimension $n$ de l'e.v.n. $E$, et soit $(e_1, \ldots, e_n)$ une base de $F$. L'application $\varphi : (\K^n, \norm{\cdot}_1) \to (F, \norm{\cdot}_E)$ définie par $\varphi(\alpha_1, \ldots, \alpha_n) = \sum_{i=1}^n \alpha_i e_i$ est un isomorphisme linéaire bicontinu (par le 
ef{thm:equiv-normes-dim-finie}). En particulier, $F$ est complet (car isomorphe à $\K^n$ complet), donc fermé dans $E$.
\end{proof}

\begin{remarque}\label{rem:dim-infinie-non-ferme}
\index{sous-espace!non fermé}
En dimension infinie, un sous-espace vectoriel n'est pas nécessairement fermé. Par exemple, l'espace $C([0,1])$ vu comme sous-espace de $L^2([0,1])$ est dense mais pas fermé (car $L^2 \neq C$).
\end{remarque}

\begin{proposition}[Critère de compacité dans $\K^n$]\label{prop:compacite-Kn}
\index{compacité!dans $\K^n$}\index{Heine-Borel, théorème de}
Dans $(\K^n, \norm{\cdot})$ (pour n'importe quelle norme), une partie est compacte si et seulement si elle est fermée et bornée (théorème de Heine-Borel).
\end{proposition}

%====================================================================
\section{Applications linéaires continues et norme d'opérateur}
%====================================================================

\begin{definition}[Application linéaire continue]\label{def:app-lin-continue}
\index{application linéaire continue}\index{opérateur!borné}
Soient $(E, \norm{\cdot}_E)$ et $(F, \norm{\cdot}_F)$ deux e.v.n. Une application linéaire $T : E \to F$ est dite \textbf{continue} (ou \textbf{bornée}) si l'une des conditions équivalentes suivantes est vérifiée :
\begin{enumerate}[label=(\roman*)]
    \item $T$ est continue en $0$ ;
    \item $T$ est continue sur $E$ ;
    \item $T$ est lipschitzienne\index{lipschitzien} ;
    \item Il existe $C \geq 0$ tel que $\norm{T(x)}_F \leq C \norm{x}_E$ pour tout $x \in E$.
\end{enumerate}
\end{definition}

\begin{proposition}\label{prop:equiv-continuite}
Les quatre conditions de la 
ef{def:app-lin-continue} sont équivalentes.
\end{proposition}

\begin{proof}
(iv) $\Rightarrow$ (iii) $\Rightarrow$ (ii) $\Rightarrow$ (i) sont immédiates. Montrons (i) $\Rightarrow$ (iv). Si $T$ est continue en $0$, pour $\varepsilon = 1$, il existe $\delta > 0$ tel que $\norm{x}_E < \delta \Rightarrow \norm{T(x)}_F < 1$. Pour $x \neq 0$, posons $y = \frac{\delta x}{2\norm{x}_E}$. Alors $\norm{y}_E = \delta/2 < \delta$, donc $\norm{T(y)}_F < 1$, c'est-à-dire $\frac{\delta}{2\norm{x}_E}\norm{T(x)}_F < 1$. D'où $\norm{T(x)}_F \leq \frac{2}{\delta}\norm{x}_E$.
\end{proof}

\begin{definition}[Norme d'opérateur]\label{def:norme-operateur}
\index{norme!d'opérateur}\index{opérateur!norme}
Soit $T : E \to F$ linéaire continue. La \textbf{norme d'opérateur} de $T$ est :
\[
    \norm{T}_{\mathcal{L}(E,F)} = \sup_{x \neq 0} \frac{\norm{T(x)}_F}{\norm{x}_E} = \sup_{\norm{x}_E \leq 1} \norm{T(x)}_F = \sup_{\norm{x}_E = 1} \norm{T(x)}_F.
\]
On note $\mathcal{L}(E,F)$\index{L(E,F)@$\mathcal{L}(E,F)$} l'espace des applications linéaires continues de $E$ dans $F$, et $\mathcal{L}(E) = \mathcal{L}(E,E)$. Le \textbf{dual topologique}\index{dual!topologique} de $E$ est $E' = \mathcal{L}(E, \K)$\index{E'@$E'$}.
\end{definition}

\begin{proposition}\label{prop:norme-op-est-norme}
$\norm{\cdot}_{\mathcal{L}(E,F)}$ est une norme sur $\mathcal{L}(E,F)$.
\end{proposition}

\begin{proof}
La séparation : si $\norm{T} = 0$, alors $T(x) = 0$ pour tout $x$, donc $T = 0$. L'homogénéité : $\norm{\lambda T} = \sup_{\norm{x}\leq 1}\norm{\lambda T(x)} = \abs{\lambda}\norm{T}$. L'inégalité triangulaire : $\norm{(T+S)(x)} \leq \norm{T(x)} + \norm{S(x)} \leq (\norm{T} + \norm{S})\norm{x}$.
\end{proof}

\begin{theoreme}[$\mathcal{L}(E,F)$ est Banach si $F$ est Banach]\label{thm:L-banach}
\index{L(E,F)@$\mathcal{L}(E,F)$!complétude}\index{complétude!de $\mathcal{L}(E,F)$}
Si $F$ est un espace de Banach, alors $\mathcal{L}(E,F)$ est un espace de Banach.
\end{theoreme}

\begin{proof}
Soit $(T_k)_{k\geq 1}$ une suite de Cauchy dans $\mathcal{L}(E,F)$.

\textbf{Étape 1 : définition de la limite.} Pour chaque $x \in E$ :
\[
    \norm{T_k(x) - T_j(x)}_F \leq \norm{T_k - T_j}\norm{x}_E \to 0.
\]
Donc $(T_k(x))_k$ est de Cauchy dans $F$ (complet). On pose $T(x) = \lim_{k\to\infty} T_k(x)$.

\textbf{Étape 2 : $T$ est linéaire.} Pour $\alpha, \beta \in \K$ et $x, y \in E$ :
\[
    T(\alpha x + \beta y) = \lim_{k\to\infty} T_k(\alpha x + \beta y) = \lim_{k\to\infty} (\alpha T_k(x) + \beta T_k(y)) = \alpha T(x) + \beta T(y).
\]

\textbf{Étape 3 : $T$ est continue et $T_k \to T$ en norme.} Soit $\varepsilon > 0$ et $K$ tel que $\norm{T_k - T_j} < \varepsilon$ pour $k, j \geq K$. Pour tout $x$ avec $\norm{x} \leq 1$ et $k \geq K$ :
\[
    \norm{T_k(x) - T_j(x)}_F \leq \varepsilon.
\]
En faisant $j \to \infty$ : $\norm{T_k(x) - T(x)}_F \leq \varepsilon$. Donc $\norm{T_k - T} \leq \varepsilon$ pour $k \geq K$. De plus, $\norm{T} \leq \norm{T - T_K} + \norm{T_K} \leq \varepsilon + \norm{T_K} < \infty$, donc $T \in \mathcal{L}(E,F)$.
\end{proof}

\begin{corollaire}\label{cor:dual-banach}
\index{dual!topologique!complétude}
Le dual topologique $E'$ de tout espace vectoriel normé $E$ est un espace de Banach.
\end{corollaire}

\begin{proof}
$E' = \mathcal{L}(E, \K)$ et $\K$ est complet.
\end{proof}

\begin{proposition}[Composition d'opérateurs]\label{prop:composition-op}
\index{opérateur!composition}
Soient $E, F, G$ des e.v.n., $T \in \mathcal{L}(E,F)$ et $S \in \mathcal{L}(F,G)$. Alors $S \circ T \in \mathcal{L}(E,G)$ et $\norm{S \circ T} \leq \norm{S}\norm{T}$.
\end{proposition}

\begin{proof}
$\norm{S(T(x))}_G \leq \norm{S}\norm{T(x)}_F \leq \norm{S}\norm{T}\norm{x}_E$.
\end{proof}

\begin{remarque}\label{rem:algebre-banach}
\index{algèbre de Banach}
En particulier, $\mathcal{L}(E)$ muni de la composition est une \textbf{algèbre de Banach}\index{algèbre de Banach} unitaire (avec $\Id_E$ comme unité) lorsque $E$ est un espace de Banach.
\end{remarque}

\begin{exemple}[Formes linéaires sur $\ell^p$]\label{ex:formes-lin-ell-p}
\index{dual!de $\ell^p$}\index{forme linéaire}
Pour $1 \leq p < \infty$, soit $y = (y_n) \in \ell^q$ (avec $\frac{1}{p} + \frac{1}{q} = 1$). L'application $\varphi_y : \ell^p \to \K$ définie par $\varphi_y(x) = \sum_{n=1}^\infty x_n y_n$ est une forme linéaire continue avec $\norm{\varphi_y} = \norm{y}_q$ (par l'inégalité de Hölder, et l'égalité se vérifie en choisissant $x_n = \abs{y_n}^{q-2}\overline{y_n}$ convenablement normalisé). On montrera plus tard que $(\ell^p)' \simeq \ell^q$ isométriquement.
\end{exemple}

\begin{exemple}[Opérateur de décalage]\label{ex:decalage}
\index{décalage à droite}\index{décalage à gauche}\index{opérateur!de décalage}
Sur $\ell^2$, le \textbf{décalage à droite} $R : (x_1, x_2, \ldots) \mapsto (0, x_1, x_2, \ldots)$ et le \textbf{décalage à gauche} $L : (x_1, x_2, \ldots) \mapsto (x_2, x_3, \ldots)$ sont des opérateurs bornés de norme $1$. On a $LR = \Id$ mais $RL \neq \Id$ (car $RL(x_1, x_2, \ldots) = (0, x_2, x_3, \ldots)$). Cela montre que l'algèbre $\mathcal{L}(\ell^2)$ n'est pas commutative.
\end{exemple}

\begin{proposition}[Série de Neumann]\label{prop:neumann}
\index{série de Neumann}\index{Neumann, série de}
Soit $E$ un espace de Banach et $T \in \mathcal{L}(E)$ avec $\norm{T} < 1$. Alors $\Id - T$ est inversible dans $\mathcal{L}(E)$ et :
\[
    (\Id - T)^{-1} = \sum_{n=0}^{\infty} T^n, \qquad \norm{(\Id - T)^{-1}} \leq \frac{1}{1 - \norm{T}}.
\]
\end{proposition}

\begin{proof}
La série $\sum T^n$ est absolument convergente car $\sum \norm{T^n} \leq \sum \norm{T}^n = \frac{1}{1-\norm{T}} < \infty$. Comme $\mathcal{L}(E)$ est un espace de Banach (
ef{thm:L-banach}), la série converge vers un $S \in \mathcal{L}(E)$. On vérifie :
\[
    (\Id - T)S_N = (\Id - T)\sum_{n=0}^N T^n = \Id - T^{N+1}.
\]
En faisant $N \to \infty$ : $(\Id - T)S = \Id$ (car $\norm{T^{N+1}} \leq \norm{T}^{N+1} \to 0$). De même $S(\Id - T) = \Id$.
\end{proof}

\begin{corollaire}[Stabilité des isomorphismes]\label{cor:stabilite-iso}
\index{isomorphisme!stabilité}\index{groupe des inversibles}
L'ensemble $GL(E) = \{T \in \mathcal{L}(E) : T \text{ inversible}\}$ est un ouvert de $\mathcal{L}(E)$, et l'application $T \mapsto T^{-1}$ est continue.
\end{corollaire}

\begin{proof}
Soit $T_0 \in GL(E)$ et $T \in \mathcal{L}(E)$ avec $\norm{T - T_0} < \norm{T_0^{-1}}^{-1}$. Alors $T = T_0(\Id - T_0^{-1}(T_0 - T))$ et $\norm{T_0^{-1}(T_0 - T)} \leq \norm{T_0^{-1}}\norm{T_0 - T} < 1$, donc $T$ est inversible par la série de Neumann. Pour la continuité de l'inversion, on utilise $T^{-1} - T_0^{-1} = T^{-1}(T_0 - T)T_0^{-1}$.
\end{proof}

%====================================================================
\section{Séries absolument convergentes et complétude}
%====================================================================

\begin{lemme}[Extraction de sous-suite de Cauchy]\label{lem:extraction-cauchy}
\index{suite de Cauchy!extraction de sous-suite}
Toute suite de Cauchy $(x_n)$ dans un e.v.n. admet une sous-suite $(x_{n_k})$ telle que $\norm{x_{n_{k+1}} - x_{n_k}} < 2^{-k}$ pour tout $k \geq 1$.
\end{lemme}

\begin{proof}
Par récurrence. On construit $(n_k)$ croissante : $n_1 = 1$, et pour $k \geq 1$, comme $(x_n)$ est de Cauchy, il existe $N_k$ tel que $\norm{x_m - x_n} < 2^{-k}$ pour $m, n \geq N_k$. On pose $n_{k+1} = \max(n_k + 1, N_k)$.
\end{proof}

\begin{remarque}\label{rem:cv-sous-suite-cauchy}
\index{suite de Cauchy!et sous-suite convergente}
Une suite de Cauchy qui admet une sous-suite convergente est elle-même convergente. En effet, si $(x_{n_k}) \to x$ et $(x_n)$ est de Cauchy, alors pour $\varepsilon > 0$, on choisit $K$ tel que $\norm{x_{n_k} - x} < \varepsilon/2$ pour $k \geq K$ et $N$ tel que $\norm{x_m - x_n} < \varepsilon/2$ pour $m, n \geq N$. Pour $n \geq \max(N, n_K)$, on choisit $k$ tel que $n_k \geq n$ et $k \geq K$ : $\norm{x_n - x} \leq \norm{x_n - x_{n_k}} + \norm{x_{n_k} - x} < \varepsilon$.
\end{remarque}

\begin{theoreme}[Caractérisation de la complétude]\label{thm:abs-cv-completude}
\index{complétude!caractérisation par les séries}\index{série!absolument convergente!et complétude}
Un espace vectoriel normé $(E, \norm{\cdot})$ est un espace de Banach si et seulement si toute série absolument convergente dans $E$ est convergente.
\end{theoreme}

\begin{proof}
$(\Rightarrow)$ Supposons $E$ complet et soit $\sum x_n$ absolument convergente. La suite des sommes partielles est de Cauchy car pour $N > M$ :
\[
    \norm{S_N - S_M} = \norm*{\sum_{n=M+1}^{N} x_n} \leq \sum_{n=M+1}^{N} \norm{x_n} \to 0
\]
puisque la série $\sum \norm{x_n}$ converge. Comme $E$ est complet, $(S_N)$ converge.

$(\Leftarrow)$ Supposons que toute série absolument convergente est convergente. Soit $(y_k)$ une suite de Cauchy dans $E$. Extrayons une sous-suite $(y_{k_j})$ telle que $\norm{y_{k_{j+1}} - y_{k_j}} < 2^{-j}$ pour tout $j$. Posons $x_1 = y_{k_1}$ et $x_j = y_{k_j} - y_{k_{j-1}}$ pour $j \geq 2$. Alors :
\[
    \sum_{j=2}^{\infty} \norm{x_j} \leq \sum_{j=1}^{\infty} 2^{-j} = 1 < \infty.
\]
Donc $\sum x_j$ est absolument convergente, et par hypothèse, elle converge vers un $y \in E$. Or $\sum_{j=1}^{N} x_j = y_{k_N}$, donc $y_{k_N} \to y$. Comme $(y_k)$ est de Cauchy et admet une sous-suite convergente, la suite entière converge vers $y$.
\end{proof}

\begin{exemple}[Application : $\ell^p$ est Banach, preuve alternative]\label{ex:ell-p-banach-alt}
\index{espace!$\ell^p$!complétude (preuve alternative)}
On peut utiliser le 
ef{thm:abs-cv-completude} pour redémontrer la complétude de $\ell^p$. Soit $\sum x^{(k)}$ absolument convergente dans $\ell^p$. Alors pour chaque $n$, $\sum_k \abs{x_n^{(k)}} \leq \sum_k \norm{x^{(k)}}_p < \infty$ par monotonie. Donc $s_n = \sum_k x_n^{(k)}$ existe. On vérifie que $s = (s_n) \in \ell^p$ par convergence dominée (discrète).
\end{exemple}

%====================================================================
\section{Espaces quotients et complétion}
%====================================================================

\begin{definition}[Espace quotient]\label{def:espace-quotient}
\index{espace quotient}\index{norme!quotient}
Soit $E$ un e.v.n. et $F$ un sous-espace vectoriel fermé de $E$. L'\textbf{espace quotient} $E/F$ est l'ensemble des classes $\dot{x} = x + F = \{x + f : f \in F\}$, muni de la norme quotient :
\[
    \norm{\dot{x}}_{E/F} = \inf_{f \in F} \norm{x + f}_E = d(x, F).
\]
\end{definition}

\begin{proposition}\label{prop:quotient-norme}
La norme quotient est bien une norme sur $E/F$ (dès que $F$ est fermé). De plus, si $E$ est un espace de Banach, alors $E/F$ est un espace de Banach.
\end{proposition}

\begin{proof}
\textbf{C'est une norme.} L'homogénéité et l'inégalité triangulaire sont immédiates. Pour la séparation : $\norm{\dot{x}} = 0$ signifie $d(x, F) = 0$, donc $x \in \overline{F} = F$ (car $F$ est fermé), donc $\dot{x} = \dot{0}$.

\textbf{Complétude.} Supposons $E$ Banach. Par le 
ef{thm:abs-cv-completude}, il suffit de montrer que toute série absolument convergente dans $E/F$ converge. Soit $\sum \dot{x}_n$ avec $\sum \norm{\dot{x}_n}_{E/F} < \infty$. Pour chaque $n$, choisissons $y_n \in \dot{x}_n$ tel que $\norm{y_n}_E \leq \norm{\dot{x}_n}_{E/F} + 2^{-n}$. Alors $\sum \norm{y_n}_E < \infty$, donc $\sum y_n$ converge vers un $y \in E$ (car $E$ est Banach). On vérifie que $\sum_{n=1}^N \dot{x}_n = \dot{S}_N$ avec $S_N = \sum_{n=1}^N y_n$, et $\norm{\dot{S}_N - \dot{y}}_{E/F} \leq \norm{S_N - y}_E \to 0$.
\end{proof}

\begin{definition}[Complétion]\label{def:completion}
\index{complétion}
Soit $(E, \norm{\cdot})$ un e.v.n. Une \textbf{complétion} de $E$ est un couple $(\hat{E}, \iota)$ où $\hat{E}$ est un espace de Banach et $\iota : E \hookrightarrow \hat{E}$ est une isométrie\index{isométrie} linéaire d'image dense dans $\hat{E}$.
\end{definition}

\begin{theoreme}[Existence et unicité de la complétion]\label{thm:completion}
\index{complétion!existence et unicité}
Tout espace vectoriel normé admet une complétion, unique à isométrie linéaire près.
\end{theoreme}

\begin{proof}[Esquisse de preuve]
\textbf{Existence.} On considère l'espace $\mathcal{C}$ des suites de Cauchy dans $E$, muni des opérations terme à terme, et on quotiente par la relation $\sim$ définie par $(x_n) \sim (y_n) \iff \norm{x_n - y_n} \to 0$. L'espace $\hat{E} = \mathcal{C}/{\sim}$ est muni de la norme $\norm{[(x_n)]} = \lim_{n\to\infty} \norm{x_n}$ (cette limite existe car $\abs{\norm{x_n} - \norm{x_m}} \leq \norm{x_n - x_m} \to 0$). L'application $\iota : x \mapsto [(x, x, x, \ldots)]$ est une isométrie linéaire d'image dense.

\textbf{Unicité.} Si $(\hat{E}_1, \iota_1)$ et $(\hat{E}_2, \iota_2)$ sont deux complétions, l'application $\iota_2 \circ \iota_1^{-1} : \iota_1(E) \to \iota_2(E)$ est une isométrie linéaire entre sous-espaces denses, qui se prolonge de manière unique en une isométrie linéaire $\hat{E}_1 \to \hat{E}_2$.
\end{proof}

\begin{proposition}[Propriété universelle de la complétion]\label{prop:completion-universelle}
\index{complétion!propriété universelle}
Soit $(\hat{E}, \iota)$ une complétion de l'e.v.n. $E$ et $F$ un espace de Banach. Toute application linéaire continue $T : E \to F$ admet un unique prolongement linéaire continu $\hat{T} : \hat{E} \to F$ tel que $\hat{T} \circ \iota = T$. De plus $\norm{\hat{T}} = \norm{T}$.
\end{proposition}

\begin{proof}
Pour $\hat{x} \in \hat{E}$, choisissons une suite $(x_n) \subset E$ telle que $\iota(x_n) \to \hat{x}$. Alors $(T(x_n))$ est de Cauchy dans $F$ car $\norm{T(x_m) - T(x_n)} \leq \norm{T}\norm{x_m - x_n}$. Comme $F$ est complet, on pose $\hat{T}(\hat{x}) = \lim T(x_n)$. La limite ne dépend pas du choix de la suite (car si $\iota(x_n'), \iota(x_n)$ convergent vers $\hat{x}$, on peut les entrelacer). La linéarité et la continuité de $\hat{T}$ sont claires, et $\norm{\hat{T}} = \norm{T}$ car $\iota(E)$ est dense.
\end{proof}

Le diagramme suivant illustre cette propriété universelle :
\begin{center}
\begin{tikzcd}
    E \arrow[r, hook, "\iota"] \arrow[dr, "T"'] & \hat{E} \arrow[d, dashed, "\hat{T}"] \\
    & F
\end{tikzcd}
\captionof{figure}{Propriété universelle de la complétion : toute application linéaire continue $T : E \to F$ ($F$ Banach) se prolonge en $\hat{T} : \hat{E} \to F$ avec $\norm{\hat{T}} = \norm{T}$.}\label{fig:completion}
\end{center}

\begin{exemple}[{Complétion de $C([0,1])$ pour différentes normes}]\label{ex:completion-C01}
\index{complétion!exemples}
L'espace $C([0,1])$ admet différentes complétions selon la norme choisie :
\begin{itemize}
    \item Pour la norme $\norm{\cdot}_\infty$ : il est déjà complet, $\hat{E} = C([0,1])$.
    \item Pour la norme $\norm{f}_p = \left(\int_0^1 \abs{f(t)}^p\,\dd t\right)^{1/p}$ ($1 \leq p < \infty$) : la complétion est $L^p([0,1])$.
    \item Pour la norme $\norm{f}_{H^1} = \left(\int_0^1 \abs{f}^2 + \abs{f'}^2\,\dd t\right)^{1/2}$ (restreinte aux fonctions $C^1$) : la complétion est l'espace de Sobolev $H^1([0,1])$\index{Sobolev, espace de}.
\end{itemize}
Cet exemple illustre que la complétion dépend crucialement de la norme.
\end{exemple}

%====================================================================
\section{Exercices}
%====================================================================

\begin{exercice}[$\star$]\label{exo:ch1-norme-verif}
\index{norme!vérification}
Montrer que les applications suivantes sont des normes sur l'espace indiqué :
\begin{enumerate}[label=(\alph*)]
    \item $\norm{f} = \int_0^1 \abs{f(t)} \, \dd t$ sur $C([0,1])$ ;
    \item $\norm{(x,y)} = \abs{x} + 2\abs{y}$ sur $\R^2$ ;
    \item $\norm{A} = \max_{1\leq i\leq n} \sum_{j=1}^n \abs{a_{ij}}$ sur $M_n(\R)$.
\end{enumerate}
\end{exercice}

\begin{exercice}[$\star$]\label{exo:ch1-non-complet}
\index{espace!non complet}
Montrer que $(C([0,1]), \norm{\cdot}_1)$ avec $\norm{f}_1 = \int_0^1 \abs{f(t)}\, \dd t$ n'est \emph{pas} un espace de Banach. \textit{Indication : considérer la suite $f_n(t) = \min(1, \max(0, n(t - 1/2)))$.}
\end{exercice}

\begin{exercice}[$\star\star$]\label{exo:ch1-sous-espace-ferme}
\index{sous-espace!fermé}
Montrer que $c_0 = \{(x_n) \in \ell^\infty : x_n \to 0\}$ est un sous-espace fermé de $\ell^\infty$. En déduire que $c_0$ est un espace de Banach.
\end{exercice}

\begin{exercice}[$\star\star$]\label{exo:ch1-normes-equiv}
\index{normes équivalentes}
Soient $\norm{\cdot}_a$ et $\norm{\cdot}_b$ deux normes sur un espace vectoriel $E$. Montrer qu'elles sont équivalentes si et seulement si elles définissent les mêmes ouverts.
\end{exercice}

\begin{exercice}[$\star\star$]\label{exo:ch1-riesz-dim-infinie}
Soit $E$ un espace vectoriel normé de dimension infinie. En utilisant le lemme de Riesz, construire une suite $(x_n)$ dans $S_E$ telle que $\norm{x_m - x_n} \geq 1/2$ pour $m \neq n$.
\end{exercice}

\begin{exercice}[$\star\star$]\label{exo:ch1-operateur-exemple}
\index{opérateur!exemple}
Soit $T : \ell^1 \to \ell^1$ défini par $T(x_1, x_2, x_3, \ldots) = (0, x_1, x_2, \ldots)$ (décalage à droite\index{décalage à droite}). Montrer que $T \in \mathcal{L}(\ell^1)$ et calculer $\norm{T}$.
\end{exercice}

\begin{exercice}[$\star\star$]\label{exo:ch1-series-operateurs}
\index{série!d'opérateurs}
Soit $E$ un espace de Banach et $T \in \mathcal{L}(E)$ avec $\norm{T} < 1$. Montrer que $\Id - T$ est inversible et que $(\Id - T)^{-1} = \sum_{n=0}^{\infty} T^n$ (série de Neumann\index{série de Neumann}).
\end{exercice}

\begin{exercice}[$\star\star\star$]\label{exo:ch1-completion-explicite}
\index{complétion!construction explicite}
Soit $E = C([0,1])$ muni de la norme $\norm{f}_1 = \int_0^1 \abs{f(t)}\,\dd t$. Montrer que la complétion de $E$ est (isométriquement isomorphe à) $L^1([0,1])$.
\end{exercice}

\begin{exercice}[$\star\star\star$]\label{exo:ch1-quotient}
Soit $E = \ell^\infty$ et $F = c_0$. Décrire l'espace quotient $\ell^\infty / c_0$. Montrer que la norme quotient de la classe de la suite constante $(1,1,1,\ldots)$ vaut $1$.
\end{exercice}

\begin{exercice}[$\star$]\label{exo:ch1-inclusion-ell-p}
\index{espace!$\ell^p$!inclusions}
Montrer que si $1 \leq p \leq q \leq \infty$, alors $\ell^p \subset \ell^q$ et $\norm{x}_q \leq \norm{x}_p$.
\end{exercice}

\begin{exercice}[$\star\star$]\label{exo:ch1-neumann-application}
\index{série de Neumann!application}
Soit $A \in M_n(\R)$ avec $\norm{A}_{\text{op}} < 1$ (norme d'opérateur subordonnée à la norme euclidienne). Montrer que la matrice $I - A$ est inversible et que $(I-A)^{-1} = \sum_{k=0}^\infty A^k$. Appliquer au cas $A = \begin{pmatrix} 0 & 1/2 \\ 1/3 & 0 \end{pmatrix}$.
\end{exercice}

\begin{exercice}[$\star\star\star$]\label{exo:ch1-dual-c0}
\index{dual!de $c_0$}
Montrer que le dual de $c_0$ est isométriquement isomorphe à $\ell^1$. \textit{Indication : pour $\varphi \in c_0'$, considérer $y_n = \varphi(e_n)$ où $e_n$ est le $n$-ième vecteur de la base canonique.}
\end{exercice}

\begin{exercice}[$\star\star$]\label{exo:ch1-compact-totborne}
\index{totalement borné}\index{compacité!et totale bornitude}
Soit $E$ un espace métrique. On dit que $E$ est \textbf{totalement borné}\index{totalement borné} si pour tout $\varepsilon > 0$, $E$ peut être recouvert par un nombre fini de boules de rayon $\varepsilon$. Montrer qu'un espace métrique est compact si et seulement s'il est complet et totalement borné.
\end{exercice}

\begin{exercice}[$\star\star$]\label{exo:ch1-densite-polynomes}
\index{Weierstrass, théorème de!approximation}
Soit $f \in C([0,1])$ et $\varepsilon > 0$. En admettant le théorème d'approximation de Weierstrass\index{Weierstrass, théorème de}, montrer qu'il existe un polynôme $P$ tel que $\norm{f - P}_\infty < \varepsilon$. En déduire que $C([0,1])$ est séparable pour la norme uniforme.
\end{exercice}

\begin{exercice}[$\star\star$]\label{exo:ch1-norme-image}
\index{norme!d'opérateur!caractérisation}
Soit $T \in \mathcal{L}(E, F)$. Montrer les caractérisations suivantes de $\norm{T}$ :
\begin{enumerate}[label=(\alph*)]
    \item $\norm{T} = \inf\{C \geq 0 : \norm{Tx} \leq C\norm{x} \text{ pour tout } x \in E\}$ ;
    \item $\norm{T} = \sup\{\norm{Tx} : x \in E, \norm{x} \leq 1\}$ ;
    \item $\norm{T} = \sup\{\norm{Tx} : x \in E, \norm{x} = 1\}$ (si $E \neq \{0\}$) ;
    \item $\norm{T} = \sup\left\{\frac{\norm{Tx}}{\norm{x}} : x \in E, x \neq 0\right\}$.
\end{enumerate}
Montrer que le supremum dans (b) n'est pas nécessairement atteint si $\dim E = \infty$.
\end{exercice}

%%%%%%%%%%%%%%%%%%%%%%%%%%%%%%%%%%%%%%%%%%%%%%%%%%%%%%%%%%%%%%%%%%%%%
%%%%%%%%%%%%%%%%%%%%%%%%%%%%%%%%%%%%%%%%%%%%%%%%%%%%%%%%%%%%%%%%%%%%%
%%  CHAPITRE 2 : ESPACES DE HILBERT
%%%%%%%%%%%%%%%%%%%%%%%%%%%%%%%%%%%%%%%%%%%%%%%%%%%%%%%%%%%%%%%%%%%%%
%%%%%%%%%%%%%%%%%%%%%%%%%%%%%%%%%%%%%%%%%%%%%%%%%%%%%%%%%%%%%%%%%%%%%

\chapter{Espaces de Hilbert}
\minitoc
\vspace{1cm}

Les espaces de Hilbert\index{espace de Hilbert} constituent une classe privilégiée d'espaces de Banach, dans laquelle la norme provient d'un produit scalaire. Cette structure supplémentaire permet de développer une théorie géométrique riche --- projection orthogonale, bases hilbertiennes, théorème de représentation de Riesz --- qui généralise la géométrie euclidienne à la dimension infinie. Les espaces de Hilbert interviennent de manière cruciale en mécanique quantique\index{mécanique quantique}, en traitement du signal\index{traitement du signal} et dans l'étude des équations aux dérivées partielles\index{équations aux dérivées partielles}.

La théorie des espaces de Hilbert a été développée par David Hilbert\index{Hilbert, David} au début du XXe siècle dans le contexte des équations intégrales\index{équation intégrale}. La formalisation abstraite est due à von Neumann\index{von Neumann, John} (1929), qui a défini la notion d'espace de Hilbert abstrait et établi les fondements de la théorie spectrale\index{théorie spectrale}. L'ouvrage de référence moderne est celui de Brezis\index{Brezis, Haïm} \emph{Analyse fonctionnelle}, dont nous suivons l'esprit dans ce chapitre.

Dans tout ce chapitre, $H$ désigne un espace vectoriel sur $\K \in \{\R, \C\}$ muni d'un produit scalaire, sauf mention contraire.

%====================================================================
\section{Produit scalaire et espaces préhilbertiens}
%====================================================================

\begin{definition}[Produit scalaire]\label{def:produit-scalaire}
\index{produit scalaire}
Soit $H$ un espace vectoriel sur $\K$. Un \textbf{produit scalaire} (ou \textbf{forme hermitienne définie positive}) sur $H$ est une application $\inner{\cdot}{\cdot} : H \times H \to \K$ vérifiant :
\begin{enumerate}[label=(\roman*)]
    \item \textbf{Linéarité à droite}\index{produit scalaire!linéarité} : pour tous $x, y, z \in H$ et $\lambda \in \K$,
    \[
        \inner{x}{\lambda y + z} = \lambda \inner{x}{y} + \inner{x}{z} ;
    \]
    \item \textbf{Symétrie hermitienne}\index{produit scalaire!symétrie hermitienne} : $\inner{x}{y} = \overline{\inner{y}{x}}$ pour tous $x, y \in H$ ;
    \item \textbf{Positivité}\index{produit scalaire!positivité} : $\inner{x}{x} \geq 0$ pour tout $x \in H$ ;
    \item \textbf{Définie}\index{produit scalaire!caractère défini} : $\inner{x}{x} = 0 \iff x = 0$.
\end{enumerate}
L'espace $(H, \inner{\cdot}{\cdot})$ est un \textbf{espace préhilbertien}\index{espace préhilbertien}.
\end{definition}

\begin{remarque}\label{rem:convention-lineaire}
\index{convention!linéarité du produit scalaire}
Notre convention est la linéarité à droite et l'antilinéarité à gauche (convention des physiciens). Dans le cas réel ($\K = \R$), la symétrie hermitienne se réduit à $\inner{x}{y} = \inner{y}{x}$ et le produit scalaire est bilinéaire. Certains auteurs (dont Brezis) adoptent la convention inverse (linéarité à gauche).
\end{remarque}

\begin{theoreme}[Inégalité de Cauchy-Schwarz]\label{thm:cauchy-schwarz}
\index{inégalité de Cauchy-Schwarz}\index{Cauchy-Schwarz, inégalité de}
Soit $(H, \inner{\cdot}{\cdot})$ un espace préhilbertien. Pour tous $x, y \in H$ :
\[
    \abs{\inner{x}{y}} \leq \sqrt{\inner{x}{x}} \cdot \sqrt{\inner{y}{y}}.
\]
L'égalité a lieu si et seulement si $x$ et $y$ sont linéairement dépendants.
\end{theoreme}

\begin{proof}
Si $y = 0$, les deux membres valent $0$ et le résultat est trivial. Supposons $y \neq 0$. Pour tout $\lambda \in \K$ :
\[
    0 \leq \inner{x - \lambda y}{x - \lambda y} = \inner{x}{x} - \lambda\inner{x}{y} - \overline{\lambda}\inner{y}{x} + \abs{\lambda}^2\inner{y}{y}.
\]
Choisissons $\lambda = \frac{\inner{y}{x}}{\inner{y}{y}}$ (qui minimise l'expression). On obtient :
\begin{align*}
    0 &\leq \inner{x}{x} - \frac{\inner{y}{x}}{\inner{y}{y}}\inner{x}{y} - \frac{\overline{\inner{y}{x}}}{\inner{y}{y}}\inner{y}{x} + \frac{\abs{\inner{y}{x}}^2}{\inner{y}{y}^2}\inner{y}{y} \\
    &= \inner{x}{x} - \frac{\abs{\inner{x}{y}}^2}{\inner{y}{y}} - \frac{\abs{\inner{x}{y}}^2}{\inner{y}{y}} + \frac{\abs{\inner{x}{y}}^2}{\inner{y}{y}} \\
    &= \inner{x}{x} - \frac{\abs{\inner{x}{y}}^2}{\inner{y}{y}}.
\end{align*}
Donc $\abs{\inner{x}{y}}^2 \leq \inner{x}{x}\inner{y}{y}$, d'où le résultat en prenant la racine carrée.

\textbf{Cas d'égalité.} L'égalité a lieu si et seulement si $\inner{x - \lambda y}{x - \lambda y} = 0$, c'est-à-dire $x = \lambda y$.
\end{proof}

\begin{corollaire}\label{cor:norme-prehilbert}
\index{norme!associée au produit scalaire}
L'application $\norm{x} = \sqrt{\inner{x}{x}}$ est une norme sur $H$.
\end{corollaire}

\begin{proof}
La séparation et l'homogénéité sont claires. Pour l'inégalité triangulaire :
\begin{align*}
    \norm{x+y}^2 &= \inner{x+y}{x+y} = \norm{x}^2 + \inner{x}{y} + \inner{y}{x} + \norm{y}^2 \\
    &= \norm{x}^2 + 2\operatorname{Re}\inner{x}{y} + \norm{y}^2 \\
    &\leq \norm{x}^2 + 2\abs{\inner{x}{y}} + \norm{y}^2 \leq \norm{x}^2 + 2\norm{x}\norm{y} + \norm{y}^2 = (\norm{x}+\norm{y})^2.
\end{align*}
\end{proof}

\begin{proposition}[Identité de polarisation]\label{prop:polarisation}
\index{identité de polarisation}\index{polarisation, identité de}
Soit $(H, \inner{\cdot}{\cdot})$ un espace préhilbertien. Le produit scalaire se retrouve à partir de la norme :
\begin{itemize}
    \item Si $\K = \R$ :
    \[
        \inner{x}{y} = \frac{1}{4}\left(\norm{x+y}^2 - \norm{x-y}^2\right) ;
    \]
    \item Si $\K = \C$ :
    \[
        \inner{x}{y} = \frac{1}{4}\sum_{k=0}^{3} i^k \norm{x + i^k y}^2.
    \]
\end{itemize}
\end{proposition}

\begin{proof}
Cas réel : $\norm{x+y}^2 = \norm{x}^2 + 2\inner{x}{y} + \norm{y}^2$ et $\norm{x-y}^2 = \norm{x}^2 - 2\inner{x}{y} + \norm{y}^2$. La différence donne $4\inner{x}{y}$.

Cas complexe : on développe $\norm{x + i^k y}^2 = \norm{x}^2 + 2\operatorname{Re}(\overline{i^k}\inner{x}{y}) + \norm{y}^2$ et on somme pour $k = 0, 1, 2, 3$. Les termes $\norm{x}^2 + \norm{y}^2$ donnent $4(\norm{x}^2 + \norm{y}^2)$. Pour la partie croisée : $\sum_{k=0}^3 i^k \cdot 2\operatorname{Re}(\overline{i^k}\inner{x}{y})$. En écrivant $\inner{x}{y} = a + ib$, un calcul direct montre que la somme vaut $4(a + ib) = 4\inner{x}{y}$.
\end{proof}

\begin{proposition}[Identité du parallélogramme]\label{prop:parallelogramme}
\index{identité du parallélogramme}\index{parallélogramme, identité du}
Dans un espace préhilbertien :
\[
    \norm{x + y}^2 + \norm{x - y}^2 = 2(\norm{x}^2 + \norm{y}^2).
\]
Réciproquement, toute norme vérifiant cette identité provient d'un produit scalaire (via l'identité de polarisation).
\end{proposition}

\begin{proof}
Le sens direct résulte du développement :
\[
    \norm{x+y}^2 + \norm{x-y}^2 = 2\norm{x}^2 + 2\operatorname{Re}\inner{x}{y} + 2\norm{y}^2 - 2\operatorname{Re}\inner{x}{y} = 2\norm{x}^2 + 2\norm{y}^2.
\]
La réciproque (théorème de Jordan-von Neumann\index{Jordan-von Neumann, théorème de}) nécessite de vérifier que la forme définie par l'identité de polarisation est bien un produit scalaire, ce qui est technique mais classique.
\end{proof}

\begin{center}
\begin{tikzpicture}[scale=1.2]
    \coordinate (O) at (0,0);
    \coordinate (X) at (3,0.5);
    \coordinate (Y) at (1,2);
    \coordinate (XY) at ($(X)+(Y)$);
    \coordinate (XmY) at ($(X)-(Y)$);
    \draw[thick,blue,-{Stealth}] (O) -- (X) node[midway,below] {$x$};
    \draw[thick,red,-{Stealth}] (O) -- (Y) node[midway,left] {$y$};
    \draw[thick,green!60!black,-{Stealth}] (O) -- (XY) node[midway,above left] {$x+y$};
    \draw[thick,orange,-{Stealth}] (Y) -- (XY);
    \draw[thick,orange,-{Stealth}] (X) -- (XY);
    \draw[dashed,purple] (O) -- (XmY) node[right] {$x-y$};
    \draw[dashed,purple] (Y) -- (XmY);
\end{tikzpicture}
\captionof{figure}{Identité du parallélogramme : la somme des carrés des diagonales égale la somme des carrés des côtés.}\label{fig:parallelogramme}
\end{center}

%====================================================================
\section{Espaces de Hilbert : définition et exemples}
%====================================================================

\begin{definition}[Espace de Hilbert]\label{def:hilbert}
\index{espace de Hilbert}\index{Hilbert, David}
Un \textbf{espace de Hilbert} est un espace préhilbertien complet, c'est-à-dire un espace de Banach dont la norme provient d'un produit scalaire.
\end{definition}

\begin{exemple}[L'espace $\ell^2$]\label{ex:ell2-hilbert}
\index{espace!$\ell^2$}\index{ell-2@$\ell^2$}
L'espace $\ell^2$ muni du produit scalaire
\[
    \inner{x}{y}_{\ell^2} = \sum_{n=1}^{\infty} \overline{x_n} y_n
\]
est un espace de Hilbert. La convergence de la série est assurée par l'inégalité de Cauchy-Schwarz.
\end{exemple}

\begin{exemple}[L'espace $L^2(\Omega, \mu)$]\label{ex:L2-hilbert}
\index{espace!$L^2$}\index{L2@$L^2$}
Soit $(\Omega, \mathcal{A}, \mu)$ un espace mesuré. L'espace $L^2(\Omega, \mu)$ muni de
\[
    \inner{f}{g}_{L^2} = \int_\Omega \overline{f} \, g \, \dd\mu
\]
est un espace de Hilbert (le théorème de Fischer-Riesz, 
ef{thm:fischer-riesz}).
\end{exemple}

\begin{exemple}[L'espace $\C^n$]\label{ex:Cn-hilbert}
\index{espace!$\C^n$!comme espace de Hilbert}
L'espace $\C^n$ muni du produit scalaire $\inner{x}{y} = \sum_{k=1}^n \overline{x_k}y_k$ est un espace de Hilbert de dimension finie.
\end{exemple}

\begin{exemple}[Aperçu des espaces de Sobolev]\label{ex:sobolev-preview}
\index{espace de Sobolev}\index{Sobolev, espace de}
L'espace de Sobolev $H^1(\Omega) = W^{1,2}(\Omega)$\index{H1@$H^1$} est l'espace des fonctions $u \in L^2(\Omega)$ dont les dérivées partielles au sens des distributions appartiennent à $L^2(\Omega)$, muni du produit scalaire :
\[
    \inner{u}{v}_{H^1} = \int_\Omega \overline{u}\, v \, \dd x + \int_\Omega \nabla \overline{u} \cdot \nabla v \, \dd x.
\]
C'est un espace de Hilbert, fondamental pour l'étude variationnelle des EDP elliptiques\index{EDP elliptique}.
\end{exemple}

\begin{remarque}\label{rem:lp-pas-hilbert}
\index{espace!$L^p$!non hilbertien pour $p \neq 2$}
Pour $p \neq 2$, l'espace $L^p$ n'est \emph{pas} un espace de Hilbert : sa norme ne satisfait pas l'identité du parallélogramme. On peut le vérifier explicitement dans $\ell^p$ avec $x = e_1$ et $y = e_2$ : $\norm{x+y}_p^2 + \norm{x-y}_p^2 = 2 \cdot 2^{2/p}$ tandis que $2(\norm{x}^2 + \norm{y}^2) = 4$, et ces quantités ne sont égales que si $p = 2$.
\end{remarque}

\begin{exemple}[L'espace $L^2$ des séries de Fourier]\label{ex:L2-fourier}
\index{série de Fourier}\index{espace!$L^2([0,2\pi])$}
L'espace $L^2([0,2\pi])$ est fondamental en analyse harmonique\index{analyse harmonique}. Toute fonction $f \in L^2([0,2\pi])$ admet un développement en série de Fourier :
\[
    f(t) = \sum_{n \in \Z} c_n(f) e^{int}, \quad c_n(f) = \frac{1}{2\pi}\int_0^{2\pi} f(t)e^{-int}\,\dd t,
\]
où la convergence est en norme $L^2$. L'identité de Parseval s'écrit :
\[
    \frac{1}{2\pi}\int_0^{2\pi}\abs{f(t)}^2\,\dd t = \sum_{n \in \Z}\abs{c_n(f)}^2.
\]
\end{exemple}

\begin{exemple}[Espace de Hardy]\label{ex:hardy}
\index{espace de Hardy}\index{Hardy, espace de}
L'espace de Hardy $H^2(\mathbb{D})$\index{H2@$H^2(\mathbb{D})$} est l'espace des fonctions holomorphes $f$ sur le disque unité $\mathbb{D} = \{z \in \C : \abs{z} < 1\}$ telles que
\[
    \norm{f}_{H^2}^2 = \sup_{0 < r < 1} \frac{1}{2\pi}\int_0^{2\pi}\abs{f(re^{i\theta})}^2\,\dd\theta < \infty.
\]
C'est un espace de Hilbert, isomorphe à $\ell^2(\N)$ via les coefficients de Taylor : si $f(z) = \sum_{n=0}^\infty a_n z^n$, alors $\norm{f}_{H^2}^2 = \sum_{n=0}^\infty \abs{a_n}^2$.
\end{exemple}

\begin{definition}[Orthogonalité]\label{def:orthogonalite}
\index{orthogonalité}
Deux vecteurs $x, y \in H$ sont \textbf{orthogonaux}, noté $x \perp y$\index{perp@$\perp$}, si $\inner{x}{y} = 0$. L'\textbf{orthogonal} d'une partie $A \subset H$ est
\[
    A^\perp = \{x \in H : \inner{x}{a} = 0 \text{ pour tout } a \in A\}.
\]
\end{definition}

\begin{proposition}\label{prop:orthogonal-ferme}
\index{orthogonal!sous-espace fermé}
Pour toute partie $A \subset H$, $A^\perp$ est un sous-espace vectoriel fermé de $H$.
\end{proposition}

\begin{proof}
La linéarité est claire. Pour la fermeture : si $(x_n) \subset A^\perp$ et $x_n \to x$, alors pour tout $a \in A$, $\inner{x}{a} = \lim_{n\to\infty}\inner{x_n}{a} = 0$ par continuité du produit scalaire\index{produit scalaire!continuité}.
\end{proof}

\begin{proposition}[Théorème de Pythagore]\label{prop:pythagore}
\index{Pythagore, théorème de}\index{théorème de Pythagore}
Si $x_1, \ldots, x_n$ sont deux à deux orthogonaux dans $H$ :
\[
    \norm{x_1 + \cdots + x_n}^2 = \norm{x_1}^2 + \cdots + \norm{x_n}^2.
\]
\end{proposition}

\begin{proof}
Par récurrence. Pour $n = 2$ : $\norm{x_1 + x_2}^2 = \norm{x_1}^2 + 2\operatorname{Re}\inner{x_1}{x_2} + \norm{x_2}^2 = \norm{x_1}^2 + \norm{x_2}^2$ car $\inner{x_1}{x_2} = 0$.
\end{proof}

%====================================================================
\section{Projection sur un convexe fermé}
%====================================================================

\begin{theoreme}[Projection sur un convexe fermé]\label{thm:proj-convexe}
\index{projection!sur un convexe fermé}\index{convexe fermé!projection}
Soit $H$ un espace de Hilbert, $C \subset H$ un sous-ensemble convexe\index{convexe} fermé non vide, et $x \in H$. Il existe un unique $p \in C$ tel que
\[
    \norm{x - p} = d(x, C) = \inf_{y \in C} \norm{x - y}.
\]
Le point $p$ est appelé la \textbf{projection} de $x$ sur $C$, noté $p = P_C(x)$\index{PC@$P_C$}.
\end{theoreme}

\begin{proof}
Posons $d = d(x, C)$.

\textbf{Existence.} Soit $(y_n) \subset C$ une suite minimisante : $\norm{x - y_n} \to d$. L'identité du parallélogramme appliquée à $x - y_m$ et $x - y_n$ donne :
\[
    \norm{(x - y_m) + (x - y_n)}^2 + \norm{y_m - y_n}^2 = 2\norm{x - y_m}^2 + 2\norm{x - y_n}^2.
\]
Or $\frac{y_m + y_n}{2} \in C$ (par convexité), donc $\norm{x - \frac{y_m + y_n}{2}} \geq d$, c'est-à-dire $\norm{(x-y_m) + (x-y_n)} \geq 2d$. En fait, on a :
\[
    \norm{y_m - y_n}^2 = 2\norm{x - y_m}^2 + 2\norm{x - y_n}^2 - \norm{2x - y_m - y_n}^2
\]
\[
    = 2\norm{x - y_m}^2 + 2\norm{x - y_n}^2 - 4\norm*{x - \frac{y_m+y_n}{2}}^2 \leq 2\norm{x - y_m}^2 + 2\norm{x - y_n}^2 - 4d^2.
\]
Comme $\norm{x - y_m}^2 \to d^2$ et $\norm{x - y_n}^2 \to d^2$, on obtient $\norm{y_m - y_n}^2 \to 0$. Donc $(y_n)$ est de Cauchy dans $H$ (complet), et converge vers un $p \in C$ ($C$ est fermé) avec $\norm{x - p} = d$.

\textbf{Unicité.} Si $p_1, p_2 \in C$ réalisent le minimum, on applique le même argument du parallélogramme à $y_m = p_1$ et $y_n = p_2$ :
\[
    \norm{p_1 - p_2}^2 \leq 2d^2 + 2d^2 - 4d^2 = 0,
\]
donc $p_1 = p_2$.
\end{proof}

\begin{proposition}[Caractérisation de la projection]\label{prop:caract-projection}
\index{projection!caractérisation}
Avec les notations du théorème précédent, $p = P_C(x)$ si et seulement si :
\[
    p \in C \quad \text{et} \quad \operatorname{Re}\inner{x - p}{y - p} \leq 0 \quad \text{pour tout } y \in C.
\]
\end{proposition}

\begin{proof}
$(\Rightarrow)$ Soit $p = P_C(x)$ et $y \in C$. Pour $t \in [0,1]$, $p + t(y-p) = (1-t)p + ty \in C$, donc :
\[
    \norm{x - p}^2 \leq \norm{x - p - t(y-p)}^2 = \norm{x-p}^2 - 2t\operatorname{Re}\inner{x-p}{y-p} + t^2\norm{y-p}^2.
\]
D'où $2t\operatorname{Re}\inner{x-p}{y-p} \leq t^2\norm{y-p}^2$. En divisant par $t > 0$ et en faisant $t \to 0^+$ : $\operatorname{Re}\inner{x-p}{y-p} \leq 0$.

$(\Leftarrow)$ Si la condition est vérifiée, pour tout $y \in C$ :
\[
    \norm{x - y}^2 = \norm{(x-p) - (y-p)}^2 = \norm{x-p}^2 - 2\operatorname{Re}\inner{x-p}{y-p} + \norm{y-p}^2 \geq \norm{x-p}^2.
\]
Donc $p$ réalise le minimum.
\end{proof}

\begin{theoreme}[Projection sur un sous-espace fermé]\label{thm:proj-sous-espace}
\index{projection!sur un sous-espace fermé}\index{sous-espace!fermé!projection}
Soit $F$ un sous-espace vectoriel fermé de $H$ et $x \in H$. Alors $p = P_F(x)$ est l'unique élément de $F$ tel que $x - p \perp F$, c'est-à-dire :
\[
    p \in F \quad \text{et} \quad \inner{x - p}{f} = 0 \quad \text{pour tout } f \in F.
\]
De plus, $H = F \oplus F^\perp$ (somme directe orthogonale\index{somme directe orthogonale}).
\end{theoreme}

\begin{proof}
\textbf{Caractérisation.} La condition de la 
ef{prop:caract-projection} pour un sous-espace $F$ s'écrit : $\operatorname{Re}\inner{x-p}{y-p} \leq 0$ pour tout $y \in F$. Comme $F$ est un sous-espace, on peut remplacer $y$ par $p + \lambda f$ pour $\lambda \in \K$ et $f \in F$ quelconques. On obtient $\operatorname{Re}(\lambda\inner{x-p}{f}) \leq 0$ pour tout $\lambda$, ce qui force $\inner{x-p}{f} = 0$.

\textbf{Décomposition.} Pour tout $x \in H$, écrivons $x = P_F(x) + (x - P_F(x))$. Le premier terme est dans $F$ et le second dans $F^\perp$. Si $x \in F \cap F^\perp$, alors $\inner{x}{x} = 0$, donc $x = 0$. D'où $H = F \oplus F^\perp$.
\end{proof}

\begin{proposition}[Propriétés de la projection orthogonale]\label{prop:proj-orth-proprietes}
\index{projection orthogonale!propriétés}
Soit $F$ un sous-espace vectoriel fermé de $H$ et $P = P_F$ la projection orthogonale sur $F$. Alors :
\begin{enumerate}[label=(\roman*)]
    \item $P$ est linéaire et continue avec $\norm{P} = 1$ (si $F \neq \{0\}$) ;
    \item $P^2 = P$ (idempotence\index{idempotent}) ;
    \item $\inner{Px}{y} = \inner{x}{Py}$ pour tous $x, y \in H$ (auto-adjoint\index{auto-adjoint}) ;
    \item $\Ker P = F^\perp$ et $\ran P = F$.
\end{enumerate}
Réciproquement, tout opérateur $P \in \mathcal{L}(H)$ vérifiant $P^2 = P$ et $P^* = P$ est la projection orthogonale sur $\ran P$.
\end{proposition}

\begin{proof}
(i) La linéarité : pour $x, y \in H$ et $\alpha, \beta \in \K$, on vérifie que $\alpha Px + \beta Py \in F$ et que $(\alpha x + \beta y) - (\alpha Px + \beta Py) = \alpha(x - Px) + \beta(y - Py) \in F^\perp$. Par unicité de la décomposition, $P(\alpha x + \beta y) = \alpha Px + \beta Py$. La continuité : $\norm{Px}^2 + \norm{x - Px}^2 = \norm{x}^2$ (Pythagore), donc $\norm{Px} \leq \norm{x}$.

(ii) Si $x \in F$, alors $Px = x$, donc $P(Px) = Px$.

(iii) $\inner{Px}{y} = \inner{Px}{Py + (y - Py)} = \inner{Px}{Py}$ (car $Px \in F$ et $y - Py \in F^\perp$). De même $\inner{x}{Py} = \inner{Px}{Py}$.

(iv) Clair par construction.
\end{proof}

\begin{corollaire}\label{cor:biorthogonal}
\index{biorthogonal}
Pour tout sous-espace vectoriel $F$ de $H$ : $(F^\perp)^\perp = \overline{F}$.
\end{corollaire}

\begin{proof}
L'inclusion $\overline{F} \subset (F^\perp)^\perp$ est claire. Pour la réciproque, $G = \overline{F}$ est un sous-espace fermé et $H = G \oplus G^\perp$. Or $G^\perp = F^\perp$. Si $x \in (F^\perp)^\perp$, on écrit $x = g + h$ avec $g \in G$ et $h \in G^\perp = F^\perp$. Alors $\inner{x}{h} = 0$ (car $x \in (F^\perp)^\perp$ et $h \in F^\perp$), et $\inner{g}{h} = 0$ (car $h \in G^\perp$), donc $\norm{h}^2 = \inner{h}{h} = \inner{x-g}{h} = 0$, d'où $h = 0$ et $x = g \in G$.
\end{proof}

%====================================================================
\section{Théorème de représentation de Riesz}
%====================================================================

\begin{theoreme}[Représentation de Riesz]\label{thm:riesz-representation}
\index{Riesz, théorème de représentation de}\index{théorème de représentation de Riesz}
Soit $H$ un espace de Hilbert. Pour toute forme linéaire continue $\varphi \in H'$, il existe un unique $u \in H$ tel que
\[
    \varphi(x) = \inner{u}{x} \quad \text{pour tout } x \in H.
\]
De plus, $\norm{\varphi}_{H'} = \norm{u}_H$. L'application $J : H' \to H$ définie par $J(\varphi) = u$ est un isomorphisme isométrique antilinéaire\index{isomorphisme isométrique antilinéaire} (linéaire dans le cas réel).
\end{theoreme}

\begin{proof}
\textbf{Unicité.} Si $\inner{u_1}{x} = \inner{u_2}{x}$ pour tout $x$, alors $\inner{u_1 - u_2}{x} = 0$ pour tout $x$. En prenant $x = u_1 - u_2$ : $\norm{u_1 - u_2}^2 = 0$, donc $u_1 = u_2$.

\textbf{Existence.} Si $\varphi = 0$, on prend $u = 0$. Sinon, $F = \Ker\varphi$\index{noyau} est un sous-espace fermé (car $\varphi$ est continue) et $F \neq H$. Par le 
ef{thm:proj-sous-espace}, $F^\perp \neq \{0\}$, donc il existe $z \in F^\perp$ avec $z \neq 0$.

Pour tout $x \in H$, l'élément $w = \varphi(x)z - \varphi(z)x$ vérifie $\varphi(w) = \varphi(x)\varphi(z) - \varphi(z)\varphi(x) = 0$, donc $w \in F$. Puisque $z \perp F$ :
\[
    0 = \inner{z}{w} = \inner{z}{\varphi(x)z - \varphi(z)x} = \varphi(x)\norm{z}^2 - \varphi(z)\inner{z}{x}.
\]
D'où
\[
    \varphi(x) = \frac{\overline{\varphi(z)}}{\norm{z}^2}\inner*{\frac{\varphi(z)}{\overline{\varphi(z)}} \cdot z}{x} \cdot \frac{\overline{\varphi(z)}}{\overline{\varphi(z)}}.
\]
Simplifions. On a $\varphi(x) = \frac{\varphi(z)}{\norm{z}^2}\inner{z}{x}$. Posons $u = \frac{\overline{\varphi(z)}}{\norm{z}^2} z$. Alors :
\[
    \inner{u}{x} = \frac{\varphi(z)}{\norm{z}^2}\inner{z}{x} = \varphi(x).
\]

\textbf{Isométrie.} D'une part, $\abs{\varphi(x)} = \abs{\inner{u}{x}} \leq \norm{u}\norm{x}$, donc $\norm{\varphi} \leq \norm{u}$. D'autre part, $\varphi(u/\norm{u}) = \inner{u}{u/\norm{u}} = \norm{u}$, donc $\norm{\varphi} \geq \norm{u}$.
\end{proof}

\begin{remarque}\label{rem:riesz-identification}
\index{identification $H \simeq H'$}
Le théorème de Riesz permet d'identifier $H$ et $H'$ via l'application antilinéaire $x \mapsto \inner{x}{\cdot}$. On écrit souvent $H \simeq H'$, mais il faut être prudent avec cette identification dans le contexte des espaces de Sobolev et des triplets de Gelfand\index{triplet de Gelfand} $V \hookrightarrow H \simeq H' \hookrightarrow V'$.
\end{remarque}

Le théorème de Riesz fournit un diagramme commutatif fondamental :
\begin{center}
\begin{tikzcd}
    H \arrow[r, "J", "\sim"'] \arrow[dr, "\Id"'] & H' \arrow[d, "\text{éval}"] \\
    & H^{**}
\end{tikzcd}
\captionof{figure}{L'identification canonique $J : H \xrightarrow{\sim} H'$ et l'injection canonique dans le bidual. Tout espace de Hilbert est réflexif.}\label{fig:riesz-diag}
\end{center}

\begin{corollaire}[Réflexivité des espaces de Hilbert]\label{cor:hilbert-reflexif}
\index{réflexivité}\index{espace de Hilbert!réflexif}
Tout espace de Hilbert est réflexif, c'est-à-dire que l'injection canonique $J : H \to H''$ est surjective.
\end{corollaire}

\begin{proof}
L'application $J : H \to H''$ définie par $J(x)(\varphi) = \varphi(x)$ pour $\varphi \in H'$ est une isométrie (Hahn-Banach). Par Riesz appliqué deux fois, $H \simeq H' \simeq H''$, et on vérifie que la composition coïncide avec $J$.
\end{proof}

\begin{exemple}[Application du théorème de Riesz]\label{ex:riesz-application}
\index{Riesz, théorème de représentation de!application}
Soit $H = L^2([0,1])$ et $\varphi : H \to \K$ définie par $\varphi(f) = \int_0^1 e^t f(t)\,\dd t$. Cette forme est linéaire continue car $\abs{\varphi(f)} \leq \norm{e^t}_{L^2}\norm{f}_{L^2}$ par Cauchy-Schwarz. Par Riesz, $\varphi(f) = \inner{u}{f}$ avec $u(t) = e^t$. On a $\norm{\varphi} = \norm{u}_{L^2} = \left(\int_0^1 e^{2t}\,\dd t\right)^{1/2} = \sqrt{(e^2-1)/2}$.
\end{exemple}

\begin{corollaire}[Théorème de Lax-Milgram]\label{cor:lax-milgram}
\index{Lax-Milgram, théorème de}\index{théorème de Lax-Milgram}
Soit $H$ un espace de Hilbert et $a : H \times H \to \K$ une forme sesquilinéaire\index{forme sesquilinéaire} vérifiant :
\begin{enumerate}[label=(\roman*)]
    \item \textbf{Continuité} : $\abs{a(u,v)} \leq M\norm{u}\norm{v}$ pour une constante $M > 0$ ;
    \item \textbf{Coercivité}\index{coercivité} : $\operatorname{Re}(a(u,u)) \geq \alpha\norm{u}^2$ pour une constante $\alpha > 0$.
\end{enumerate}
Alors pour tout $\varphi \in H'$, il existe un unique $u \in H$ tel que $a(u,v) = \varphi(v)$ pour tout $v \in H$.
\end{corollaire}

\begin{proof}
Par Riesz, $\varphi(v) = \inner{f}{v}$ pour un unique $f \in H$. Pour $u$ fixé, $v \mapsto a(u,v)$ est continue, donc par Riesz il existe un unique $Au \in H$ tel que $a(u,v) = \inner{Au}{v}$. L'application $A : H \to H$ est linéaire avec $\norm{Au} \leq M\norm{u}$ (par (i)) et $\alpha\norm{u}^2 \leq \operatorname{Re}\inner{Au}{u} \leq \norm{Au}\norm{u}$, d'où $\norm{Au} \geq \alpha\norm{u}$ (par (ii)).

Ainsi $A$ est injectif et $\ran(A)$ est fermé (car si $Au_n \to y$, alors $(u_n)$ est de Cauchy par $\norm{u_n - u_m} \leq \alpha^{-1}\norm{Au_n - Au_m}$). De plus, $\ran(A)^\perp = \{0\}$ : si $\inner{Au}{w} = 0$ pour tout $u$, en prenant $u = w$ on obtient $\alpha\norm{w}^2 \leq \operatorname{Re}\inner{Aw}{w} = 0$, donc $w = 0$. Par le 
ef{cor:biorthogonal}, $\ran(A) = H$. Donc $A$ est bijectif et $u = A^{-1}f$ est l'unique solution.
\end{proof}

%====================================================================
\section{Bases hilbertiennes}
%====================================================================

\begin{definition}[Système orthonormal]\label{def:systeme-orthonormal}
\index{système orthonormal}\index{famille orthonormale}
Une famille $(e_i)_{i \in I}$ dans $H$ est un \textbf{système orthonormal} (ou \textbf{famille orthonormale}) si :
\[
    \inner{e_i}{e_j} = \delta_{ij} = \begin{cases} 1 & \text{si } i = j, \\ 0 & \text{si } i \neq j. \end{cases}
\]
\end{definition}

\begin{definition}[Base hilbertienne]\label{def:base-hilbertienne}
\index{base hilbertienne}\index{base orthonormale complète}
Un système orthonormal $(e_i)_{i \in I}$ est une \textbf{base hilbertienne} (ou \textbf{base orthonormale complète}) de $H$ si $\overline{\spn\{e_i : i \in I\}} = H$, c'est-à-dire si l'adhérence de l'espace engendré par les $e_i$ est $H$ tout entier.
\end{definition}

\begin{proposition}[Coefficients de Fourier]\label{prop:coeff-fourier}
\index{coefficients de Fourier}
Soit $(e_i)_{i \in I}$ un système orthonormal dans $H$ et $x \in H$. Les \textbf{coefficients de Fourier}\index{Fourier, coefficients de} de $x$ sont $\hat{x}_i = \inner{e_i}{x}$ pour $i \in I$. Pour tout sous-ensemble fini $J \subset I$, la meilleure approximation de $x$ dans $\spn\{e_j : j \in J\}$ est :
\[
    P_J(x) = \sum_{j \in J} \hat{x}_j e_j,
\]
et $\norm{x - P_J(x)}^2 = \norm{x}^2 - \sum_{j \in J} \abs{\hat{x}_j}^2$.
\end{proposition}

\begin{proof}
Pour $y = \sum_{j \in J} \alpha_j e_j$ :
\begin{align*}
    \norm{x - y}^2 &= \norm{x}^2 - \sum_{j \in J}\overline{\alpha_j}\inner{e_j}{x} - \sum_{j \in J}\alpha_j\inner{x}{e_j} + \sum_{j \in J}\abs{\alpha_j}^2 \\
    &= \norm{x}^2 - \sum_{j \in J}\abs{\hat{x}_j}^2 + \sum_{j \in J}\abs{\alpha_j - \hat{x}_j}^2.
\end{align*}
Ce minimum est atteint pour $\alpha_j = \hat{x}_j$, et vaut $\norm{x}^2 - \sum_{j \in J}\abs{\hat{x}_j}^2$.
\end{proof}

\begin{theoreme}[Inégalité de Bessel]\label{thm:bessel}
\index{inégalité de Bessel}\index{Bessel, inégalité de}
Soit $(e_i)_{i \in I}$ un système orthonormal dans $H$. Pour tout $x \in H$ :
\[
    \sum_{i \in I} \abs{\inner{e_i}{x}}^2 \leq \norm{x}^2.
\]
En particulier, l'ensemble $\{i \in I : \inner{e_i}{x} \neq 0\}$ est au plus dénombrable.
\end{theoreme}

\begin{proof}
Pour tout sous-ensemble fini $J \subset I$, la 
ef{prop:coeff-fourier} donne $\sum_{j \in J}\abs{\inner{e_j}{x}}^2 \leq \norm{x}^2$. En prenant le supremum sur tous les $J$ finis, on obtient $\sum_{i \in I}\abs{\inner{e_i}{x}}^2 \leq \norm{x}^2$.

Pour la dénombrabilité : pour $n \geq 1$, l'ensemble $I_n = \{i \in I : \abs{\inner{e_i}{x}} > 1/n\}$ est fini (car $\sum_{i \in I_n} 1/n^2 \leq \norm{x}^2$, donc $\abs{I_n} \leq n^2\norm{x}^2$). L'ensemble $\{i : \inner{e_i}{x} \neq 0\} = \bigcup_{n\geq 1} I_n$ est dénombrable.
\end{proof}

\begin{theoreme}[Théorème de Parseval]\label{thm:parseval}
\index{théorème de Parseval}\index{Parseval, théorème de}\index{identité de Parseval}
Soit $(e_n)_{n \geq 1}$ une base hilbertienne dénombrable de $H$. Pour tout $x \in H$ :
\begin{enumerate}[label=(\roman*)]
    \item $x = \sum_{n=1}^{\infty} \inner{e_n}{x} e_n$ \quad (convergence dans $H$) ;
    \item $\norm{x}^2 = \sum_{n=1}^{\infty} \abs{\inner{e_n}{x}}^2$ \quad (\textbf{identité de Parseval}).
\end{enumerate}
\end{theoreme}

\begin{proof}
\textbf{(i)} Soit $S_N = \sum_{n=1}^{N}\inner{e_n}{x}e_n = P_N(x)$ la projection sur $\spn\{e_1, \ldots, e_N\}$. Par Bessel, la série $\sum\abs{\inner{e_n}{x}}^2$ converge, donc pour $N > M$ :
\[
    \norm{S_N - S_M}^2 = \sum_{n=M+1}^{N}\abs{\inner{e_n}{x}}^2 \to 0.
\]
Donc $(S_N)$ est de Cauchy dans $H$ (complet), et converge vers un $y \in H$. Pour tout $m$ : $\inner{e_m}{y} = \lim_N \inner{e_m}{S_N} = \inner{e_m}{x}$. Donc $\inner{e_m}{x - y} = 0$ pour tout $m$, et $x - y \in \{e_n\}^\perp$. Comme $(e_n)$ est une base, $\{e_n\}^\perp = \{0\}$, donc $x = y$.

\textbf{(ii)} Par la 
ef{prop:coeff-fourier}, $\norm{x - S_N}^2 = \norm{x}^2 - \sum_{n=1}^N \abs{\inner{e_n}{x}}^2$. Comme $S_N \to x$, on fait $N \to \infty$.
\end{proof}

\begin{remarque}\label{rem:parseval-espace-non-separable}
\index{Parseval, théorème de!cas non séparable}
Le théorème de Parseval se généralise aux systèmes orthonormaux indexés par un ensemble $I$ quelconque (non nécessairement dénombrable). Pour $x \in H$, on a $x = \sum_{i \in I}\inner{e_i}{x}e_i$ où la somme est définie comme limite des sommes partielles finies (elle n'a qu'un nombre dénombrable de termes non nuls, par l'inégalité de Bessel). L'identité de Parseval reste valable : $\norm{x}^2 = \sum_{i \in I}\abs{\inner{e_i}{x}}^2$.
\end{remarque}

\begin{corollaire}[Isomorphisme avec $\ell^2$]\label{cor:iso-ell2}
\index{isomorphisme!avec $\ell^2$}
Si $H$ admet une base hilbertienne dénombrable $(e_n)_{n\geq 1}$, alors l'application $U : H \to \ell^2$ définie par $U(x) = (\inner{e_n}{x})_{n\geq 1}$ est un isomorphisme isométrique. En particulier, tous les espaces de Hilbert séparables de dimension infinie sont isomorphes.
\end{corollaire}

\begin{proof}
$U$ est linéaire (antilinéarité de la première variable, mais ici on prend les coefficients). Par Parseval, $\norm{U(x)}_{\ell^2} = \norm{x}_H$ : c'est une isométrie, donc injective. Surjectivité : pour $(\alpha_n) \in \ell^2$, posons $x = \sum_n \alpha_n e_n$ (converge car $\sum \abs{\alpha_n}^2 < \infty$). Alors $U(x) = (\alpha_n)$.
\end{proof}

%====================================================================
\section{Critères de complétude d'un système orthonormal}
%====================================================================

\begin{theoreme}[Critères de complétude]\label{thm:criteres-completude-son}
\index{base hilbertienne!critères de complétude}\index{système orthonormal!complétude}
Soit $(e_i)_{i \in I}$ un système orthonormal dans un espace de Hilbert $H$. Les conditions suivantes sont équivalentes :
\begin{enumerate}[label=(\roman*)]
    \item $(e_i)_{i \in I}$ est une base hilbertienne (c'est-à-dire $\overline{\spn\{e_i\}} = H$) ;
    \item $\{e_i\}^\perp = \{0\}$ (le système est \textbf{total}\index{système total}) ;
    \item Pour tout $x \in H$ : $x = \sum_{i \in I} \inner{e_i}{x} e_i$ ;
    \item Pour tout $x \in H$ : $\norm{x}^2 = \sum_{i \in I}\abs{\inner{e_i}{x}}^2$ (Parseval) ;
    \item Pour tous $x, y \in H$ : $\inner{x}{y} = \sum_{i \in I}\overline{\inner{e_i}{x}}\inner{e_i}{y}$ (Parseval généralisée\index{Parseval, identité de!généralisée}).
\end{enumerate}
\end{theoreme}

\begin{proof}
(i) $\Leftrightarrow$ (ii) : $\overline{\spn\{e_i\}} = H$ si et seulement si $(\spn\{e_i\})^\perp = \{0\}$ par le 
ef{cor:biorthogonal}, et $(\spn\{e_i\})^\perp = \{e_i\}^\perp$.

(ii) $\Rightarrow$ (iii) : C'est la preuve du 
ef{thm:parseval}(i).

(iii) $\Rightarrow$ (iv) : On prend la norme des deux côtés et on utilise Pythagore.

(iv) $\Rightarrow$ (v) : Par l'identité de polarisation appliquée dans $\ell^2$ via l'application $x \mapsto (\inner{e_i}{x})_{i \in I}$.

(v) $\Rightarrow$ (ii) : Si $\inner{e_i}{x} = 0$ pour tout $i$, alors par (v) avec $y = x$ : $\norm{x}^2 = 0$, donc $x = 0$.
\end{proof}

%====================================================================
\section{Espaces séparables et bases dénombrables}
%====================================================================

\begin{definition}[Espace séparable]\label{def:separable}
\index{espace séparable}\index{séparable}
Un espace topologique est \textbf{séparable} s'il contient un sous-ensemble dénombrable dense.
\end{definition}

\begin{theoreme}\label{thm:hilbert-separable-base}
\index{base hilbertienne!dénombrable}\index{espace de Hilbert!séparable}
Un espace de Hilbert $H$ est séparable si et seulement s'il admet une base hilbertienne dénombrable.
\end{theoreme}

\begin{proof}
$(\Leftarrow)$ Si $(e_n)_{n\geq 1}$ est une base hilbertienne, alors $D = \{\sum_{n=1}^N q_n e_n : N \in \N, q_n \in \Q + i\Q\}$ est dénombrable et dense dans $H$.

$(\Rightarrow)$ Soit $\{d_n\}$ un sous-ensemble dénombrable dense. Par le procédé de Gram-Schmidt\index{Gram-Schmidt, procédé de}, on peut en extraire un système orthonormal $(e_n)$. Montrons qu'il est total. Si $x \perp e_n$ pour tout $n$, alors $x \perp \spn\{d_n\}$ (car chaque $e_n$ est combinaison linéaire des $d_k$, et réciproquement chaque $d_k$ est combinaison linéaire des $e_n$). Mais $\spn\{d_n\}$ est dense dans $H$, donc $x = 0$.
\end{proof}

\begin{exemple}[{Base de Fourier de $L^2([0, 2\pi])$}]\label{ex:base-fourier}
\index{base de Fourier}\index{Fourier, base de}
Le système $\left(\frac{e^{int}}{\sqrt{2\pi}}\right)_{n \in \Z}$ est une base hilbertienne de $L^2([0, 2\pi])$. L'orthonormalité se vérifie par calcul direct :
\[
    \frac{1}{2\pi}\int_0^{2\pi} e^{-imt}e^{int}\,\dd t = \frac{1}{2\pi}\int_0^{2\pi}e^{i(n-m)t}\,\dd t = \delta_{mn}.
\]
La totalité résulte du théorème de Weierstrass trigonométrique\index{Weierstrass, théorème de!trigonométrique}.
\end{exemple}

\begin{exemple}[Polynômes de Hermite]\label{ex:hermite}
\index{polynômes de Hermite}\index{Hermite, polynômes de}
Dans $L^2(\R, e^{-x^2}\dd x)$, les polynômes de Hermite (convenablement normalisés) forment une base hilbertienne. Ceci est lié à l'oscillateur harmonique quantique\index{oscillateur harmonique}.
\end{exemple}

\begin{proposition}[Procédé de Gram-Schmidt]\label{prop:gram-schmidt}
\index{Gram-Schmidt, procédé de}
Soit $(x_n)_{n \geq 1}$ une famille libre dans un espace préhilbertien $H$. Il existe un unique système orthonormal $(e_n)_{n\geq 1}$ tel que pour tout $N \geq 1$ :
\[
    \spn\{e_1, \ldots, e_N\} = \spn\{x_1, \ldots, x_N\}
\]
et $\inner{e_n}{x_n} > 0$ pour tout $n$. On le construit par la récurrence :
\[
    \tilde{e}_1 = x_1, \quad \tilde{e}_n = x_n - \sum_{k=1}^{n-1}\inner{e_k}{x_n}e_k, \quad e_n = \frac{\tilde{e}_n}{\norm{\tilde{e}_n}}.
\]
\end{proposition}

\begin{proof}
Par récurrence. L'initialisation $e_1 = x_1/\norm{x_1}$ est claire. Si $e_1, \ldots, e_{n-1}$ sont construits, l'élément $\tilde{e}_n = x_n - \sum_{k=1}^{n-1}\inner{e_k}{x_n}e_k$ est la composante de $x_n$ orthogonale à $\spn\{e_1, \ldots, e_{n-1}\}$. Comme $x_n \notin \spn\{x_1, \ldots, x_{n-1}\} = \spn\{e_1, \ldots, e_{n-1}\}$ (famille libre), on a $\tilde{e}_n \neq 0$. L'orthonormalité se vérifie : pour $k < n$, $\inner{e_k}{\tilde{e}_n} = \inner{e_k}{x_n} - \inner{e_k}{x_n} = 0$.
\end{proof}

\begin{exemple}[Polynômes de Legendre]\label{ex:legendre}
\index{polynômes de Legendre}\index{Legendre, polynômes de}
En appliquant Gram-Schmidt à $(1, t, t^2, \ldots)$ dans $L^2([-1,1])$, on obtient les polynômes de Legendre\index{Legendre, polynômes de} (à normalisation près). Les premiers sont :
\[
    P_0(t) = 1, \quad P_1(t) = t, \quad P_2(t) = \frac{1}{2}(3t^2 - 1), \quad P_3(t) = \frac{1}{2}(5t^3 - 3t).
\]
Ils satisfont l'orthogonalité $\int_{-1}^1 P_m(t)P_n(t)\,\dd t = \frac{2}{2n+1}\delta_{mn}$ et forment une base hilbertienne de $L^2([-1,1])$.
\end{exemple}

%====================================================================
\section{Théorème de Fischer-Riesz}
%====================================================================

\begin{theoreme}[Fischer-Riesz]\label{thm:fischer-riesz}
\index{Fischer-Riesz, théorème de}\index{théorème de Fischer-Riesz}
Soit $(\Omega, \mathcal{A}, \mu)$ un espace mesuré. L'espace $L^2(\Omega, \mu)$ est un espace de Hilbert.
\end{theoreme}

\begin{proof}
Le produit scalaire $\inner{f}{g} = \int_\Omega \overline{f}\, g\,\dd\mu$ est bien défini sur $L^2$ (par l'inégalité de Cauchy-Schwarz intégrale\index{inégalité de Cauchy-Schwarz!intégrale}). Il reste à montrer la complétude.

Soit $(f_n)$ une suite de Cauchy dans $L^2$. Par le 
ef{thm:abs-cv-completude}, il suffit de montrer que toute série absolument convergente converge. Soit $\sum g_k$ avec $\sum \norm{g_k}_{L^2} < \infty$.

\textbf{Étape 1.} Posons $G_N = \sum_{k=1}^N \abs{g_k}$ et $G = \sum_{k=1}^\infty \abs{g_k}$ (limite ponctuelle, éventuellement $+\infty$). Par l'inégalité de Minkowski :
\[
    \norm{G_N}_{L^2} \leq \sum_{k=1}^N \norm{g_k}_{L^2} \leq \sum_{k=1}^\infty \norm{g_k}_{L^2} = M < \infty.
\]
Par le théorème de convergence monotone\index{convergence monotone, théorème de} :
\[
    \int_\Omega G^2 \,\dd\mu = \lim_{N\to\infty}\int_\Omega G_N^2\,\dd\mu \leq M^2.
\]
Donc $G \in L^2$, et en particulier $G(x) < \infty$ pour $\mu$-presque tout $x$.

\textbf{Étape 2.} Pour $\mu$-presque tout $x$, la série $\sum g_k(x)$ est absolument convergente (car $\sum \abs{g_k(x)} = G(x) < \infty$). Posons $f(x) = \sum_{k=1}^\infty g_k(x)$ (définie p.p.).

\textbf{Étape 3.} On a $\abs{\sum_{k=1}^N g_k(x)} \leq G(x)$ avec $G \in L^2$. Par le théorème de convergence dominée\index{convergence dominée, théorème de} :
\[
    \norm*{\sum_{k=1}^N g_k - f}_{L^2}^2 = \int_\Omega \abs*{\sum_{k=1}^N g_k - f}^2 \,\dd\mu \to 0.
\]
Donc $\sum g_k$ converge vers $f$ dans $L^2$.
\end{proof}

\begin{remarque}\label{rem:fischer-riesz-general}
\index{Fischer-Riesz, théorème de!cas $L^p$}
Le théorème de Fischer-Riesz se généralise : pour $1 \leq p < \infty$, $L^p(\Omega, \mu)$ est un espace de Banach. La preuve suit la même stratégie, mais en utilisant la convergence monotone pour le cas $p = 1$ et l'inégalité de Minkowski pour le cas général. Nous renvoyons au cours de théorie de la mesure pour les détails.
\end{remarque}

\begin{proposition}[Densité dans $L^2$]\label{prop:densite-L2}
\index{densité!dans $L^2$}
Soit $\Omega \subset \R^n$ un ouvert. Les espaces suivants sont denses dans $L^2(\Omega)$ :
\begin{enumerate}[label=(\roman*)]
    \item L'espace $C_c(\Omega)$ des fonctions continues à support compact\index{support compact} ;
    \item L'espace $C_c^\infty(\Omega)$ des fonctions lisses à support compact\index{fonction lisse à support compact} ;
    \item L'espace des fonctions en escalier\index{fonction en escalier} (si $\Omega$ est un intervalle borné).
\end{enumerate}
\end{proposition}

\begin{proof}[Esquisse]
(i) Par approximation : pour $f \in L^2$, on approche d'abord $f$ par des fonctions bornées à support compact (troncature), puis par des fonctions continues (régularisation ou théorème de Lusin\index{Lusin, théorème de}).

(ii) Par convolution\index{convolution} avec un noyau régularisant $\rho_\varepsilon$\index{noyau régularisant}.

(iii) Par approximation des fonctions continues par des fonctions en escalier.
\end{proof}

%====================================================================
\section{Somme directe hilbertienne}
%====================================================================

\begin{definition}[Somme directe hilbertienne]\label{def:somme-directe-hilbertienne}
\index{somme directe hilbertienne}
Soit $(H_i)_{i \in I}$ une famille d'espaces de Hilbert. La \textbf{somme directe hilbertienne} est :
\[
    \bigoplus_{i \in I} H_i = \left\{ (x_i)_{i \in I} : x_i \in H_i, \; \sum_{i \in I} \norm{x_i}_{H_i}^2 < \infty \right\},
\]
munie du produit scalaire $\inner{(x_i)}{(y_i)} = \sum_{i \in I} \inner{x_i}{y_i}_{H_i}$.
\end{definition}

\begin{proposition}\label{prop:somme-directe-hilbert}
La somme directe hilbertienne $\bigoplus_{i \in I} H_i$ est un espace de Hilbert.
\end{proposition}

\begin{proof}
Le produit scalaire est bien défini par l'inégalité de Cauchy-Schwarz (la série converge absolument). La complétude se démontre comme pour $\ell^2$ : une suite de Cauchy converge coordonnée par coordonnée, et la limite est dans l'espace.
\end{proof}

\begin{exemple}[Somme directe de copies de $\C$]\label{ex:somme-directe-C}
\index{somme directe hilbertienne!de copies de $\C$}
La somme directe $\bigoplus_{n \geq 1} \C$ (chaque $H_n = \C$ muni du produit scalaire usuel) est isomorphe à $\ell^2$. Plus généralement, si $H_n = \C^{d_n}$, alors $\bigoplus_{n\geq 1} H_n = \{(x_n) : x_n \in \C^{d_n}, \sum_n \norm{x_n}^2 < \infty\}$.
\end{exemple}

\begin{theoreme}[Décomposition en somme directe]\label{thm:decomp-somme-directe}
\index{décomposition orthogonale}
Soit $H$ un espace de Hilbert et $(F_i)_{i \in I}$ une famille de sous-espaces fermés mutuellement orthogonaux ($F_i \perp F_j$ pour $i \neq j$) tels que $\overline{\bigoplus_{\text{alg}} F_i} = H$. Alors $H$ est isomorphe (isométriquement) à la somme directe hilbertienne $\bigoplus_{i \in I} F_i$.
\end{theoreme}

\begin{proof}
L'application $U : \bigoplus_{i \in I} F_i \to H$ définie par $U((x_i)) = \sum_{i \in I} x_i$ (la somme converge car $\sum \norm{x_i}^2 < \infty$ et les termes sont orthogonaux) est une isométrie par le théorème de Pythagore. Son image contient $\bigoplus_{\text{alg}} F_i$ qui est dense dans $H$, et l'image d'une isométrie depuis un espace complet est fermée. Donc $U$ est surjective.
\end{proof}

%====================================================================
\section{Opérateurs dans les espaces de Hilbert : introduction}
%====================================================================

La théorie des opérateurs\index{théorie des opérateurs} dans les espaces de Hilbert est considérablement plus riche que dans les espaces de Banach généraux, grâce à l'existence du produit scalaire. En particulier, la notion d'adjoint joue un rôle central et permet de classifier les opérateurs en familles aux propriétés remarquables : opérateurs auto-adjoints, normaux, unitaires et positifs.

\begin{definition}[Opérateur adjoint]\label{def:adjoint}
\index{opérateur!adjoint}\index{adjoint}
Soit $T \in \mathcal{L}(H)$ un opérateur borné sur un espace de Hilbert $H$. L'\textbf{opérateur adjoint} de $T$ est l'unique opérateur $T^* \in \mathcal{L}(H)$ tel que
\[
    \inner{T^*x}{y} = \inner{x}{Ty} \quad \text{pour tous } x, y \in H.
\]
\end{definition}

\begin{proof}[Justification de l'existence]
Pour $x$ fixé, l'application $y \mapsto \inner{x}{Ty}$ est une forme linéaire continue sur $H$ (de norme $\leq \norm{x}\norm{T}$). Par le théorème de Riesz (
ef{thm:riesz-representation}), il existe un unique $z \in H$ tel que $\inner{x}{Ty} = \inner{z}{y}$ pour tout $y$. On pose $T^*x = z$. La linéarité de $T^*$ résulte de l'unicité dans le théorème de Riesz.
\end{proof}

\begin{proposition}[Propriétés de l'adjoint]\label{prop:proprietes-adjoint}
\index{adjoint!propriétés}
Soient $S, T \in \mathcal{L}(H)$ et $\lambda \in \K$. Alors :
\begin{enumerate}[label=(\roman*)]
    \item $(T^*)^* = T$ ;
    \item $(\lambda T)^* = \overline{\lambda} T^*$ ;
    \item $(S + T)^* = S^* + T^*$ ;
    \item $(ST)^* = T^*S^*$ ;
    \item $\norm{T^*} = \norm{T}$ ;
    \item $\norm{T^*T} = \norm{T}^2$.
\end{enumerate}
\end{proposition}

\begin{proof}
Les points (i)--(iv) se vérifient directement à partir de la définition.

(v) On a $\norm{T^*x}^2 = \inner{T^*x}{T^*x} = \inner{x}{TT^*x} \leq \norm{x}\norm{TT^*x} \leq \norm{x}\norm{T}\norm{T^*x}$, d'où $\norm{T^*x} \leq \norm{T}\norm{x}$ et $\norm{T^*} \leq \norm{T}$. Par symétrie (en appliquant le même argument à $T^*$ et en utilisant (i)), $\norm{T} = \norm{(T^*)^*} \leq \norm{T^*}$.

(vi) On a $\norm{T^*T} \leq \norm{T^*}\norm{T} = \norm{T}^2$. D'autre part, $\norm{Tx}^2 = \inner{Tx}{Tx} = \inner{x}{T^*Tx} \leq \norm{x}\norm{T^*Tx} \leq \norm{x}\norm{T^*T}\norm{x}$. Donc $\norm{T}^2 \leq \norm{T^*T}$.
\end{proof}

\begin{exemple}[Adjoint d'un opérateur de multiplication]\label{ex:adjoint-multiplication}
\index{opérateur!de multiplication}\index{adjoint!de l'opérateur de multiplication}
Soit $\varphi \in L^\infty(\Omega, \mu)$ et $M_\varphi : L^2 \to L^2$ l'opérateur de multiplication\index{opérateur!de multiplication} défini par $(M_\varphi f)(x) = \varphi(x)f(x)$. Alors $M_\varphi \in \mathcal{L}(L^2)$ avec $\norm{M_\varphi} = \norm{\varphi}_{L^\infty}$, et $M_\varphi^* = M_{\overline{\varphi}}$. En effet :
\[
    \inner{M_\varphi^* f}{g} = \inner{f}{M_\varphi g} = \int \overline{f}\,\varphi g\,\dd\mu = \int \overline{\overline{\varphi}f}\,g\,\dd\mu = \inner{M_{\overline{\varphi}}f}{g}.
\]
L'opérateur $M_\varphi$ est auto-adjoint si et seulement si $\varphi$ est à valeurs réelles (p.p.), et unitaire si et seulement si $\abs{\varphi} = 1$ p.p.
\end{exemple}

\begin{exemple}[Opérateur intégral]\label{ex:operateur-integral}
\index{opérateur!intégral}\index{noyau intégral}
Soit $K \in L^2([0,1]^2)$ un noyau intégral. L'opérateur $T : L^2([0,1]) \to L^2([0,1])$ défini par
\[
    (Tf)(x) = \int_0^1 K(x,t)f(t)\,\dd t
\]
est borné avec $\norm{T} \leq \norm{K}_{L^2([0,1]^2)}$ (opérateur de Hilbert-Schmidt\index{Hilbert-Schmidt, opérateur de}\index{opérateur!de Hilbert-Schmidt}). Son adjoint est $T^*g(x) = \int_0^1 \overline{K(t,x)}\,g(t)\,\dd t$.
\end{exemple}

\begin{definition}[Opérateurs spéciaux]\label{def:op-speciaux}
\index{opérateur!auto-adjoint}\index{opérateur!normal}\index{opérateur!unitaire}\index{opérateur!positif}
Un opérateur $T \in \mathcal{L}(H)$ est dit :
\begin{enumerate}[label=(\roman*)]
    \item \textbf{auto-adjoint}\index{auto-adjoint} (ou hermitien\index{hermitien}) si $T^* = T$ ;
    \item \textbf{normal}\index{normal} si $T^*T = TT^*$ ;
    \item \textbf{unitaire}\index{unitaire} si $T^*T = TT^* = \Id$ ;
    \item \textbf{positif}\index{positif} si $\inner{Tx}{x} \geq 0$ pour tout $x \in H$ (on écrit $T \geq 0$).
\end{enumerate}
\end{definition}

\begin{proposition}\label{prop:auto-adjoint-reel}
\index{auto-adjoint!spectre réel}
Si $T \in \mathcal{L}(H)$ est auto-adjoint, alors $\inner{Tx}{x} \in \R$ pour tout $x \in H$. De plus :
\[
    \norm{T} = \sup_{\norm{x}=1} \abs{\inner{Tx}{x}}.
\]
\end{proposition}

\begin{proof}
$\overline{\inner{Tx}{x}} = \inner{x}{Tx} = \inner{T^*x}{x} = \inner{Tx}{x}$, donc $\inner{Tx}{x} \in \R$.

Posons $M = \sup_{\norm{x}=1}\abs{\inner{Tx}{x}}$. Clairement $M \leq \norm{T}$ (par Cauchy-Schwarz). Pour l'autre sens, on utilise l'identité de polarisation : pour $x, y$ de norme $1$,
\[
    4\operatorname{Re}\inner{Tx}{y} = \inner{T(x+y)}{x+y} - \inner{T(x-y)}{x-y}.
\]
D'où $4\abs{\operatorname{Re}\inner{Tx}{y}} \leq M(\norm{x+y}^2 + \norm{x-y}^2) = 2M(\norm{x}^2 + \norm{y}^2) = 4M$, et en optimisant sur le choix de phase de $y$ : $\abs{\inner{Tx}{y}} \leq M$. En prenant $y = Tx/\norm{Tx}$ : $\norm{Tx} \leq M$, donc $\norm{T} \leq M$.
\end{proof}

\begin{remarque}\label{rem:ker-ran-adjoint}
\index{noyau!et image de l'adjoint}
Pour $T \in \mathcal{L}(H)$, on a les relations fondamentales :
\[
    \Ker(T^*) = (\ran T)^\perp, \qquad \overline{\ran(T)} = (\Ker T^*)^\perp.
\]
Ces relations sont immédiates : $x \in \Ker(T^*)$ ssi $T^*x = 0$ ssi $\inner{T^*x}{y} = 0$ pour tout $y$ ssi $\inner{x}{Ty} = 0$ pour tout $y$ ssi $x \perp \ran T$.
\end{remarque}

Le diagramme suivant résume les relations entre noyaux et images :
\begin{center}
\begin{tikzcd}[column sep=large]
    H \arrow[r, "T"] & H \\
    \Ker T \arrow[u, hook] \arrow[r, phantom, "\perp"] & \overline{\ran T^*} \arrow[u, hook] \\
    (\ran T^*)^\perp \arrow[u, equal] \arrow[r, phantom, "\oplus"] & \overline{\ran T^*} \arrow[u, equal]
\end{tikzcd}
\captionof{figure}{Décomposition $H = \Ker T \oplus \overline{\ran T^*}$ pour un opérateur $T \in \mathcal{L}(H)$.}\label{fig:ker-ran}
\end{center}

\begin{proposition}[Spectre d'un opérateur auto-adjoint]\label{prop:spectre-autoadjoint}
\index{spectre!d'un opérateur auto-adjoint}\index{auto-adjoint!spectre}
Soit $T \in \mathcal{L}(H)$ un opérateur auto-adjoint. Alors :
\begin{enumerate}[label=(\roman*)]
    \item Toute valeur propre\index{valeur propre} de $T$ est réelle ;
    \item Les espaces propres\index{espace propre} correspondant à des valeurs propres distinctes sont orthogonaux ;
    \item $\norm{T} = \sup_{\norm{x}=1}\abs{\inner{Tx}{x}}$.
\end{enumerate}
\end{proposition}

\begin{proof}
(i) Si $Tx = \lambda x$ avec $x \neq 0$ : $\lambda\norm{x}^2 = \inner{Tx}{x} \in \R$ (car $T = T^*$), donc $\lambda \in \R$.

(ii) Si $Tx = \lambda x$ et $Ty = \mu y$ avec $\lambda \neq \mu$ : $\lambda\inner{x}{y} = \inner{Tx}{y} = \inner{x}{Ty} = \mu\inner{x}{y}$, donc $(\lambda - \mu)\inner{x}{y} = 0$ et $\inner{x}{y} = 0$.

(iii) C'est la 
ef{prop:auto-adjoint-reel}.
\end{proof}

\begin{remarque}[Vers la théorie spectrale]\label{rem:vers-spectrale}
\index{théorie spectrale}
La théorie spectrale des opérateurs auto-adjoints (ou normaux) dans les espaces de Hilbert est l'un des sujets majeurs de l'analyse fonctionnelle. Le théorème spectral\index{théorème spectral} affirme que tout opérateur auto-adjoint borné admet une représentation intégrale $T = \int_{\sigma(T)} \lambda\,\dd E(\lambda)$ où $E$ est une mesure spectrale\index{mesure spectrale}. Cette théorie sera développée dans un chapitre ultérieur.
\end{remarque}

%====================================================================
\section{Exercices}
%====================================================================

\begin{exercice}[$\star$]\label{exo:ch2-verif-ps}
\index{produit scalaire!vérification}
Montrer que $\inner{f}{g} = \int_0^1 f(t)\overline{g(t)} \, \dd t$ définit un produit scalaire sur $C([0,1], \C)$.
\end{exercice}

\begin{exercice}[$\star$]\label{exo:ch2-parallelogramme}
\index{identité du parallélogramme!contre-exemple}
Montrer que la norme $\norm{\cdot}_\infty$ sur $C([0,1])$ ne provient pas d'un produit scalaire. \textit{Indication : vérifier que l'identité du parallélogramme échoue.}
\end{exercice}

\begin{exercice}[$\star$]\label{exo:ch2-orthogonal}
Soit $H = L^2([0,1])$ et $F = \{f \in L^2 : \int_0^1 f(t)\,\dd t = 0\}$. Déterminer $F^\perp$.
\end{exercice}

\begin{exercice}[$\star\star$]\label{exo:ch2-proj-calcul}
\index{projection!calcul explicite}
Dans $H = L^2([0,1])$, soit $F = \spn\{1, t\}$. Calculer la projection orthogonale de $f(t) = t^2$ sur $F$.
\end{exercice}

\begin{exercice}[$\star\star$]\label{exo:ch2-cs-egalite}
Soit $H$ un espace de Hilbert et $x, y \in H$ avec $\norm{x} = \norm{y} = 1$. Montrer que $\abs{\inner{x}{y}} = 1$ si et seulement si $x = \lambda y$ pour un $\lambda \in \K$ avec $\abs{\lambda} = 1$.
\end{exercice}

\begin{exercice}[$\star\star$]\label{exo:ch2-proj-contraction}
\index{projection!contraction}
Montrer que la projection $P_C$ sur un convexe fermé $C$ d'un espace de Hilbert est une contraction\index{contraction} : $\norm{P_C(x) - P_C(y)} \leq \norm{x - y}$.
\end{exercice}

\begin{exercice}[$\star\star$]\label{exo:ch2-gram-schmidt}
\index{Gram-Schmidt, procédé de}
Appliquer le procédé de Gram-Schmidt aux polynômes $1, t, t^2, t^3$ dans $L^2([-1,1])$ pour retrouver les premiers polynômes de Legendre\index{Legendre, polynômes de} (à normalisation près).
\end{exercice}

\begin{exercice}[$\star\star$]\label{exo:ch2-bessel-optimal}
Montrer que l'inégalité de Bessel peut être stricte si le système orthonormal n'est pas complet. Donner un exemple.
\end{exercice}

\begin{exercice}[$\star\star$]\label{exo:ch2-adjoint-calcul}
\index{adjoint!calcul}
Soit $T : \ell^2 \to \ell^2$ défini par $T(x_1, x_2, x_3, \ldots) = (0, x_1, x_2, \ldots)$ (décalage à droite\index{décalage à droite}). Calculer $T^*$ et vérifier que $\norm{T^*T} = \norm{T}^2$.
\end{exercice}

\begin{exercice}[$\star\star\star$]\label{exo:ch2-lax-milgram-app}
\index{Lax-Milgram, théorème de!application}
Soit $\Omega \subset \R^n$ un ouvert borné. En utilisant le théorème de Lax-Milgram, montrer l'existence et l'unicité de la solution faible $u \in H^1_0(\Omega)$ de $-\Delta u + u = f$ pour $f \in L^2(\Omega)$.
\end{exercice}

\begin{exercice}[$\star\star\star$]\label{exo:ch2-hilbert-non-separable}
\index{espace de Hilbert!non séparable}
Montrer que $\ell^2(I)$ pour un ensemble $I$ non dénombrable est un espace de Hilbert non séparable. \textit{Indication : montrer que la base canonique $(e_i)_{i\in I}$ est un système orthonormal tel que $\norm{e_i - e_j} = \sqrt{2}$ pour $i \neq j$.}
\end{exercice}

\begin{exercice}[$\star\star\star$]\label{exo:ch2-operateur-compact}
\index{opérateur!compact}
Soit $T : \ell^2 \to \ell^2$ défini par $T(x_n) = (x_n/n)$. Montrer que $T$ est auto-adjoint, positif, injectif, et que $T$ est compact\index{compact!opérateur} (c'est-à-dire que $T(B_{\ell^2})$ est relativement compact).
\end{exercice}

\begin{exercice}[$\star$]\label{exo:ch2-pythagore-infini}
\index{Pythagore, théorème de!séries}
Soit $(e_n)_{n\geq 1}$ un système orthonormal dans un espace de Hilbert $H$ et $(\alpha_n) \in \ell^2$. Montrer que la série $\sum_{n=1}^\infty \alpha_n e_n$ converge dans $H$ et que $\norm{\sum_{n=1}^\infty \alpha_n e_n}^2 = \sum_{n=1}^\infty \abs{\alpha_n}^2$.
\end{exercice}

\begin{exercice}[$\star\star$]\label{exo:ch2-riesz-app}
\index{Riesz, théorème de représentation de!application}
Soit $H = L^2([0,1])$ et $\varphi : H \to \R$ définie par $\varphi(f) = \int_0^1 t f(t)\,\dd t$. Identifier le représentant de Riesz de $\varphi$ et calculer $\norm{\varphi}$.
\end{exercice}

\begin{exercice}[$\star\star$]\label{exo:ch2-proj-formule}
\index{projection!formule explicite}
Soit $F = \spn\{e_1, \ldots, e_n\}$ un sous-espace de dimension finie d'un espace de Hilbert $H$, où $(e_1, \ldots, e_n)$ est un système orthonormal. Montrer que la projection orthogonale sur $F$ est donnée par :
\[
    P_F(x) = \sum_{k=1}^n \inner{e_k}{x}\, e_k.
\]
Montrer que si $(f_1, \ldots, f_n)$ est une base quelconque (non nécessairement orthonormale) de $F$, la projection s'exprime via la matrice de Gram\index{matrice de Gram} $G = (\inner{f_i}{f_j})_{1\leq i,j\leq n}$ :
\[
    P_F(x) = \sum_{i,j=1}^n (G^{-1})_{ij}\, \inner{f_i}{x}\, f_j.
\]
\end{exercice}

\begin{exercice}[$\star\star\star$]\label{exo:ch2-noyau-reproduisant}
\index{noyau reproduisant}\index{espace de Hilbert à noyau reproduisant}
Soit $H$ un espace de Hilbert de fonctions sur un ensemble $X$ tel que pour tout $x \in X$, l'évaluation $\delta_x : f \mapsto f(x)$ est une forme linéaire continue sur $H$. Par le théorème de Riesz, il existe $K_x \in H$ tel que $f(x) = \inner{K_x}{f}$ pour tout $f \in H$.
\begin{enumerate}[label=(\alph*)]
    \item Montrer que $K(x,y) = \inner{K_x}{K_y} = K_y(x)$ est un noyau défini positif\index{noyau défini positif}.
    \item Montrer que $K$ a la \emph{propriété de reproduction} : $\inner{K(\cdot, y)}{f} = f(y)$.
    \item Vérifier que l'espace de Hardy $H^2(\mathbb{D})$ est un espace de Hilbert à noyau reproduisant avec $K(z,w) = \frac{1}{1-\overline{w}z}$ (noyau de Szeg\H{o}\index{Szego, noyau de@Szeg\H{o}, noyau de}).
\end{enumerate}
\end{exercice}

\begin{exercice}[$\star\star$]\label{exo:ch2-somme-directe}
\index{somme directe hilbertienne!exercice}
Soit $H = L^2([0,2\pi])$. Posons $H_n = \spn\{e^{int}\}$ pour $n \in \Z$.
\begin{enumerate}[label=(\alph*)]
    \item Montrer que $H_n \perp H_m$ pour $n \neq m$.
    \item Montrer que $H = \overline{\bigoplus_{n \in \Z} H_n}$ (somme directe hilbertienne).
    \item En déduire une preuve du théorème de Parseval pour les séries de Fourier.
\end{enumerate}
\end{exercice}

\begin{exercice}[$\star\star$]\label{exo:ch2-operateurs-borneinf}
\index{opérateur!borné inférieurement}
Soit $T \in \mathcal{L}(H)$. On dit que $T$ est \textbf{borné inférieurement}\index{borné inférieurement} s'il existe $c > 0$ tel que $\norm{Tx} \geq c\norm{x}$ pour tout $x \in H$.
\begin{enumerate}[label=(\alph*)]
    \item Montrer que si $T$ est borné inférieurement, alors $T$ est injectif et $\ran T$ est fermé.
    \item Montrer que si $T$ est auto-adjoint et borné inférieurement, alors $T$ est bijectif.
    \item Donner un exemple d'opérateur borné inférieurement non surjectif.
\end{enumerate}
\end{exercice}

\begin{exercice}[$\star\star\star$]\label{exo:ch2-projection-non-lineaire}
\index{projection!sur un convexe!propriétés}
Soit $C$ un convexe fermé non vide dans un espace de Hilbert $H$ et $P_C : H \to C$ la projection sur $C$.
\begin{enumerate}[label=(\alph*)]
    \item Montrer que $P_C$ est une contraction : $\norm{P_C(x) - P_C(y)} \leq \norm{x - y}$ pour tous $x, y \in H$.
    \item Montrer que $P_C$ est un opérateur monotone\index{opérateur monotone} : $\operatorname{Re}\inner{P_C(x) - P_C(y)}{x - y} \geq 0$.
    \item En déduire que $P_C$ est \emph{fermement non expansif}\index{fermement non expansif} :
    \[
        \norm{P_C(x) - P_C(y)}^2 \leq \operatorname{Re}\inner{P_C(x) - P_C(y)}{x - y}.
    \]
\end{enumerate}
\end{exercice}

\begin{exercice}[$\star$]\label{exo:ch2-cont-prodscal}
\index{produit scalaire!continuité}
Montrer que le produit scalaire est une application continue de $H \times H$ dans $\K$. Plus précisément, si $x_n \to x$ et $y_n \to y$ dans $H$, alors $\inner{x_n}{y_n} \to \inner{x}{y}$.
\end{exercice}

\begin{exercice}[$\star\star$]\label{exo:ch2-sous-espace-dense}
\index{sous-espace!dense}
Soit $F$ un sous-espace vectoriel d'un espace de Hilbert $H$. Montrer que les conditions suivantes sont équivalentes :
\begin{enumerate}[label=(\roman*)]
    \item $F$ est dense dans $H$ ;
    \item $F^\perp = \{0\}$ ;
    \item Pour tout $x \in H$, si $\inner{x}{f} = 0$ pour tout $f \in F$, alors $x = 0$.
\end{enumerate}
\end{exercice}

\begin{exercice}[$\star\star$]\label{exo:ch2-adjoint-matrice}
\index{adjoint!et matrice}
Soit $A \in M_n(\C)$ vue comme opérateur sur $(\C^n, \inner{\cdot}{\cdot})$ (produit scalaire canonique). Montrer que la matrice de $A^*$ dans la base canonique est $\overline{A}^T$ (transposée conjuguée\index{transposée conjuguée}). En déduire qu'un opérateur sur $\C^n$ est auto-adjoint si et seulement si sa matrice est hermitienne\index{matrice hermitienne}.
\end{exercice}

\begin{exercice}[$\star\star$]\label{exo:ch2-unitaire}
\index{opérateur!unitaire!caractérisation}
Soit $U \in \mathcal{L}(H)$. Montrer que les conditions suivantes sont équivalentes :
\begin{enumerate}[label=(\roman*)]
    \item $U$ est unitaire ($U^*U = UU^* = \Id$) ;
    \item $U$ est surjectif et isométrique ($\norm{Ux} = \norm{x}$ pour tout $x$) ;
    \item $U$ préserve le produit scalaire : $\inner{Ux}{Uy} = \inner{x}{y}$ pour tous $x, y$.
\end{enumerate}
\end{exercice}

\begin{exercice}[$\star\star\star$]\label{exo:ch2-parseval-explicite}
\index{Parseval, identité de!calcul explicite}
En utilisant la base de Fourier de $L^2([0,2\pi])$ et l'identité de Parseval, calculer les sommes suivantes :
\begin{enumerate}[label=(\alph*)]
    \item $\sum_{n=1}^\infty \frac{1}{n^2} = \frac{\pi^2}{6}$ \quad (en développant $f(t) = t$ sur $[0, 2\pi]$) ;
    \item $\sum_{n=1}^\infty \frac{1}{n^4} = \frac{\pi^4}{90}$ \quad (en développant $f(t) = t^2$ sur $[0, 2\pi]$).
\end{enumerate}
\textit{Indication pour (a) : les coefficients de Fourier de $f(t) = t$ sur $[0,2\pi]$ sont $c_0 = \pi$ et $c_n = \frac{i}{n}$ pour $n \neq 0$ (à un facteur $2\pi$ près selon la convention). Appliquer Parseval et identifier.}
\end{exercice}

\begin{exercice}[$\star\star\star$]\label{exo:ch2-alternative-fredholm}
\index{alternative de Fredholm}\index{Fredholm, alternative de}
Soit $T \in \mathcal{L}(H)$ un opérateur compact\index{opérateur!compact} auto-adjoint sur un espace de Hilbert $H$ et $\lambda \neq 0$. Montrer que :
\begin{enumerate}[label=(\alph*)]
    \item $\Ker(T - \lambda\Id)$ est de dimension finie ;
    \item $\ran(T - \lambda\Id)$ est fermé ;
    \item $\ran(T - \lambda\Id) = \Ker(T - \lambda\Id)^\perp$ ;
    \item En déduire l'\textbf{alternative de Fredholm} : soit $T - \lambda\Id$ est bijectif, soit $\lambda$ est valeur propre de $T$.
\end{enumerate}
\textit{Cet exercice anticipe le chapitre sur la théorie spectrale des opérateurs compacts.}
\end{exercice}
%%%%%%%%%%%%%%%%%%%%%%%%%%%%%%%%%%%%%%%%%%%%%%%%%%%%%%%%%%%%%%%%%%%%
%  Analyse Fonctionnelle — Chapitre 3
%  Opérateurs Linéaires Bornés
%%%%%%%%%%%%%%%%%%%%%%%%%%%%%%%%%%%%%%%%%%%%%%%%%%%%%%%%%%%%%%%%%%%%

\chapter{Opérateurs Linéaires Bornés}
\label{chap:operateurs-bornes}
\minitoc

\vspace{1em}

L'étude des \emph{opérateurs linéaires continus}\index{opérateur linéaire continu}
entre espaces de Banach constitue le c\oe ur de l'analyse fonctionnelle.
Ce chapitre développe la théorie de l'espace $\mathcal{L}(E,F)$\index{L(E,F)@$\mathcal{L}(E,F)$}
des opérateurs bornés, introduit les opérateurs compacts\index{opérateur compact}, puis
aborde la théorie spectrale\index{théorie spectrale} élémentaire. Dans le cadre hilbertien,
nous étudions l'opérateur adjoint\index{opérateur!adjoint} et les classes d'opérateurs
normaux, auto-adjoints et unitaires.

Dans tout ce chapitre, $E$, $F$ et $G$ désignent des espaces vectoriels normés sur
le corps $\K$ ($\R$ ou $\C$), sauf mention contraire. Lorsque nous abordons la théorie
spectrale, nous travaillons sur $\C$.

%% ══════════════════════════════════════════════════════════════════
\section{L'espace \texorpdfstring{$\mathcal{L}(E,F)$}{L(E,F)}}
\label{sec:espace-LEF}
%% ══════════════════════════════════════════════════════════════════

\index{opérateur linéaire borné|(}

\begin{definition}[Opérateur linéaire borné]\label{def:operateur-borne}
\index{opérateur linéaire borné!définition}
Soient $E$ et $F$ deux espaces vectoriels normés. Une application linéaire
$T : E \to F$ est dite \textbf{bornée} s'il existe une constante $C \geq 0$ telle que
\[
    \norm{Tx}_F \leq C \norm{x}_E \quad \text{pour tout } x \in E.
\]
On note $\mathcal{L}(E,F)$ l'ensemble de toutes les applications linéaires bornées de
$E$ dans $F$, et $\mathcal{L}(E) := \mathcal{L}(E,E)$.
\end{definition}

\begin{proposition}[Continuité et bornitude]\label{prop:continu-borne}
\index{opérateur linéaire borné!continuité}
Soit $T : E \to F$ une application linéaire. Les assertions suivantes sont équivalentes :
\begin{enumerate}[label=(\roman*)]
    \item $T$ est continue sur $E$ ;
    \item $T$ est continue en $0$ ;
    \item $T$ est bornée ;
    \item $T$ est lipschitzienne\index{application lipschitzienne}.
\end{enumerate}
\end{proposition}

\begin{proof}
$(iv) \Rightarrow (i) \Rightarrow (ii)$ est immédiat.

$(ii) \Rightarrow (iii)$ : Supposons $T$ continue en $0$. Pour $\varepsilon = 1$, il
existe $\delta > 0$ tel que $\norm{x}_E \leq \delta \Rightarrow \norm{Tx}_F \leq 1$.
Pour $x \neq 0$, posons $y = \delta x / \norm{x}_E$. Alors $\norm{y}_E = \delta$, donc
$\norm{Ty}_F \leq 1$, c'est-à-dire $\norm{Tx}_F \leq \norm{x}_E / \delta$.
On prend $C = 1/\delta$.

$(iii) \Rightarrow (iv)$ : Si $\norm{Tx}_F \leq C\norm{x}_E$ pour tout $x$, alors par
linéarité $\norm{Tx - Ty}_F = \norm{T(x-y)}_F \leq C\norm{x-y}_E$.
\end{proof}

\begin{definition}[Norme d'opérateur]\label{def:norme-operateur}
\index{norme!d'opérateur}
Pour $T \in \mathcal{L}(E,F)$, on définit la \textbf{norme d'opérateur} par
\[
    \norm{T}_{\mathcal{L}(E,F)} := \sup_{x \neq 0} \frac{\norm{Tx}_F}{\norm{x}_E}
    = \sup_{\norm{x}_E \leq 1} \norm{Tx}_F
    = \sup_{\norm{x}_E = 1} \norm{Tx}_F.
\]
\end{definition}

\begin{theoreme}[$\mathcal{L}(E,F)$ est un espace de Banach]\label{thm:LEF-banach}
\index{L(E,F)@$\mathcal{L}(E,F)$!complétude}
Si $F$ est un espace de Banach, alors $\mathcal{L}(E,F)$ muni de la norme d'opérateur
est un espace de Banach.
\end{theoreme}

\begin{proof}
Soit $(T_n)$ une suite de Cauchy dans $\mathcal{L}(E,F)$. Pour tout $x \in E$,
\[
    \norm{T_n x - T_m x}_F \leq \norm{T_n - T_m} \cdot \norm{x}_E \to 0
    \quad \text{quand } n,m \to \infty.
\]
Donc $(T_n x)$ est de Cauchy dans $F$, qui est complet. On pose
$Tx := \lim_{n\to\infty} T_n x$. Par passage à la limite, $T$ est linéaire.

Soit $\varepsilon > 0$. Il existe $N$ tel que $\norm{T_n - T_m} < \varepsilon$ pour
$n,m \geq N$. Pour $\norm{x}_E \leq 1$ et $n,m \geq N$ :
$\norm{T_n x - T_m x}_F \leq \varepsilon$.
En faisant $m \to \infty$ : $\norm{T_n x - Tx}_F \leq \varepsilon$. Donc
$\norm{T_n - T} \leq \varepsilon$ pour $n \geq N$. En particulier
$T = T_N + (T - T_N) \in \mathcal{L}(E,F)$ et $T_n \to T$ en norme.
\end{proof}

\begin{corollaire}[Dual est toujours complet]\label{cor:dual-banach}
\index{espace dual!complétude}
Pour tout espace vectoriel normé $E$, le dual topologique
$E' = \mathcal{L}(E, \K)$\index{E'@$E'$} est un espace de Banach.
\end{corollaire}

\begin{proposition}[Composition d'opérateurs]\label{prop:composition-operateurs}
\index{opérateur linéaire borné!composition}
Soient $T \in \mathcal{L}(E,F)$ et $S \in \mathcal{L}(F,G)$. Alors
$ST \in \mathcal{L}(E,G)$ et $\norm{ST} \leq \norm{S}\cdot\norm{T}$.
\end{proposition}

\begin{proof}
Pour tout $x \in E$, $\norm{STx}_G \leq \norm{S}\cdot\norm{Tx}_F \leq
\norm{S}\cdot\norm{T}\cdot\norm{x}_E$.
\end{proof}

\begin{corollaire}\label{cor:LEE-algebre}
\index{algèbre de Banach}
Si $E$ est un espace de Banach, alors $\mathcal{L}(E)$ est une algèbre de Banach
unitaire\index{algèbre de Banach!unitaire} (l'unité étant $\Id_E$), c'est-à-dire
$\norm{ST} \leq \norm{S}\cdot\norm{T}$ et $\norm{\Id} = 1$ (si $E \neq \{0\}$).
\end{corollaire}

\begin{notation}\label{not:GL-E}
\index{GL(E)@$GL(E)$}
On note $GL(E)$ l'ensemble des opérateurs inversibles de $\mathcal{L}(E)$, c'est-à-dire
les $T \in \mathcal{L}(E)$ tels qu'il existe $T^{-1} \in \mathcal{L}(E)$ avec
$TT^{-1} = T^{-1}T = \Id_E$.
\end{notation}

\index{opérateur linéaire borné|)}

%% ══════════════════════════════════════════════════════════════════
\section{Opérateurs compacts}
\label{sec:operateurs-compacts}
%% ══════════════════════════════════════════════════════════════════

\index{opérateur compact|(}

\begin{definition}[Opérateur compact]\label{def:operateur-compact}
\index{opérateur compact!définition}
Un opérateur $T \in \mathcal{L}(E,F)$ est dit \textbf{compact} si l'image de la boule
unité fermée $\overline{B}_E$ par $T$ est relativement compacte dans $F$, c'est-à-dire
si $\overline{T(\overline{B}_E)}$ est compact dans $F$.

On note $\mathcal{K}(E,F)$\index{K(E,F)@$\mathcal{K}(E,F)$} l'ensemble des opérateurs
compacts de $E$ dans $F$, et $\mathcal{K}(E) := \mathcal{K}(E,E)$.
\end{definition}

\begin{remarque}\label{rem:compact-suites}
\index{opérateur compact!caractérisation séquentielle}
De manière équivalente, $T$ est compact si et seulement si pour toute suite bornée
$(x_n)$ de $E$, la suite $(Tx_n)$ admet une sous-suite convergente dans $F$.
\end{remarque}

\begin{proposition}[Premières propriétés des opérateurs compacts]%
\label{prop:proprietes-compacts}
\index{opérateur compact!propriétés}
\begin{enumerate}[label=(\roman*)]
    \item Tout opérateur compact est borné : $\mathcal{K}(E,F) \subset \mathcal{L}(E,F)$.
    \item Tout opérateur de rang fini\index{opérateur!de rang fini} est compact.
    \item Si $T \in \mathcal{K}(E,F)$, $S \in \mathcal{L}(F,G)$ et
          $R \in \mathcal{L}(G,E)$, alors $ST \in \mathcal{K}(E,G)$ et
          $TR \in \mathcal{K}(G,F)$.
    \item $\mathcal{K}(E,F)$ est un sous-espace vectoriel fermé de $\mathcal{L}(E,F)$
          lorsque $F$ est complet.
\end{enumerate}
\end{proposition}

\begin{proof}
(i) L'image de la boule unité est relativement compacte, donc bornée.

(ii) Si $\dim(\ran T) < \infty$, alors $T(\overline{B}_E)$ est bornée dans un espace
de dimension finie, donc relativement compacte par le théorème de Bolzano-Weierstrass%
\index{Bolzano-Weierstrass, théorème de}.

(iii) Soit $(x_n)$ bornée dans $E$. Comme $T$ est compact, il existe une sous-suite
$(x_{n_k})$ telle que $Tx_{n_k} \to y$ dans $F$. Alors $STx_{n_k} \to Sy$ dans $G$
par continuité de $S$. Pour $TR$, si $(z_n)$ est bornée dans $G$, alors $(Rz_n)$ est
bornée dans $E$, donc $(TRz_n)$ admet une sous-suite convergente.

(iv) Soient $T, S \in \mathcal{K}(E,F)$ et $\lambda \in \K$. Pour toute suite bornée
$(x_n)$, on extrait d'abord une sous-suite telle que $(Tx_{n_k})$ converge, puis une
sous-sous-suite telle que $(Sx_{n_{k_l}})$ converge. Alors
$((T + \lambda S)x_{n_{k_l}})$ converge, donc $T + \lambda S$ est compact.

Pour la fermeture, soit $(T_k)$ une suite dans $\mathcal{K}(E,F)$ avec $T_k \to T$
en norme. Soit $(x_n)$ bornée dans $E$ avec $\norm{x_n} \leq M$. Par un procédé
diagonal\index{procédé diagonal}, on construit une sous-suite $(x_{n_j})$ telle que
$(T_k x_{n_j})_j$ converge pour tout $k$. Pour $\varepsilon > 0$, choisissons $k$
tel que $\norm{T - T_k} < \varepsilon / (3M)$, puis $J$ tel que
$\norm{T_k x_{n_i} - T_k x_{n_j}} < \varepsilon/3$ pour $i, j \geq J$. Alors
\[
    \norm{Tx_{n_i} - Tx_{n_j}} \leq \norm{(T-T_k)x_{n_i}} +
    \norm{T_k x_{n_i} - T_k x_{n_j}} + \norm{(T_k-T)x_{n_j}} < \varepsilon.
\]
Donc $(Tx_{n_j})$ est de Cauchy dans $F$ complet, d'où converge.
\end{proof}

\begin{corollaire}\label{cor:compact-ideal}
\index{opérateur compact!idéal bilatère}
L'espace $\mathcal{K}(E)$ est un idéal bilatère fermé de l'algèbre $\mathcal{L}(E)$
lorsque $E$ est un espace de Banach.
\end{corollaire}

\begin{exemple}\label{ex:operateur-integral-compact}
\index{opérateur intégral!compact}
Soit $k \in C([0,1]^2)$. L'opérateur intégral $T : C([0,1]) \to C([0,1])$ défini par
\[
    (Tf)(x) = \int_0^1 k(x,t) f(t) \, \dd t
\]
est compact. En effet, on vérifie que $T(\overline{B})$ est équicontinue et uniformément
bornée, puis on applique le théorème d'Ascoli\index{Ascoli, théorème d'}.
\end{exemple}

\begin{exemple}\label{ex:compact-hilbert}
\index{opérateur diagonal}
Sur $\ell^2$, soit $T$ défini par $T(e_n) = \lambda_n e_n$ où $(\lambda_n)$ est une
suite de scalaires. Alors $T \in \mathcal{K}(\ell^2)$ si et seulement si
$\lambda_n \to 0$.
\end{exemple}

%% ══════════════════════════════════════════════════════════════════
\section{Série de Neumann et inversibilité}
\label{sec:neumann}
%% ══════════════════════════════════════════════════════════════════

\index{série de Neumann|(}

\begin{theoreme}[Série de Neumann]\label{thm:neumann}
\index{série de Neumann}
Soit $E$ un espace de Banach et $T \in \mathcal{L}(E)$ avec $\norm{T} < 1$. Alors
$\Id - T$ est inversible dans $\mathcal{L}(E)$ et
\[
    (\Id - T)^{-1} = \sum_{n=0}^{\infty} T^n,
\]
avec $\norm{(\Id - T)^{-1}} \leq \dfrac{1}{1 - \norm{T}}$.
\end{theoreme}

\begin{proof}
Posons $S_N = \sum_{n=0}^{N} T^n$. Comme $\norm{T^n} \leq \norm{T}^n$ et
$\norm{T} < 1$, la série $\sum \norm{T^n}$ converge. Donc $(S_N)$ est de Cauchy dans
$\mathcal{L}(E)$ qui est complet (car $E$ est de Banach). Soit
$S = \sum_{n=0}^{\infty} T^n$ sa limite.

On a $S_N(\Id - T) = \Id - T^{N+1}$ et $(\Id - T)S_N = \Id - T^{N+1}$.
Comme $\norm{T^{N+1}} \leq \norm{T}^{N+1} \to 0$, on obtient en passant à la limite
$S(\Id - T) = (\Id - T)S = \Id$. Enfin,
\[
    \norm{S} \leq \sum_{n=0}^{\infty} \norm{T}^n = \frac{1}{1 - \norm{T}}.
    \qedhere
\]
\end{proof}

\begin{corollaire}[Ouverture de $GL(E)$]\label{cor:GL-ouvert}
\index{GL(E)@$GL(E)$!ouverture}
L'ensemble $GL(E)$ des opérateurs inversibles est un ouvert de $\mathcal{L}(E)$, et
l'application $T \mapsto T^{-1}$ est continue sur $GL(E)$.
\end{corollaire}

\begin{proof}
Soit $T_0 \in GL(E)$ et $T \in \mathcal{L}(E)$ avec
$\norm{T - T_0} < \norm{T_0^{-1}}^{-1}$. Écrivons
$T = T_0(\Id - T_0^{-1}(T_0 - T))$.
On a $\norm{T_0^{-1}(T_0 - T)} \leq \norm{T_0^{-1}}\cdot\norm{T_0 - T} < 1$,
donc $\Id - T_0^{-1}(T_0 - T)$ est inversible par le théorème de Neumann. Par suite
$T$ est inversible, et
\[
    T^{-1} = \left(\sum_{n=0}^{\infty} (T_0^{-1}(T_0 - T))^n\right) T_0^{-1}.
\]
La continuité de l'inversion\index{inversion!continuité} découle de l'estimation
$\norm{T^{-1} - T_0^{-1}} \leq \norm{T_0^{-1}}^2 \norm{T - T_0} / (1 -
\norm{T_0^{-1}}\norm{T-T_0})$.
\end{proof}

\begin{corollaire}[Perturbation de l'identité]\label{cor:perturbation-id}
\index{perturbation!de l'identité}
Soit $E$ un espace de Banach et $T \in \mathcal{L}(E)$ avec $\norm{T} < 1$. Alors
$\Id + T$ est inversible et
\[
    (\Id + T)^{-1} = \sum_{n=0}^{\infty} (-1)^n T^n.
\]
\end{corollaire}

\index{série de Neumann|)}

%% ══════════════════════════════════════════════════════════════════
\section{Spectre et résolvante}
\label{sec:spectre-resolvante}
%% ══════════════════════════════════════════════════════════════════

\index{spectre|(}
\index{résolvante|(}

Dans toute cette section, $E$ est un espace de Banach \emph{complexe}
($\K = \C$).

\begin{definition}[Spectre, résolvante]\label{def:spectre}
\index{spectre!définition}
\index{résolvante!ensemble}
Soit $T \in \mathcal{L}(E)$. L'\textbf{ensemble résolvant}\index{ensemble résolvant}
de $T$ est
\[
    \rho(T) := \{\lambda \in \C : \lambda\Id - T \in GL(E)\}.
\]
Le \textbf{spectre} de $T$ est $\sigma(T) := \C \setminus \rho(T)$.

Pour $\lambda \in \rho(T)$, l'opérateur
$R(\lambda, T) := (\lambda\Id - T)^{-1} \in \mathcal{L}(E)$
est appelé la \textbf{résolvante}\index{résolvante!opérateur} de $T$ en $\lambda$.
\end{definition}

\begin{definition}[Décomposition du spectre]\label{def:decomposition-spectre}
\index{spectre!ponctuel}\index{spectre!continu}\index{spectre!résiduel}
Le spectre se décompose en :
\begin{enumerate}[label=(\roman*)]
    \item le \textbf{spectre ponctuel} $\sigma_p(T) := \{\lambda \in \C :
          \lambda\Id - T \text{ n'est pas injectif}\}$ (ensemble des valeurs propres\index{valeur propre}) ;
    \item le \textbf{spectre continu} $\sigma_c(T) := \{\lambda \in \sigma(T) \setminus
          \sigma_p(T) : \overline{\ran(\lambda\Id - T)} = E\}$ ;
    \item le \textbf{spectre résiduel} $\sigma_r(T) := \sigma(T) \setminus
          (\sigma_p(T) \cup \sigma_c(T))$.
\end{enumerate}
\end{definition}

\begin{proposition}[Propriétés de la résolvante]\label{prop:resolvante-proprietes}
\index{résolvante!propriétés}
\index{identité de la résolvante}
Soit $T \in \mathcal{L}(E)$.
\begin{enumerate}[label=(\roman*)]
    \item $\rho(T)$ est un ouvert de $\C$ et $\sigma(T)$ est fermé.
    \item \textbf{Identité de la résolvante} : pour $\lambda, \mu \in \rho(T)$,
    \[
        R(\lambda,T) - R(\mu,T) = (\mu - \lambda) R(\lambda,T) R(\mu,T).
    \]
    \item L'application $\lambda \mapsto R(\lambda,T)$ est analytique\index{résolvante!analyticité}
          sur $\rho(T)$ (au sens de la norme d'opérateur).
\end{enumerate}
\end{proposition}

\begin{proof}
(i) Découle du corollaire~\ref{cor:GL-ouvert} : pour $\lambda_0 \in \rho(T)$,
l'opérateur $\lambda\Id - T$ reste inversible pour $\lambda$ proche de $\lambda_0$.

(ii) On écrit
$R(\lambda,T) - R(\mu,T) = R(\lambda,T)[(\mu\Id - T) - (\lambda\Id - T)]R(\mu,T)
= (\mu - \lambda) R(\lambda,T)R(\mu,T)$.

(iii) L'identité de la résolvante donne
$\frac{R(\lambda,T) - R(\mu,T)}{\lambda - \mu} = -R(\lambda,T)R(\mu,T) \to -R(\mu,T)^2$
quand $\lambda \to \mu$, ce qui montre la dérivabilité complexe.
\end{proof}

\begin{theoreme}[Le spectre est compact et non vide]\label{thm:spectre-compact}
\index{spectre!compacité}
\index{spectre!non-vacuité}
Soit $T \in \mathcal{L}(E)$ avec $E \neq \{0\}$. Alors :
\begin{enumerate}[label=(\roman*)]
    \item $\sigma(T) \subset \overline{B}(0, \norm{T})$ (le spectre est borné) ;
    \item $\sigma(T) \neq \emptyset$.
\end{enumerate}
\end{theoreme}

\begin{proof}
(i) Si $\abs{\lambda} > \norm{T}$, alors $\lambda\Id - T = \lambda(\Id - T/\lambda)$.
Comme $\norm{T/\lambda} < 1$, la série de Neumann montre que $\Id - T/\lambda$ est
inversible, donc $\lambda \in \rho(T)$.

(ii) Supposons $\sigma(T) = \emptyset$. Alors $R(\cdot, T)$ est une fonction entière
à valeurs dans $\mathcal{L}(E)$. Pour $\abs{\lambda} > \norm{T}$,
\[
    R(\lambda, T) = \frac{1}{\lambda}\left(\Id - \frac{T}{\lambda}\right)^{-1}
    = \frac{1}{\lambda}\sum_{n=0}^{\infty} \frac{T^n}{\lambda^n},
\]
donc $\norm{R(\lambda,T)} \leq \frac{1}{\abs{\lambda} - \norm{T}} \to 0$ quand
$\abs{\lambda} \to \infty$. Par le théorème de Liouville\index{Liouville, théorème de}
vectoriel, $R(\cdot,T) = 0$, ce qui est absurde car $R(\lambda,T)$ est inversible.
\end{proof}

\begin{definition}[Rayon spectral]\label{def:rayon-spectral}
\index{rayon spectral!définition}
Le \textbf{rayon spectral} de $T \in \mathcal{L}(E)$ est
\[
    r(T) := \sup\{\abs{\lambda} : \lambda \in \sigma(T)\}.
\]
\end{definition}

\begin{theoreme}[Formule du rayon spectral (Gelfand)]\label{thm:rayon-spectral}
\index{rayon spectral!formule de Gelfand}
\index{Gelfand, formule de}
Soit $E$ un espace de Banach complexe et $T \in \mathcal{L}(E)$. Alors
\[
    r(T) = \lim_{n \to \infty} \norm{T^n}^{1/n} = \inf_{n \geq 1} \norm{T^n}^{1/n}.
\]
\end{theoreme}

\begin{proof}
\emph{Étape 1 : $r(T) \leq \inf_n \norm{T^n}^{1/n}$.}
Soit $\lambda \in \sigma(T)$. On montre que $\lambda^n \in \sigma(T^n)$ pour tout
$n \geq 1$, car
$\lambda^n\Id - T^n = (\lambda\Id - T)(\lambda^{n-1}\Id + \lambda^{n-2}T + \cdots + T^{n-1})$.
Si $\lambda^n \notin \sigma(T^n)$, alors $\lambda^n\Id - T^n$ serait inversible, et
comme $\lambda\Id - T$ est un facteur gauche, $\lambda\Id - T$ aurait un inverse à
gauche ; un argument similaire avec la factorisation à droite donne un inverse à droite,
d'où $\lambda \notin \sigma(T)$, contradiction.
Donc $\abs{\lambda}^n \leq \norm{T^n}$ par le théorème~\ref{thm:spectre-compact}(i)
appliqué à $T^n$, et $\abs{\lambda} \leq \norm{T^n}^{1/n}$ pour tout $n$.

\emph{Étape 2 : $\limsup_n \norm{T^n}^{1/n} \leq r(T)$.}
Pour $\abs{\lambda} > r(T)$, on a $\lambda \in \rho(T)$ et la série de Laurent de
$R(\lambda, T)$ en $\infty$ s'écrit $R(\lambda,T) = \sum_{n=0}^{\infty} T^n / \lambda^{n+1}$.
Cette série converge pour $\abs{\lambda} > r(T)$ (car $R(\cdot,T)$ est analytique sur
$\rho(T)$ qui contient $\{\abs{\lambda} > r(T)\}$).

Par le lemme de Hadamard\index{Hadamard, formule de} pour les séries de Laurent, le
rayon de convergence de $\sum T^n z^{n+1}$ est au moins $1/r(T)$.
Rappelons que le rayon de convergence est $1/\limsup_n \norm{T^n}^{1/n}$.
Donc $\limsup_n \norm{T^n}^{1/n} \leq r(T)$.

\emph{Conclusion.}
Des étapes 1 et 2, $r(T) \leq \inf_n \norm{T^n}^{1/n}
\leq \liminf_n \norm{T^n}^{1/n} \leq \limsup_n \norm{T^n}^{1/n} \leq r(T)$.
Donc la limite existe et toutes les quantités sont égales.
\end{proof}

\begin{remarque}\label{rem:rayon-spectral-sous-multiplicatif}
\index{suite sous-multiplicative}
L'existence de la limite $\lim \norm{T^n}^{1/n}$ peut aussi se déduire du lemme de
Fekete\index{Fekete, lemme de} : si $(a_n)$ est sous-additive (i.e.\ $a_{n+m} \leq
a_n + a_m$), alors $a_n/n \to \inf a_n/n$. On applique ceci à $a_n = \log\norm{T^n}$.
\end{remarque}

\begin{exemple}\label{ex:spectre-shift}
\index{opérateur de décalage!spectre}
Considérons le shift à droite\index{shift à droite} $S : \ell^2 \to \ell^2$ défini par
$S(x_1, x_2, \ldots) = (0, x_1, x_2, \ldots)$. On a $\norm{S} = 1$,
$\sigma_p(S) = \emptyset$ et $\sigma(S) = \overline{B}(0,1)$.
Le shift à gauche\index{shift à gauche}
$S^*(x_1, x_2, \ldots) = (x_2, x_3, \ldots)$ a $\sigma_p(S^*) = B(0,1)$ (le disque
ouvert) et $\sigma(S^*) = \overline{B}(0,1)$.
\end{exemple}

\index{spectre|)}
\index{résolvante|)}

%% ══════════════════════════════════════════════════════════════════
\section{Adjoint sur les espaces de Hilbert}
\label{sec:adjoint-hilbert}
%% ══════════════════════════════════════════════════════════════════

\index{opérateur!adjoint|(}

Dans toute cette section, $H$ et $K$ désignent des espaces de Hilbert\index{espace de Hilbert}
sur $\K$ ($\R$ ou $\C$).

\begin{theoreme}[Existence de l'adjoint]\label{thm:adjoint-existence}
\index{opérateur!adjoint!existence}
\index{Riesz, théorème de représentation de}
Soit $T \in \mathcal{L}(H, K)$. Il existe un unique opérateur
$T^* \in \mathcal{L}(K, H)$\index{T*@$T^*$} tel que
\[
    \inner{Tx}{y}_K = \inner{x}{T^*y}_H \quad
    \text{pour tous } x \in H,\, y \in K.
\]
De plus, $\norm{T^*} = \norm{T}$.
\end{theoreme}

\begin{proof}
Pour $y \in K$ fixé, l'application $\varphi_y : x \mapsto \inner{Tx}{y}_K$ est une
forme linéaire continue sur $H$ (si $\K = \C$, elle est antilinéaire en $y$ et linéaire
en $x$; on considère $\varphi_y$ comme forme linéaire en $x$). On a
$\abs{\varphi_y(x)} \leq \norm{T}\norm{x}\norm{y}$, donc $\varphi_y \in H'$ avec
$\norm{\varphi_y} \leq \norm{T}\norm{y}$.

Par le théorème de représentation de Riesz\index{Riesz, théorème de représentation de},
il existe un unique $z_y \in H$ tel que $\varphi_y(x) = \inner{x}{z_y}_H$ pour tout
$x \in H$, avec $\norm{z_y} = \norm{\varphi_y} \leq \norm{T}\norm{y}$.

On pose $T^*y := z_y$. L'unicité dans le théorème de Riesz assure que $T^*$ est bien
défini. Vérifions la linéarité (dans le cas $\K = \C$) : pour $y_1, y_2 \in K$ et
$\alpha \in \C$,
\[
    \inner{x}{T^*(\alpha y_1 + y_2)}_H = \inner{Tx}{\alpha y_1 + y_2}_K
    = \bar{\alpha}\inner{Tx}{y_1}_K + \inner{Tx}{y_2}_K
    = \inner{x}{\alpha T^*y_1 + T^*y_2}_H
\]
pour tout $x$, donc $T^*(\alpha y_1 + y_2) = \alpha T^*y_1 + T^*y_2$.
La bornitude est assurée par $\norm{T^*y} \leq \norm{T}\norm{y}$, d'où
$\norm{T^*} \leq \norm{T}$.

Réciproquement, pour tout $x \in H$,
$\norm{Tx}^2 = \inner{Tx}{Tx} = \inner{x}{T^*Tx} \leq \norm{x}\norm{T^*Tx}
\leq \norm{x}\norm{T^*}\norm{T}\norm{x}$,
donc $\norm{T}^2 \leq \norm{T^*}\norm{T}$, d'où $\norm{T} \leq \norm{T^*}$.
\end{proof}

\begin{proposition}[Propriétés de l'adjoint]\label{prop:proprietes-adjoint}
\index{opérateur!adjoint!propriétés}
Soient $S, T \in \mathcal{L}(H,K)$, $R \in \mathcal{L}(K,G)$ (où $G$ est un Hilbert)
et $\alpha \in \K$. Alors :
\begin{enumerate}[label=(\roman*)]
    \item $(S + T)^* = S^* + T^*$ et $(\alpha T)^* = \bar{\alpha} T^*$ ;
    \item $(T^*)^* = T$ ;
    \item $(RT)^* = T^* R^*$\index{opérateur!adjoint!composition} ;
    \item $\norm{T^*T} = \norm{TT^*} = \norm{T}^2$\index{identité $C^*$} ;
    \item $\Ker T^* = (\ran T)^\perp$\index{noyau!et image de l'adjoint} et
          $\overline{\ran T} = (\Ker T^*)^\perp$.
\end{enumerate}
\end{proposition}

\begin{proof}
(i) et (ii) découlent directement de la définition.

(iii) Pour tous $x \in H$ et $z \in G$ :
\[
    \inner{RTx}{z}_G = \inner{Tx}{R^*z}_K = \inner{x}{T^*R^*z}_H,
\]
donc $(RT)^* = T^*R^*$ par unicité de l'adjoint.

(iv) On a $\norm{T^*T} \leq \norm{T^*}\norm{T} = \norm{T}^2$.
Inversement, $\norm{Tx}^2 = \inner{T^*Tx}{x} \leq \norm{T^*Tx}\norm{x}
\leq \norm{T^*T}\norm{x}^2$,
donc $\norm{T}^2 \leq \norm{T^*T}$. L'identité $\norm{T^*T} = \norm{T}^2$ est
appelée l'\textbf{identité $C^*$}\index{identité $C^*$}.
L'égalité $\norm{TT^*} = \norm{T}^2$ s'obtient en appliquant le résultat à $T^*$
et en utilisant $(T^*)^* = T$.

(v) On a $y \in \Ker T^*$ ssi $T^*y = 0$ ssi $\inner{x}{T^*y} = 0$ pour tout $x$
ssi $\inner{Tx}{y} = 0$ pour tout $x$ ssi $y \in (\ran T)^\perp$.
La seconde égalité en découle par passage à l'orthogonal.
\end{proof}

\index{opérateur!adjoint|)}

%% ══════════════════════════════════════════════════════════════════
\section{Opérateurs normaux, auto-adjoints et unitaires}
\label{sec:normaux-autoadjoints-unitaires}
%% ══════════════════════════════════════════════════════════════════

\index{opérateur!normal|(}
\index{opérateur!auto-adjoint|(}
\index{opérateur!unitaire|(}

On travaille sur un espace de Hilbert $H$ (sur $\C$ sauf mention contraire).

\begin{definition}[Opérateurs normaux, auto-adjoints, unitaires]\label{def:normal-sa-unitaire}
Soit $T \in \mathcal{L}(H)$.
\begin{enumerate}[label=(\roman*)]
    \item $T$ est \textbf{normal}\index{opérateur!normal!définition} si $T^*T = TT^*$.
    \item $T$ est \textbf{auto-adjoint} (ou \textbf{hermitien})\index{opérateur!auto-adjoint!définition}%
          \index{opérateur hermitien|see{opérateur auto-adjoint}} si $T^* = T$.
    \item $T$ est \textbf{unitaire}\index{opérateur!unitaire!définition} si
          $T^*T = TT^* = \Id$.
    \item $T$ est une \textbf{projection orthogonale}\index{projection orthogonale}
          si $T^2 = T$ et $T^* = T$.
\end{enumerate}
\end{definition}

\begin{remarque}\label{rem:inclusions-normaux}
On a les inclusions : auto-adjoint $\Rightarrow$ normal, et unitaire $\Rightarrow$ normal.
Tout opérateur $T \in \mathcal{L}(H)$ s'écrit de manière unique
$T = A + iB$ où $A = \frac{T+T^*}{2}$ et $B = \frac{T-T^*}{2i}$ sont auto-adjoints%
\index{décomposition en parties réelle et imaginaire}.
$T$ est normal si et seulement si $AB = BA$.
\end{remarque}

\begin{proposition}[Caractérisation des opérateurs normaux]\label{prop:caract-normal}
\index{opérateur!normal!caractérisation}
$T \in \mathcal{L}(H)$ est normal si et seulement si $\norm{Tx} = \norm{T^*x}$ pour
tout $x \in H$.
\end{proposition}

\begin{proof}
On a $\norm{Tx}^2 = \inner{T^*Tx}{x}$ et $\norm{T^*x}^2 = \inner{TT^*x}{x}$. Si
$T$ est normal, ces quantités sont égales. Réciproquement, si
$\inner{T^*Tx}{x} = \inner{TT^*x}{x}$ pour tout $x$, alors
$\inner{(T^*T - TT^*)x}{x} = 0$ pour tout $x$. Par l'identité de polarisation%
\index{identité de polarisation}, l'opérateur auto-adjoint $T^*T - TT^*$
est nul.
\end{proof}

\begin{proposition}[Valeurs propres d'un opérateur auto-adjoint]\label{prop:vp-autoadjoint}
\index{opérateur!auto-adjoint!valeurs propres}
Soit $T \in \mathcal{L}(H)$ auto-adjoint.
\begin{enumerate}[label=(\roman*)]
    \item Les valeurs propres de $T$ sont réelles.
    \item Les sous-espaces propres correspondant à des valeurs propres distinctes
          sont orthogonaux.
\end{enumerate}
\end{proposition}

\begin{proof}
(i) Soit $Tx = \lambda x$ avec $x \neq 0$. Alors
$\lambda\norm{x}^2 = \inner{\lambda x}{x} = \inner{Tx}{x} = \inner{x}{Tx}
= \overline{\inner{Tx}{x}} = \bar{\lambda}\norm{x}^2$,
donc $\lambda = \bar{\lambda}$, i.e.\ $\lambda \in \R$.

(ii) Soient $Tx = \lambda x$ et $Ty = \mu y$ avec $\lambda \neq \mu$. Alors
$\lambda\inner{x}{y} = \inner{Tx}{y} = \inner{x}{Ty} = \mu\inner{x}{y}$
(car $\bar{\mu} = \mu$ par (i)), donc $(\lambda - \mu)\inner{x}{y} = 0$ et
$\inner{x}{y} = 0$.
\end{proof}

\begin{theoreme}[Spectre d'un opérateur auto-adjoint]\label{thm:spectre-autoadjoint}
\index{opérateur!auto-adjoint!spectre}
\index{spectre!d'un opérateur auto-adjoint}
Soit $T \in \mathcal{L}(H)$ auto-adjoint. Alors $\sigma(T) \subset \R$.
Plus précisément, $\sigma(T) \subset [m, M]$ où
$m = \inf_{\norm{x}=1} \inner{Tx}{x}$ et $M = \sup_{\norm{x}=1} \inner{Tx}{x}$,
et $\{m, M\} \subset \sigma(T)$.
\end{theoreme}

\begin{proof}
\emph{$\sigma(T) \subset \R$.}
Soit $\lambda = \alpha + i\beta$ avec $\beta \neq 0$. Pour tout $x \in H$,
\begin{align*}
    \norm{(\lambda\Id - T)x}^2
    &= \norm{(\alpha\Id - T)x + i\beta x}^2 \\
    &= \norm{(\alpha\Id - T)x}^2 + \beta^2\norm{x}^2
       + 2\beta\, \im\inner{(\alpha\Id - T)x}{x}.
\end{align*}
Or $\inner{(\alpha\Id - T)x}{x} = \alpha\norm{x}^2 - \inner{Tx}{x} \in \R$ car $T$
est auto-adjoint. Donc $\im\inner{(\alpha\Id-T)x}{x} = 0$ et
\[
    \norm{(\lambda\Id - T)x}^2 \geq \beta^2\norm{x}^2.
\]
Ceci montre que $\lambda\Id - T$ est injectif et d'image fermée. De même,
$\bar{\lambda}\Id - T = (\lambda\Id - T)^*$ est injectif, donc
$\ran(\lambda\Id - T)$ est dense (car $\Ker(\lambda\Id-T)^* = (\ran(\lambda\Id-T))^\perp = \{0\}$).
Fermée et dense, l'image est $H$ tout entier, donc $\lambda \in \rho(T)$.

\emph{$\sigma(T) \subset [m, M]$.}
Soit $\lambda > M$. Alors pour tout $x$ avec $\norm{x} = 1$,
$\inner{(\lambda\Id - T)x}{x} = \lambda - \inner{Tx}{x} \geq \lambda - M > 0$.
L'opérateur $\lambda\Id - T$ est auto-adjoint et vérifie
$\inner{(\lambda\Id - T)x}{x} \geq (\lambda - M)\norm{x}^2$, d'où il est inversible
(par un argument analogue au cas $\beta \neq 0$, en utilisant la minoration directe).
De même pour $\lambda < m$.

\emph{$M \in \sigma(T)$.}
Par définition du supremum, il existe une suite $(x_n)$ avec $\norm{x_n} = 1$ et
$\inner{Tx_n}{x_n} \to M$. On a
$\norm{(T - M\Id)x_n}^2 = \norm{Tx_n}^2 - 2M\inner{Tx_n}{x_n} + M^2
\leq \norm{T}^2 - 2M\inner{Tx_n}{x_n} + M^2$.
Comme $\norm{T} = \max(\abs{m}, \abs{M})$ pour un opérateur auto-adjoint (on admet ce
point), un calcul montre que $\norm{(T-M\Id)x_n}^2 \to M^2 - M^2 = 0$
lorsqu'on raffine l'argument en utilisant $\norm{Tx_n}^2 \leq \norm{T}\inner{Tx_n}{x_n}$
(qui découle de $\norm{Tx_n}^2 = \inner{T^2 x_n}{x_n} \leq \norm{T}\inner{Tx_n}{x_n}$
puisque $\inner{T^2x_n}{x_n} \leq \norm{T}\inner{Tx_n}{x_n}$ par l'inégalité
$T^2 \leq \norm{T}T$ au sens des formes, car $\inner{Tx}{x} \leq M \leq \norm{T}$).
Donc $T - M\Id$ n'est pas inversible par en dessous, ce qui donne $M \in \sigma(T)$.
L'argument pour $m$ est analogue en considérant $-T$.
\end{proof}

\begin{corollaire}\label{cor:sa-norme-spectre}
\index{opérateur!auto-adjoint!norme et spectre}
Si $T$ est auto-adjoint, $\norm{T} = r(T) = \sup_{\norm{x}=1} \abs{\inner{Tx}{x}}$.
\end{corollaire}

\begin{proposition}[Opérateur auto-adjoint positif]\label{prop:operateur-positif}
\index{opérateur!positif}
Soit $T \in \mathcal{L}(H)$ auto-adjoint. On dit que $T$ est \textbf{positif}
(noté $T \geq 0$) si $\inner{Tx}{x} \geq 0$ pour tout $x \in H$.
Alors $\sigma(T) \subset [0, +\infty[$.
\end{proposition}

\begin{proposition}[Caractérisation des opérateurs unitaires]\label{prop:caract-unitaire}
\index{opérateur!unitaire!caractérisation}
Soit $U \in \mathcal{L}(H)$. Les assertions suivantes sont équivalentes :
\begin{enumerate}[label=(\roman*)]
    \item $U$ est unitaire ;
    \item $U$ est surjectif et isométrique ($\norm{Ux} = \norm{x}$ pour tout $x$) ;
    \item $U$ est surjectif et $\inner{Ux}{Uy} = \inner{x}{y}$ pour tous $x, y$.
\end{enumerate}
De plus, $\sigma(U) \subset \{z \in \C : \abs{z} = 1\}$ (le cercle unité).
\end{proposition}

\begin{proof}
$(i) \Rightarrow (iii)$ : $\inner{Ux}{Uy} = \inner{x}{U^*Uy} = \inner{x}{y}$, et $U$
est inversible donc surjectif.

$(iii) \Rightarrow (ii)$ : prendre $y = x$.

$(ii) \Rightarrow (i)$ : $U$ isométrique donne $\inner{U^*Ux}{x} = \norm{Ux}^2 =
\norm{x}^2 = \inner{x}{x}$, donc $U^*U = \Id$ par polarisation.
Comme $U$ est surjectif et isométrique, $U$ est bijectif, et $U^{-1} = U^*$, d'où
$UU^* = \Id$.

Pour le spectre : si $\lambda \in \sigma_p(U)$, alors $Ux = \lambda x$ avec $x \neq 0$,
et $\norm{x} = \norm{Ux} = \abs{\lambda}\norm{x}$, donc $\abs{\lambda} = 1$.
L'inclusion $\sigma(U) \subset \{\abs{z}=1\}$ découle de $\norm{U} = 1$ (d'où
$\sigma(U) \subset \overline{B}(0,1)$) et $\norm{U^{-1}} = \norm{U^*} = 1$ (d'où
$\sigma(U^{-1}) \subset \overline{B}(0,1)$, et comme
$\sigma(U^{-1}) = \{1/\lambda : \lambda \in \sigma(U)\}$, on obtient $\abs{\lambda} \geq 1$).
\end{proof}

\index{opérateur!normal|)}
\index{opérateur!auto-adjoint|)}
\index{opérateur!unitaire|)}

%% ══════════════════════════════════════════════════════════════════
\section{Exercices}
\label{sec:exos-chap3}
%% ══════════════════════════════════════════════════════════════════

\begin{exercice}[Norme d'opérateur]\label{exo:norme-operateur}
\index{norme!d'opérateur!calcul}
Soit $T : \ell^1 \to \ell^1$ défini par $T(x_n) = (x_n / n)_{n \geq 1}$.
\begin{enumerate}[label=(\alph*)]
    \item Montrer que $T \in \mathcal{L}(\ell^1)$ et calculer $\norm{T}$.
    \item $T$ est-il compact ?
\end{enumerate}
\end{exercice}

\begin{exercice}[Opérateur de Volterra]\label{exo:volterra}
\index{opérateur de Volterra}
Soit $V : L^2([0,1]) \to L^2([0,1])$ l'opérateur de Volterra défini par
$(Vf)(x) = \int_0^x f(t)\, \dd t$.
\begin{enumerate}[label=(\alph*)]
    \item Montrer que $V$ est borné et calculer $\norm{V^n}$ (on pourra montrer que
          $\norm{V^n} \leq 1/n!$).
    \item En déduire $r(V) = 0$, puis $\sigma(V) = \{0\}$.
    \item $V$ est-il compact ? normal ? auto-adjoint ?
    \item Calculer $V^*$.
\end{enumerate}
\end{exercice}

\begin{exercice}[Spectre du shift pondéré]\label{exo:shift-pondere}
\index{opérateur de décalage!pondéré}
Soit $(w_n)_{n \geq 1}$ une suite bornée de scalaires non nuls. On définit le shift
pondéré $T_w : \ell^2 \to \ell^2$ par $T_w(x_1, x_2, \ldots) = (0, w_1 x_1, w_2 x_2, \ldots)$.
\begin{enumerate}[label=(\alph*)]
    \item Montrer que $\norm{T_w} = \sup_n \abs{w_n}$.
    \item Calculer $\sigma_p(T_w)$ et $\sigma_p(T_w^*)$.
    \item Calculer $r(T_w)$ en fonction des poids $(w_n)$.
\end{enumerate}
\end{exercice}

\begin{exercice}[Inversibilité par perturbation]\label{exo:perturbation}
\index{perturbation!d'un opérateur inversible}
Soient $E$ un espace de Banach, $T \in GL(E)$ et $S \in \mathcal{L}(E)$ avec
$\norm{S} < \norm{T^{-1}}^{-1}$.
\begin{enumerate}[label=(\alph*)]
    \item Montrer que $T + S \in GL(E)$.
    \item Établir l'estimation $\norm{(T+S)^{-1} - T^{-1}} \leq
          \dfrac{\norm{T^{-1}}^2 \norm{S}}{1 - \norm{T^{-1}}\norm{S}}$.
\end{enumerate}
\end{exercice}

\begin{exercice}[Opérateurs compacts sur $\ell^2$]\label{exo:compact-l2}
\index{opérateur compact!sur $\ell^2$}
Soit $T \in \mathcal{L}(\ell^2)$ défini par $T(e_n) = a_n e_{\sigma(n)}$ où $\sigma$
est une permutation de $\N^*$ et $(a_n)$ est bornée.
\begin{enumerate}[label=(\alph*)]
    \item Montrer que $T$ est borné et calculer $\norm{T}$.
    \item Montrer que $T$ est compact si et seulement si $a_n \to 0$.
\end{enumerate}
\end{exercice}

\begin{exercice}[Opérateur normal]\label{exo:normal}
\index{opérateur!normal!exercice}
Soit $T \in \mathcal{L}(H)$ normal.
\begin{enumerate}[label=(\alph*)]
    \item Montrer que $\Ker T = \Ker T^*$.
    \item Montrer que $\norm{T^n} = \norm{T}^n$ pour tout $n \geq 1$.
    \item En déduire $r(T) = \norm{T}$.
    \item Montrer que si $T$ est compact et normal, alors il existe une valeur propre
          $\lambda$ de $T$ avec $\abs{\lambda} = \norm{T}$ (si $T \neq 0$).
\end{enumerate}
\end{exercice}

\begin{exercice}[Spectre d'un opérateur de multiplication]\label{exo:multiplication}
\index{opérateur de multiplication!spectre}
Soit $(X, \mu)$ un espace mesuré $\sigma$-fini et $\varphi \in L^\infty(X, \mu)$.
On définit $M_\varphi : L^2(X) \to L^2(X)$ par $M_\varphi f = \varphi f$.
\begin{enumerate}[label=(\alph*)]
    \item Montrer que $\norm{M_\varphi} = \norm{\varphi}_\infty$.
    \item Montrer que $M_\varphi^* = M_{\bar{\varphi}}$.
    \item Déterminer quand $M_\varphi$ est normal, auto-adjoint, unitaire, positif.
    \item Montrer que $\sigma(M_\varphi) = \operatorname{ess\,ran}(\varphi)$
          (l'image essentielle de $\varphi$).
\end{enumerate}
\end{exercice}

\begin{exercice}[Identité $C^*$]\label{exo:identite-cstar}
\index{identité $C^*$!exercice}
Soit $T \in \mathcal{L}(H)$.
\begin{enumerate}[label=(\alph*)]
    \item Montrer que $\norm{T^*T} = \norm{T}^2$ sans utiliser le théorème de
          l'adjoint, en prouvant directement l'inégalité
          $\norm{T}^2 \leq \norm{T^*T}$ via $\norm{Tx}^2 = \inner{T^*Tx}{x}$.
    \item En déduire que $\norm{T^{2^n}} = \norm{T}^{2^n}$ pour les opérateurs normaux,
          puis retrouver $r(T) = \norm{T}$ par la formule du rayon spectral.
\end{enumerate}
\end{exercice}

\begin{exercice}[Projection orthogonale et sous-espaces]\label{exo:projection-ortho}
\index{projection orthogonale!exercice}
Soient $P, Q \in \mathcal{L}(H)$ deux projections orthogonales.
\begin{enumerate}[label=(\alph*)]
    \item Montrer que $\sigma(P) \subset \{0, 1\}$.
    \item Montrer que $PQ$ est une projection orthogonale si et seulement si $PQ = QP$.
    \item Montrer que $P + Q$ est une projection orthogonale si et seulement si
          $PQ = 0$.
    \item Montrer que $\norm{P - Q} \leq 1$, avec égalité si et seulement si
          $\ran P \not\subset \ran Q$ ou $\ran Q \not\subset \ran P$.
\end{enumerate}
\end{exercice}


%%%%%%%%%%%%%%%%%%%%%%%%%%%%%%%%%%%%%%%%%%%%%%%%%%%%%%%%%%%%%%%%%%%%
%%%%%%%%%%%%%%%%%%%%%%%%%%%%%%%%%%%%%%%%%%%%%%%%%%%%%%%%%%%%%%%%%%%%
%%  CHAPITRE 4 : THÉORÈME DE HAHN-BANACH
%%%%%%%%%%%%%%%%%%%%%%%%%%%%%%%%%%%%%%%%%%%%%%%%%%%%%%%%%%%%%%%%%%%%
%%%%%%%%%%%%%%%%%%%%%%%%%%%%%%%%%%%%%%%%%%%%%%%%%%%%%%%%%%%%%%%%%%%%

\chapter{Théorème de Hahn-Banach}
\label{chap:hahn-banach}
\minitoc

\vspace{1em}

Le théorème de Hahn-Banach\index{Hahn-Banach, théorème de|(} est l'un des piliers de
l'analyse fonctionnelle. Sous sa forme analytique, il permet de prolonger des formes
linéaires continues définies sur un sous-espace à l'espace tout entier, en préservant
une majoration par une fonctionnelle sous-linéaire. Sous sa forme géométrique, il
fournit des résultats de séparation de convexes\index{séparation de convexes} par des
hyperplans\index{hyperplan}, qui jouent un rôle fondamental en optimisation et en théorie
des distributions.

Ce chapitre développe les deux formes du théorème, en explore les conséquences pour la
dualité\index{dualité} des espaces normés, et introduit les espaces réflexifs%
\index{espace réflexif}. L'axiome de Zorn\index{Zorn, lemme de} est utilisé de manière
essentielle dans les preuves.

%% ══════════════════════════════════════════════════════════════════
\section{Forme analytique réelle}
\label{sec:hb-reel}
%% ══════════════════════════════════════════════════════════════════

\index{Hahn-Banach, théorème de!forme analytique réelle|(}

\begin{definition}[Fonctionnelle sous-linéaire]\label{def:sous-lineaire}
\index{fonctionnelle sous-linéaire}
Soit $E$ un espace vectoriel réel. Une application $p : E \to \R$ est dite
\textbf{sous-linéaire} si :
\begin{enumerate}[label=(\roman*)]
    \item $p(\lambda x) = \lambda p(x)$ pour tout $\lambda > 0$ et tout $x \in E$
          (homogénéité positive) ;
    \item $p(x + y) \leq p(x) + p(y)$ pour tous $x, y \in E$ (sous-additivité).
\end{enumerate}
\end{definition}

\begin{remarque}\label{rem:semi-norme-sous-lineaire}
\index{semi-norme}
Toute semi-norme est sous-linéaire. Réciproquement, une fonctionnelle sous-linéaire
n'est pas nécessairement paire : on peut avoir $p(-x) \neq p(x)$.
\end{remarque}

\begin{theoreme}[Hahn-Banach analytique réel]\label{thm:hahn-banach-reel}
\index{Hahn-Banach, théorème de!forme analytique réelle}
Soit $E$ un espace vectoriel réel, $p : E \to \R$ une fonctionnelle sous-linéaire, $F$
un sous-espace vectoriel de $E$, et $f : F \to \R$ une forme linéaire telle que
$f(x) \leq p(x)$ pour tout $x \in F$. Alors il existe une forme linéaire
$\tilde{f} : E \to \R$ qui prolonge $f$ (i.e.\ $\tilde{f}|_F = f$) et vérifie
\[
    \tilde{f}(x) \leq p(x) \quad \text{pour tout } x \in E.
\]
\end{theoreme}

\begin{proof}
La preuve procède en deux étapes.

\emph{Étape 1 : Extension d'une dimension.}
Supposons $F \subsetneq E$ et soit $x_0 \in E \setminus F$. Posons
$G = F \oplus \R x_0 = \{x + t x_0 : x \in F,\, t \in \R\}$.
On cherche $g : G \to \R$ linéaire prolongeant $f$ avec $g \leq p$ sur $G$.
Nécessairement $g(x + t x_0) = f(x) + t\alpha$ où $\alpha = g(x_0)$ est à déterminer.

Pour $t > 0$, la condition $g(x + tx_0) \leq p(x + tx_0)$ s'écrit
$f(x) + t\alpha \leq p(x + tx_0)$, soit $\alpha \leq \frac{p(x+tx_0) - f(x)}{t}$.
En substituant $x/t \to x$, on obtient
$\alpha \leq p(y + x_0) - f(y)$ pour tout $y \in F$.

Pour $t < 0$, la condition s'écrit (en posant $t = -s$, $s > 0$)
$f(x) - s\alpha \leq p(x - sx_0)$, soit
$\alpha \geq f(z) - p(z - x_0)$ pour tout $z \in F$.

Il faut donc trouver $\alpha$ tel que
\[
    \sup_{z \in F} [f(z) - p(z - x_0)] \leq \alpha \leq
    \inf_{y \in F} [p(y + x_0) - f(y)].
\]
Ceci est possible car pour tous $y, z \in F$,
$f(y) + f(z) = f(y+z) \leq p(y+z) = p((y+x_0) + (z-x_0)) \leq p(y+x_0) + p(z-x_0)$,
d'où $f(z) - p(z-x_0) \leq p(y+x_0) - f(y)$.

\emph{Étape 2 : Application du lemme de Zorn\index{Zorn, lemme de}.}
Considérons l'ensemble $\mathcal{E}$ des couples $(G, g)$ où $G$ est un sous-espace
de $E$ contenant $F$, et $g : G \to \R$ est linéaire, prolonge $f$, et vérifie
$g \leq p$ sur $G$. On munit $\mathcal{E}$ de l'ordre $(G_1, g_1) \preceq (G_2, g_2)$
si $G_1 \subset G_2$ et $g_2|_{G_1} = g_1$.

$\mathcal{E}$ est non vide (car $(F, f) \in \mathcal{E}$). Si $\{(G_i, g_i)\}_{i \in I}$
est une chaîne totalement ordonnée, alors $(G_\infty, g_\infty)$ avec
$G_\infty = \bigcup_i G_i$ et $g_\infty|_{G_i} = g_i$ est un majorant.
Par le lemme de Zorn, $\mathcal{E}$ admet un élément maximal $(\tilde{E}, \tilde{f})$.

Si $\tilde{E} \neq E$, l'étape~1 permettrait d'étendre $\tilde{f}$ à un sous-espace
strictement plus grand, contredisant la maximalité. Donc $\tilde{E} = E$.
\end{proof}

\index{Hahn-Banach, théorème de!forme analytique réelle|)}

%% ══════════════════════════════════════════════════════════════════
\section{Forme analytique complexe}
\label{sec:hb-complexe}
%% ══════════════════════════════════════════════════════════════════

\index{Hahn-Banach, théorème de!forme analytique complexe|(}

\begin{lemme}[Lemme de passage réel-complexe]\label{lem:reel-complexe}
\index{forme linéaire!passage réel-complexe}
Soit $E$ un espace vectoriel complexe et $f : E \to \C$ une forme $\C$-linéaire.
Posons $u = \operatorname{Re} f$. Alors $u$ est $\R$-linéaire et
$f(x) = u(x) - i\, u(ix)$ pour tout $x \in E$.

Réciproquement, si $u : E \to \R$ est $\R$-linéaire et vérifie $u(ix) = -u(x)$
(condition automatiquement satisfaite si $u = \operatorname{Re} f$ pour une forme
$\C$-linéaire $f$), alors $f(x) := u(x) - i\, u(ix)$ est $\C$-linéaire.
\end{lemme}

\begin{proof}
Si $f$ est $\C$-linéaire et $u = \operatorname{Re} f$, alors
$f(x) = u(x) + i\, \operatorname{Im} f(x)$. Or
$f(ix) = if(x) = i\, u(x) - \operatorname{Im} f(x)$, donc
$u(ix) = \operatorname{Re} f(ix) = -\operatorname{Im} f(x)$, d'où
$f(x) = u(x) - i\, u(ix)$.

Réciproquement, soit $g(x) = u(x) - iu(ix)$. Alors $g$ est $\R$-linéaire. Vérifions
la $\C$-linéarité : $g(ix) = u(ix) - iu(i^2x) = u(ix) + iu(x) = i(u(x) - iu(ix)) = ig(x)$.
\end{proof}

\begin{theoreme}[Hahn-Banach analytique complexe]\label{thm:hahn-banach-complexe}
\index{Hahn-Banach, théorème de!forme analytique complexe}
Soit $E$ un espace vectoriel complexe, $p : E \to [0,+\infty[$ une semi-norme, $F$
un sous-espace vectoriel de $E$, et $f : F \to \C$ une forme $\C$-linéaire telle que
$\abs{f(x)} \leq p(x)$ pour tout $x \in F$. Alors il existe une forme $\C$-linéaire
$\tilde{f} : E \to \C$ qui prolonge $f$ et vérifie
$\abs{\tilde{f}(x)} \leq p(x)$ pour tout $x \in E$.
\end{theoreme}

\begin{proof}
Considérons $E$ comme espace vectoriel réel et posons $u = \operatorname{Re} f : F \to \R$.
Alors $u$ est $\R$-linéaire et $u(x) \leq \abs{f(x)} \leq p(x)$ pour tout $x \in F$.

Par le théorème de Hahn-Banach réel (théorème~\ref{thm:hahn-banach-reel}), il existe
$\tilde{u} : E \to \R$, $\R$-linéaire, prolongeant $u$, avec $\tilde{u}(x) \leq p(x)$
pour tout $x \in E$.

Posons $\tilde{f}(x) := \tilde{u}(x) - i\,\tilde{u}(ix)$. Par le
lemme~\ref{lem:reel-complexe}, $\tilde{f}$ est $\C$-linéaire. Sur $F$, on retrouve
$\tilde{f}|_F = u - i\,u(i\cdot) = f$ par le même lemme. Il reste à montrer la
majoration.

Pour $x \in E$, écrivons $\tilde{f}(x) = \abs{\tilde{f}(x)} e^{i\theta}$. Alors
\[
    \abs{\tilde{f}(x)} = e^{-i\theta}\tilde{f}(x) = \tilde{f}(e^{-i\theta}x)
    = \tilde{u}(e^{-i\theta}x) \leq p(e^{-i\theta}x) = p(x),
\]
la dernière égalité utilisant $p(\alpha x) = \abs{\alpha}p(x)$ pour la semi-norme.
\end{proof}

\begin{corollaire}[Prolongement de formes linéaires continues]\label{cor:prolongement-norme}
\index{Hahn-Banach, théorème de!prolongement isométrique}
\index{prolongement!de formes linéaires}
Soit $E$ un espace vectoriel normé, $F$ un sous-espace de $E$, et $f \in F'$.
Alors il existe $\tilde{f} \in E'$ tel que $\tilde{f}|_F = f$ et
$\norm{\tilde{f}}_{E'} = \norm{f}_{F'}$.
\end{corollaire}

\begin{proof}
On applique le théorème~\ref{thm:hahn-banach-complexe} (ou \ref{thm:hahn-banach-reel}
dans le cas réel) avec la semi-norme $p(x) = \norm{f}_{F'} \cdot \norm{x}_E$.
On obtient $\tilde{f}$ avec $\abs{\tilde{f}(x)} \leq \norm{f}\norm{x}$, donc
$\norm{\tilde{f}} \leq \norm{f}$. Comme $\tilde{f}$ prolonge $f$,
$\norm{\tilde{f}} \geq \norm{f}$.
\end{proof}

\begin{corollaire}[Séparation des points par le dual]\label{cor:separation-points}
\index{dual!séparation des points}
Soit $E$ un espace vectoriel normé et $x_0 \in E$ avec $x_0 \neq 0$. Alors il existe
$f \in E'$ tel que $\norm{f} = 1$ et $f(x_0) = \norm{x_0}$.
\end{corollaire}

\begin{proof}
Sur $F = \K x_0$, définissons $g(\lambda x_0) = \lambda \norm{x_0}$. Alors
$\norm{g}_{F'} = 1$ et $g(x_0) = \norm{x_0}$. Le corollaire~\ref{cor:prolongement-norme}
fournit $f \in E'$ prolongeant $g$ avec $\norm{f} = 1$.
\end{proof}

\begin{corollaire}\label{cor:norme-dual}
\index{norme!expression duale}
Pour tout $x \in E$,
$\norm{x} = \sup\{\abs{f(x)} : f \in E',\, \norm{f} \leq 1\}$.
\end{corollaire}

\index{Hahn-Banach, théorème de!forme analytique complexe|)}

%% ══════════════════════════════════════════════════════════════════
\section{Formes géométriques : séparation de convexes}
\label{sec:separation}
%% ══════════════════════════════════════════════════════════════════

\index{séparation de convexes|(}
\index{Hahn-Banach, théorème de!forme géométrique|(}

\begin{definition}[Hyperplan]\label{def:hyperplan}
\index{hyperplan!définition}
Un \textbf{hyperplan} de $E$ est un sous-espace affine de codimension~$1$, c'est-à-dire
un ensemble de la forme $H = \{x \in E : f(x) = \alpha\}$ pour une forme linéaire
$f \neq 0$ et un scalaire $\alpha \in \R$ (dans le cas réel). L'hyperplan est dit
\textbf{fermé}\index{hyperplan!fermé} si $f$ est continue.
\end{definition}

\begin{definition}[Fonctionnelle de Minkowski]\label{def:jauge-minkowski}
\index{fonctionnelle de Minkowski}\index{jauge|see{fonctionnelle de Minkowski}}
Soit $C$ un sous-ensemble convexe\index{ensemble convexe} de $E$ avec $0 \in
\operatorname{int}(C)$ (intérieur non vide contenant l'origine). La \textbf{fonctionnelle
de Minkowski} (ou \textbf{jauge}) de $C$ est
\[
    p_C(x) := \inf\{t > 0 : x/t \in C\} = \inf\{t > 0 : x \in tC\}.
\]
\end{definition}

\begin{proposition}[Propriétés de la jauge]\label{prop:proprietes-jauge}
\index{fonctionnelle de Minkowski!propriétés}
Soit $C$ un convexe ouvert\index{ensemble convexe!ouvert} avec $0 \in C$ dans un e.v.n.\
réel $E$. Alors :
\begin{enumerate}[label=(\roman*)]
    \item $p_C$ est sous-linéaire et continue ;
    \item $C = \{x \in E : p_C(x) < 1\}$ ;
    \item si $C$ est de plus symétrique ($C = -C$) et borné, alors $p_C$ est une norme
          équivalente à la norme de $E$.
\end{enumerate}
\end{proposition}

\begin{proof}
(i) \emph{Homogénéité positive :} pour $\lambda > 0$,
$p_C(\lambda x) = \inf\{t > 0 : \lambda x \in tC\} = \lambda \inf\{s > 0 : x \in sC\}
= \lambda p_C(x)$.

\emph{Sous-additivité :} Soient $s > p_C(x)$ et $t > p_C(y)$. Alors $x \in sC$ et
$y \in tC$. Par convexité,
$\frac{x+y}{s+t} = \frac{s}{s+t}\cdot\frac{x}{s} + \frac{t}{s+t}\cdot\frac{y}{t} \in C$,
donc $x + y \in (s+t)C$ et $p_C(x+y) \leq s+t$. En passant à la borne inférieure,
$p_C(x+y) \leq p_C(x) + p_C(y)$.

\emph{Continuité :} Comme $0 \in C$ ouvert, il existe $r > 0$ tel que $B(0,r) \subset C$.
Pour $x \in E$, $x \in (\norm{x}/r)B(0,r) \subset (\norm{x}/r)C$, donc
$p_C(x) \leq \norm{x}/r$. La sous-linéarité et cette majoration donnent la continuité.

(ii) Si $x \in C$ (ouvert), alors par continuité en $0$ de $t \mapsto x/t$, il existe
$t < 1$ avec $x \in tC$, i.e.\ $x/t \in C$, d'où $p_C(x) \leq t < 1$.
Réciproquement, si $p_C(x) < 1$, il existe $t < 1$ avec $x \in tC$.
Alors $x = t(x/t) + (1-t)\cdot 0 \in C$ par convexité.
\end{proof}

\begin{theoreme}[Séparation d'un point et d'un convexe]\label{thm:separation-point-convexe}
\index{séparation!point et convexe}
Soit $E$ un espace vectoriel normé réel, $C \subset E$ un convexe ouvert non vide, et
$x_0 \in E \setminus C$. Alors il existe $f \in E' \setminus \{0\}$ et $\alpha \in \R$
tels que
\[
    f(x) < \alpha \leq f(x_0) \quad \text{pour tout } x \in C.
\]
\end{theoreme}

\begin{proof}
Quitte à translater, on peut supposer $0 \in C$ (remplacer $C$ par $C - a$ pour un
$a \in C$ et $x_0$ par $x_0 - a$). Soit $p = p_C$ la jauge de $C$.
Comme $x_0 \notin C$, on a $p(x_0) \geq 1$ (par la
proposition~\ref{prop:proprietes-jauge}(ii)).

Sur $F = \R x_0$, définissons $g(tx_0) = t\, p(x_0)$ pour $t \in \R$.
Pour $t \geq 0$, $g(tx_0) = tp(x_0) = p(tx_0)$, donc $g \leq p$ sur $F$ (pour $t<0$,
$g(tx_0) = tp(x_0) < 0 \leq p(tx_0)$).

Par Hahn-Banach (théorème~\ref{thm:hahn-banach-reel}), il existe $f : E \to \R$
linéaire avec $f|_F = g$ et $f \leq p$ sur $E$.

Comme $p$ est continue et $f \leq p$, pour tout $x \in C$, $f(x) \leq p(x) < 1$.
D'autre part, $f(x_0) = g(x_0) = p(x_0) \geq 1$.
On prend $\alpha = 1$.

La continuité de $f$ découle de $f(x) \leq p(x) \leq \norm{x}/r$ et
$f(-x) \leq p(-x) \leq \norm{x}/r'$ (pour des constantes $r, r' > 0$), d'où
$\abs{f(x)} \leq C'\norm{x}$.
\end{proof}

\begin{theoreme}[Séparation stricte de deux convexes]\label{thm:separation-stricte}
\index{séparation!stricte}
\index{séparation!de deux convexes}
Soient $E$ un espace vectoriel normé réel, $A \subset E$ un convexe compact non vide et
$B \subset E$ un convexe fermé non vide, avec $A \cap B = \emptyset$. Alors il existe
$f \in E' \setminus \{0\}$ et $\alpha \in \R$ tels que
\[
    f(a) < \alpha < f(b) \quad \text{pour tous } a \in A,\, b \in B.
\]
\end{theoreme}

\begin{proof}
\emph{Étape 1.} Posons $C = B - A = \{b - a : b \in B,\, a \in A\}$.
L'ensemble $C$ est convexe. Montrons qu'il est fermé. Soit $(c_n) = (b_n - a_n)$ une
suite dans $C$ avec $c_n \to c$. Comme $A$ est compact, on extrait $(a_{n_k})$
convergente : $a_{n_k} \to a \in A$. Alors $b_{n_k} = c_{n_k} + a_{n_k} \to c + a$.
Comme $B$ est fermé, $c + a \in B$, donc $c = (c+a) - a \in B - A = C$.

Comme $A \cap B = \emptyset$, on a $0 \notin C$. Comme $A$ est compact et $B$ fermé,
$d(A, B) = \inf\{\norm{a-b} : a \in A,\, b \in B\} > 0$, ce qui implique que $0$
est à distance positive de $C$.

\emph{Étape 2.} L'ensemble $U = C + B(0, \delta)$ (où $\delta = d(A,B)/2$) est un
convexe ouvert ne contenant pas $0$. Par le
théorème~\ref{thm:separation-point-convexe}, il existe $f \in E' \setminus \{0\}$ avec
$f(x) < 1 \leq f(0) = 0$ pour tout $x \in U$... On ajuste l'argument :
appliquons directement le théorème de séparation à $U$ et $\{0\}$, on obtient
$f \in E'$ non nulle et $\beta$ tel que $f(x) < \beta \leq 0$ pour tout $x \in U$.

En particulier, pour $b \in B$ et $a \in A$ : $f(b-a) \leq f(b-a+u)$ pour de petits $u$,
et la borne donne $\sup_A f(a) + \delta' < \inf_B f(b)$ pour un certain $\delta' > 0$
(en exploitant la compacité de $A$ et l'ouverture de $U$).

Plus précisément : comme $f \leq 0$ sur $U \supset C + B(0,\delta)$, pour tout
$b - a \in C$ et $\norm{u} < \delta$ : $f(b-a+u) \leq 0$, donc
$f(b) - f(a) \leq -f(u)$ pour tout $\norm{u} < \delta$, d'où
$f(b) - f(a) \leq -\sup_{\norm{u}<\delta} f(u)$.
Comme $f \neq 0$, $\sup_{\norm{u}<\delta} f(u) = \delta\norm{f} > 0$, donc
$f(b) - f(a) \leq -\delta\norm{f} < 0$... Cela donne $f(b) < f(a)$.

Reconsidérons en prenant $f$ séparant $0$ de $U$ dans le sens $f(u) < 0 \leq f(0)$
pour $u \in U$. Alors pour tout $b \in B$, $a \in A$ : $f(b - a) < 0$, i.e.\
$f(b) < f(a)$ pour tous $b, a$. Donc $\sup_B f \leq \inf_A f$.

En raffinant avec le terme $B(0,\delta)$ : pour $\norm{v} < \delta$,
$f(b-a+v) < 0$, d'où $f(b) - f(a) < -f(v)$ pour tout tel $v$. En prenant $v$ tel
que $f(v)$ est proche de $\delta\norm{f}$ : $f(b) - f(a) < -\delta\norm{f} + \varepsilon$.
Donc $\sup_B f < \inf_A f - \delta\norm{f}/2$, et on choisit $\alpha$ entre ces deux
valeurs. On remplace $f$ par $-f$ pour obtenir la conclusion dans la forme souhaitée
$f(a) < \alpha < f(b)$.
\end{proof}

\begin{remarque}\label{rem:separation-complexe}
\index{séparation!cas complexe}
Les théorèmes de séparation s'étendent au cas complexe. Pour un espace vectoriel normé
complexe $E$, un convexe ouvert $C$ et $x_0 \notin C$, il existe $f \in E'$ (au sens
complexe) et $\alpha \in \R$ avec $\operatorname{Re} f(x) < \alpha \leq
\operatorname{Re} f(x_0)$ pour tout $x \in C$.
\end{remarque}

\index{séparation de convexes|)}
\index{Hahn-Banach, théorème de!forme géométrique|)}

%% ══════════════════════════════════════════════════════════════════
\section{Dual des espaces \texorpdfstring{$\ell^p$}{lp} et bidual}
\label{sec:dual-lp-bidual}
%% ══════════════════════════════════════════════════════════════════

\index{dual!de $\ell^p$|(}

\begin{theoreme}[Dual de $\ell^p$]\label{thm:dual-lp}
\index{dual!de $\ell^p$}
Pour $1 \leq p < \infty$, le dual $(\ell^p)' $ est isométriquement isomorphe à
$\ell^q$ où $\frac{1}{p} + \frac{1}{q} = 1$ (avec $q = \infty$ si $p = 1$).
L'isomorphisme est donné par
\[
    \Phi : \ell^q \to (\ell^p)', \quad
    \Phi(y)(x) = \sum_{n=1}^{\infty} x_n y_n
    \quad \text{pour } y = (y_n) \in \ell^q,\, x = (x_n) \in \ell^p.
\]
\end{theoreme}

\begin{proof}
\emph{$\Phi$ est bien défini et isométrique.}
Par l'inégalité de Hölder\index{Hölder, inégalité de},
$\abs{\Phi(y)(x)} \leq \norm{x}_p \norm{y}_q$,
donc $\Phi(y) \in (\ell^p)'$ avec $\norm{\Phi(y)} \leq \norm{y}_q$.

Pour l'égalité, si $1 < p < \infty$, on prend
$x_n = \abs{y_n}^{q-2} \bar{y}_n$ (avec $x_n = 0$ si $y_n = 0$).
Alors $\abs{x_n}^p = \abs{y_n}^{p(q-1)} = \abs{y_n}^q$, donc
$\norm{x}_p = \norm{y}_q^{q/p}$ et
$\Phi(y)(x) = \sum \abs{y_n}^q = \norm{y}_q^q$.
Ainsi $\norm{\Phi(y)} \geq \norm{y}_q^q / \norm{y}_q^{q/p} = \norm{y}_q$.
Le cas $p = 1$ est analogue en prenant $x = e_{n_0}$ où $\abs{y_{n_0}}$ est proche
de $\norm{y}_\infty$.

\emph{$\Phi$ est surjectif.}
Soit $f \in (\ell^p)'$. Posons $y_n = f(e_n)$ où $(e_n)$ est la base canonique.
Pour tout $N$, considérons $x^{(N)} = \sum_{n=1}^N x_n e_n$ avec les $x_n$ comme
ci-dessus. Alors $f(x^{(N)}) = \sum_{n=1}^N x_n y_n$ et
$\abs{f(x^{(N)})} \leq \norm{f}\norm{x^{(N)}}_p$.

Pour $1 < p < \infty$, en prenant $x_n = \abs{y_n}^{q-2}\bar{y}_n$ :
$\sum_{n=1}^N \abs{y_n}^q \leq \norm{f} \left(\sum_{n=1}^N
\abs{y_n}^q\right)^{1/p}$,
d'où $\left(\sum_{n=1}^N \abs{y_n}^q\right)^{1/q} \leq \norm{f}$ pour tout $N$.
Donc $y \in \ell^q$ avec $\norm{y}_q \leq \norm{f}$.

Par densité des suites à support fini dans $\ell^p$, $f = \Phi(y)$.
\end{proof}

\begin{remarque}\label{rem:dual-l-infty}
\index{dual!de $\ell^\infty$}
Le dual de $\ell^\infty$ est strictement plus grand que $\ell^1$ : il s'identifie à
l'espace $\operatorname{ba}(\N)$\index{ba(N)@$\operatorname{ba}(\N)$} des mesures
bornées finiment additives sur $\N$, qui contient des éléments non représentables par
des suites sommables (comme les limites de Banach\index{limite de Banach}).
\end{remarque}

\subsection{Bidual et injection canonique}

\index{bidual|(}
\index{injection canonique|(}

\begin{definition}[Bidual, injection canonique]\label{def:bidual}
\index{bidual!définition}
\index{injection canonique!définition}
Soit $E$ un espace vectoriel normé. Le \textbf{bidual} de $E$ est l'espace
$E'' := (E')'$\index{E''@$E''$}. L'\textbf{injection canonique} est l'application
\[
    J : E \to E'', \quad J(x)(f) = f(x) \quad \text{pour } x \in E,\, f \in E'.
\]
\end{definition}

\begin{theoreme}[L'injection canonique est isométrique]\label{thm:J-isometrique}
\index{injection canonique!isométrie}
L'application $J : E \to E''$ est une isométrie linéaire : $\norm{J(x)}_{E''} = \norm{x}_E$
pour tout $x \in E$.
\end{theoreme}

\begin{proof}
Pour tout $x \in E$,
\[
    \norm{J(x)}_{E''} = \sup_{\norm{f}_{E'} \leq 1} \abs{J(x)(f)}
    = \sup_{\norm{f} \leq 1} \abs{f(x)}.
\]
D'une part, $\abs{f(x)} \leq \norm{f}\norm{x} \leq \norm{x}$, donc
$\norm{J(x)} \leq \norm{x}$.

D'autre part, par le corollaire~\ref{cor:separation-points}, il existe $f_0 \in E'$
avec $\norm{f_0} = 1$ et $f_0(x) = \norm{x}$. Donc
$\norm{J(x)} \geq \abs{f_0(x)} = \norm{x}$.
\end{proof}

\begin{corollaire}\label{cor:J-injectif}
L'injection canonique $J$ est injective. En particulier, $E$ s'identifie
isométriquement à un sous-espace de $E''$.
\end{corollaire}

\index{injection canonique|)}

%% ══════════════════════════════════════════════════════════════════
\section{Espaces réflexifs et théorème de Banach-Alaoglu}
\label{sec:reflexifs-alaoglu}
%% ══════════════════════════════════════════════════════════════════

\index{espace réflexif|(}

\begin{definition}[Espace réflexif]\label{def:reflexif}
\index{espace réflexif!définition}
Un espace de Banach $E$ est dit \textbf{réflexif} si l'injection canonique
$J : E \to E''$ est surjective (et donc un isomorphisme isométrique).
\end{definition}

\begin{exemple}\label{ex:reflexifs}
\index{espace réflexif!exemples}
\begin{enumerate}[label=(\roman*)]
    \item Les espaces de Hilbert sont réflexifs (par le théorème de Riesz).
    \item $\ell^p$ est réflexif pour $1 < p < \infty$, car
          $(\ell^p)'' \cong (\ell^q)' \cong \ell^p$.
    \item $\ell^1$ n'est pas réflexif : $(\ell^1)' = \ell^\infty$ et
          $(\ell^\infty)' \supsetneq \ell^1$.
    \item $L^p(\mu)$ est réflexif pour $1 < p < \infty$.
    \item $c_0$\index{c0@$c_0$}, $\ell^1$, $\ell^\infty$, $L^1$, $L^\infty$,
          $C([0,1])$ ne sont pas réflexifs.
    \item Tout espace de Banach de dimension finie est réflexif.
\end{enumerate}
\end{exemple}

\begin{proposition}[Propriétés des espaces réflexifs]\label{prop:proprietes-reflexifs}
\index{espace réflexif!propriétés}
\begin{enumerate}[label=(\roman*)]
    \item Un espace de Banach réflexif est séparable si et seulement si son dual l'est.
    \item Un sous-espace fermé d'un espace réflexif est réflexif.
    \item Un espace de Banach est réflexif si et seulement si son dual l'est.
\end{enumerate}
\end{proposition}

\begin{definition}[Topologie faible-*]\label{def:topologie-faible-etoile}
\index{topologie faible-*}
\index{faible-*|see{topologie faible-*}}
Soit $E$ un espace vectoriel normé. La \textbf{topologie faible-$*$} sur $E'$,
notée $\sigma(E', E)$\index{sigma(E',E)@$\sigma(E',E)$}, est la topologie la moins
fine rendant continues toutes les applications $f \mapsto f(x)$ pour $x \in E$.
Une suite $(f_n)$ dans $E'$ converge faible-$*$ vers $f$ (noté $f_n \weakstarto f$)
si et seulement si $f_n(x) \to f(x)$ pour tout $x \in E$.
\end{definition}

\begin{remarque}\label{rem:faible-etoile-moins-fine}
\index{topologie faible-*!comparaison}
La topologie faible-$*$ est moins fine que la topologie faible $\sigma(E', E'')$, qui
est elle-même moins fine que la topologie forte (normique) de $E'$.
En dimension infinie, ces topologies sont toutes distinctes.
\end{remarque}

\begin{theoreme}[Banach-Alaoglu]\label{thm:banach-alaoglu}
\index{Banach-Alaoglu, théorème de}
\index{boule unité!compacité faible-*}
Soit $E$ un espace vectoriel normé. La boule unité fermée
$\overline{B}_{E'} = \{f \in E' : \norm{f} \leq 1\}$
est compacte pour la topologie faible-$*$.
\end{theoreme}

\begin{proof}[Idée de la preuve]
Pour chaque $x \in E$, soit $D_x = \{\alpha \in \K : \abs{\alpha} \leq \norm{x}\}$.
Chaque $D_x$ est compact. L'application
\[
    \Phi : \overline{B}_{E'} \to \prod_{x \in E} D_x, \quad f \mapsto (f(x))_{x \in E}
\]
est bien définie (car $\abs{f(x)} \leq \norm{f}\norm{x} \leq \norm{x}$), injective et
continue pour la topologie faible-$*$ sur $\overline{B}_{E'}$ et la topologie produit
sur $\prod D_x$.

Par le théorème de Tychonov\index{Tychonov, théorème de}, $\prod_{x \in E} D_x$ est
compact. Il suffit de montrer que $\Phi(\overline{B}_{E'})$ est fermé dans le produit.
Si un réseau $\Phi(f_\alpha)$ converge vers $(\xi_x)_{x \in E}$ dans le produit, alors
$f_\alpha(x) \to \xi_x$ pour tout $x$. On vérifie que $g : x \mapsto \xi_x$ est
linéaire (par passage à la limite) et bornée ($\abs{\xi_x} \leq \norm{x}$), donc
$g \in \overline{B}_{E'}$ et $\Phi(g) = (\xi_x)$. Ainsi $\Phi(\overline{B}_{E'})$ est
fermé, donc compact, et $\overline{B}_{E'}$ est compacte pour la topologie faible-$*$.
\end{proof}

\begin{remarque}\label{rem:alaoglu-reflexif}
\index{Banach-Alaoglu, théorème de!et réflexivité}
Si $E$ est réflexif, la topologie faible-$*$ sur $E' = E''$ coïncide avec la topologie
faible $\sigma(E', E'')$ via l'identification $J$. Le théorème de Banach-Alaoglu montre
alors que la boule unité de $E$ est faiblement compacte.
Réciproquement, un résultat profond (théorème d'Eberlein-\v{S}mulian%
\index{Eberlein-Smulian, théorème d'@Eberlein-\v{S}mulian, théorème d'}) affirme que
la boule unité d'un espace de Banach est faiblement compacte si et seulement si l'espace
est réflexif.
\end{remarque}

\index{espace réflexif|)}
\index{bidual|)}

%% ══════════════════════════════════════════════════════════════════
\section{Exercices}
\label{sec:exos-chap4}
%% ══════════════════════════════════════════════════════════════════

\begin{exercice}[Prolongement concret]\label{exo:prolongement-concret}
\index{Hahn-Banach, théorème de!exercice}
Soit $F = \{(x_n) \in \ell^\infty : \lim_{n\to\infty} x_n \text{ existe}\}$ et
$f : F \to \R$ définie par $f((x_n)) = \lim_{n\to\infty} x_n$.
\begin{enumerate}[label=(\alph*)]
    \item Montrer que $f \in F'$ et calculer $\norm{f}$.
    \item Par Hahn-Banach, $f$ se prolonge en $\tilde{f} \in (\ell^\infty)'$ avec
          $\norm{\tilde{f}} = 1$. Un tel prolongement est-il unique ?
    \item Montrer qu'un prolongement positif (i.e.\ $\tilde{f}(x) \geq 0$ si $x_n \geq 0$
          pour tout $n$) est invariant par translation\index{limite de Banach!invariance}
          : $\tilde{f}(Sx) = \tilde{f}(x)$ où $S$ est le shift.
\end{enumerate}
\end{exercice}

\begin{exercice}[Annihilateur]\label{exo:annihilateur}
\index{annihilateur}
Soit $E$ un espace vectoriel normé et $M \subset E$ un sous-espace. On définit
l'annihilateur $M^\perp = \{f \in E' : f(x) = 0 \text{ pour tout } x \in M\}$.
\begin{enumerate}[label=(\alph*)]
    \item Montrer que $M^\perp$ est un sous-espace fermé de $E'$.
    \item Montrer que $(\overline{M})^\perp = M^\perp$.
    \item Montrer que $E'/M^\perp$ est isométriquement isomorphe à $\overline{M}'$.
    \item Si $M$ est fermé, montrer que $M' \cong E'/M^\perp$ (isométriquement).
\end{enumerate}
\end{exercice}

\begin{exercice}[Hahn-Banach et sous-espaces de codimension finie]\label{exo:codim-finie}
\index{sous-espace!de codimension finie}
Soit $E$ un espace vectoriel normé de dimension infinie.
\begin{enumerate}[label=(\alph*)]
    \item Montrer que tout sous-espace fermé de codimension finie est l'intersection
          de noyaux de formes linéaires continues.
    \item Montrer que tout sous-espace de codimension~$1$ est soit fermé soit dense.
\end{enumerate}
\end{exercice}

\begin{exercice}[Séparation et enveloppe convexe fermée]\label{exo:conv-fermee}
\index{enveloppe convexe!fermée}
Soit $E$ un espace vectoriel normé et $A \subset E$.
\begin{enumerate}[label=(\alph*)]
    \item Montrer que $\overline{\conv(A)} = \bigcap\{C : C \text{ convexe fermé},\,
          A \subset C\}$.
    \item En déduire que $x \in \overline{\conv(A)}$ si et seulement si $f(x) \leq
          \sup_{a \in A} f(a)$ pour toute $f \in E'$.
\end{enumerate}
\end{exercice}

\begin{exercice}[Non-réflexivité de $c_0$]\label{exo:c0-non-reflexif}
\index{c0@$c_0$!non-réflexivité}
\begin{enumerate}[label=(\alph*)]
    \item Montrer que $(c_0)' \cong \ell^1$.
    \item En déduire $(c_0)'' \cong \ell^\infty$.
    \item Montrer que l'injection canonique $J : c_0 \to \ell^\infty$ est l'inclusion
          naturelle, et conclure que $c_0$ n'est pas réflexif.
\end{enumerate}
\end{exercice}

\begin{exercice}[Jauge et convexes]\label{exo:jauge}
\index{fonctionnelle de Minkowski!exercice}
Soit $C$ un convexe ouvert d'un espace vectoriel normé réel $E$ avec $0 \in C$.
\begin{enumerate}[label=(\alph*)]
    \item Montrer que $p_C$ est la plus grande fonctionnelle sous-linéaire $q$ vérifiant
          $q \leq 1$ sur $C$.
    \item Si $C$ est de plus borné et symétrique, montrer que $p_C$ est une norme et
          que $C = \{x : p_C(x) < 1\}$.
    \item Calculer $p_C$ lorsque $C = \{(x,y) \in \R^2 : \abs{x} + \abs{y} < 1\}$.
\end{enumerate}
\end{exercice}

\begin{exercice}[Topologie faible-$*$ et métrisabilité]\label{exo:faible-etoile-metrisable}
\index{topologie faible-*!métrisabilité}
Soit $E$ un espace de Banach séparable de dual $E'$.
\begin{enumerate}[label=(\alph*)]
    \item Soit $(x_n)_{n \geq 1}$ une suite dense dans la boule unité de $E$. Montrer que
          \[
              d(f,g) = \sum_{n=1}^{\infty} \frac{1}{2^n} \abs{f(x_n) - g(x_n)}
          \]
          définit une distance sur $\overline{B}_{E'}$ dont la topologie coïncide avec
          la topologie faible-$*$ restreinte à $\overline{B}_{E'}$.
    \item En déduire que $\overline{B}_{E'}$ est faible-$*$ séquentiellement compacte.
    \item Ceci est-il vrai sans l'hypothèse de séparabilité de $E$ ?
\end{enumerate}
\end{exercice}

\begin{exercice}[Dual de $L^p$]\label{exo:dual-Lp}
\index{dual!de $L^p$}
Soit $(X, \mathcal{A}, \mu)$ un espace mesuré $\sigma$-fini et $1 \leq p < \infty$.
\begin{enumerate}[label=(\alph*)]
    \item Montrer que l'application $\Phi : L^q \to (L^p)'$ définie par
          $\Phi(g)(f) = \int fg\, \dd\mu$ est une isométrie.
    \item Montrer la surjectivité de $\Phi$ dans le cas $p = 2$ en utilisant le
          théorème de Riesz.
    \item Admettre le cas général (théorème de Radon-Nikodym\index{Radon-Nikodym, théorème de})
          et en déduire la réflexivité de $L^p$ pour $1 < p < \infty$.
\end{enumerate}
\end{exercice}

\begin{exercice}[Applications du théorème de séparation]\label{exo:applications-separation}
\index{séparation de convexes!applications}
\begin{enumerate}[label=(\alph*)]
    \item Soit $E$ un e.v.n.\ réel et $C \subset E$ un convexe fermé non vide.
          Montrer que $C = \bigcap\{H^+ : H^+ \text{ demi-espace fermé contenant } C\}$.
    \item Soit $C \subset E$ convexe fermé et $x_0 \notin C$. Montrer qu'il existe
          $f \in E'$ telle que $f(x_0) > \sup_{x \in C} f(x)$.
    \item (\textbf{Théorème de Mazur}\index{Mazur, théorème de}) Montrer que la fermeture
          faible d'un convexe coïncide avec sa fermeture forte.
\end{enumerate}
\end{exercice}

\begin{exercice}[Bidual et réflexivité]\label{exo:bidual-reflexivite}
\index{bidual!exercice}
\index{espace réflexif!exercice}
\begin{enumerate}[label=(\alph*)]
    \item Soit $E$ un espace de Banach réflexif et $F$ un sous-espace fermé.
          Montrer que $F$ est réflexif.
          \emph{Indication : utiliser l'exercice~\ref{exo:annihilateur}(d).}
    \item Montrer qu'un espace de Banach $E$ est réflexif si et seulement si $E'$ est
          réflexif.
    \item Montrer que si $E$ est réflexif, alors toute forme linéaire continue sur $E$
          atteint sa norme, c'est-à-dire : pour tout $f \in E'$, il existe $x_0 \in E$
          avec $\norm{x_0} = 1$ et $f(x_0) = \norm{f}$.
          (\textbf{Théorème de James}\index{James, théorème de}, direction facile.)
\end{enumerate}
\end{exercice}

\index{Hahn-Banach, théorème de|)}
%%%%%%%%%%%%%%%%%%%%%%%%%%%%%%%%%%%%%%%%%%%%%%%%%%%%%%%%%%%%%%%%%%%%
%  Analyse Fonctionnelle — Chapitre 5
%  Théorème de Banach-Steinhaus et Applications
%%%%%%%%%%%%%%%%%%%%%%%%%%%%%%%%%%%%%%%%%%%%%%%%%%%%%%%%%%%%%%%%%%%%

\chapter{Théorème de Banach-Steinhaus et Applications}
\label{chap:banach-steinhaus}
\minitoc

\vspace{1em}

Ce chapitre est consacré au \emph{principe de borne uniforme}\index{principe de borne uniforme},
l'un des quatre grands théorèmes de l'analyse fonctionnelle. Ce résultat, dû à
Stefan Banach\index{Banach, Stefan} et Hugo Steinhaus\index{Steinhaus, Hugo} (1927),
établit qu'une famille d'opérateurs linéaires continus qui est bornée ponctuellement
est nécessairement bornée uniformément en norme. La preuve repose de manière essentielle
sur le \emph{théorème de Baire}\index{théorème!de Baire}, que nous démontrons en premier lieu.

%% ══════════════════════════════════════════════════════════════════
\section{Le théorème de Baire}
\label{sec:baire}
%% ══════════════════════════════════════════════════════════════════

\index{théorème!de Baire|(}
\index{espace!de Baire}

Le théorème de Baire est un résultat fondamental de topologie générale qui fournit
des informations sur la «~taille~» des intersections dénombrables d'ouverts denses.
Il admet deux formulations équivalentes.

\begin{definition}[Ensemble maigre et ensemble résiduel]
\label{def:maigre-residuel}
Soit $(X,d)$ un espace métrique\index{espace!métrique}.
\begin{enumerate}[label=(\roman*)]
  \item Un sous-ensemble $A \subset X$ est dit \emph{nulle part dense}\index{ensemble!nulle part dense}
    (ou \emph{rare}\index{ensemble!rare}) si $\overline{A}$ est d'intérieur vide :
    $\mathring{\overline{A}} = \varnothing$.
  \item Un sous-ensemble $A \subset X$ est dit \emph{maigre}\index{ensemble!maigre}
    (ou \emph{de première catégorie}\index{première catégorie}) s'il est réunion
    dénombrable d'ensembles nulle part denses.
  \item Un sous-ensemble $A \subset X$ est dit \emph{résiduel}\index{ensemble!résiduel}
    (ou \emph{de deuxième catégorie en lui-même}\index{deuxième catégorie}) si son
    complémentaire $X \setminus A$ est maigre.
\end{enumerate}
\end{definition}

\begin{exemple}
\label{ex:Q-maigre}
L'ensemble $\Q$\index{rationnels} est maigre dans $\R$ : en effet, $\Q = \bigcup_{q \in \Q} \{q\}$,
et chaque singleton $\{q\}$ est fermé d'intérieur vide, donc nulle part dense.
\end{exemple}

\begin{exemple}
\label{ex:cantor-maigre}
L'ensemble de Cantor\index{ensemble!de Cantor} $\mathcal{C} \subset [0,1]$ est fermé
d'intérieur vide (car il ne contient aucun intervalle ouvert), donc nulle part dense,
et a fortiori maigre. Pourtant, $\mathcal{C}$ est non dénombrable.
\end{exemple}

\begin{theoreme}[Théorème de Baire]
\label{thm:baire}
\index{théorème!de Baire}
Soit $(X,d)$ un espace métrique complet\index{espace!métrique complet}. Alors :
\begin{enumerate}[label=(\roman*)]
  \item Toute intersection dénombrable d'ouverts denses\index{ouvert dense} est dense dans $X$.
  \item $X$ n'est pas maigre en lui-même : $X$ ne peut pas s'écrire comme réunion
    dénombrable d'ensembles nulle part denses\index{ensemble!nulle part dense}.
\end{enumerate}
\end{theoreme}

\begin{proof}
Montrons d'abord (i). Soient $(U_n)_{n \geq 1}$ des ouverts denses de $X$.
Nous devons montrer que $G = \bigcap_{n=1}^{\infty} U_n$ est dense dans $X$,
c'est-à-dire que pour tout ouvert non vide $V \subset X$, on a $G \cap V \neq \varnothing$.

Soit $V \subset X$ un ouvert non vide. Comme $U_1$ est dense et ouvert,
$V \cap U_1$ est un ouvert non vide. Il existe donc une boule ouverte
$B(x_1, r_1) \subset V \cap U_1$ avec $0 < r_1 < 1$.

Par récurrence, supposons construite une boule $B(x_n, r_n)$ avec $0 < r_n < 1/n$
et $\overline{B(x_n, r_n)} \subset U_n \cap B(x_{n-1}, r_{n-1})$ (avec la convention
$B(x_0, r_0) = V$). Comme $U_{n+1}$ est dense et ouvert, $B(x_n, r_n) \cap U_{n+1}$
est un ouvert non vide. On peut donc choisir $x_{n+1} \in X$ et $0 < r_{n+1} < \frac{1}{n+1}$
tels que
\[
  \overline{B(x_{n+1}, r_{n+1})} \subset B(x_n, r_n) \cap U_{n+1}.
\]

La suite $(x_n)_{n \geq 1}$ est de Cauchy\index{suite!de Cauchy} : pour $m > n$, on a
$x_m \in B(x_n, r_n)$, donc $d(x_n, x_m) < r_n < 1/n \to 0$. Comme $X$ est complet,
$(x_n)$ converge vers un point $x^* \in X$.

Pour tout $n \geq 1$, et pour tout $m > n$, on a $x_m \in \overline{B(x_n, r_n)}$.
Par passage à la limite, $x^* \in \overline{B(x_n, r_n)} \subset U_n$. De plus,
$x^* \in \overline{B(x_1, r_1)} \subset V$. Donc $x^* \in V \cap \bigcap_{n=1}^{\infty} U_n = V \cap G$.

Montrons l'équivalence avec (ii). Si $X = \bigcup_{n=1}^{\infty} F_n$ avec chaque $F_n$
nulle part dense, alors $X \setminus \overline{F_n}$ est un ouvert dense pour tout $n$, et
\[
  \bigcap_{n=1}^{\infty} (X \setminus \overline{F_n})
  = X \setminus \bigcup_{n=1}^{\infty} \overline{F_n}
  \subset X \setminus \bigcup_{n=1}^{\infty} F_n = \varnothing,
\]
ce qui contredit (i) puisqu'une intersection dénombrable d'ouverts denses
devrait être dense, donc non vide (car $X \neq \varnothing$).
\end{proof}

\begin{remarque}
\label{rem:baire-localement-compact}
Le théorème de Baire reste valable pour les espaces localement compacts
de Hausdorff\index{espace!localement compact}\index{Hausdorff}, même sans
métrisabilité. La preuve est analogue, en utilisant la compacité locale
pour extraire des boules fermées emboîtées.
\end{remarque}

\begin{corollaire}
\label{cor:baire-ouvert}
Dans un espace métrique complet, tout ouvert non vide est de deuxième
catégorie\index{deuxième catégorie} en lui-même.
\end{corollaire}

\begin{proof}
Un ouvert d'un espace métrique complet est un espace de Baire
(car il est homéomorphe à un espace métrique complet par le
théorème d'Alexandrov\index{théorème!d'Alexandrov}).
\end{proof}

\begin{exercice}[$\star$]
\label{exo:Gdelta}
\index{$G_\delta$}
Soit $(X,d)$ un espace métrique complet et $A \subset X$ un sous-ensemble
$G_\delta$\index{ensemble!$G_\delta$} dense (i.e.\ intersection dénombrable
d'ouverts denses). Montrer que $A$ est de deuxième catégorie en lui-même.
\end{exercice}

\begin{exercice}[$\star$]
\label{exo:irrationnels-residuels}
Montrer que l'ensemble des irrationnels $\R \setminus \Q$ est résiduel
dans $\R$. En déduire que $\R \setminus \Q$ est non dénombrable.
\end{exercice}

\begin{exercice}[$\star\star$]
\label{exo:nowhere-diff}
\index{fonction!nulle part dérivable}
En utilisant le théorème de Baire, montrer que l'ensemble des fonctions
continues sur $[0,1]$ qui sont dérivables en au moins un point est
maigre dans $\mathcal{C}([0,1])$. En déduire l'existence de fonctions
continues nulle part dérivables.
\end{exercice}


%% ══════════════════════════════════════════════════════════════════
\section{Le théorème de Banach-Steinhaus}
\label{sec:banach-steinhaus}
%% ══════════════════════════════════════════════════════════════════

\index{théorème!de Banach-Steinhaus|(}

Nous arrivons maintenant au résultat principal de ce chapitre.

\begin{theoreme}[Banach-Steinhaus / Principe de borne uniforme]
\label{thm:banach-steinhaus}
\index{principe de borne uniforme}
\index{Banach-Steinhaus}
Soient $X$ un espace de Banach\index{espace!de Banach}, $Y$ un espace
vectoriel normé\index{espace!vectoriel normé}, et $(T_i)_{i \in I}$ une
famille d'opérateurs linéaires continus $T_i : X \to Y$. Si cette famille
est \emph{bornée ponctuellement}\index{borné!ponctuellement}, c'est-à-dire si
\[
  \forall x \in X, \quad \sup_{i \in I} \norm{T_i x}_Y < +\infty,
\]
alors elle est \emph{bornée uniformément}\index{borné!uniformément} :
\[
  \sup_{i \in I} \norm{T_i}_{\mathcal{L}(X,Y)} < +\infty.
\]
\end{theoreme}

\begin{proof}
Pour chaque $n \in \N$, posons
\[
  F_n = \bigl\{ x \in X : \sup_{i \in I} \norm{T_i x}_Y \leq n \bigr\}
  = \bigcap_{i \in I} \bigl\{ x \in X : \norm{T_i x}_Y \leq n \bigr\}.
\]
Chaque $F_n$ est une intersection de fermés (car $x \mapsto \norm{T_i x}$ est
continue pour chaque $i$), donc $F_n$ est fermé\index{fermé}.

Par hypothèse de bornitude ponctuelle, pour tout $x \in X$, il existe $n \in \N$
tel que $\sup_{i \in I} \norm{T_i x} \leq n$, c'est-à-dire $x \in F_n$. Donc
$X = \bigcup_{n=1}^{\infty} F_n$.

Comme $X$ est un espace de Banach (donc complet), le théorème de Baire
(Théorème~\ref{thm:baire}) implique qu'il existe $n_0 \in \N$ tel que $F_{n_0}$
est d'intérieur non vide : il existe $x_0 \in X$ et $r > 0$ tels que
$B(x_0, r) \subset F_{n_0}$.

Soit $x \in X$ avec $\norm{x} \leq 1$. Alors $x_0 + rx \in B(x_0, r) \subset F_{n_0}$
et $x_0 \in F_{n_0}$, donc pour tout $i \in I$ :
\[
  \norm{T_i(rx)}_Y = \norm{T_i(x_0 + rx) - T_i(x_0)}_Y
  \leq \norm{T_i(x_0 + rx)}_Y + \norm{T_i(x_0)}_Y
  \leq n_0 + n_0 = 2n_0.
\]
Par conséquent, pour tout $x \in X$ avec $\norm{x} \leq 1$ et tout $i \in I$ :
\[
  \norm{T_i x}_Y \leq \frac{2n_0}{r}.
\]
On en déduit $\sup_{i \in I} \norm{T_i}_{\mathcal{L}(X,Y)} \leq \frac{2n_0}{r} < +\infty$.
\end{proof}

\begin{remarque}
\label{rem:BS-hypotheses}
\index{complétude!nécessité dans Banach-Steinhaus}
L'hypothèse de complétude de $X$ est \textbf{essentielle}. Considérons
$X = c_{00}$\index{$c_{00}$} (suites à support fini) muni de la norme $\norm{\cdot}_\infty$,
et les fonctionnelles linéaires $f_n : X \to \K$ définies par
$f_n(x) = n\, x_n$. Pour tout $x \in c_{00}$ fixé, $x_n = 0$ pour $n$ assez grand,
donc $\sup_n \abs{f_n(x)} < +\infty$. Pourtant $\norm{f_n} = n \to +\infty$.
\end{remarque}

\begin{remarque}
\label{rem:BS-Y-norme}
L'espace $Y$ n'a pas besoin d'être complet. Seule la complétude du
domaine $X$ est requise.
\end{remarque}

\begin{corollaire}
\label{cor:BS-convergence}
\index{convergence!forte d'opérateurs}
Soient $X$ un espace de Banach, $Y$ un espace vectoriel normé, et
$(T_n)_{n \geq 1} \subset \mathcal{L}(X,Y)$ une suite d'opérateurs.
Si pour tout $x \in X$, la suite $(T_n x)_{n \geq 1}$ converge dans $Y$,
alors :
\begin{enumerate}[label=(\roman*)]
  \item $\sup_{n \geq 1} \norm{T_n} < +\infty$ ;
  \item L'application $T : X \to Y$ définie par $Tx = \lim_{n \to \infty} T_n x$
    est linéaire continue ;
  \item $\norm{T} \leq \liminf_{n \to \infty} \norm{T_n}$.
\end{enumerate}
\end{corollaire}

\begin{proof}
(i) est une conséquence immédiate du théorème de Banach-Steinhaus, car une suite
convergente est bornée.

(ii) La linéarité de $T$ est claire par passage à la limite. Pour la continuité,
soit $M = \sup_n \norm{T_n}$. Pour tout $x \in X$, $\norm{T_n x} \leq M \norm{x}$
pour tout $n$, donc par passage à la limite :
\[
  \norm{Tx} = \lim_{n \to \infty} \norm{T_n x} \leq M \norm{x}.
\]
Ainsi $T \in \mathcal{L}(X,Y)$ et $\norm{T} \leq M$.

(iii) Pour tout $x \in X$ avec $\norm{x} \leq 1$,
$\norm{Tx} = \lim_n \norm{T_n x} \leq \liminf_n \norm{T_n}$, d'où le résultat.
\end{proof}


%% ══════════════════════════════════════════════════════════════════
\section{Applications aux suites de fonctionnelles}
\label{sec:suites-fonctionnelles}
%% ══════════════════════════════════════════════════════════════════

\index{fonctionnelle!linéaire|(}

\begin{proposition}[Convergence faible-$*$ dans le dual]
\label{prop:weak-star-bounded}
\index{convergence!faible-$*$}
\index{topologie!faible-$*$}
Soient $X$ un espace de Banach et $(f_n)_{n \geq 1} \subset X^*$ une suite
de fonctionnelles linéaires continues. Si $(f_n)$ converge dans la topologie
faible-$*$\index{faible-$*$}, c'est-à-dire si pour tout $x \in X$,
$\lim_{n \to \infty} f_n(x)$ existe dans $\K$, alors :
\begin{enumerate}[label=(\roman*)]
  \item $\sup_n \norm{f_n}_{X^*} < +\infty$ ;
  \item La limite $f(x) = \lim_n f_n(x)$ définit un élément $f \in X^*$ ;
  \item $\norm{f}_{X^*} \leq \liminf_n \norm{f_n}_{X^*}$.
\end{enumerate}
\end{proposition}

\begin{proof}
C'est un cas particulier du Corollaire~\ref{cor:BS-convergence} avec $Y = \K$.
\end{proof}

\begin{exemple}[Suite faible-$*$ convergente dans $\ell^\infty$]
\label{ex:weak-star-linfty}
\index{$\ell^\infty$}
Considérons $X = \ell^1$ et $X^* = \ell^\infty$.
Soit $e_n = (0,\ldots,0,1,0,\ldots)$ le $n$-ième vecteur canonique de $\ell^\infty$.
Pour tout $x = (x_k) \in \ell^1$, on a $e_n(x) = x_n \to 0$ car $\sum \abs{x_k} < \infty$.
Donc $e_n \weakstarto 0$. Pourtant $\norm{e_n}_{\ell^\infty} = 1$ pour tout $n$:
la convergence faible-$*$ n'entraîne pas la convergence en norme.
\end{exemple}

\begin{proposition}[Principe de fonctionnelles bornées ponctuellement]
\label{prop:fonctionnelles-bornees}
\index{borné!ponctuellement}
Soit $X$ un espace de Banach et $A \subset X^*$ un sous-ensemble du dual.
Si $A$ est ponctuellement borné\index{borné!ponctuellement},
c'est-à-dire si $\sup_{f \in A} \abs{f(x)} < +\infty$ pour tout $x \in X$,
alors $A$ est borné en norme : $\sup_{f \in A} \norm{f}_{X^*} < +\infty$.
\end{proposition}

\begin{proof}
Application directe du Théorème~\ref{thm:banach-steinhaus} avec
$Y = \K$ et la famille $(T_f)_{f \in A}$ où $T_f = f$.
\end{proof}


%% ══════════════════════════════════════════════════════════════════
\section{Convergence forte versus convergence uniforme d'opérateurs}
\label{sec:convergence-operateurs}
%% ══════════════════════════════════════════════════════════════════

\index{convergence!forte d'opérateurs|(}
\index{convergence!uniforme d'opérateurs|(}

\begin{definition}[Modes de convergence d'opérateurs]
\label{def:modes-convergence-op}
Soient $X$ et $Y$ des espaces vectoriels normés et $(T_n)_{n \geq 1} \subset \mathcal{L}(X,Y)$.
\begin{enumerate}[label=(\roman*)]
  \item On dit que $(T_n)$ \emph{converge en norme}\index{convergence!en norme d'opérateurs}
    (ou \emph{uniformément}) vers $T \in \mathcal{L}(X,Y)$ si $\norm{T_n - T} \to 0$.
  \item On dit que $(T_n)$ \emph{converge fortement}\index{convergence!forte d'opérateurs}
    (ou \emph{ponctuellement}) vers $T$ si $\norm{T_n x - Tx} \to 0$ pour tout $x \in X$.
\end{enumerate}
\end{definition}

\begin{proposition}
\label{prop:uniform-implies-strong}
La convergence en norme entraîne la convergence forte, mais la réciproque
est fausse en général.
\end{proposition}

\begin{proof}
Si $\norm{T_n - T} \to 0$, alors pour tout $x \in X$ :
$\norm{T_n x - Tx} \leq \norm{T_n - T} \norm{x} \to 0$.
La réciproque est fausse : considérer par exemple les projections
$P_n : \ell^2 \to \ell^2$ définies par $P_n(x_1, x_2, \ldots) = (x_1, \ldots, x_n, 0, 0, \ldots)$.
On a $P_n x \to x$ pour tout $x \in \ell^2$, mais $\norm{P_n - \Id} = 1$ pour tout $n$.
\end{proof}

\begin{theoreme}[Principe de Banach-Steinhaus pour les suites d'opérateurs]
\label{thm:BS-suites-operateurs}
\index{Banach-Steinhaus!pour les suites d'opérateurs}
Soient $X$ un espace de Banach, $Y$ un espace vectoriel normé et
$(T_n)_{n \geq 1} \subset \mathcal{L}(X,Y)$. Les assertions suivantes
sont équivalentes :
\begin{enumerate}[label=(\roman*)]
  \item $(T_n)$ converge fortement ;
  \item $(T_n)$ est bornée en norme \textbf{et} converge sur un sous-ensemble
    dense de $X$.
\end{enumerate}
\end{theoreme}

\begin{proof}
(i) $\Rightarrow$ (ii) : Si $(T_n)$ converge fortement, le
Corollaire~\ref{cor:BS-convergence} donne $\sup_n \norm{T_n} < +\infty$,
et $(T_n x)$ converge a fortiori sur tout sous-ensemble dense.

(ii) $\Rightarrow$ (i) : Soit $D \subset X$ un sous-ensemble dense sur lequel
$(T_n)$ converge, et soit $M = \sup_n \norm{T_n} < +\infty$. Soit $x \in X$ et
$\varepsilon > 0$. Il existe $y \in D$ tel que $\norm{x - y} < \varepsilon$.
Comme $(T_n y)$ converge, il existe $N$ tel que pour $m, n \geq N$,
$\norm{T_n y - T_m y} < \varepsilon$. Alors pour $m, n \geq N$ :
\begin{align*}
  \norm{T_n x - T_m x}
  &\leq \norm{T_n(x-y)} + \norm{T_n y - T_m y} + \norm{T_m(y-x)} \\
  &\leq M\norm{x-y} + \varepsilon + M\norm{x-y} \\
  &\leq (2M+1)\varepsilon.
\end{align*}
Donc $(T_n x)$ est de Cauchy dans $Y$. Si $Y$ est complet, $(T_n x)$ converge.
Si $Y$ n'est pas complet, on utilise le fait que la suite est de Cauchy et que
la complétude de $Y$ n'est pas nécessaire si on sait a priori que la limite existe
(par exemple sur $D$, et on prolonge par continuité uniforme).

En fait, une formulation plus précise : on suppose que $(T_n y)$ converge pour
tout $y \in D$, et si $Y$ est un espace de Banach, alors $(T_n x)$ converge
pour tout $x \in X$. Si $Y$ n'est pas complet, la conclusion est que $(T_n x)$
est de Cauchy pour tout $x \in X$.
\end{proof}


%% ══════════════════════════════════════════════════════════════════
\section{Application aux séries de Fourier : théorème de du Bois-Reymond}
\label{sec:du-bois-reymond}
%% ══════════════════════════════════════════════════════════════════

\index{séries de Fourier|(}
\index{du Bois-Reymond, théorème de|(}
\index{noyau!de Dirichlet|(}

L'une des applications les plus frappantes du théorème de Banach-Steinhaus
est la démonstration de l'existence de fonctions continues dont la série de
Fourier diverge en un point.

\subsection{Noyau de Dirichlet et sommes partielles de Fourier}

\begin{definition}[Coefficients et sommes partielles de Fourier]
\label{def:fourier}
\index{coefficients de Fourier}
\index{sommes partielles de Fourier}
Soit $f \in L^1(\mathbb{T})$, où $\mathbb{T} = \R / 2\pi\Z$
désigne le tore\index{tore}. Les \emph{coefficients de Fourier} de $f$ sont
\[
  \hat{f}(n) = \frac{1}{2\pi} \int_{-\pi}^{\pi} f(t)\, e^{-int}\, \dd t,
  \quad n \in \Z.
\]
La \emph{$N$-ième somme partielle de Fourier}\index{somme partielle de Fourier} est
\[
  S_N f(x) = \sum_{n=-N}^{N} \hat{f}(n)\, e^{inx}.
\]
\end{definition}

\begin{definition}[Noyau de Dirichlet]
\label{def:noyau-dirichlet}
\index{noyau!de Dirichlet}
Le \emph{noyau de Dirichlet} d'ordre $N$ est la fonction
\[
  D_N(t) = \sum_{n=-N}^{N} e^{int} = \frac{\sin\bigl((N + \tfrac{1}{2})t\bigr)}{\sin(t/2)},
  \quad t \neq 0.
\]
On a alors la formule de convolution :
\[
  S_N f(x) = \frac{1}{2\pi} \int_{-\pi}^{\pi} f(x-t)\, D_N(t)\, \dd t
  = (f * D_N)(x).
\]
\end{definition}

\begin{lemme}[Norme $L^1$ du noyau de Dirichlet]
\label{lem:norme-dirichlet}
\index{noyau!de Dirichlet!norme $L^1$}
On a
\[
  \frac{1}{2\pi} \int_{-\pi}^{\pi} \abs{D_N(t)}\, \dd t
  = \frac{4}{\pi^2} \ln N + O(1) \quad \text{quand } N \to \infty.
\]
En particulier, $\norm{D_N}_{L^1(\mathbb{T})} \to +\infty$.
\end{lemme}

\begin{proof}
On utilise $\abs{\sin(t/2)} \geq \abs{t}/\pi$ pour $t \in [-\pi, \pi]$ et on majore/minore
par des sommes harmoniques\index{somme harmonique} :
\begin{align*}
  \frac{1}{2\pi}\int_{-\pi}^{\pi} \abs{D_N(t)}\, \dd t
  &= \frac{1}{\pi}\int_0^{\pi} \frac{\abs{\sin((N+\tfrac{1}{2})t)}}{\sin(t/2)}\, \dd t \\
  &\geq \frac{1}{\pi}\int_0^{\pi} \frac{\abs{\sin((N+\tfrac{1}{2})t)}}{t/2}\, \dd t \\
  &= \frac{2}{\pi}\int_0^{(N+1/2)\pi} \frac{\abs{\sin u}}{u}\, \dd u \\
  &\geq \frac{2}{\pi} \sum_{k=1}^{N} \int_{(k-1)\pi}^{k\pi} \frac{\abs{\sin u}}{u}\, \dd u \\
  &\geq \frac{2}{\pi} \sum_{k=1}^{N} \frac{1}{k\pi} \int_{(k-1)\pi}^{k\pi} \abs{\sin u}\, \dd u \\
  &= \frac{2}{\pi} \sum_{k=1}^{N} \frac{2}{k\pi}
  = \frac{4}{\pi^2} \sum_{k=1}^{N} \frac{1}{k}
  = \frac{4}{\pi^2}\ln N + O(1).
\end{align*}
La majoration supérieure est analogue.
\end{proof}

\begin{theoreme}[du Bois-Reymond, 1876]
\label{thm:du-bois-reymond}
\index{du Bois-Reymond, théorème de}
\index{divergence!des séries de Fourier}
Il existe une fonction $f \in \mathcal{C}(\mathbb{T})$\index{$\mathcal{C}(\mathbb{T})$}
dont la série de Fourier diverge en un point. Plus précisément,
l'ensemble des fonctions $f \in \mathcal{C}(\mathbb{T})$ telles que
$(S_N f(0))_{N \geq 0}$ est borné est maigre dans $\mathcal{C}(\mathbb{T})$.
\end{theoreme}

\begin{proof}
Considérons les fonctionnelles linéaires $\Lambda_N : \mathcal{C}(\mathbb{T}) \to \K$
définies par
\[
  \Lambda_N(f) = S_N f(0) = \frac{1}{2\pi} \int_{-\pi}^{\pi} f(-t)\, D_N(t)\, \dd t.
\]
Chaque $\Lambda_N$ est linéaire continue sur $\mathcal{C}(\mathbb{T})$ muni de la norme
$\norm{\cdot}_\infty$, et
\[
  \norm{\Lambda_N} = \frac{1}{2\pi} \int_{-\pi}^{\pi} \abs{D_N(t)}\, \dd t.
\]
En effet, l'inégalité $\leq$ est immédiate. Pour l'égalité, on approche $\sgn(D_N)$
par des fonctions continues.

Par le Lemme~\ref{lem:norme-dirichlet}, $\norm{\Lambda_N} \to +\infty$.

Or $\mathcal{C}(\mathbb{T})$ est un espace de Banach\index{espace!de Banach!$\mathcal{C}(\mathbb{T})$}.
Si pour toute $f \in \mathcal{C}(\mathbb{T})$, la suite $(S_N f(0))_N$ était bornée,
le théorème de Banach-Steinhaus donnerait $\sup_N \norm{\Lambda_N} < +\infty$,
ce qui est une contradiction.

Donc il existe $f \in \mathcal{C}(\mathbb{T})$ telle que $\sup_N \abs{S_N f(0)} = +\infty$.

Pour l'affirmation plus forte, notons $A = \{f \in \mathcal{C}(\mathbb{T}) :
\sup_N \abs{S_N f(0)} < +\infty\}$. Alors $A = \bigcup_{k=1}^{\infty} A_k$ avec
$A_k = \{f : \abs{S_N f(0)} \leq k \text{ pour tout } N\}$. Chaque $A_k$ est fermé
(par continuité des $\Lambda_N$) et d'intérieur vide (car sinon, en raisonnant comme
dans la preuve du théorème de Banach-Steinhaus, on obtiendrait $\sup_N \norm{\Lambda_N} < \infty$).
Donc $A$ est maigre.
\end{proof}

\begin{remarque}
\label{rem:carleson}
\index{Carleson, théorème de}
Le théorème de Carleson (1966) montre que pour $f \in L^2(\mathbb{T})$,
la série de Fourier converge presque partout. Le théorème de du Bois-Reymond
montre que la convergence \emph{ponctuelle en tout point} ne peut pas être
garantie, même pour les fonctions continues.
\end{remarque}


%% ══════════════════════════════════════════════════════════════════
\section{Théorème de condensation de Banach-Steinhaus}
\label{sec:condensation}
%% ══════════════════════════════════════════════════════════════════

\index{condensation de Banach-Steinhaus|(}

Le théorème de condensation est un renforcement du principe de borne uniforme
qui décrit la structure de l'ensemble des points où la bornitude ponctuelle échoue.

\begin{theoreme}[Condensation de Banach-Steinhaus]
\label{thm:condensation}
\index{théorème!de condensation}
Soient $X$ un espace de Banach, $Y$ un espace vectoriel normé et
$(T_n)_{n \geq 1} \subset \mathcal{L}(X,Y)$ une suite d'opérateurs
telle que $\sup_n \norm{T_n} = +\infty$. Alors l'ensemble
\[
  \mathcal{R} = \bigl\{ x \in X : \sup_{n \geq 1} \norm{T_n x}_Y = +\infty \bigr\}
\]
est un $G_\delta$ dense\index{ensemble!$G_\delta$ dense} dans $X$. En particulier,
$\mathcal{R}$ est résiduel\index{ensemble!résiduel}.
\end{theoreme}

\begin{proof}
Pour chaque $k \in \N$, posons
\[
  U_k = \bigl\{ x \in X : \exists\, n \geq 1,\ \norm{T_n x}_Y > k \bigr\}
  = \bigcup_{n=1}^{\infty} \bigl\{ x \in X : \norm{T_n x}_Y > k \bigr\}.
\]
Chaque ensemble $\{x : \norm{T_n x} > k\}$ est ouvert (par continuité de $T_n$),
donc $U_k$ est ouvert. De plus,
\[
  \mathcal{R} = \bigcap_{k=1}^{\infty} U_k,
\]
donc $\mathcal{R}$ est un $G_\delta$.

Il reste à montrer que chaque $U_k$ est dense. Supposons par l'absurde que
$U_k$ n'est pas dense pour un certain $k$. Alors il existe une boule ouverte
$B(x_0, r)$ telle que $B(x_0, r) \cap U_k = \varnothing$, ce qui signifie
que pour tout $x \in B(x_0, r)$ et tout $n$, $\norm{T_n x} \leq k$.

Par le même argument que dans la preuve du Théorème~\ref{thm:banach-steinhaus},
pour tout $x$ avec $\norm{x} \leq 1$ :
\[
  \norm{T_n x} \leq \frac{2k}{r} \quad \text{pour tout } n.
\]
Donc $\sup_n \norm{T_n} \leq 2k/r < +\infty$, ce qui contredit l'hypothèse
$\sup_n \norm{T_n} = +\infty$. Donc $U_k$ est dense.

Par le théorème de Baire, $\mathcal{R} = \bigcap_{k} U_k$ est un $G_\delta$ dense.
\end{proof}

\begin{corollaire}
\label{cor:condensation-fourier}
\index{séries de Fourier!divergence générique}
L'ensemble des fonctions $f \in \mathcal{C}(\mathbb{T})$ telles que
$(S_N f(0))_N$ est non bornée est un $G_\delta$ dense dans $\mathcal{C}(\mathbb{T})$.
En d'autres termes, la divergence de la série de Fourier en $0$ est
\emph{générique}\index{générique}.
\end{corollaire}

\begin{proof}
Les fonctionnelles $\Lambda_N$ de la preuve du Théorème~\ref{thm:du-bois-reymond}
satisfont $\sup_N \norm{\Lambda_N} = +\infty$. Le résultat découle directement
du Théorème~\ref{thm:condensation}.
\end{proof}


%% ══════════════════════════════════════════════════════════════════
\section{Principe de borne uniforme pour les familles d'opérateurs}
\label{sec:PBU-familles}
%% ══════════════════════════════════════════════════════════════════

\index{principe de borne uniforme!pour les familles|(}

\begin{theoreme}[Principe de borne uniforme, version bilinéaire]
\label{thm:PBU-bilineaire}
\index{forme bilinéaire!borne uniforme}
Soient $X$ et $Y$ des espaces de Banach et $B : X \times Y \to \K$
une forme bilinéaire\index{forme bilinéaire} séparément continue
(i.e.\ $B(\cdot, y)$ est continue pour tout $y$ fixé, et $B(x, \cdot)$
est continue pour tout $x$ fixé). Alors $B$ est continue (conjointement) :
\[
  \exists\, C > 0, \quad \forall (x,y) \in X \times Y, \quad
  \abs{B(x,y)} \leq C \norm{x}_X \norm{y}_Y.
\]
\end{theoreme}

\begin{proof}
Pour chaque $x \in X$, l'application $y \mapsto B(x,y)$ est une fonctionnelle
linéaire continue sur $Y$. Notons-la $T_x \in Y^*$. La famille $(T_x)_{x \in B_X}$
(où $B_X$ est la boule unité fermée de $X$) est ponctuellement bornée :
pour tout $y \in Y$,
\[
  \sup_{\norm{x} \leq 1} \abs{T_x(y)} = \sup_{\norm{x} \leq 1} \abs{B(x,y)}
  \leq \norm{B(\cdot, y)}_{X^*} \norm{x} < +\infty,
\]
car $B(\cdot, y)$ est continue donc bornée sur $B_X$.

Par le théorème de Banach-Steinhaus appliqué à la famille $(T_x)_{\norm{x} \leq 1}$
dans $\mathcal{L}(Y, \K)$ avec $Y$ espace de Banach :
\[
  C := \sup_{\norm{x} \leq 1} \norm{T_x}_{Y^*} < +\infty.
\]
Pour tout $(x,y)$ avec $\norm{x} \leq 1$ et $\norm{y} \leq 1$:
$\abs{B(x,y)} = \abs{T_x(y)} \leq \norm{T_x} \leq C$.
Le cas général suit par homogénéité.
\end{proof}

\begin{exemple}[Application à l'espace $\ell^2$]
\label{ex:bilinear-l2}
\index{$\ell^2$}
Soit $A = (a_{ij})_{i,j \geq 1}$ une matrice infinie telle que pour
tout $x \in \ell^2$ et tout $y \in \ell^2$, la série
$\sum_{i,j} a_{ij} x_i y_j$ converge. Alors il existe $C > 0$ tel que
\[
  \biggl\lvert \sum_{i,j=1}^{\infty} a_{ij} x_i y_j \biggr\rvert
  \leq C \norm{x}_{\ell^2} \norm{y}_{\ell^2}.
\]
\end{exemple}

\begin{proposition}[Principe de borne uniforme pour les opérateurs bornés]
\label{prop:PBU-operateurs}
\index{principe de borne uniforme!pour les opérateurs bornés}
Soient $X$, $Y$, $Z$ des espaces de Banach et $(T_i)_{i \in I} \subset \mathcal{L}(X,Y)$,
$(S_j)_{j \in J} \subset \mathcal{L}(Y,Z)$ des familles d'opérateurs. Si
\[
  \forall (i,j) \in I \times J, \quad \norm{S_j T_i}_{\mathcal{L}(X,Z)} \leq M
\]
pour une constante $M > 0$, alors
\[
  \sup_{j \in J} \norm{S_j} < +\infty \quad \text{ou} \quad
  \sup_{i \in I} \norm{T_i} < +\infty.
\]
\end{proposition}

\begin{proof}
Supposons $\sup_i \norm{T_i} = +\infty$. Pour tout $y$ dans l'image
d'un $T_i$, on a un contrôle. Plus précisément, pour tout $x \in X$ et
tout $j \in J$ : $\norm{S_j T_i x} \leq M \norm{x}$. En prenant le sup
sur $\norm{x} \leq 1$ : $\norm{S_j} \cdot \norm{T_i} \leq M / \norm{x}$...

En fait, on raisonne directement : pour tout $x \in X$ fixé,
$\sup_j \norm{S_j(T_i x)} \leq M\norm{x}$ pour tout $i$. Pour tout
$y \in \overline{\bigcup_i T_i(B_X)}$, on a $\sup_j \norm{S_j y} < +\infty$.
Si $\overline{\spn(\bigcup_i \im T_i)} = Y$, alors Banach-Steinhaus donne
$\sup_j \norm{S_j} < +\infty$.
\end{proof}


%% ══════════════════════════════════════════════════════════════════
\section{Le phénomène de Gibbs et le noyau de Dirichlet}
\label{sec:gibbs}
%% ══════════════════════════════════════════════════════════════════

\index{phénomène de Gibbs|(}
\index{noyau!de Dirichlet}

\begin{definition}[Noyau de Fejér]
\label{def:noyau-fejer}
\index{noyau!de Fejér}
Le \emph{noyau de Fejér}\index{Fejér, noyau de} d'ordre $N$ est défini par
\[
  K_N(t) = \frac{1}{N+1} \sum_{n=0}^{N} D_n(t)
  = \frac{1}{N+1} \cdot \frac{\sin^2\bigl(\frac{(N+1)t}{2}\bigr)}{\sin^2(t/2)}.
\]
\end{definition}

\begin{proposition}[Propriétés du noyau de Fejér]
\label{prop:fejer-proprietes}
\index{noyau!de Fejér!propriétés}
Le noyau de Fejér satisfait :
\begin{enumerate}[label=(\roman*)]
  \item $K_N(t) \geq 0$ pour tout $t$ ;
  \item $\frac{1}{2\pi}\int_{-\pi}^{\pi} K_N(t)\, \dd t = 1$ ;
  \item Pour tout $\delta > 0$, $\sup_{\delta \leq \abs{t} \leq \pi} K_N(t) \to 0$ quand $N \to \infty$.
\end{enumerate}
Autrement dit, $(K_N)$ est une \emph{approximation de l'identité}\index{approximation de l'identité}.
\end{proposition}

\begin{proof}
(i) est immédiat par la formule explicite (carré au numérateur et au dénominateur).

(ii) Par linéarité et $\frac{1}{2\pi}\int D_n = 1$ :
$\frac{1}{2\pi}\int K_N = \frac{1}{N+1}\sum_{n=0}^{N} 1 = 1$.

(iii) Pour $\abs{t} \geq \delta$, $\sin^2(t/2) \geq \sin^2(\delta/2) > 0$, donc
$K_N(t) \leq \frac{1}{(N+1)\sin^2(\delta/2)} \to 0$.
\end{proof}

\begin{theoreme}[Fejér]
\label{thm:fejer}
\index{théorème!de Fejér}
Si $f \in \mathcal{C}(\mathbb{T})$, alors les moyennes de Cesàro\index{moyennes de Cesàro}
des sommes partielles de Fourier,
\[
  \sigma_N f(x) = \frac{1}{N+1} \sum_{n=0}^{N} S_n f(x) = (f * K_N)(x),
\]
convergent uniformément vers $f$.
\end{theoreme}

\begin{proof}
Soit $\varepsilon > 0$. Comme $f$ est uniformément continue sur $\mathbb{T}$ (compact),
il existe $\delta > 0$ tel que $\abs{t} < \delta \Rightarrow \abs{f(x-t) - f(x)} < \varepsilon$
pour tout $x$. Alors
\begin{align*}
  \abs{\sigma_N f(x) - f(x)}
  &= \biggl\lvert \frac{1}{2\pi}\int_{-\pi}^{\pi} (f(x-t) - f(x)) K_N(t)\, \dd t \biggr\rvert \\
  &\leq \frac{1}{2\pi}\int_{\abs{t} < \delta} \abs{f(x-t)-f(x)} K_N(t)\, \dd t
    + \frac{1}{2\pi}\int_{\abs{t} \geq \delta} \abs{f(x-t)-f(x)} K_N(t)\, \dd t \\
  &\leq \varepsilon \cdot 1 + 2\norm{f}_\infty \cdot \frac{\pi}{(N+1)\sin^2(\delta/2)}.
\end{align*}
Pour $N$ assez grand, le second terme est $< \varepsilon$, d'où la convergence uniforme.
\end{proof}

Le phénomène de Gibbs décrit le comportement des sommes partielles de Fourier
près d'une discontinuité de saut.

\begin{theoreme}[Phénomène de Gibbs]
\label{thm:gibbs}
\index{phénomène de Gibbs}
Considérons la fonction $f : [-\pi, \pi] \to \R$ définie par $f(t) = \sgn(t)$
pour $t \neq 0$ (et étendue par $2\pi$-périodicité). Alors :
\begin{enumerate}[label=(\roman*)]
  \item $S_N f(0) = 0$ pour tout $N$ ;
  \item $\displaystyle\lim_{N \to \infty} S_N f\!\left(\frac{\pi}{N}\right)
    = \frac{2}{\pi}\int_0^{\pi} \frac{\sin t}{t}\, \dd t \approx 1{,}1790$,
    ce qui représente un dépassement\index{dépassement de Gibbs} d'environ $9\%$
    par rapport à la valeur $1$ de $f$ en $0^+$.
\end{enumerate}
\end{theoreme}

\begin{proof}
(i) Par imparité de $f$ et de $\sin(nt)$, on a $\hat{f}(n) = -\hat{f}(-n)$
pour $n > 0$, et $\hat{f}(0) = 0$. Donc $S_N f(0) = \sum_{n=-N}^{N} \hat{f}(n) = 0$.

(ii) On calcule explicitement :
\[
  S_N f(x) = \frac{2}{\pi} \sum_{k=0}^{N-1} \frac{\sin((2k+1)x)}{2k+1}.
\]
En posant $x = \pi/N$ et en passant à la limite (somme de Riemann) :
\begin{align*}
  \lim_{N \to \infty} S_N f\!\left(\frac{\pi}{N}\right)
  &= \lim_{N \to \infty} \frac{1}{\pi}\int_{-\pi}^{\pi} f\!\left(\frac{\pi}{N}-t\right) D_N(t)\, \dd t \\
  &= \frac{2}{\pi} \int_0^{\pi} \frac{\sin t}{t}\, \dd t,
\end{align*}
par un changement de variables et un passage à la limite dominé.
L'intégrale $\int_0^{\pi} \frac{\sin t}{t}\, \dd t \approx 1{,}8519$ donne
le résultat $\frac{2}{\pi} \cdot 1{,}8519 \approx 1{,}1790$.
\end{proof}


%% ══════════════════════════════════════════════════════════════════
\section{Exercices}
\label{sec:exos-ch5}
%% ══════════════════════════════════════════════════════════════════

\begin{exercice}[$\star$]
\label{exo:baire-fermé}
\index{théorème!de Baire!exercice}
Soit $(X,d)$ un espace métrique complet et $(F_n)_{n \geq 1}$ une suite
de fermés telle que $X = \bigcup_{n=1}^{\infty} F_n$. Montrer qu'au moins
un $F_n$ est d'intérieur non vide.
\end{exercice}

\begin{exercice}[$\star$]
\label{exo:semicontinue-continue}
\index{fonction!semi-continue}
Soit $f : X \to \R$ une fonction sur un espace métrique complet, qui est
à la fois semi-continue inférieurement et semi-continue supérieurement.
Montrer que $f$ est continue (cet exercice est trivial, c'est un rappel).
Montrer ensuite que si $f$ est limite simple d'une suite de fonctions continues,
alors l'ensemble des points de continuité de $f$ est un $G_\delta$ dense.
\end{exercice}

\begin{exercice}[$\star$]
\label{exo:BS-dual}
\index{dual!borne uniforme}
Soit $X$ un espace de Banach et $A \subset X$. Montrer que $A$ est borné si et
seulement si pour tout $f \in X^*$, $\sup_{x \in A} \abs{f(x)} < +\infty$.
\textit{(Indication : utiliser le plongement canonique $J : X \to X^{**}$ et
le principe de borne uniforme.)}
\end{exercice}

\begin{exercice}[$\star\star$]
\label{exo:divergence-fourier-tout-point}
\index{séries de Fourier!divergence}
En utilisant le théorème de condensation (Théorème~\ref{thm:condensation}),
montrer qu'il existe $f \in \mathcal{C}(\mathbb{T})$ dont la série de Fourier
diverge en tout point rationnel (i.e.\ en tout point de la forme $2\pi p/q$,
$p/q \in \Q$).
\end{exercice}

\begin{exercice}[$\star\star$]
\label{exo:BS-lp}
\index{$\ell^p$!borne uniforme}
Soit $1 \leq p < \infty$ et $(a^{(n)})_{n \geq 1}$ une suite dans $\ell^p$
telle que pour tout $x \in \ell^q$ (où $1/p + 1/q = 1$), la série
$\sum_k a^{(n)}_k x_k$ converge et $\sup_n \abs{\sum_k a^{(n)}_k x_k} < \infty$.
Montrer que $\sup_n \norm{a^{(n)}}_{\ell^p} < \infty$.
\end{exercice}

\begin{exercice}[$\star\star$]
\label{exo:matrice-infinie}
\index{matrice!infinie}
Soit $A = (a_{ij})_{i,j \geq 1}$ une matrice infinie telle que pour tout
$x = (x_j) \in \ell^2$, la suite $(y_i)$ définie par $y_i = \sum_{j=1}^{\infty} a_{ij} x_j$
est dans $\ell^2$. Montrer que l'opérateur $T : \ell^2 \to \ell^2$ ainsi
défini est borné. \textit{(Utiliser le théorème du graphe fermé, ou
Banach-Steinhaus avec les projections.)}
\end{exercice}

\begin{exercice}[$\star\star$]
\label{exo:convergence-faible-forte}
\index{convergence!faible vs forte}
Soit $X$ un espace de Banach et $(x_n) \subset X$ telle que $x_n \weakto x$
(convergence faible\index{convergence!faible}). Montrer que $(x_n)$ est bornée
et que $\norm{x} \leq \liminf_n \norm{x_n}$. \textit{(Utiliser $J(x_n) \in X^{**}$
et Banach-Steinhaus.)}
\end{exercice}

\begin{exercice}[$\star\star\star$]
\label{exo:BS-non-separable}
\index{Banach-Steinhaus!espace non séparable}
Construire un exemple d'un espace vectoriel normé \emph{non complet} $X$ et
d'une suite $(f_n) \subset X^*$ telle que $\sup_n \abs{f_n(x)} < \infty$ pour
tout $x \in X$, mais $\sup_n \norm{f_n} = +\infty$. Pourquoi cela ne contredit-il
pas le théorème de Banach-Steinhaus ?
\end{exercice}

\begin{exercice}[$\star\star\star$]
\label{exo:polynomes-bernstein}
\index{polynômes de Bernstein}
Soit $f \in \mathcal{C}([0,1])$. Les \emph{polynômes de Bernstein} de $f$ sont
\[
  B_n f(x) = \sum_{k=0}^{n} f\!\left(\frac{k}{n}\right) \binom{n}{k} x^k (1-x)^{n-k}.
\]
\begin{enumerate}[label=(\alph*)]
  \item Montrer que $B_n f \to f$ uniformément.
  \item Montrer que les opérateurs $B_n : \mathcal{C}([0,1]) \to \mathcal{C}([0,1])$
    satisfont $\norm{B_n} = 1$.
  \item En déduire que la convergence forte $B_n \to \Id$ n'implique pas
    la convergence en norme.
\end{enumerate}
\end{exercice}

\begin{exercice}[$\star\star\star$]
\label{exo:non-convergence-uniforme}
\index{convergence!non uniforme}
Soit $X$ un espace de Banach de dimension infinie et $(T_n) \subset \mathcal{L}(X,X)$
une suite telle que $T_n \to \Id$ fortement (i.e.\ $T_n x \to x$ pour tout $x$).
Montrer que si chaque $T_n$ est de rang fini\index{opérateur!de rang fini},
alors $\norm{T_n - \Id} \geq 1$ pour tout $n$ assez grand
(en utilisant le fait qu'un opérateur de rang fini ne peut pas être surjectif
en dimension infinie, ou une majoration de norme).
\end{exercice}

\index{théorème!de Banach-Steinhaus|)}
\index{théorème!de Baire|)}
\index{séries de Fourier|)}
\index{du Bois-Reymond, théorème de|)}
\index{noyau!de Dirichlet|)}
\index{fonctionnelle!linéaire|)}
\index{convergence!forte d'opérateurs|)}
\index{convergence!uniforme d'opérateurs|)}
\index{condensation de Banach-Steinhaus|)}
\index{phénomène de Gibbs|)}
\index{principe de borne uniforme!pour les familles|)}


%%%%%%%%%%%%%%%%%%%%%%%%%%%%%%%%%%%%%%%%%%%%%%%%%%%%%%%%%%%%%%%%%%%%
%  Analyse Fonctionnelle — Chapitre 6
%  Théorème de l'Application Ouverte et du Graphe Fermé
%%%%%%%%%%%%%%%%%%%%%%%%%%%%%%%%%%%%%%%%%%%%%%%%%%%%%%%%%%%%%%%%%%%%

\chapter{Théorème de l'Application Ouverte et du Graphe Fermé}
\label{chap:app-ouverte-graphe-ferme}
\minitoc

\vspace{1em}

Ce chapitre est consacré à deux autres piliers de l'analyse fonctionnelle :
le \emph{théorème de l'application ouverte}\index{théorème!de l'application ouverte}
(aussi appelé théorème de Banach-Schauder\index{Banach-Schauder, théorème de})
et le \emph{théorème du graphe fermé}\index{théorème!du graphe fermé}.
Ces résultats, comme le théorème de Banach-Steinhaus du chapitre précédent,
reposent de manière fondamentale sur le théorème de Baire\index{théorème!de Baire}.
Ils constituent des outils indispensables pour l'étude des opérateurs
linéaires entre espaces de Banach.

%% ══════════════════════════════════════════════════════════════════
\section{Le théorème de l'application ouverte}
\label{sec:application-ouverte}
%% ══════════════════════════════════════════════════════════════════

\index{théorème!de l'application ouverte|(}
\index{application!ouverte|(}

\begin{definition}[Application ouverte]
\label{def:application-ouverte}
\index{application!ouverte}
Une application $f : X \to Y$ entre espaces topologiques est dite
\emph{ouverte}\index{ouverte, application} si l'image de tout ouvert
de $X$ est un ouvert de $Y$.
\end{definition}

\begin{remarque}
\label{rem:ouverte-vs-continue}
Attention : une application ouverte n'est pas nécessairement continue,
et une application continue n'est pas nécessairement ouverte.
Par exemple, $x \mapsto x^2$ de $\R$ dans $\R$ est continue mais pas ouverte
(l'image de $(-1,1)$ est $[0,1)$, qui n'est pas ouvert).
\end{remarque}

Le lemme suivant est l'étape clé de la démonstration.

\begin{lemme}[Lemme d'ouverture approchée]
\label{lem:ouverture-approchee}
\index{lemme!d'ouverture approchée}
Soient $X$ et $Y$ des espaces de Banach et $T \in \mathcal{L}(X,Y)$
un opérateur surjectif. Alors il existe $\delta > 0$ tel que
\[
  B_Y(0, \delta) \subset \overline{T(B_X(0,1))},
\]
où $B_X(0,1)$ et $B_Y(0,\delta)$ désignent les boules ouvertes.
\end{lemme}

\begin{proof}
Comme $T$ est surjectif, $Y = T(X) = \bigcup_{n=1}^{\infty} T(B_X(0,n))$
(car tout $x \in X$ est dans $B_X(0,n)$ pour $n$ assez grand). Donc
\[
  Y = \bigcup_{n=1}^{\infty} \overline{T(B_X(0,n))}.
\]
Par le théorème de Baire\index{théorème!de Baire} (car $Y$ est complet),
il existe $n_0$ tel que $\overline{T(B_X(0,n_0))}$ est d'intérieur non vide :
il existe $y_0 \in Y$ et $r > 0$ tels que
\[
  B_Y(y_0, r) \subset \overline{T(B_X(0,n_0))}.
\]

Comme $y_0 \in \overline{T(B_X(0,n_0))}$, par symétrie et convexité de
$\overline{T(B_X(0,n_0))}$, on obtient
$B_Y(-y_0, r) \subset \overline{T(B_X(0,n_0))}$
(car $T(B_X(0,n_0))$ est symétrique : $-T(B_X(0,n_0)) = T(B_X(0,n_0))$).

Pour tout $y \in B_Y(0,r)$, on peut écrire $y = \frac{1}{2}(y_0+y) + \frac{1}{2}(-y_0+y)$
avec $y_0+y \in B_Y(y_0,r)$ et $-y_0+y \in B_Y(-y_0,r)$.
Donc $y_0 + y$ et $-y_0 + y$ sont dans $\overline{T(B_X(0,n_0))}$, et par convexité
de l'adhérence d'un ensemble convexe :
\[
  y = \tfrac{1}{2}(y_0+y) + \tfrac{1}{2}(-y_0+y) \in \overline{T(B_X(0,n_0))}.
\]
Donc $B_Y(0,r) \subset \overline{T(B_X(0,n_0))}$.

Par homogénéité ($T$ est linéaire), $\overline{T(B_X(0,n_0))} = n_0 \overline{T(B_X(0,1))}$,
donc
\[
  B_Y\!\left(0, \frac{r}{n_0}\right) \subset \overline{T(B_X(0,1))}.
\]
On pose $\delta = r/n_0 > 0$.
\end{proof}

\begin{theoreme}[Théorème de l'application ouverte / Banach-Schauder]
\label{thm:application-ouverte}
\index{théorème!de l'application ouverte}
\index{Banach-Schauder, théorème de}
Soient $X$ et $Y$ des espaces de Banach et $T \in \mathcal{L}(X,Y)$
surjectif\index{surjectif}. Alors $T$ est une application ouverte.

Plus précisément, il existe $\delta > 0$ tel que
\[
  B_Y(0, \delta) \subset T(B_X(0,1)).
\]
\end{theoreme}

\begin{proof}
Par le Lemme~\ref{lem:ouverture-approchee}, il existe $\delta > 0$ tel que
$B_Y(0, 2\delta) \subset \overline{T(B_X(0,1))}$. Par homogénéité, pour tout
$n \geq 1$ :
\begin{equation}
\label{eq:inclusion-approchee}
  B_Y(0, 2\delta \cdot 2^{-n}) \subset \overline{T(B_X(0, 2^{-n}))}.
\end{equation}

Montrons que $B_Y(0, \delta) \subset T(B_X(0,1))$.

Soit $y \in B_Y(0, \delta)$, donc $\norm{y} < \delta$.
Comme $\norm{y} < \delta < 2\delta$, par \eqref{eq:inclusion-approchee} avec $n=0$,
$y \in \overline{T(B_X(0,1))}$, donc il existe $x_0 \in B_X(0, 1)$ tel que
$\norm{y - Tx_0} < \delta/2$.

Par récurrence : supposons construits $x_0, x_1, \ldots, x_{n-1}$ avec
$\norm{x_k} < 2^{-k}$ et
\[
  \biggl\| y - T\biggl(\sum_{k=0}^{n-1} x_k\biggr) \biggr\| < \delta \cdot 2^{-n}.
\]
Posons $y_n = y - T\bigl(\sum_{k=0}^{n-1} x_k\bigr)$. Comme $\norm{y_n} < \delta \cdot 2^{-n}
< 2\delta \cdot 2^{-n}$, par \eqref{eq:inclusion-approchee}, il existe
$x_n \in B_X(0, 2^{-n})$ tel que $\norm{y_n - Tx_n} < \delta \cdot 2^{-(n+1)}$.

La série $\sum_{k=0}^{\infty} x_k$ converge absolument dans $X$ car
$\sum_{k=0}^{\infty} \norm{x_k} < \sum_{k=0}^{\infty} 2^{-k} = 2$.
Comme $X$ est complet, soit $x = \sum_{k=0}^{\infty} x_k \in X$.
On a $\norm{x} \leq \sum_k \norm{x_k} < 2$, mais en fait
$\norm{x_0} < 1$ et $\sum_{k=1}^{\infty} \norm{x_k} < 1$, donc $\norm{x} < 2$.

Par continuité de $T$ :
\[
  Tx = T\biggl(\sum_{k=0}^{\infty} x_k\biggr)
  = \lim_{n \to \infty} T\biggl(\sum_{k=0}^{n} x_k\biggr) = y,
\]
car $\norm{y - T(\sum_{k=0}^{n} x_k)} < \delta \cdot 2^{-(n+1)} \to 0$.

Nous avons montré que pour tout $y \in B_Y(0, \delta)$, il existe $x \in X$
avec $\norm{x} < 2$ et $Tx = y$. En remplaçant $\delta$ par $\delta/2$, on
obtient : pour tout $y \in B_Y(0, \delta/2)$, il existe $x$ avec $\norm{x} < 1$
et $Tx = y$. Autrement dit, $B_Y(0, \delta/2) \subset T(B_X(0,1))$.

Enfin, montrons que $T$ est ouverte. Soit $U \subset X$ un ouvert et
$y_0 = Tx_0 \in T(U)$. Comme $U$ est ouvert, il existe $r > 0$ tel que
$B_X(x_0, r) \subset U$, i.e.\ $x_0 + B_X(0,r) \subset U$. Alors
\[
  T(U) \supset T(x_0 + B_X(0,r)) = Tx_0 + T(B_X(0,r))
  \supset y_0 + B_Y(0, r\delta/2).
\]
Donc $y_0$ est intérieur à $T(U)$, ce qui prouve que $T(U)$ est ouvert.
\end{proof}

\begin{remarque}
\label{rem:app-ouverte-hypotheses}
\index{complétude!nécessité dans l'application ouverte}
Les hypothèses de complétude de $X$ \textbf{et} de $Y$ sont toutes
deux nécessaires, ainsi que la surjectivité de $T$.
\begin{itemize}
  \item Si $X$ n'est pas complet : la série $\sum x_k$ pourrait ne pas converger.
  \item Si $Y$ n'est pas complet : le théorème de Baire ne s'applique pas.
  \item Si $T$ n'est pas surjectif : l'image de $T$ pourrait être un
    sous-espace propre fermé, qui ne contient aucune boule ouverte de $Y$.
\end{itemize}
\end{remarque}


%% ══════════════════════════════════════════════════════════════════
\section{Corollaire : isomorphisme de Banach}
\label{sec:isomorphisme-banach}
%% ══════════════════════════════════════════════════════════════════

\index{isomorphisme!de Banach|(}

\begin{corollaire}[Isomorphisme de Banach]
\label{cor:isomorphisme-banach}
\index{isomorphisme!de Banach}
\index{théorème!de l'isomorphisme de Banach}
Soient $X$ et $Y$ des espaces de Banach et $T \in \mathcal{L}(X,Y)$
un opérateur bijectif. Alors $T^{-1} \in \mathcal{L}(Y,X)$, c'est-à-dire
que $T^{-1}$ est automatiquement continu.
\end{corollaire}

\begin{proof}
L'application $T^{-1} : Y \to X$ est linéaire (vérification immédiate).
Pour montrer qu'elle est continue, il suffit de montrer que $T$ est ouverte,
car alors pour tout ouvert $U \subset X$, $(T^{-1})^{-1}(U) = T(U)$ est
ouvert dans $Y$.

Or $T$ est surjectif (car bijectif), donc par le Théorème~\ref{thm:application-ouverte},
$T$ est une application ouverte. Donc $T^{-1}$ est continue.
\end{proof}

\begin{remarque}
\label{rem:iso-banach-important}
Ce résultat est remarquable : pour montrer qu'un opérateur linéaire continu
bijectif entre espaces de Banach est un isomorphisme topologique, il suffit
de vérifier la bijectivité. La continuité de l'inverse est \emph{gratuite}.
Cela n'est pas vrai en général pour les espaces vectoriels normés non complets.
\end{remarque}

\begin{corollaire}[Équivalence des normes]
\label{cor:equivalence-normes}
\index{normes!équivalentes}
Soit $X$ un espace vectoriel muni de deux normes $\norm{\cdot}_1$ et $\norm{\cdot}_2$
telles que $(X, \norm{\cdot}_1)$ et $(X, \norm{\cdot}_2)$ soient des espaces de
Banach. Si l'une des normes est plus fine que l'autre, i.e.\ s'il existe $C > 0$
tel que $\norm{x}_2 \leq C \norm{x}_1$ pour tout $x$, alors les deux normes sont
équivalentes : il existe $c > 0$ tel que $c\norm{x}_1 \leq \norm{x}_2$ pour tout $x$.
\end{corollaire}

\begin{proof}
L'application $\Id : (X, \norm{\cdot}_1) \to (X, \norm{\cdot}_2)$ est linéaire,
continue (par hypothèse) et bijective. Par le Corollaire~\ref{cor:isomorphisme-banach},
$\Id^{-1} : (X, \norm{\cdot}_2) \to (X, \norm{\cdot}_1)$ est continue, ce qui donne
l'inégalité inverse.
\end{proof}

\begin{exemple}
\label{ex:normes-Lp}
\index{$L^p$!comparaison des normes}
Sur l'espace $X = L^1([0,1]) \cap L^\infty([0,1])$, muni de la norme
$\norm{f}_1 = \norm{f}_{L^1} + \norm{f}_{L^\infty}$, $(X, \norm{\cdot}_1)$ n'est
pas complet. La norme $\norm{f}_{L^\infty}$ satisfait $\norm{f}_{L^1} \leq \norm{f}_{L^\infty}$,
mais les deux normes ne sont pas équivalentes sur $X$ --- ce qui ne contredit pas
le corollaire précédent, puisque $(X, \norm{\cdot}_1)$ n'est pas un espace de Banach.
\end{exemple}

\begin{proposition}[Supplémentaire topologique]
\label{prop:supplementaire-topo}
\index{supplémentaire!topologique}
Soit $X$ un espace de Banach et $X = M \oplus N$ une décomposition en somme
directe algébrique, où $M$ et $N$ sont des sous-espaces fermés\index{sous-espace!fermé}.
Alors la projection\index{projection} $P : X \to M$ parallèlement à $N$ est continue.
\end{proposition}

\begin{proof}
L'espace $M \times N$ muni de la norme $\norm{(m,n)} = \norm{m}_X + \norm{n}_X$
est un espace de Banach (produit de Banach). L'application
$\Phi : M \times N \to X$, $(m,n) \mapsto m+n$ est linéaire, continue (par
l'inégalité triangulaire) et bijective (par hypothèse de somme directe).
Par le Corollaire~\ref{cor:isomorphisme-banach}, $\Phi^{-1}$ est continue.
Or $P = \pi_1 \circ \Phi^{-1}$ où $\pi_1 : M \times N \to M$ est la première
projection (continue). Donc $P$ est continue.
\end{proof}


%% ══════════════════════════════════════════════════════════════════
\section{Le théorème du graphe fermé}
\label{sec:graphe-ferme}
%% ══════════════════════════════════════════════════════════════════

\index{théorème!du graphe fermé|(}
\index{graphe!fermé|(}

\begin{definition}[Graphe d'un opérateur]
\label{def:graphe}
\index{graphe!d'un opérateur}
Soient $X$ et $Y$ des espaces vectoriels normés et $T : D(T) \subset X \to Y$
un opérateur linéaire (éventuellement non borné\index{opérateur!non borné}),
de domaine $D(T)$\index{domaine}. Le \emph{graphe}\index{graphe} de $T$ est
le sous-ensemble de $X \times Y$ défini par
\[
  \Gamma(T) = \bigl\{ (x, Tx) : x \in D(T) \bigr\} \subset X \times Y.
\]
L'espace $X \times Y$ est muni de la norme produit
$\norm{(x,y)} = \norm{x}_X + \norm{y}_Y$.
\end{definition}

\begin{remarque}
\label{rem:graphe-sous-espace}
Le graphe $\Gamma(T)$ est toujours un sous-espace vectoriel de $X \times Y$
(car $T$ est linéaire). La question est de savoir s'il est fermé.
\end{remarque}

\begin{definition}[Opérateur fermé]
\label{def:operateur-ferme}
\index{opérateur!fermé}
Un opérateur linéaire $T : D(T) \subset X \to Y$ est dit \emph{fermé}
si son graphe $\Gamma(T)$ est fermé dans $X \times Y$. De manière
équivalente : si $(x_n) \subset D(T)$, $x_n \to x$ dans $X$ et $Tx_n \to y$
dans $Y$, alors $x \in D(T)$ et $Tx = y$.
\end{definition}

\begin{exemple}
\label{ex:derivation-fermee}
\index{opérateur!de dérivation}
L'opérateur de dérivation $T : D(T) \subset \mathcal{C}([0,1]) \to \mathcal{C}([0,1])$,
$Tf = f'$, avec $D(T) = \mathcal{C}^1([0,1])$, est un opérateur fermé mais non borné.

En effet, si $f_n \to f$ uniformément et $f_n' \to g$ uniformément, alors
$f$ est dérivable et $f' = g$ (théorème classique d'analyse). Donc $T$ est fermé.

Il n'est pas borné : prendre $f_n(x) = x^n / n$, alors $\norm{f_n}_\infty = 1/n \to 0$
mais $\norm{f_n'}_\infty = \norm{x^{n-1}}_\infty = 1$.
\end{exemple}

\begin{exemple}
\label{ex:operateur-non-ferme}
\index{opérateur!non fermé}
L'opérateur $T : c_{00} \subset \ell^2 \to \ell^2$ défini par
$T(e_n) = n\, e_n$ (étendu linéairement) n'est pas fermé. En effet,
considérons $x_m = \sum_{n=1}^{m} \frac{1}{n^2} e_n$. On a $x_m \to x = \sum_{n=1}^{\infty} \frac{1}{n^2} e_n \in \ell^2$
et $Tx_m = \sum_{n=1}^{m} \frac{1}{n} e_n$. La suite $(Tx_m)$ ne converge pas dans $\ell^2$
car la série harmonique diverge. De plus, $x \notin c_{00}$, donc $x \notin D(T)$.
\end{exemple}

\begin{theoreme}[Théorème du graphe fermé]
\label{thm:graphe-ferme}
\index{théorème!du graphe fermé}
Soient $X$ et $Y$ des espaces de Banach et $T : X \to Y$ un opérateur linéaire
(défini sur tout $X$). Alors $T$ est continu si et seulement si son graphe
$\Gamma(T)$ est fermé dans $X \times Y$.
\end{theoreme}

\begin{proof}
$(\Rightarrow)$ Si $T$ est continu et $(x_n, Tx_n) \to (x,y)$ dans $X \times Y$,
alors $x_n \to x$ et $Tx_n \to y$. Par continuité de $T$, $Tx_n \to Tx$, donc
$y = Tx$ par unicité de la limite. Ainsi $(x,y) = (x,Tx) \in \Gamma(T)$.

$(\Leftarrow)$ Supposons que $\Gamma(T)$ est fermé dans $X \times Y$.
L'espace $X \times Y$ muni de la norme $\norm{(x,y)} = \norm{x}_X + \norm{y}_Y$
est un espace de Banach\index{espace!de Banach!produit}. Comme $\Gamma(T)$ est
un sous-espace vectoriel fermé d'un espace de Banach, $\Gamma(T)$ est lui-même
un espace de Banach\index{sous-espace!fermé}.

Considérons les projections\index{projection} :
\begin{align*}
  \pi_1 &: \Gamma(T) \to X, \quad (x, Tx) \mapsto x, \\
  \pi_2 &: \Gamma(T) \to Y, \quad (x, Tx) \mapsto Tx.
\end{align*}
Les deux projections sont linéaires continues (car $\norm{x} \leq \norm{(x,Tx)}$
et $\norm{Tx} \leq \norm{(x,Tx)}$). De plus, $\pi_1$ est bijective
(car $T$ est défini sur tout $X$).

Par le Corollaire~\ref{cor:isomorphisme-banach} (isomorphisme de Banach),
$\pi_1^{-1} : X \to \Gamma(T)$ est continue. Donc $T = \pi_2 \circ \pi_1^{-1}$
est continue, comme composée d'applications continues.
\end{proof}

\begin{remarque}
\label{rem:graphe-ferme-domaine}
\index{graphe!fermé!domaine}
L'hypothèse $D(T) = X$ est cruciale dans le sens $(\Leftarrow)$.
L'Exemple~\ref{ex:derivation-fermee} montre un opérateur à graphe fermé
mais non borné, car son domaine $\mathcal{C}^1([0,1])$ est un sous-espace
propre de $\mathcal{C}([0,1])$ (non fermé pour la norme $\norm{\cdot}_\infty$).
\end{remarque}

\begin{corollaire}[Critère séquentiel du graphe fermé]
\label{cor:critere-graphe-ferme}
\index{graphe!fermé!critère séquentiel}
Soient $X$ et $Y$ des espaces de Banach et $T : X \to Y$ linéaire.
$T$ est continu si et seulement si : pour toute suite $(x_n)$ dans $X$
telle que $x_n \to 0$ et $Tx_n \to y$, on a $y = 0$.
\end{corollaire}

\begin{proof}
$(\Rightarrow)$ est évident.

$(\Leftarrow)$ Montrons que $\Gamma(T)$ est fermé. Soit $(x_n, Tx_n) \to (x,y)$.
Alors $x_n \to x$ et $Tx_n \to y$. Posons $z_n = x_n - x$. Alors $z_n \to 0$
et $Tz_n = Tx_n - Tx \to y - Tx$. Par hypothèse, $y - Tx = 0$, donc $y = Tx$.
Le Théorème~\ref{thm:graphe-ferme} conclut.
\end{proof}

\begin{exemple}[Application du critère du graphe fermé]
\label{ex:application-critere}
\index{opérateur!de multiplication}
Soit $g \in L^\infty([0,1])$ et $M_g : L^2([0,1]) \to L^2([0,1])$\index{opérateur!de multiplication}
l'opérateur de multiplication $M_g f = gf$. Montrons que $M_g$ est borné
en utilisant le critère du graphe fermé.

Si $f_n \to 0$ dans $L^2$ et $gf_n \to h$ dans $L^2$, alors il existe une
sous-suite $f_{n_k} \to 0$ p.p.\ et $gf_{n_k} \to h$ p.p.\ Comme $g$ est
bornée, $gf_{n_k} \to 0$ p.p., donc $h = 0$ p.p.
\end{exemple}


%% ══════════════════════════════════════════════════════════════════
\section{Applications : opérateurs non bornés}
\label{sec:operateurs-non-bornes}
%% ══════════════════════════════════════════════════════════════════

\index{opérateur!non borné|(}

Le théorème du graphe fermé a des conséquences importantes pour l'étude des
opérateurs non bornés, qui apparaissent naturellement en théorie spectrale
et en équations aux dérivées partielles.

\begin{definition}[Opérateur fermable]
\label{def:fermable}
\index{opérateur!fermable}
\index{fermable}
Un opérateur $T : D(T) \subset X \to Y$ est dit \emph{fermable} si
l'adhérence de son graphe $\overline{\Gamma(T)}$ dans $X \times Y$ est
encore le graphe d'un opérateur (linéaire). Cet opérateur est alors noté
$\overline{T}$ et appelé la \emph{fermeture}\index{fermeture!d'un opérateur} de $T$.
\end{definition}

\begin{proposition}[Caractérisation des opérateurs fermables]
\label{prop:caract-fermable}
\index{opérateur!fermable!caractérisation}
Un opérateur $T : D(T) \subset X \to Y$ est fermable si et seulement si
la condition suivante est satisfaite :
\[
  \bigl[ x_n \in D(T),\; x_n \to 0,\; Tx_n \to y \bigr]
  \implies y = 0.
\]
\end{proposition}

\begin{proof}
$(\Rightarrow)$ Si $T$ est fermable, $\overline{\Gamma(T)} = \Gamma(\overline{T})$.
Si $x_n \to 0$ et $Tx_n \to y$, alors $(x_n, Tx_n) \to (0,y) \in \overline{\Gamma(T)} = \Gamma(\overline{T})$,
donc $\overline{T}(0) = y$, d'où $y = 0$.

$(\Leftarrow)$ Supposons la condition satisfaite. Il faut vérifier que
$\overline{\Gamma(T)}$ est le graphe d'une application, i.e.\ que si $(0,y) \in \overline{\Gamma(T)}$,
alors $y = 0$. C'est exactement la condition donnée : si $(0,y) \in \overline{\Gamma(T)}$,
il existe $(x_n, Tx_n) \to (0,y)$, donc $x_n \to 0$ et $Tx_n \to y$, d'où $y = 0$.
\end{proof}

\begin{proposition}[Norme du graphe]
\label{prop:norme-graphe}
\index{norme!du graphe}
Soit $T : D(T) \subset X \to Y$ un opérateur fermé. Alors $D(T)$ muni de
la \emph{norme du graphe}
\[
  \norm{x}_T = \norm{x}_X + \norm{Tx}_Y
\]
est un espace de Banach. De plus, $T : (D(T), \norm{\cdot}_T) \to Y$ est
continu de norme $\leq 1$.
\end{proposition}

\begin{proof}
L'application $\Phi : (D(T), \norm{\cdot}_T) \to (\Gamma(T), \norm{\cdot}_{X \times Y})$,
$x \mapsto (x, Tx)$ est une isométrie surjective\index{isométrie}. Comme $\Gamma(T)$
est fermé dans l'espace de Banach $X \times Y$, c'est un espace de Banach,
et donc $(D(T), \norm{\cdot}_T)$ aussi.

Enfin, $\norm{Tx}_Y \leq \norm{x}_X + \norm{Tx}_Y = \norm{x}_T$, donc
$\norm{T}_{(D(T),\norm{\cdot}_T) \to Y} \leq 1$.
\end{proof}

\begin{corollaire}
\label{cor:domaine-ferme-borne}
\index{opérateur!fermé!à domaine fermé}
Soit $T : D(T) \subset X \to Y$ un opérateur fermé entre espaces de Banach.
Si $D(T)$ est fermé dans $X$, alors $T$ est borné.
\end{corollaire}

\begin{proof}
Si $D(T)$ est fermé dans l'espace de Banach $X$, alors $(D(T), \norm{\cdot}_X)$
est un espace de Banach. L'identité $\Id : (D(T), \norm{\cdot}_T) \to (D(T), \norm{\cdot}_X)$
est continue (car $\norm{x}_X \leq \norm{x}_T$). Par le
Corollaire~\ref{cor:isomorphisme-banach}, $\Id^{-1}$ est continue :
il existe $C > 0$ tel que $\norm{x}_T \leq C \norm{x}_X$, d'où
$\norm{Tx}_Y \leq \norm{x}_T \leq C \norm{x}_X$.
\end{proof}


%% ══════════════════════════════════════════════════════════════════
\section{Théorème de l'image fermée}
\label{sec:image-fermee}
%% ══════════════════════════════════════════════════════════════════

\index{théorème!de l'image fermée|(}
\index{image!fermée|(}

\begin{definition}[Image et conoyau]
\label{def:image-conoyau}
Soit $T \in \mathcal{L}(X,Y)$. L'\emph{image}\index{image} de $T$ est
$\im T = \{Tx : x \in X\}$, et le \emph{conoyau}\index{conoyau} est
$\coker T = Y / \im T$.
\end{definition}

\begin{theoreme}[Théorème de l'image fermée]
\label{thm:image-fermee}
\index{théorème!de l'image fermée}
Soient $X$ et $Y$ des espaces de Banach et $T \in \mathcal{L}(X,Y)$.
Les assertions suivantes sont équivalentes :
\begin{enumerate}[label=(\roman*)]
  \item $\im T$ est fermée dans $Y$ ;
  \item $\im T^*$ est fermée dans $X^*$ ;
  \item $\im T = (\Ker T^*)^\perp := \{y \in Y : f(y) = 0\ \forall f \in \Ker T^*\}$ ;
  \item $\im T^* = (\Ker T)^\perp := \{f \in X^* : f(x) = 0\ \forall x \in \Ker T\}$.
\end{enumerate}
\end{theoreme}

\begin{proof}
Nous démontrons les implications principales.

\textbf{(i) $\Rightarrow$ (iii) :}
L'inclusion $\im T \subset (\Ker T^*)^\perp$ est toujours vraie : si $y = Tx$
et $f \in \Ker T^*$, alors $f(y) = f(Tx) = (T^*f)(x) = 0$.

Pour l'inclusion inverse, supposons $\im T$ fermée et soit $y_0 \notin \im T$.
Comme $\im T$ est un sous-espace fermé, par le théorème de Hahn-Banach\index{Hahn-Banach, théorème de},
il existe $f \in Y^*$ tel que $f|_{\im T} = 0$ et $f(y_0) \neq 0$.
Alors $T^*f = f \circ T = 0$, donc $f \in \Ker T^*$, et $f(y_0) \neq 0$,
donc $y_0 \notin (\Ker T^*)^\perp$.

\textbf{(i) $\Rightarrow$ (iv) :}
L'inclusion $\im T^* \subset (\Ker T)^\perp$ est toujours vraie.
Pour l'inverse, supposons $\im T$ fermée. L'opérateur
$\tilde{T} : X / \Ker T \to \im T$, $\tilde{T}([x]) = Tx$ est linéaire,
continu, bijectif entre espaces de Banach ($X/\Ker T$ est Banach car
$\Ker T$ est fermé, et $\im T$ est Banach par hypothèse).
Par le Corollaire~\ref{cor:isomorphisme-banach}, $\tilde{T}$ est un isomorphisme.

Soit $g \in (\Ker T)^\perp$. Alors $g$ s'annule sur $\Ker T$, donc
$g$ passe au quotient : il existe $\tilde{g} \in (X/\Ker T)^*$ tel que
$g = \tilde{g} \circ \pi$ (où $\pi : X \to X/\Ker T$ est la projection
canonique). Posons $h = \tilde{g} \circ \tilde{T}^{-1} \in (\im T)^*$.
En prolongeant $h$ à $Y$ par Hahn-Banach, on obtient $f \in Y^*$, et
$T^*f(x) = f(Tx) = h(Tx) = \tilde{g}(\tilde{T}^{-1}(Tx)) = \tilde{g}([x]) = g(x)$.
Donc $g = T^*f \in \im T^*$.

\textbf{(iv) $\Rightarrow$ (ii) :} Trivial.

\textbf{(ii) $\Rightarrow$ (i) :}
On admet cette implication, dont la preuve repose sur des arguments de
dualité plus fins et le théorème bipolaire.
\end{proof}

\begin{proposition}[Caractérisation par la constante de surjection]
\label{prop:constante-surjection}
\index{constante de surjection}
\index{minoré!opérateur}
Soient $X$ et $Y$ des espaces de Banach et $T \in \mathcal{L}(X,Y)$. Alors
$\im T$ est fermée si et seulement s'il existe $c > 0$ tel que
\[
  \forall y \in \im T, \quad \inf_{x \in T^{-1}(\{y\})} \norm{x}_X \leq c\, \norm{y}_Y,
\]
autrement dit, $T$ induit un isomorphisme de $X / \Ker T$ sur $\im T$.
\end{proposition}

\begin{proof}
$(\Leftarrow)$ Soit $(y_n) \subset \im T$ avec $y_n \to y$. Pour chaque $n$,
choisissons $x_n$ tel que $Tx_n = y_n$ et $\norm{x_n} \leq c\norm{y_n} + 1/n$.
La suite $(x_n)$ est bornée (car $(y_n)$ converge donc est bornée). En passant
au quotient $X/\Ker T$ et en utilisant la complétude, on montre que $(x_n + \Ker T)$
a une sous-suite convergente vers $[x] \in X/\Ker T$. Alors $y = Tx \in \im T$.

Plus directement : l'opérateur induit $\tilde{T} : X/\Ker T \to Y$ est
injectif et $\norm{y} \geq c^{-1}\norm{[x]}_{X/\Ker T}$ pour $y = Tx$.
Donc $\tilde{T}$ est un isomorphisme sur son image. L'image d'un espace de Banach
par un opérateur minoré est complète, donc fermée.

$(\Rightarrow)$ Si $\im T$ est fermée, c'est un espace de Banach.
$\tilde{T} : X/\Ker T \to \im T$ est bijectif continu entre Banach,
donc c'est un isomorphisme par le Corollaire~\ref{cor:isomorphisme-banach}.
En particulier $\norm{[x]} \leq \norm{\tilde{T}^{-1}}\,\norm{Tx}$, ce qui donne
la constante $c = \norm{\tilde{T}^{-1}}$.
\end{proof}


%% ══════════════════════════════════════════════════════════════════
\section{Principe de factorisation}
\label{sec:factorisation}
%% ══════════════════════════════════════════════════════════════════

\index{principe de factorisation|(}
\index{factorisation|(}

\begin{theoreme}[Principe de factorisation]
\label{thm:factorisation}
\index{factorisation!principe de}
Soient $X$, $Y$, $Z$ des espaces de Banach, $S \in \mathcal{L}(X,Z)$ et
$T \in \mathcal{L}(X,Y)$. Supposons que $\Ker T \subset \Ker S$.
Alors il existe un unique opérateur $\tilde{S} : \im T \to Z$ tel que
$S = \tilde{S} \circ T$. Si de plus $\im T$ est fermée, alors
$\tilde{S}$ est continu.
\end{theoreme}

\begin{proof}
\textbf{Unicité :} Si $\tilde{S}$ existe, alors pour $y = Tx \in \im T$,
$\tilde{S}(y) = \tilde{S}(Tx) = Sx$, ce qui détermine $\tilde{S}$ de manière unique.
Il faut vérifier la bonne définition : si $Tx = Tx'$, alors $x - x' \in \Ker T \subset \Ker S$,
donc $Sx = Sx'$.

\textbf{Existence :} Posons $\tilde{S}(Tx) = Sx$. Par ce qui précède,
$\tilde{S}$ est bien défini et linéaire.

\textbf{Continuité :} Si $\im T$ est fermée, l'opérateur induit
$\bar{T} : X/\Ker T \to \im T$ est un isomorphisme de Banach
(par le Corollaire~\ref{cor:isomorphisme-banach}). On a
$\tilde{S} = \bar{S} \circ \bar{T}^{-1}$, où $\bar{S} : X/\Ker T \to Z$
est induit par $S$ (bien défini car $\Ker T \subset \Ker S$). Comme
$\bar{S}$ et $\bar{T}^{-1}$ sont continus, $\tilde{S}$ l'est aussi.
\end{proof}

\begin{corollaire}
\label{cor:factorisation-noyaux}
Sous les hypothèses du théorème, si $\im T$ est fermée, alors
\[
  \norm{\tilde{S}}_{\mathcal{L}(\im T, Z)} \leq \norm{S} \cdot \norm{\bar{T}^{-1}}.
\]
\end{corollaire}


%% ══════════════════════════════════════════════════════════════════
\section{Espaces de Fréchet et théorème du graphe fermé généralisé}
\label{sec:frechet}
%% ══════════════════════════════════════════════════════════════════

\index{espace!de Fréchet|(}
\index{théorème!du graphe fermé!généralisé|(}

Les espaces de Fréchet\index{Fréchet, espace de} généralisent les espaces de Banach
en autorisant une topologie définie par une \emph{famille dénombrable} de semi-normes,
avec complétude.

\begin{definition}[Espace de Fréchet]
\label{def:frechet}
\index{espace!de Fréchet}
Un espace de Fréchet est un espace vectoriel topologique $X$ dont la topologie
est définie par une famille dénombrable de semi-normes\index{semi-norme}
$(p_n)_{n \geq 1}$ séparante (i.e.\ $p_n(x) = 0$ pour tout $n$ implique $x = 0$),
et qui est complet pour la métrique\index{métrique!d'un espace de Fréchet}
\[
  d(x,y) = \sum_{n=1}^{\infty} \frac{1}{2^n} \frac{p_n(x-y)}{1 + p_n(x-y)}.
\]
\end{definition}

\begin{exemple}
\label{ex:frechet-exemples}
\begin{enumerate}[label=(\roman*)]
  \item Tout espace de Banach est un espace de Fréchet (une seule semi-norme,
    qui est une norme).
  \item L'espace $\mathcal{C}^\infty([0,1])$\index{$\mathcal{C}^\infty$} muni des semi-normes
    $p_n(f) = \sup_{x \in [0,1]} \abs{f^{(n)}(x)}$ est un espace de Fréchet.
  \item L'espace $\mathcal{C}(\R)$\index{$\mathcal{C}(\R)$} muni des semi-normes
    $p_n(f) = \sup_{\abs{x} \leq n} \abs{f(x)}$ (convergence uniforme
    sur les compacts\index{convergence!uniforme sur les compacts})
    est un espace de Fréchet.
  \item L'espace des suites $\omega = \K^{\N}$\index{$\omega$} muni des semi-normes
    $p_n((x_k)) = \abs{x_n}$ est un espace de Fréchet.
\end{enumerate}
\end{exemple}

\begin{remarque}
\label{rem:frechet-localement-convexe}
Un espace de Fréchet est un espace localement convexe\index{localement convexe}
métrisable complet. La métrique $d$ ci-dessus est invariante par translation.
\end{remarque}

\begin{theoreme}[Application ouverte pour les espaces de Fréchet]
\label{thm:app-ouverte-frechet}
\index{théorème!de l'application ouverte!pour les espaces de Fréchet}
Soient $X$ et $Y$ des espaces de Fréchet et $T : X \to Y$ un opérateur
linéaire continu surjectif. Alors $T$ est ouvert.
\end{theoreme}

\begin{proof}
La preuve suit le même schéma que pour les espaces de Banach. Le théorème
de Baire s'applique car un espace de Fréchet est un espace métrique complet.
L'argument de «~suite géométrique convergente~» de la preuve du
Théorème~\ref{thm:application-ouverte} se généralise en utilisant la
famille de semi-normes au lieu de la norme unique.
\end{proof}

\begin{theoreme}[Graphe fermé pour les espaces de Fréchet]
\label{thm:graphe-ferme-frechet}
\index{théorème!du graphe fermé!pour les espaces de Fréchet}
Soient $X$ et $Y$ des espaces de Fréchet et $T : X \to Y$ un opérateur
linéaire à graphe fermé. Alors $T$ est continu.
\end{theoreme}

\begin{proof}
$\Gamma(T)$ est un sous-espace fermé de l'espace de Fréchet $X \times Y$
(muni des semi-normes $q_n(x,y) = p_n^X(x) + p_n^Y(y)$), donc c'est un
espace de Fréchet. La projection $\pi_1 : \Gamma(T) \to X$ est linéaire,
continue, bijective et surjective. Par le Théorème~\ref{thm:app-ouverte-frechet},
$\pi_1$ est ouverte, donc $\pi_1^{-1}$ est continue.
Alors $T = \pi_2 \circ \pi_1^{-1}$ est continue.
\end{proof}

\begin{remarque}[Limites du théorème du graphe fermé]
\label{rem:limites-graphe-ferme}
Le théorème du graphe fermé est faux pour des espaces vectoriels topologiques
généraux. Il existe des extensions au-delà des espaces de Fréchet
(espaces \emph{webbed}\index{espace!webbed} de De Wilde\index{De Wilde},
espaces ultrabornologiques\index{ultrabornologique}), mais ces extensions
nécessitent des hypothèses structurelles sur les espaces.
\end{remarque}


%% ══════════════════════════════════════════════════════════════════
\section{Exercices}
\label{sec:exos-ch6}
%% ══════════════════════════════════════════════════════════════════

\begin{exercice}[$\star$]
\label{exo:ouverte-quotient}
\index{application!ouverte!quotient}
Soit $X$ un espace de Banach et $M \subset X$ un sous-espace fermé.
Montrer que la projection canonique $\pi : X \to X/M$ est une application ouverte.
\textit{(Ne pas utiliser le théorème de l'application ouverte ; le démontrer directement.)}
\end{exercice}

\begin{exercice}[$\star$]
\label{exo:graphe-ferme-simple}
\index{graphe!fermé!exercice}
Soit $T : \ell^1 \to \ell^1$ l'opérateur défini par
$T(x_1, x_2, x_3, \ldots) = (x_1, 2x_2, 3x_3, \ldots)$, avec
$D(T) = \{x \in \ell^1 : \sum n\abs{x_n} < \infty\}$.
\begin{enumerate}[label=(\alph*)]
  \item Montrer que $T$ est un opérateur fermé.
  \item Montrer que $T$ n'est pas borné.
  \item Pourquoi cela ne contredit-il pas le théorème du graphe fermé ?
\end{enumerate}
\end{exercice}

\begin{exercice}[$\star$]
\label{exo:isomorphisme-lp}
\index{$\ell^p$!isomorphisme}
Montrer que si $1 \leq p < q \leq \infty$, alors $\ell^p$ et $\ell^q$
ne sont pas isomorphes comme espaces de Banach.
\textit{(Indication : considérer les propriétés de réflexivité, ou utiliser
le fait que l'identité $\ell^p \hookrightarrow \ell^q$ n'est pas surjective.)}
\end{exercice}

\begin{exercice}[$\star\star$]
\label{exo:app-ouverte-non-surjective}
\index{application!ouverte!image d'un opérateur}
Soit $T \in \mathcal{L}(X,Y)$ entre espaces de Banach, avec $\im T$ fermée.
Montrer que $T$ est une application ouverte de $X$ sur $\im T$ (muni de
la topologie induite).
\end{exercice}

\begin{exercice}[$\star\star$]
\label{exo:graphe-ferme-convergence}
\index{graphe!fermé!convergence de suites}
Soient $X$ et $Y$ des espaces de Banach et $T_n : X \to Y$ des opérateurs
linéaires continus tels que pour tout $x \in X$, $(T_n x)$ converge vers
$Tx$. On sait par le Corollaire~\ref{cor:BS-convergence} que $T$ est borné.
Donner une preuve alternative en utilisant le théorème du graphe fermé.
\end{exercice}

\begin{exercice}[$\star\star$]
\label{exo:operateur-involutif}
\index{opérateur!involutif}
Soit $T : X \to X$ un opérateur linéaire sur un espace de Banach $X$
tel que $T^2 = \Id$. Montrer que $T$ est borné.
\textit{(Utiliser le théorème du graphe fermé : si $x_n \to 0$ et $Tx_n \to y$,
appliquer $T$ à nouveau.)}
\end{exercice}

\begin{exercice}[$\star\star$]
\label{exo:image-fermee-codimension}
\index{image!fermée!de codimension finie}
Soit $T \in \mathcal{L}(X,Y)$ entre espaces de Banach.
Montrer que si $\im T$ est de codimension finie\index{codimension!finie}
dans $Y$ (i.e.\ $\dim(Y/\im T) < \infty$), alors $\im T$ est fermée.
\textit{(Écrire $Y = \im T \oplus F$ avec $\dim F < \infty$, et utiliser
le fait que les sous-espaces de dimension finie sont fermés.)}
\end{exercice}

\begin{exercice}[$\star\star$]
\label{exo:Hellinger-Toeplitz}
\index{théorème!de Hellinger-Toeplitz}
\index{Hellinger-Toeplitz}
(\textbf{Théorème de Hellinger-Toeplitz})
Soit $H$ un espace de Hilbert\index{espace!de Hilbert} et $T : H \to H$
un opérateur linéaire \emph{symétrique}\index{opérateur!symétrique}, i.e.\
$\inner{Tx}{y} = \inner{x}{Ty}$ pour tous $x, y \in H$. Montrer que $T$
est automatiquement borné.
\textit{(Utiliser le théorème du graphe fermé.)}
\end{exercice}

\begin{exercice}[$\star\star\star$]
\label{exo:inverse-automatique}
\index{inverse!automatique}
Soient $X$ et $Y$ des espaces de Banach et $T \in \mathcal{L}(X,Y)$.
Montrer les équivalences suivantes :
\begin{enumerate}[label=(\alph*)]
  \item $T$ est surjectif $\iff$ il existe $c > 0$ tel que $\norm{T^*f}_{X^*} \geq c\norm{f}_{Y^*}$ pour tout $f \in Y^*$.
  \item $T$ est injectif et $\im T$ est fermée $\iff$ il existe $c > 0$ tel que $\norm{Tx}_Y \geq c\norm{x}_X$ pour tout $x \in X$.
\end{enumerate}
\end{exercice}

\begin{exercice}[$\star\star\star$]
\label{exo:factorisation-application}
\index{factorisation!application}
Soient $X$, $Y$ des espaces de Banach et $S, T \in \mathcal{L}(X,Y)$
tels que $\norm{Sx} \leq C \norm{Tx}$ pour tout $x \in X$ et une constante
$C > 0$. Montrer qu'il existe $\tilde{S} \in \mathcal{L}(\overline{\im T}, Y)$
tel que $S = \tilde{S} \circ T$.
\end{exercice}

\begin{exercice}[$\star\star\star$]
\label{exo:trois-espaces}
\index{problème des trois espaces}
(\textbf{Problème des trois espaces})
Soit $X$ un espace de Banach et $M \subset X$ un sous-espace fermé.
\begin{enumerate}[label=(\alph*)]
  \item Montrer que si $M$ et $X/M$ sont des espaces de Hilbert, alors $X$
    est isomorphe (comme espace de Banach) à $M \oplus X/M$.
  \item En utilisant l'application ouverte, montrer que la suite exacte courte
    $0 \to M \xrightarrow{\iota} X \xrightarrow{\pi} X/M \to 0$
    admet une section continue à droite $s : X/M \to X$ (i.e.\ $\pi \circ s = \Id$)
    si et seulement si $M$ admet un supplémentaire topologique fermé dans $X$.
\end{enumerate}
\end{exercice}

\begin{exercice}[$\star\star\star$]
\label{exo:frechet-quotient}
\index{espace!de Fréchet!quotient}
Soit $X$ un espace de Fréchet et $M \subset X$ un sous-espace fermé.
Montrer que $X/M$ est un espace de Fréchet.
\end{exercice}

\begin{exercice}[$\star\star\star$]
\label{exo:espaces-webbed}
\index{espace!webbed}
\index{De Wilde}
Montrer que tout espace de Banach est un espace webbed au sens de
De Wilde. \textit{(Construire un réseau ($=$~web) en utilisant les boules
$B(0, 2^{-n})$.)}
\end{exercice}

\index{théorème!de l'application ouverte|)}
\index{application!ouverte|)}
\index{isomorphisme!de Banach|)}
\index{théorème!du graphe fermé|)}
\index{graphe!fermé|)}
\index{opérateur!non borné|)}
\index{théorème!de l'image fermée|)}
\index{image!fermée|)}
\index{principe de factorisation|)}
\index{factorisation|)}
\index{espace!de Fréchet|)}
\index{théorème!du graphe fermé!généralisé|)}
%%%%%%%%%%%%%%%%%%%%%%%%%%%%%%%%%%%%%%%%%%%%%%%%%%%%%%%%%%%%%%%%%%%%
%  Analyse Fonctionnelle — Chapitre 7
%  Espaces Duaux et Réflexivité
%%%%%%%%%%%%%%%%%%%%%%%%%%%%%%%%%%%%%%%%%%%%%%%%%%%%%%%%%%%%%%%%%%%%

\chapter{Espaces Duaux et Réflexivité}
\label{chap:dualite-reflexivite}
\minitoc

\vspace{1em}

Ce chapitre approfondit l'étude des espaces duaux\index{espace dual!topologique}
amorcée dans les chapitres précédents. Nous identifions explicitement les duaux
des espaces $\ell^p$\index{lp@$\ell^p$} et $c_0$\index{c0@$c_0$}, introduisons
l'injection canonique\index{injection canonique} dans le bidual\index{bidual}, et
développons la théorie de la réflexivité\index{réflexivité}. Les résultats
profonds de James\index{James, Robert C.}, Eberlein--\v{S}mulian\index{Eberlein--Smulian@Eberlein--\v{S}mulian}
et Milman--Pettis\index{Milman--Pettis} sont présentés, certains avec preuve complète.

%% ══════════════════════════════════════════════════════════════════
\section{Le dual topologique}
\label{sec:dual-topologique}
%% ══════════════════════════════════════════════════════════════════

\index{dual topologique|(}

\begin{definition}[Dual topologique]
\label{def:dual-topologique}
Soit $E$ un espace vectoriel normé\index{espace vectoriel normé} sur $\K$ ($= \R$ ou $\C$).
Le \emph{dual topologique} de $E$, noté $E'$\index{E'@$E'$}, est l'espace
$\mathcal{L}(E, \K)$ des formes linéaires continues\index{forme linéaire!continue} sur $E$,
muni de la norme d'opérateur :
\[
  \norm{f}_{E'} = \sup_{\substack{x \in E \\ x \neq 0}} \frac{\abs{f(x)}}{\norm{x}}
  = \sup_{\norm{x} \le 1} \abs{f(x)}.
\]
\end{definition}

\begin{remarque}
\label{rem:dual-banach}
Le dual $E'$ est toujours un espace de Banach\index{espace de Banach}, même si $E$
ne l'est pas, car $\K$ est complet.
\end{remarque}

\begin{notation}
\label{not:crochet-dualite}
Pour $f \in E'$ et $x \in E$, on note souvent
$\inner{f}{x}_{E',E}$\index{crochet de dualité} ou simplement $\inner{f}{x}$
la valeur $f(x)$, appelée \emph{crochet de dualité}\index{crochet de dualité}.
\end{notation}

\subsection{Rappels : le théorème de Hahn-Banach}
\label{subsec:rappel-hahn-banach}
\index{théorème!de Hahn-Banach|(}

Le théorème de Hahn-Banach\index{Hahn-Banach|see{théorème de Hahn-Banach}},
démontré au chapitre~4, est l'outil fondamental pour construire des formes
linéaires continues. Rappelons ses conséquences les plus utiles.

\begin{theoreme}[Hahn-Banach analytique --- rappel]
\label{thm:hahn-banach-rappel}
\index{théorème!de Hahn-Banach!forme analytique}
Soit $E$ un espace vectoriel normé, $F \subset E$ un sous-espace vectoriel,
et $g \in F'$. Il existe $f \in E'$ tel que $f|_F = g$ et $\norm{f}_{E'} = \norm{g}_{F'}$.
\end{theoreme}

\begin{corollaire}[Séparation d'un point]
\label{cor:separation-point}
\index{séparation!d'un point}
Pour tout $x_0 \in E$ avec $x_0 \neq 0$, il existe $f \in E'$ tel que
$\norm{f} = 1$ et $f(x_0) = \norm{x_0}$.
\end{corollaire}

\begin{corollaire}[Le dual sépare les points]
\label{cor:dual-separe}
\index{dual topologique!sépare les points}
Si $x, y \in E$ sont tels que $f(x) = f(y)$ pour tout $f \in E'$, alors $x = y$.
En particulier, $E'$ sépare les points de $E$.
\end{corollaire}

\begin{corollaire}[Norme via le dual]
\label{cor:norme-via-dual}
\index{norme!via le dual}
Pour tout $x \in E$,
\[
  \norm{x} = \sup_{\substack{f \in E' \\ \norm{f} \le 1}} \abs{f(x)}
  = \max_{\substack{f \in E' \\ \norm{f} \le 1}} \abs{f(x)}.
\]
\end{corollaire}

\index{théorème!de Hahn-Banach|)}

%% ══════════════════════════════════════════════════════════════════
\section{Dual de $\ell^p$ pour $1 \le p < \infty$}
\label{sec:dual-lp}
%% ════════════════════════════════════════════════════════════════════

\index{lp@$\ell^p$!dual de|(}

Nous identifions maintenant les duaux des espaces de suites classiques. On note
$p'$ l'exposant conjugué\index{exposant conjugué} de $p$, défini par
$\frac{1}{p} + \frac{1}{p'} = 1$ (avec la convention $1' = \infty$).

\begin{theoreme}[{Dual de $\ell^p$, $1 \le p < \infty$}]
\label{thm:dual-lp}
\index{lp@$\ell^p$!dual de}
\index{isométrie!isomorphe}
Soit $1 \le p < \infty$ et $q = p'$ l'exposant conjugué de $p$.
L'application
\[
  \Phi : \ell^q \longrightarrow (\ell^p)', \quad
  y = (y_n) \longmapsto \Phi_y, \quad
  \Phi_y(x) = \sum_{n=1}^{\infty} x_n y_n,
\]
est une isométrie isomorphe : $(\ell^p)' \cong \ell^q$.
\end{theoreme}

\begin{proof}
La preuve se décompose en trois étapes.

\medskip
\noindent\textbf{Étape 1 : $\Phi$ est bien définie et $\norm{\Phi_y} \le \norm{y}_q$.}

Soit $y \in \ell^q$. Pour $x \in \ell^p$, l'inégalité de Hölder\index{inégalité!de Hölder} donne
\[
  \abs{\Phi_y(x)} = \abs[\bigg]{\sum_{n=1}^{\infty} x_n y_n}
  \le \sum_{n=1}^{\infty} \abs{x_n} \abs{y_n}
  \le \norm{x}_p \, \norm{y}_q.
\]
Donc $\Phi_y \in (\ell^p)'$ et $\norm{\Phi_y}_{(\ell^p)'} \le \norm{y}_q$.

\medskip
\noindent\textbf{Étape 2 : $\norm{\Phi_y} = \norm{y}_q$.}

Il faut exhiber un $x \in \ell^p$ avec $\norm{x}_p = 1$ tel que
$\abs{\Phi_y(x)} = \norm{y}_q$.

\emph{Cas $1 < p < \infty$} (donc $1 < q < \infty$). Si $y = 0$, c'est trivial.
Sinon, posons
\[
  x_n = \frac{\abs{y_n}^{q-1} \sgn(\overline{y_n})}{\norm{y}_q^{q/p}},
  \quad n \ge 1.
\]
Comme $q - 1 = q/p$, on a $\abs{x_n}^p = \abs{y_n}^q / \norm{y}_q^q$, d'où
$\norm{x}_p^p = \norm{y}_q^q / \norm{y}_q^q = 1$. De plus,
\[
  \Phi_y(x) = \sum_{n=1}^{\infty} x_n y_n
  = \frac{1}{\norm{y}_q^{q/p}} \sum_{n=1}^{\infty} \abs{y_n}^q
  = \frac{\norm{y}_q^q}{\norm{y}_q^{q/p}} = \norm{y}_q^{q - q/p} = \norm{y}_q.
\]

\emph{Cas $p = 1$} (donc $q = \infty$). Pour tout $\varepsilon > 0$, il existe $n_0$
tel que $\abs{y_{n_0}} > \norm{y}_\infty - \varepsilon$. On prend $x = \sgn(\overline{y_{n_0}}) \, e_{n_0}$,
où $e_{n_0}$ est le $n_0$-ième vecteur canonique\index{vecteur canonique}, de sorte
que $\norm{x}_1 = 1$ et $\Phi_y(x) = \abs{y_{n_0}} > \norm{y}_\infty - \varepsilon$.
Comme $\varepsilon$ est arbitraire, $\norm{\Phi_y} = \norm{y}_\infty$.

\medskip
\noindent\textbf{Étape 3 : $\Phi$ est surjective.}

Soit $f \in (\ell^p)'$. Posons $y_n = f(e_n)$\index{vecteur canonique} pour $n \ge 1$.
Montrons que $y = (y_n) \in \ell^q$ et que $f = \Phi_y$.

Pour tout $x = (x_n) \in \ell^p$, on a $x = \sum_{n=1}^{\infty} x_n e_n$ (convergence
dans $\ell^p$), donc par continuité de $f$,
\[
  f(x) = \sum_{n=1}^{\infty} x_n f(e_n) = \sum_{n=1}^{\infty} x_n y_n.
\]
Il reste à montrer $y \in \ell^q$.

\emph{Cas $1 < p < \infty$.} Pour $N \in \N^*$, posons
$x_n^{(N)} = \abs{y_n}^{q-1} \sgn(\overline{y_n})$ pour $1 \le n \le N$ et
$x_n^{(N)} = 0$ pour $n > N$. Alors $x^{(N)} \in \ell^p$ et
\[
  \norm{x^{(N)}}_p^p = \sum_{n=1}^N \abs{y_n}^{(q-1)p} = \sum_{n=1}^N \abs{y_n}^q,
\]
car $(q-1)p = q$. De plus,
\[
  f(x^{(N)}) = \sum_{n=1}^N \abs{y_n}^q \le \norm{f} \, \norm{x^{(N)}}_p
  = \norm{f} \left( \sum_{n=1}^N \abs{y_n}^q \right)^{1/p}.
\]
Donc $\left(\sum_{n=1}^N \abs{y_n}^q\right)^{1 - 1/p} \le \norm{f}$,
soit $\left(\sum_{n=1}^N \abs{y_n}^q\right)^{1/q} \le \norm{f}$.
En passant à la limite $N \to \infty$, on obtient $\norm{y}_q \le \norm{f} < \infty$.

\emph{Cas $p = 1$.} On a $\abs{y_n} = \abs{f(e_n)} \le \norm{f} \norm{e_n}_1 = \norm{f}$,
donc $\norm{y}_\infty \le \norm{f} < \infty$ et $y \in \ell^\infty$.

Ainsi $f = \Phi_y$ avec $y \in \ell^q$, ce qui achève la preuve de la surjectivité,
et donc de l'isométrie isomorphe.
\end{proof}

\index{lp@$\ell^p$!dual de|)}

%% ══════════════════════════════════════════════════════════════════
\section{Dual de $c_0$}
\label{sec:dual-c0}
%% ══════════════════════════════════════════════════════════════════

\index{c0@$c_0$!dual de|(}

\begin{definition}[L'espace $c_0$]
\label{def:c0}
\index{c0@$c_0$!définition}
L'espace $c_0$ est le sous-espace fermé de $\ell^\infty$ formé des suites
convergentes vers $0$ :
\[
  c_0 = \Bigl\{ x = (x_n)_{n \ge 1} \in \K^\N \,:\, \lim_{n \to \infty} x_n = 0 \Bigr\},
\]
muni de la norme $\norm{x}_\infty = \sup_{n \ge 1} \abs{x_n}$.
\end{definition}

\begin{theoreme}[{Dual de $c_0$}]
\label{thm:dual-c0}
\index{c0@$c_0$!dual de}
L'application
\[
  \Phi : \ell^1 \longrightarrow (c_0)', \quad
  y = (y_n) \longmapsto \Phi_y, \quad
  \Phi_y(x) = \sum_{n=1}^{\infty} x_n y_n,
\]
est une isométrie isomorphe : $(c_0)' \cong \ell^1$.
\end{theoreme}

\begin{proof}
\noindent\textbf{Étape 1 : $\Phi$ est bien définie et isométrique.}

Pour $y \in \ell^1$ et $x \in c_0 \subset \ell^\infty$,
\[
  \abs{\Phi_y(x)} \le \sum_{n=1}^{\infty} \abs{x_n}\abs{y_n}
  \le \norm{x}_\infty \norm{y}_1,
\]
donc $\Phi_y \in (c_0)'$ et $\norm{\Phi_y} \le \norm{y}_1$.

Pour l'inégalité inverse, soit $\varepsilon > 0$. Pour $N$ assez grand, considérons
\[
  x^{(N)} = \bigl(\sgn(\overline{y_1}), \sgn(\overline{y_2}), \ldots, \sgn(\overline{y_N}), 0, 0, \ldots\bigr) \in c_0.
\]
Alors $\norm{x^{(N)}}_\infty \le 1$ et
$\Phi_y(x^{(N)}) = \sum_{n=1}^N \abs{y_n}$.
En passant à la limite,
$\norm{\Phi_y} \ge \sup_N \sum_{n=1}^N \abs{y_n} = \norm{y}_1$.
D'où $\norm{\Phi_y} = \norm{y}_1$.

\medskip
\noindent\textbf{Étape 2 : $\Phi$ est surjective.}

Soit $f \in (c_0)'$. On pose $y_n = f(e_n)$ pour tout $n \ge 1$, où $(e_n)$ est la
base canonique\index{base canonique}. Comme pour $\ell^p$, tout $x \in c_0$ s'écrit
$x = \sum_{n=1}^{\infty} x_n e_n$ (convergence dans $c_0$), donc
$f(x) = \sum_{n=1}^{\infty} x_n y_n$.

Il reste à montrer $y \in \ell^1$. Pour $N \ge 1$, posons
$x^{(N)} = (\sgn(\overline{y_1}), \ldots, \sgn(\overline{y_N}), 0, 0, \ldots)$.
Alors $\norm{x^{(N)}}_\infty \le 1$ et
\[
  \sum_{n=1}^N \abs{y_n} = f(x^{(N)}) \le \norm{f}.
\]
En passant à la limite, $\norm{y}_1 \le \norm{f} < \infty$, donc $y \in \ell^1$.
\end{proof}

\begin{remarque}
\label{rem:dual-linfty}
\index{linfty@$\ell^\infty$!dual de}
Le dual de $\ell^\infty$ est bien plus compliqué : il s'identifie à l'espace
$\mathrm{ba}(\N)$\index{ba(N)@$\mathrm{ba}(\N)$} des mesures de Banach
finiment additives sur $\N$, qui contient $\ell^1$ comme sous-espace propre.
En particulier, $(\ell^\infty)' \supsetneq \ell^1$.
\end{remarque}

\index{c0@$c_0$!dual de|)}

%% ══════════════════════════════════════════════════════════════════
\section{Dual de $L^p$}
\label{sec:dual-Lp}
%% ══════════════════════════════════════════════════════════════════

\index{Lp@$L^p$!dual de|(}

Le résultat suivant généralise le théorème~\ref{thm:dual-lp} au cadre des espaces
de Lebesgue.

\begin{theoreme}[{Dual de $L^p$, théorème de Riesz}]
\label{thm:dual-Lp}
\index{théorème!de Riesz (représentation dans $L^p$)}
\index{Riesz, Frigyes}
Soit $(\Omega, \mathcal{A}, \mu)$ un espace mesuré $\sigma$-fini\index{sigma-fini@$\sigma$-fini}
et $1 \le p < \infty$, avec $q = p'$.
L'application
\[
  \Phi : L^q(\mu) \longrightarrow (L^p(\mu))', \quad
  g \longmapsto \Phi_g, \quad
  \Phi_g(f) = \int_\Omega f \, g \, \dd\mu,
\]
est une isométrie isomorphe : $(L^p(\mu))' \cong L^q(\mu)$.
\end{theoreme}

\begin{remarque}
La preuve de ce théorème repose sur le théorème de Radon-Nikodym\index{théorème!de Radon-Nikodym}
et sera donnée au chapitre~11. L'hypothèse de $\sigma$-finitude\index{sigma-finitude@$\sigma$-finitude}
est essentielle.
\end{remarque}

\index{Lp@$L^p$!dual de|)}

%% ══════════════════════════════════════════════════════════════════
\section{L'injection canonique et le bidual}
\label{sec:injection-canonique-bidual}
%% ══════════════════════════════════════════════════════════════════

\index{bidual|(}
\index{injection canonique|(}

\begin{definition}[Bidual]
\label{def:bidual}
\index{bidual!définition}
Le \emph{bidual} de $E$ est le dual du dual : $E'' = (E')'$.
C'est un espace de Banach pour la norme d'opérateur.
\end{definition}

\begin{definition}[Injection canonique]
\label{def:injection-canonique}
\index{injection canonique!définition}
L'\emph{injection canonique} est l'application
\[
  J : E \longrightarrow E'', \quad
  J(x)(f) = f(x), \quad \forall f \in E'.
\]
Autrement dit, $J(x)$ est l'\emph{évaluation en $x$}\index{évaluation}.
\end{definition}

\begin{theoreme}[L'injection canonique est une isométrie]
\label{thm:J-isometrie}
\index{injection canonique!isométrie}
L'application $J : E \to E''$ est linéaire et isométrique. En particulier,
$J$ est injective et $J(E)$ est un sous-espace fermé de $E''$ si $E$ est complet.
\end{theoreme}

\begin{proof}
La linéarité est immédiate : pour $x, y \in E$, $\lambda \in \K$ et $f \in E'$,
\[
  J(\lambda x + y)(f) = f(\lambda x + y) = \lambda f(x) + f(y) = (\lambda J(x) + J(y))(f).
\]
Pour l'isométrie, on calcule
\[
  \norm{J(x)}_{E''} = \sup_{\substack{f \in E' \\ \norm{f} \le 1}} \abs{J(x)(f)}
  = \sup_{\substack{f \in E' \\ \norm{f} \le 1}} \abs{f(x)} = \norm{x},
\]
où la dernière égalité découle du corollaire~\ref{cor:norme-via-dual} (conséquence de
Hahn-Banach\index{théorème!de Hahn-Banach}).

Comme $J$ est isométrique, elle est injective. Si $E$ est un espace de Banach, alors
$J(E)$ est un sous-espace complet de $E''$, donc fermé.
\end{proof}

\begin{remarque}
\label{rem:J-naturelle}
\index{injection canonique!caractère naturel}
L'injection $J$ est \emph{canonique} au sens où elle ne dépend d'aucun choix
(pas de base, pas de forme linéaire particulière). C'est un morphisme naturel
au sens de la théorie des catégories\index{catégorie, théorie des}. On identifie
souvent $E$ à $J(E) \subset E''$.
\end{remarque}

\index{injection canonique|)}
\index{bidual|)}

%% ══════════════════════════════════════════════════════════════════
\section{Espaces réflexifs}
\label{sec:reflexifs}
%% ══════════════════════════════════════════════════════════════════

\index{réflexivité|(}
\index{espace réflexif|(}

\begin{definition}[Espace réflexif]
\label{def:reflexif}
\index{espace réflexif!définition}
Un espace de Banach $E$ est dit \emph{réflexif} si l'injection canonique
$J : E \to E''$ est surjective, c'est-à-dire si $J(E) = E''$.
\end{definition}

\begin{remarque}
\label{rem:reflexif-surjectif}
\index{espace réflexif!versus isomorphe au bidual}
La réflexivité exige que $J$ elle-même soit surjective, et pas seulement qu'il
existe un isomorphisme isométrique entre $E$ et $E''$. L'espace de James\index{espace de James}
fournit un exemple d'espace de Banach isométriquement isomorphe à son bidual sans être réflexif.
\end{remarque}

\begin{exemple}[Espaces réflexifs classiques]
\label{ex:reflexifs}
\begin{enumerate}[label=(\roman*)]
  \item\index{espace!de dimension finie!réflexif} Tout espace de Banach de dimension finie
    est réflexif (car $\dim E = \dim E' = \dim E''$).
  \item\index{espace de Hilbert!réflexif} Tout espace de Hilbert $H$ est réflexif, par le
    théorème de représentation de Riesz\index{théorème!de Riesz (représentation hilbertienne)}.
  \item\index{lp@$\ell^p$!réflexif} Pour $1 < p < \infty$, $\ell^p$ est réflexif.
  \item\index{Lp@$L^p$!réflexif} Pour $1 < p < \infty$ et $\mu$ $\sigma$-finie, $L^p(\mu)$ est réflexif.
\end{enumerate}
\end{exemple}

\begin{proposition}[{$\ell^p$ est réflexif pour $1 < p < \infty$}]
\label{prop:lp-reflexif}
\index{lp@$\ell^p$!réflexif}
Pour $1 < p < \infty$, l'espace $\ell^p$ est réflexif.
\end{proposition}

\begin{proof}
Soit $q = p'$ l'exposant conjugué de $p$. Par le théorème~\ref{thm:dual-lp},
$(\ell^p)' \cong \ell^q$ via $\Phi$, et $(\ell^q)' \cong \ell^p$ via $\Psi$
(car $q' = p$). Plus précisément, notons $\Phi : \ell^q \to (\ell^p)'$
l'isométrie $y \mapsto \Phi_y$ et $\Psi : \ell^p \to (\ell^q)'$ l'isométrie
$x \mapsto \Psi_x$. L'injection canonique $J : \ell^p \to (\ell^p)''$
satisfait, pour $x \in \ell^p$ et $f \in (\ell^p)'$,
\[
  J(x)(f) = f(x).
\]
Si $f = \Phi_y$ pour $y \in \ell^q$, alors
\[
  J(x)(\Phi_y) = \Phi_y(x) = \sum_{n=1}^{\infty} x_n y_n = \Psi_x(y).
\]
Ainsi $J(x) = \Psi_x \circ \Phi^{-1}$, et la surjectivité de $J$ découle de celles
de $\Phi$ et $\Psi$. Formellement, tout $\xi \in (\ell^p)''$ donne par composition
$\xi \circ \Phi \in (\ell^q)'$, donc il existe $x \in \ell^p$ tel que
$\xi \circ \Phi = \Psi_x$, i.e.\ $\xi(\Phi_y) = \sum_{n} x_n y_n$ pour tout
$y \in \ell^q$, ce qui donne $\xi = J(x)$.
\end{proof}

\begin{proposition}[$\ell^1$ n'est pas réflexif]
\label{prop:l1-non-reflexif}
\index{l1@$\ell^1$!non réflexif}
L'espace $\ell^1$ n'est pas réflexif.
\end{proposition}

\begin{proof}
Par le théorème~\ref{thm:dual-lp} (cas $p = 1$), $(\ell^1)' \cong \ell^\infty$.
Donc $(\ell^1)'' \cong (\ell^\infty)'$. Si $\ell^1$ était réflexif, on aurait
$(\ell^\infty)' \cong \ell^1$, ce qui impliquerait que $\ell^\infty$ est séparable
(car le dual d'un espace séparable est séparable lorsque le prédual est séparable...
en fait, c'est l'inverse : si $E'$ est séparable, $E$ l'est). Or on sait
que $\ell^\infty$ n'est pas séparable\index{linfty@$\ell^\infty$!non séparable}.

Plus directement, la boule unité fermée de $\ell^1$ n'est pas faiblement
compacte\index{compacité!faible}. En effet, la suite des vecteurs canoniques
$(e_n)$ vérifie $\norm{e_n}_1 = 1$ et $e_n \weakto 0$ faiblement (car pour tout
$y \in \ell^\infty = (\ell^1)'$, $\inner{y}{e_n} = y_n \to 0$ est faux en général,
par exemple pour $y = (1,1,1,\ldots)$, $\inner{y}{e_n} = 1 \not\to 0$).
Ainsi $(e_n)$ n'admet aucune sous-suite faiblement convergente dans $\ell^1$.
Le théorème d'Eberlein--\v{S}mulian (théorème~\ref{thm:eberlein-smulian})
implique que la boule unité fermée n'est pas faiblement compacte, donc $\ell^1$
n'est pas réflexif.
\end{proof}

\begin{proposition}[$c_0$ n'est pas réflexif]
\label{prop:c0-non-reflexif}
\index{c0@$c_0$!non réflexif}
L'espace $c_0$ n'est pas réflexif.
\end{proposition}

\begin{proof}
On a $(c_0)' \cong \ell^1$ et $(\ell^1)' \cong \ell^\infty$, donc
$(c_0)'' \cong \ell^\infty$. L'injection canonique $J : c_0 \hookrightarrow \ell^\infty$
s'identifie à l'inclusion naturelle, qui n'est pas surjective (la suite
constante $(1,1,1,\ldots) \in \ell^\infty \setminus c_0$).
\end{proof}

%% ══════════════════════════════════════════════════════════════════
\section{Théorèmes fondamentaux de réflexivité}
\label{sec:thm-reflexivite}
%% ══════════════════════════════════════════════════════════════════

\subsection{Le théorème de James}
\label{subsec:james}

\begin{theoreme}[Théorème de James]
\label{thm:james}
\index{théorème!de James}
\index{James, Robert C.}
Soit $E$ un espace de Banach. Les assertions suivantes sont équivalentes :
\begin{enumerate}[label=(\roman*)]
  \item $E$ est réflexif.
  \item Toute forme linéaire continue $f \in E'$ atteint sa norme, c'est-à-dire
    qu'il existe $x_0 \in E$ avec $\norm{x_0} = 1$ tel que $\abs{f(x_0)} = \norm{f}$.
  \item La boule unité fermée $\overline{B}_E$ est faiblement compacte\index{compacité!faible}.
\end{enumerate}
\end{theoreme}

\begin{remarque}
L'implication (i) $\Rightarrow$ (ii) est un exercice (exercice~\ref{exo:reflexif-atteint}).
L'implication (i) $\Rightarrow$ (iii) résulte du théorème de Banach-Alaoglu appliqué à
$E''$ via $J$, et sera démontrée au chapitre~8. L'implication (ii) $\Rightarrow$ (i)
est le résultat profond de R.C.~James (1964)\index{James, Robert C.!théorème de};
sa preuve dépasse le cadre de ce cours.
\end{remarque}

\subsection{Le théorème d'Eberlein--\v{S}mulian}
\label{subsec:eberlein-smulian}

\begin{theoreme}[Eberlein--\v{S}mulian]
\label{thm:eberlein-smulian}
\index{théorème!d'Eberlein--Smulian@d'Eberlein--\v{S}mulian}
\index{Eberlein, William Frederick}
\index{Smulian@\v{S}mulian, Vitold}
Soit $E$ un espace de Banach et $A \subset E$. Les assertions suivantes sont
équivalentes :
\begin{enumerate}[label=(\roman*)]
  \item $A$ est relativement faiblement compact\index{compacité!faible!relative}
    (l'adhérence de $A$ dans la topologie $\sigma(E,E')$ est compacte).
  \item $A$ est relativement faiblement séquentiellement compact (toute suite de $A$
    admet une sous-suite faiblement convergente dans $E$).
  \item $A$ est faiblement séquentiellement compact
    dans $A$ même\index{compacité!séquentielle}.
\end{enumerate}
\end{theoreme}

\begin{remarque}
Ce théorème est remarquable car la topologie faible n'est en général pas métrisable.
Pour les espaces métriques, compacité et compacité séquentielle coïncident, mais
ce n'est pas le cas en général pour les espaces topologiques. Le théorème
d'Eberlein--\v{S}mulian affirme que cette équivalence subsiste pour la topologie
faible des espaces de Banach. La preuve complète, technique, peut être trouvée
dans \cite{Megginson1998}\index{Megginson}.
\end{remarque}

\subsection{Sous-espaces, quotients, et réflexivité}
\label{subsec:sous-espaces-quotients-reflexifs}

\begin{proposition}[Sous-espace fermé d'un espace réflexif]
\label{prop:sous-espace-reflexif}
\index{espace réflexif!sous-espace fermé}
Tout sous-espace fermé d'un espace de Banach réflexif est réflexif.
\end{proposition}

\begin{proof}
Soit $E$ un espace de Banach réflexif et $F \subset E$ un sous-espace fermé.
Notons $J_E : E \to E''$ et $J_F : F \to F''$ les injections canoniques.
Soit $\xi \in F''$. Nous devons montrer qu'il existe $y \in F$ tel que $\xi = J_F(y)$.

Considérons l'application restriction $R : E' \to F'$ définie par $R(f) = f|_F$.
Par le théorème de Hahn-Banach, $R$ est surjective. Son transposée
$R' : F'' \to E''$ est définie par
\[
  R'(\xi)(f) = \xi(R(f)) = \xi(f|_F), \quad \forall f \in E'.
\]
Comme $E$ est réflexif, il existe $x \in E$ tel que $J_E(x) = R'(\xi)$, c'est-à-dire
\[
  f(x) = \xi(f|_F), \quad \forall f \in E'.
\]
Montrons que $x \in F$. Si $x \notin F$, par Hahn-Banach il existe $f_0 \in E'$ tel que
$f_0|_F = 0$ et $f_0(x) \neq 0$. Mais alors
$f_0(x) = \xi(f_0|_F) = \xi(0) = 0$, contradiction. Donc $x \in F$.

Pour tout $g \in F'$, en choisissant $f \in E'$ tel que $f|_F = g$ (Hahn-Banach),
\[
  J_F(x)(g) = g(x) = f(x) = \xi(f|_F) = \xi(g).
\]
Donc $J_F(x) = \xi$, ce qui prouve que $J_F$ est surjective.
\end{proof}

\begin{proposition}[Quotient d'un espace réflexif]
\label{prop:quotient-reflexif}
\index{espace réflexif!quotient}
\index{espace quotient!réflexif}
Si $E$ est un espace de Banach réflexif et $F \subset E$ un sous-espace fermé,
alors $E/F$ est réflexif.
\end{proposition}

\begin{proof}
Notons $\pi : E \to E/F$ la projection canonique\index{projection canonique} et
$J : E/F \to (E/F)''$ l'injection canonique. Soit $\xi \in (E/F)''$.

L'application $\pi' : (E/F)' \to E'$ (transposée de $\pi$) est une isométrie de
$(E/F)'$ sur $F^\perp = \{ f \in E' : f|_F = 0 \}$\index{annulateur}.
On a aussi $\pi'' : E'' \to (E/F)''$.

Par la réflexivité de $E$, l'application $\pi'' \circ J_E : E \to (E/F)''$ est
surjective (car $J_E$ est surjective et $\pi''$ l'est aussi). Donc il existe
$x \in E$ tel que $\pi''(J_E(x)) = \xi$. Or pour $g \in (E/F)'$,
\[
  \pi''(J_E(x))(g) = J_E(x)(\pi'(g)) = (\pi'(g))(x) = g(\pi(x)) = J(\pi(x))(g).
\]
Donc $\xi = J(\pi(x))$, et $J$ est surjective.
\end{proof}

\begin{corollaire}
\label{cor:reflexif-dual}
\index{espace réflexif!dual réflexif}
Un espace de Banach $E$ est réflexif si et seulement si $E'$ est réflexif.
\end{corollaire}

\begin{proof}
($\Rightarrow$) Supposons $E$ réflexif. Soit $\Phi \in E'''$. Considérons
$\Phi \circ J_E^{-1} \in E'$ (bien défini car $J_E$ est un isomorphisme isométrique
de $E$ sur $E''$). Plus précisément, définissons $f \in E'$ par
$f(x) = \Phi(J_E(x))$ pour tout $x \in E$. Alors pour tout $\xi \in E''$,
il existe $x \in E$ tel que $\xi = J_E(x)$ (réflexivité), et
\[
  J_{E'}(f)(\xi) = \xi(f) = J_E(x)(f) = f(x) = \Phi(J_E(x)) = \Phi(\xi).
\]
Donc $J_{E'}(f) = \Phi$, et $E'$ est réflexif.

($\Leftarrow$) Si $E'$ est réflexif, alors $E'' = J_{E'}(E')$. Comme $J_E(E)$ est
un sous-espace fermé de $E''$ (car $E$ est complet), il suffit de montrer que
$J_E(E)$ est dense dans $E''$ pour la topologie de la norme. Supposons par l'absurde que
$J_E(E) \subsetneq E''$. Par Hahn-Banach, il existe $\Phi \in E''' \setminus \{0\}$
tel que $\Phi|_{J_E(E)} = 0$. Comme $E'$ est réflexif, $\Phi = J_{E'}(f)$ pour un
$f \in E'$. Alors pour tout $x \in E$,
\[
  0 = \Phi(J_E(x)) = J_{E'}(f)(J_E(x)) = J_E(x)(f) = f(x).
\]
Donc $f = 0$, d'où $\Phi = 0$, contradiction.
\end{proof}

\subsection{Le théorème de Milman-Pettis}
\label{subsec:milman-pettis}

\begin{theoreme}[Milman-Pettis]
\label{thm:milman-pettis}
\index{théorème!de Milman-Pettis}
\index{Milman, David}
\index{Pettis, Billy James}
\index{espace!uniformément convexe}
Tout espace de Banach uniformément convexe\index{convexité uniforme} est réflexif.
\end{theoreme}

Rappelons d'abord la définition.

\begin{definition}[Convexité uniforme]
\label{def:unif-convexe}
\index{convexité uniforme!définition}
Un espace de Banach $(E, \norm{\cdot})$ est \emph{uniformément convexe} si pour tout
$\varepsilon > 0$, il existe $\delta > 0$ tel que
\[
  \norm{x} \le 1, \; \norm{y} \le 1, \; \norm{x - y} \ge \varepsilon
  \implies \norm[\bigg]{\frac{x+y}{2}} \le 1 - \delta.
\]
\end{definition}

\begin{exemple}
\label{ex:lp-unif-convexe}
\index{lp@$\ell^p$!uniformément convexe}
Pour $1 < p < \infty$, l'espace $\ell^p$ (et plus généralement $L^p$) est
uniformément convexe (inégalités de Clarkson\index{inégalités!de Clarkson}).
Les espaces $\ell^1$, $\ell^\infty$, $c_0$ ne le sont pas.
\end{exemple}

\begin{proof}[Preuve du théorème de Milman-Pettis]
Soit $E$ un espace de Banach uniformément convexe. Par le théorème de James
(théorème~\ref{thm:james}), il suffit de montrer que la boule unité fermée
$\overline{B}_E$ est faiblement compacte. Nous allons plutôt montrer directement
que $J(E) = E''$.

Soit $\xi \in E''$ avec $\norm{\xi}_{E''} = 1$. Il suffit de trouver $x \in E$ tel
que $J(x) = \xi$. Par le théorème de Goldstine (théorème~\ref{thm:goldstine},
démontré au chapitre~8), $J(\overline{B}_E)$ est dense dans $\overline{B}_{E''}$
pour la topologie faible-*\index{topologie!faible-*} $\sigma(E'', E')$.
Donc il existe un filet $(x_\alpha)$ dans $\overline{B}_E$ tel que
$J(x_\alpha) \weakstarto \xi$ dans $(E'', \sigma(E'', E'))$.

Montrons que $(x_\alpha)$ possède un sous-filet convergent en norme, ce qui impliquera
$\xi \in J(E)$.

Comme $\norm{\xi} = 1$, il existe $f \in E'$ avec $\norm{f} = 1$ et
$\xi(f) = \norm{\xi} = 1$ (le dual d'un espace uniformément convexe est
\emph{strictement convexe}\index{convexité stricte}, et la norme y est lisse,
ce qui garantit l'existence d'un tel $f$; alternativement, on applique
le lemme de Bishop-Phelps\index{lemme!de Bishop-Phelps}).
Comme $J(x_\alpha)(f) = f(x_\alpha) \to \xi(f) = 1$ et $\norm{x_\alpha} \le 1$,
on obtient $\mathrm{Re}\, f(x_\alpha) \to 1$.

Montrons que $(x_\alpha)$ est de Cauchy en norme. Par l'absurde, supposons qu'il
existe $\varepsilon > 0$ et des sous-filets $x_{\alpha_i}, x_{\beta_i}$ avec
$\norm{x_{\alpha_i} - x_{\beta_i}} \ge \varepsilon$. Par convexité uniforme, il
existe $\delta > 0$ tel que $\norm{(x_{\alpha_i} + x_{\beta_i})/2} \le 1 - \delta$.
Mais alors
\[
  \mathrm{Re}\, f\!\left(\frac{x_{\alpha_i} + x_{\beta_i}}{2}\right) \le
  \norm{f} \norm[\bigg]{\frac{x_{\alpha_i} + x_{\beta_i}}{2}} \le 1 - \delta.
\]
Or le membre de gauche tend vers $1$, contradiction. Donc $(x_\alpha)$ est de Cauchy
dans $E$ (complet), donc convergente vers un $x \in E$.
Ainsi $J(x) = \xi$, et $E$ est réflexif.
\end{proof}

\begin{corollaire}
\label{cor:lp-reflexif-milman}
\index{lp@$\ell^p$!réflexif}
Pour $1 < p < \infty$, $\ell^p$ et $L^p(\mu)$ ($\mu$ $\sigma$-finie) sont réflexifs.
\end{corollaire}

\begin{proof}
Ces espaces sont uniformément convexes (inégalités de Clarkson\index{inégalités!de Clarkson}),
donc réflexifs par le théorème de Milman-Pettis.
\end{proof}

\index{espace réflexif|)}
\index{réflexivité|)}

%% ══════════════════════════════════════════════════════════════════
\section{Exercices}
\label{sec:exo-ch7}
%% ══════════════════════════════════════════════════════════════════

\begin{exercice}[$\star$]
\label{exo:dual-separe}
\index{exercice!dual sépare les points}
Soit $E$ un espace vectoriel normé et $x \in E$.
Montrer que $\norm{x} = \sup\{ \abs{f(x)} : f \in E',\, \norm{f} \le 1 \}$
et que le supremum est atteint.
\end{exercice}

\begin{exercice}[$\star$]
\label{exo:dual-sous-espace}
\index{exercice!dual d'un sous-espace}
Soit $F$ un sous-espace fermé de $E$. Montrer que $F' \cong E'/F^\perp$ isométriquement,
où $F^\perp = \{ f \in E' : f|_F = 0 \}$\index{annulateur}.
\end{exercice}

\begin{exercice}[$\star\star$]
\label{exo:reflexif-atteint}
\index{exercice!réflexif implique norme atteinte}
Soit $E$ un espace de Banach réflexif. Montrer que toute forme linéaire continue
$f \in E'$ atteint sa norme, c'est-à-dire qu'il existe $x_0 \in E$ avec
$\norm{x_0} = 1$ et $\abs{f(x_0)} = \norm{f}$.

\textit{Indication :} Considérer une suite maximisante et utiliser la compacité faible
de la boule unité (théorème de James/Banach-Alaoglu).
\end{exercice}

\begin{exercice}[$\star\star$]
\label{exo:c0-non-dual}
\index{exercice!$c_0$ n'est dual d'aucun espace}
Montrer que $c_0$ n'est isométriquement isomorphe au dual d'aucun espace de Banach.

\textit{Indication :} Si $c_0 \cong X'$, alors $X' \cong c_0$ serait séparable, donc
$X$ serait séparable, donc la boule unité de $X'$ serait faiblement-* séquentiellement compacte
(par Banach-Alaoglu et métrisabilité). Examiner les vecteurs canoniques dans $c_0$.
\end{exercice}

\begin{exercice}[$\star\star$]
\label{exo:reflexif-separable}
\index{exercice!réflexif et séparable}
Montrer qu'un espace de Banach réflexif $E$ est séparable si et seulement si $E'$
est séparable.
\end{exercice}

\begin{exercice}[$\star\star$]
\label{exo:lp-dual-calcul}
\index{exercice!calcul de normes duales}
Pour $p = 4/3$, calculer la norme de la forme linéaire
$f : \ell^{4/3} \to \R$ définie par
$f(x) = \sum_{n=1}^{\infty} x_n / n^3$.
Identifier $f$ comme élément de $\ell^4$ et calculer $\norm{f}_4$.
\end{exercice}

\begin{exercice}[$\star\star\star$]
\label{exo:l1-quotient-c0}
\index{exercice!quotient de $\ell^1$}
Soit $T : \ell^1 \to c_0$ un opérateur linéaire continu surjectif.
\begin{enumerate}[label=(\alph*)]
  \item Montrer qu'un tel opérateur existe (considérer $\ell^1 \to \ell^1/\Ker(T)$
    et utiliser le théorème de l'application ouverte).
  \item En déduire que $c_0$ est isomorphe à un quotient de $\ell^1$.
  \item Montrer que $\Ker(T)$ n'est pas complémenté dans $\ell^1$.
\end{enumerate}
\end{exercice}

\begin{exercice}[$\star\star$]
\label{exo:unif-convexe-strict}
\index{exercice!convexité uniforme}
Montrer que tout espace uniformément convexe est strictement convexe\index{convexité stricte}
(c'est-à-dire que si $\norm{x} = \norm{y} = 1$ et $x \neq y$, alors
$\norm{(x+y)/2} < 1$). La réciproque est-elle vraie ?
\end{exercice}

\begin{exercice}[$\star\star\star$]
\label{exo:james-l1}
\index{exercice!théorème de James appliqué à $\ell^1$}
Donner un exemple de forme linéaire continue $f \in (\ell^1)'$ qui n'atteint pas
sa norme. En déduire que $\ell^1$ n'est pas réflexif (via le théorème de James).

\textit{Indication :} Considérer $y = (1 - 1/n)_{n \ge 1} \in \ell^\infty$ et la forme
$\Phi_y(x) = \sum_n (1-1/n) x_n$.
\end{exercice}

\begin{exercice}[$\star\star$]
\label{exo:reflexif-bidual-quotient}
\index{exercice!bidual et réflexivité}
Soit $E$ un espace de Banach. Montrer que si $E''$ est séparable, alors $E$
est réflexif.

\textit{Indication :} Montrer que $E'/J(E)^\perp \cong E'$, donc $E'$ est séparable,
puis $E$ est séparable, et utiliser le fait que $J(E)$ est fermé et de codimension
nulle.
\end{exercice}

%%%%%%%%%%%%%%%%%%%%%%%%%%%%%%%%%%%%%%%%%%%%%%%%%%%%%%%%%%%%%%%%%%%%
%  Analyse Fonctionnelle — Chapitre 8
%  Topologies Faibles
%%%%%%%%%%%%%%%%%%%%%%%%%%%%%%%%%%%%%%%%%%%%%%%%%%%%%%%%%%%%%%%%%%%%

\chapter{Topologies Faibles}
\label{chap:topologies-faibles}
\minitoc

\vspace{1em}

La topologie de la norme d'un espace de Banach est souvent «~trop fine~» pour obtenir
des résultats de compacité utiles en analyse. Ce chapitre introduit les
\emph{topologies faibles}\index{topologie!faible|(}, qui fournissent des
notions de convergence plus souples et sont essentielles en théorie des
EDP\index{EDP}, en calcul des variations\index{calcul des variations},
et en optimisation\index{optimisation}. Les théorèmes centraux ---
Banach-Alaoglu\index{théorème!de Banach-Alaoglu},
Goldstine\index{théorème!de Goldstine},
Kakutani\index{théorème!de Kakutani},
Mazur\index{théorème!de Mazur} --- sont démontrés en détail.

%% ══════════════════════════════════════════════════════════════════
\section{La topologie faible $\sigma(E, E')$}
\label{sec:topo-faible}
%% ══════════════════════════════════════════════════════════════════

\index{topologie!faible!définition|(}

\begin{definition}[Topologie faible]
\label{def:topo-faible}
\index{topologie!faible}
\index{sigma(E,E')@$\sigma(E,E')$}
Soit $E$ un espace vectoriel normé. La \emph{topologie faible} sur $E$, notée
$\sigma(E, E')$, est la topologie la moins fine rendant continues toutes les formes
linéaires $f \in E'$. C'est la topologie initiale\index{topologie!initiale}
engendrée par la famille $(f)_{f \in E'}$.
\end{definition}

\begin{proposition}[Base de voisinages]
\label{prop:base-voisinages-faible}
\index{topologie!faible!base de voisinages}
Une base de voisinages de $x_0 \in E$ pour $\sigma(E, E')$ est constituée des
ensembles de la forme
\[
  V(x_0; f_1, \ldots, f_n; \varepsilon) = \bigl\{ x \in E :
  \abs{f_i(x) - f_i(x_0)} < \varepsilon, \; i = 1, \ldots, n \bigr\},
\]
où $n \ge 1$, $f_1, \ldots, f_n \in E'$, et $\varepsilon > 0$.
\end{proposition}

\begin{proof}
Par définition de la topologie initiale, les sous-ensembles ouverts de $\sigma(E,E')$
sont les réunions d'intersections finies d'ensembles de la forme
$f^{-1}(U)$ avec $f \in E'$ et $U$ ouvert dans $\K$. Les ensembles
$V(x_0; f_1, \ldots, f_n; \varepsilon)$ s'écrivent
$\bigcap_{i=1}^n f_i^{-1}(B(f_i(x_0), \varepsilon))$, qui sont bien des
intersections finies de pré-images d'ouverts. Réciproquement, toute intersection
finie $\bigcap_{i=1}^n f_i^{-1}(U_i)$ contenant $x_0$ contient un
$V(x_0; f_1, \ldots, f_n; \varepsilon)$ pour $\varepsilon$ assez petit.
\end{proof}

\begin{proposition}[La topologie faible est séparée]
\label{prop:faible-separee}
\index{topologie!faible!séparée}
\index{Hausdorff}
La topologie $\sigma(E,E')$ est séparée (Hausdorff).
\end{proposition}

\begin{proof}
Soient $x \neq y$ dans $E$. Par le corollaire~\ref{cor:dual-separe} (Hahn-Banach),
il existe $f \in E'$ tel que $f(x) \neq f(y)$. Posons
$\varepsilon = \abs{f(x) - f(y)}/3 > 0$. Alors
$V(x; f; \varepsilon) \cap V(y; f; \varepsilon) = \varnothing$,
ce qui sépare $x$ et $y$.
\end{proof}

\begin{proposition}[Comparaison avec la topologie forte]
\label{prop:faible-moins-fine}
\index{topologie!faible!comparaison avec la norme}
\index{topologie!forte}
La topologie $\sigma(E,E')$ est moins fine que la topologie de la norme. Si $E$ est
de dimension infinie, elle est strictement moins fine.
\end{proposition}

\begin{proof}
Tout ouvert faible est ouvert pour la norme, car chaque $f \in E'$ est continue
pour la norme. Réciproquement, si $\dim E = \infty$, la boule ouverte unité
$B_E = \{ x : \norm{x} < 1 \}$ n'est pas ouverte pour $\sigma(E,E')$.
En effet, tout voisinage faible de $0$ contient un sous-espace de codimension finie
$\bigcap_{i=1}^n \Ker(f_i)$, qui est de dimension infinie, donc non borné; il ne peut
donc pas être contenu dans $B_E$.
\end{proof}

%% ══════════════════════════════════════════════════════════════════
\section{Convergence faible}
\label{sec:conv-faible}
%% ══════════════════════════════════════════════════════════════════

\index{convergence faible|(}

\begin{definition}[Convergence faible]
\label{def:conv-faible}
\index{convergence faible!définition}
\index{weakto@$\weakto$}
Soit $(x_n)$ une suite dans $E$. On dit que $(x_n)$ \emph{converge faiblement}
vers $x \in E$, noté $x_n \weakto x$, si
\[
  f(x_n) \to f(x), \quad \forall f \in E'.
\]
\end{definition}

\begin{remarque}
\label{rem:conv-forte-implique-faible}
\index{convergence faible!versus convergence forte}
La convergence forte (en norme) implique la convergence faible, mais la réciproque
est fausse en dimension infinie. Par exemple, dans $\ell^2$, la base orthonormée
$(e_n)$ converge faiblement vers $0$ (car pour $y \in \ell^2$,
$\inner{y}{e_n} = y_n \to 0$), mais $\norm{e_n} = 1 \not\to 0$.
\end{remarque}

\begin{proposition}[Unicité de la limite faible]
\label{prop:unicite-limite-faible}
\index{convergence faible!unicité}
Si $x_n \weakto x$ et $x_n \weakto y$, alors $x = y$.
\end{proposition}

\begin{proof}
Pour tout $f \in E'$, $f(x) = \lim f(x_n) = f(y)$. Comme $E'$ sépare les points
(corollaire~\ref{cor:dual-separe}), $x = y$.
\end{proof}

\begin{theoreme}[Les suites faiblement convergentes sont bornées]
\label{thm:faible-borne}
\index{convergence faible!bornée}
\index{théorème!de Banach-Steinhaus}
Soit $E$ un espace de Banach. Si $x_n \weakto x$, alors la suite $(x_n)$ est
bornée et $\norm{x} \le \liminf_{n \to \infty} \norm{x_n}$.
\end{theoreme}

\begin{proof}
Considérons la famille $(J(x_n))_{n \ge 1}$ dans $E''$. Pour tout $f \in E'$,
\[
  \sup_{n \ge 1} \abs{J(x_n)(f)} = \sup_{n \ge 1} \abs{f(x_n)} < \infty,
\]
car $(f(x_n))$ converge et est donc bornée. Par le théorème de Banach-Steinhaus\index{théorème!de Banach-Steinhaus}
(principe de borne uniforme\index{principe de borne uniforme}) appliqué à la famille
$(J(x_n))_{n \ge 1} \subset E'' = \mathcal{L}(E', \K)$
et à l'espace de Banach $E'$,
\[
  \sup_{n \ge 1} \norm{J(x_n)}_{E''} < \infty.
\]
Comme $J$ est isométrique, $\sup_{n \ge 1} \norm{x_n} < \infty$.

Pour la seconde inégalité, soit $f \in E'$ avec $\norm{f} \le 1$. Alors
$\abs{f(x)} = \lim_n \abs{f(x_n)} \le \liminf_n \norm{x_n}$.
En prenant le sup sur $f$, $\norm{x} \le \liminf_n \norm{x_n}$.
\end{proof}

\begin{remarque}
\label{rem:norme-sci-faible}
\index{norme!semi-continuité inférieure faible}
La seconde inégalité signifie que la norme est \emph{séquentiellement semi-continue
inférieurement}\index{semi-continuité inférieure} pour la topologie faible.
En fait, elle est s.c.i.\ pour $\sigma(E,E')$ (pas seulement séquentiellement).
\end{remarque}

\begin{exemple}[Convergence faible dans $\ell^p$]
\label{ex:conv-faible-lp}
\index{convergence faible!dans $\ell^p$}
Soit $1 < p < \infty$. Alors $x_n \weakto x$ dans $\ell^p$ si et seulement si
\begin{enumerate}[label=(\roman*)]
  \item $\sup_n \norm{x_n}_p < \infty$, et
  \item pour tout $k \ge 1$, $x_n^{(k)} \to x^{(k)}$ (convergence coordonnée par coordonnée).
\end{enumerate}
En effet, les $e_k^*$ (évaluation de la $k$-ième coordonnée) sont dans $(\ell^p)' \cong \ell^q$,
et ces formes linéaires engendrent un sous-espace dense de $\ell^q$.
\end{exemple}

\begin{exemple}[Convergence faible dans $L^p$]
\label{ex:conv-faible-Lp}
\index{convergence faible!dans $L^p$}
Soit $1 < p < \infty$ et $\mu$ une mesure $\sigma$-finie. Alors $f_n \weakto f$ dans
$L^p(\mu)$ si et seulement si $\sup_n \norm{f_n}_p < \infty$ et
$\int_A f_n \, \dd\mu \to \int_A f \, \dd\mu$ pour tout ensemble mesurable $A$ de
mesure finie.
\end{exemple}

\index{convergence faible|)}
\index{topologie!faible!définition|)}

%% ══════════════════════════════════════════════════════════════════
\section{La topologie faible-* $\sigma(E', E)$}
\label{sec:topo-faible-etoile}
%% ══════════════════════════════════════════════════════════════════

\index{topologie!faible-*|(}

\begin{definition}[Topologie faible-*]
\label{def:topo-faible-etoile}
\index{topologie!faible-*!définition}
\index{sigma(E',E)@$\sigma(E',E)$}
Soit $E$ un espace de Banach. La \emph{topologie faible-*} sur $E'$, notée
$\sigma(E', E)$, est la topologie la moins fine rendant continues toutes les applications
d'évaluation $\hat{x} : E' \to \K$, $\hat{x}(f) = f(x)$, pour $x \in E$.
\end{definition}

\begin{proposition}[Base de voisinages faible-*]
\label{prop:base-voisinages-faible-etoile}
\index{topologie!faible-*!base de voisinages}
Une base de voisinages de $f_0 \in E'$ pour $\sigma(E', E)$ est constituée des
ensembles
\[
  V^*(f_0; x_1, \ldots, x_n; \varepsilon) = \bigl\{ f \in E' :
  \abs{f(x_i) - f_0(x_i)} < \varepsilon, \; i = 1, \ldots, n \bigr\},
\]
où $n \ge 1$, $x_1, \ldots, x_n \in E$, et $\varepsilon > 0$.
\end{proposition}

\begin{definition}[Convergence faible-*]
\label{def:conv-faible-etoile}
\index{convergence faible-*}
\index{weakstarto@$\weakstarto$}
On dit que $(f_n)$ \emph{converge faible-*} vers $f$ dans $E'$, noté
$f_n \weakstarto f$, si
\[
  f_n(x) \to f(x), \quad \forall x \in E.
\]
\end{definition}

\begin{proposition}[Relations entre topologies]
\label{prop:relations-topologies}
\index{topologie!faible-*!comparaison}
Sur $E'$, on a les inclusions
\[
  \sigma(E', E) \subset \sigma(E', E'') \subset \tau_{\text{norme}},
\]
où $\sigma(E', E'')$ est la topologie faible de $E'$ (vu comme espace de Banach)
et $\tau_{\text{norme}}$ la topologie de la norme. Si $E$ est réflexif, les deux
premières coïncident : $\sigma(E', E) = \sigma(E', E'')$.
\end{proposition}

\begin{proof}
L'inclusion $\sigma(E',E) \subset \sigma(E',E'')$ est claire : les évaluations par les
éléments de $J(E) \subset E''$ forment un sous-ensemble de toutes les évaluations par
$E''$, et les évaluations par $E''$ définissent $\sigma(E',E'')$. Si $E$ est réflexif,
$J(E) = E''$, d'où l'égalité.
\end{proof}

\begin{remarque}
\label{rem:faible-etoile-vs-faible}
\index{topologie!faible-*!versus topologie faible}
La topologie $\sigma(E',E)$ est en général strictement moins fine que $\sigma(E',E'')$.
C'est précisément lorsque $E$ n'est pas réflexif qu'elles diffèrent.
\end{remarque}

%% ══════════════════════════════════════════════════════════════════
\section{Le théorème de Banach-Alaoglu}
\label{sec:banach-alaoglu}
%% ══════════════════════════════════════════════════════════════════

\index{théorème!de Banach-Alaoglu|(}

Le théorème de Banach-Alaoglu est l'un des résultats les plus importants de l'analyse
fonctionnelle. Il fournit la compacité de la boule unité du dual pour la topologie
faible-*, compensant l'absence de compacité pour la topologie forte en dimension infinie.

\begin{theoreme}[Banach-Alaoglu]
\label{thm:banach-alaoglu}
\index{théorème!de Banach-Alaoglu}
\index{Banach, Stefan}
\index{Alaoglu, Leonidas}
Soit $E$ un espace vectoriel normé. La boule unité fermée du dual
\[
  \overline{B}_{E'} = \{ f \in E' : \norm{f} \le 1 \}
\]
est compacte pour la topologie faible-* $\sigma(E', E)$.
\end{theoreme}

\begin{proof}
Pour chaque $x \in E$, posons $D_x = \{ \lambda \in \K : \abs{\lambda} \le \norm{x} \}$.
C'est un compact de $\K$. Considérons l'espace produit\index{espace produit}
\[
  P = \prod_{x \in E} D_x,
\]
muni de la topologie produit\index{topologie!produit}. Par le théorème de
Tychonoff\index{théorème!de Tychonoff}, $P$ est compact.

Un élément de $P$ est une application $\varphi : E \to \K$ telle que
$\abs{\varphi(x)} \le \norm{x}$ pour tout $x \in E$. Considérons l'application
\[
  \Psi : \overline{B}_{E'} \longrightarrow P, \quad f \longmapsto (f(x))_{x \in E}.
\]
Cette application est bien définie car si $\norm{f} \le 1$, alors
$\abs{f(x)} \le \norm{f} \norm{x} \le \norm{x}$, donc $f(x) \in D_x$.

\medskip
\noindent\textbf{Étape 1 : $\Psi$ est un homéomorphisme sur son image.}

L'application $\Psi$ est injective : si $\Psi(f) = \Psi(g)$, alors $f(x) = g(x)$
pour tout $x$, donc $f = g$.

La topologie produit sur $P$ est la topologie de la convergence simple\index{convergence!simple}
(chaque projection $\pi_x : P \to D_x$ est continue). Or la topologie $\sigma(E',E)$
est aussi la topologie de la convergence simple sur $E'$ restreinte à $E$. Donc
$\Psi$ est continue et son inverse $\Psi^{-1} : \Psi(\overline{B}_{E'}) \to \overline{B}_{E'}$
est continue. Ainsi $\Psi$ est un homéomorphisme de $(\overline{B}_{E'}, \sigma(E',E))$
sur $\Psi(\overline{B}_{E'}) \subset P$.

\medskip
\noindent\textbf{Étape 2 : $\Psi(\overline{B}_{E'})$ est fermé dans $P$.}

Soit $(\varphi_\alpha)$ un filet dans $\Psi(\overline{B}_{E'})$ convergeant vers
$\varphi \in P$. Pour chaque $\alpha$, soit $f_\alpha \in \overline{B}_{E'}$ tel que
$\varphi_\alpha = \Psi(f_\alpha)$, c'est-à-dire $\varphi_\alpha(x) = f_\alpha(x)$
pour tout $x \in E$.

La convergence dans $P$ signifie $f_\alpha(x) \to \varphi(x)$ pour tout $x \in E$.
Montrons que $\varphi$ est linéaire. Pour $x, y \in E$ et $\lambda \in \K$,
\begin{align*}
  \varphi(\lambda x + y) &= \lim_\alpha f_\alpha(\lambda x + y)
  = \lim_\alpha (\lambda f_\alpha(x) + f_\alpha(y)) \\
  &= \lambda \varphi(x) + \varphi(y).
\end{align*}
De plus, $\abs{\varphi(x)} \le \norm{x}$ pour tout $x$ (car $\varphi \in P$), donc
$\varphi$ est une forme linéaire continue de norme $\le 1$. Ainsi
$\varphi \in \Psi(\overline{B}_{E'})$.

\medskip
\noindent\textbf{Conclusion.}
Comme $\Psi(\overline{B}_{E'})$ est un fermé d'un compact ($P$), il est compact.
Par l'homéomorphisme $\Psi$, $\overline{B}_{E'}$ est compact pour $\sigma(E',E)$.
\end{proof}

\begin{remarque}
\label{rem:alaoglu-separable}
\index{théorème!de Banach-Alaoglu!cas séparable}
\index{métrisabilité!de la topologie faible-*}
Si $E$ est séparable, alors $(\overline{B}_{E'}, \sigma(E',E))$ est métrisable.
En effet, soit $(x_n)$ une suite dense dans $\overline{B}_E$. La métrique
\[
  d(f,g) = \sum_{n=1}^{\infty} 2^{-n} \abs{f(x_n) - g(x_n)}
\]
engendre la topologie $\sigma(E',E)$ sur $\overline{B}_{E'}$. Dans ce cas,
la compacité séquentielle et la compacité coïncident, et on peut remplacer
Tychonoff par un argument diagonal\index{argument diagonal} plus élémentaire
(théorème de Banach-Alaoglu séquentiel\index{théorème!de Banach-Alaoglu!séquentiel},
aussi appelé théorème de Banach\index{théorème!de Banach (Alaoglu séquentiel)}).
\end{remarque}

\index{théorème!de Banach-Alaoglu|)}

%% ══════════════════════════════════════════════════════════════════
\section{Le théorème de Goldstine}
\label{sec:goldstine}
%% ══════════════════════════════════════════════════════════════════

\index{théorème!de Goldstine|(}

\begin{theoreme}[Goldstine]
\label{thm:goldstine}
\index{Goldstine, Herman Heine}
Soit $E$ un espace de Banach et $J : E \to E''$ l'injection canonique. Alors
$J(\overline{B}_E)$ est dense dans $\overline{B}_{E''}$ pour la topologie
faible-* $\sigma(E'', E')$.
\end{theoreme}

\begin{proof}
Notons $K = \overline{J(\overline{B}_E)}^{\sigma(E'',E')}$ l'adhérence faible-*
de $J(\overline{B}_E)$ dans $E''$. Comme $J(\overline{B}_E) \subset \overline{B}_{E''}$
et que $\overline{B}_{E''}$ est faible-* compact (Banach-Alaoglu appliqué à $E'$)
donc faible-* fermé, on a $K \subset \overline{B}_{E''}$.

Il faut montrer $K = \overline{B}_{E''}$. Raisonnons par contraposée : supposons
qu'il existe $\xi_0 \in \overline{B}_{E''} \setminus K$. Comme $K$ est un convexe
faible-* fermé\index{convexe!fermé faible-*} (adhérence d'un convexe, puisque
$J(\overline{B}_E)$ est convexe par linéarité de $J$), le théorème de séparation
de Hahn-Banach\index{théorème!de Hahn-Banach!forme géométrique} dans l'espace
$(E'', \sigma(E'',E'))$ permet de séparer $\xi_0$ de $K$ par une forme linéaire
continue pour $\sigma(E'',E')$.

Or les formes linéaires continues sur $(E'', \sigma(E'',E'))$ sont exactement les
éléments de $E'$ (vus comme évaluations). Il existe donc $f \in E'$ et $\alpha \in \R$
tels que
\[
  \mathrm{Re}\, \xi(f) \le \alpha < \mathrm{Re}\, \xi_0(f), \quad \forall \xi \in K.
\]
En particulier, pour tout $x \in \overline{B}_E$,
\[
  \mathrm{Re}\, f(x) = \mathrm{Re}\, J(x)(f) \le \alpha.
\]
Or $\sup_{\norm{x} \le 1} \mathrm{Re}\, f(x) = \norm{f}$
(en choisissant $x$ convenablement). Donc $\norm{f} \le \alpha$. Mais alors
\[
  \mathrm{Re}\, \xi_0(f) > \alpha \ge \norm{f},
\]
ce qui contredit $\abs{\xi_0(f)} \le \norm{\xi_0} \norm{f} \le \norm{f}$
(car $\norm{\xi_0} \le 1$). Cette contradiction achève la preuve.
\end{proof}

\begin{corollaire}
\label{cor:goldstine-dense}
\index{injection canonique!image dense}
$J(E)$ est dense dans $E''$ pour $\sigma(E'', E')$. En particulier, $E$ est réflexif
si et seulement si $\overline{B}_E$ est faiblement compacte.
\end{corollaire}

\begin{proof}
La densité de $J(E)$ suit de Goldstine par homogénéité (appliquer Goldstine à
$r \overline{B}_E$ pour tout $r > 0$). Pour la seconde assertion :

($\Rightarrow$) Si $E$ est réflexif, $J(\overline{B}_E) = \overline{B}_{E''}$ est
faible-* compact dans $E''$. Via $J^{-1}$, $\overline{B}_E$ est faiblement compacte
dans $E$ (car $\sigma(E,E') = J^{-1}(\sigma(E'',E')|_{J(E)})$ quand $J$ est surjectif).

($\Leftarrow$) Si $\overline{B}_E$ est faiblement compacte, alors $J(\overline{B}_E)$
est faible-* compact dans $E''$ (car $J$ est continue de $\sigma(E,E')$ vers
$\sigma(E'',E')$), donc faible-* fermé. Par Goldstine, $J(\overline{B}_E)$ est
faible-* dense dans $\overline{B}_{E''}$. Fermé et dense implique
$J(\overline{B}_E) = \overline{B}_{E''}$, donc $J$ est surjective.
\end{proof}

\index{théorème!de Goldstine|)}

%% ══════════════════════════════════════════════════════════════════
\section{Le théorème de Kakutani}
\label{sec:kakutani}
%% ══════════════════════════════════════════════════════════════════

\index{théorème!de Kakutani|(}

\begin{theoreme}[Kakutani]
\label{thm:kakutani}
\index{Kakutani, Shizuo}
\index{espace réflexif!caractérisation de Kakutani}
Soit $E$ un espace de Banach. Les assertions suivantes sont équivalentes :
\begin{enumerate}[label=(\roman*)]
  \item $E$ est réflexif.
  \item La boule unité fermée $\overline{B}_E$ est compacte pour $\sigma(E,E')$.
\end{enumerate}
\end{theoreme}

\begin{proof}
C'est exactement le contenu du corollaire~\ref{cor:goldstine-dense}. Redonnons
l'argument pour la commodité du lecteur.

\medskip
\noindent (i) $\Rightarrow$ (ii). Si $E$ est réflexif, $J : E \to E''$ est un
isomorphisme isométrique. Par Banach-Alaoglu, $\overline{B}_{E''}$ est
$\sigma(E'',E')$-compacte. Comme $J(\overline{B}_E) = \overline{B}_{E''}$,
et que $J$ est un homéomorphisme de $(E, \sigma(E,E'))$ sur $(J(E), \sigma(E'',E')|_{J(E)})$
(car $f \circ J^{-1}$ parcourt $E'$ quand $f$ le fait, via la surjectivité de la
transposée), $\overline{B}_E$ est $\sigma(E,E')$-compacte.

\medskip
\noindent (ii) $\Rightarrow$ (i). Supposons $\overline{B}_E$ faiblement compacte.
L'application $J$ est continue de $(E, \sigma(E,E'))$ vers $(E'', \sigma(E'',E'))$
(car pour $f \in E'$, la composée $\xi \mapsto \xi(f)$ avec $J$ donne $x \mapsto f(x)$,
qui est $\sigma(E,E')$-continue). L'image continue d'un compact est compacte, donc
$J(\overline{B}_E)$ est $\sigma(E'',E')$-compacte, en particulier fermée. Par le
théorème de Goldstine, $J(\overline{B}_E)$ est $\sigma(E'',E')$-dense dans
$\overline{B}_{E''}$. Un ensemble à la fois dense et fermé est égal à l'espace entier :
$J(\overline{B}_E) = \overline{B}_{E''}$. Donc $J$ est surjective et $E$ est réflexif.
\end{proof}

\index{théorème!de Kakutani|)}

%% ══════════════════════════════════════════════════════════════════
\section{Le théorème de Mazur}
\label{sec:mazur}
%% ══════════════════════════════════════════════════════════════════

\index{théorème!de Mazur|(}

\begin{theoreme}[Mazur]
\label{thm:mazur}
\index{Mazur, Stanisław}
\index{convexe!fermé = faiblement fermé}
Soit $E$ un espace vectoriel normé et $C \subset E$ un ensemble convexe.
Alors $C$ est fermé pour la topologie de la norme si et seulement si $C$ est fermé
pour la topologie faible $\sigma(E, E')$.
\end{theoreme}

\begin{proof}
\noindent($\Leftarrow$) Tout faiblement fermé est fortement fermé, car la topologie
faible est moins fine que la topologie forte (proposition~\ref{prop:faible-moins-fine}).

\medskip
\noindent($\Rightarrow$) Supposons $C$ convexe et fermé (pour la norme). Montrons
que $C$ est faiblement fermé, c'est-à-dire que $E \setminus C$ est faiblement ouvert.

Soit $x_0 \in E \setminus C$. Comme $C$ est convexe et fermé, le théorème de séparation
de Hahn-Banach\index{théorème!de Hahn-Banach!forme géométrique (séparation stricte)}
(sous sa forme géométrique, séparation stricte d'un point et d'un convexe fermé)
fournit $f \in E'$ et $\alpha \in \R$ tels que
\[
  \mathrm{Re}\, f(x_0) > \alpha \ge \sup_{x \in C} \mathrm{Re}\, f(x).
\]
L'ensemble $U = \{ x \in E : \mathrm{Re}\, f(x) > \alpha \}$ est un ouvert faible
(pré-image d'un ouvert par $\mathrm{Re} \circ f$, qui est $\sigma(E,E')$-continue).
De plus, $x_0 \in U$ et $U \cap C = \varnothing$. Donc $x_0$ n'est pas dans
l'adhérence faible de $C$.

Comme $x_0 \in E \setminus C$ est arbitraire, l'adhérence faible de $C$ est contenue
dans $C$, d'où $C$ est faiblement fermé.
\end{proof}

\begin{corollaire}[Mazur pour les suites]
\label{cor:mazur-suites}
\index{théorème!de Mazur!pour les suites}
Soit $(x_n)$ une suite dans un espace de Banach $E$ telle que $x_n \weakto x$.
Alors $x$ appartient à l'adhérence forte de l'enveloppe convexe\index{enveloppe convexe}
de $\{x_n : n \ge 1\}$ :
\[
  x \in \overline{\conv}\{x_n : n \ge 1\}.
\]
Autrement dit, il existe une suite de combinaisons convexes des $x_n$ qui converge
fortement vers $x$.
\end{corollaire}

\begin{proof}
Posons $C = \overline{\conv}\{x_n : n \ge 1\}$. C'est un convexe fermé, donc faiblement
fermé par le théorème de Mazur. Comme $x_n \in C$ pour tout $n$ et $x_n \weakto x$,
la limite faible $x$ appartient à l'adhérence faible de $C$, qui est $C$ lui-même.
\end{proof}

\begin{remarque}
\label{rem:mazur-constructif}
\index{théorème!de Mazur!aspect constructif}
Le corollaire de Mazur est souvent utilisé de la manière suivante : si $(x_n)$ est une
suite faiblement convergente dans un espace de Banach, on peut en extraire des
combinaisons convexes qui convergent fortement. Cette observation est cruciale en
théorie des EDP\index{EDP} et en calcul des variations\index{calcul des variations}.
\end{remarque}

\index{théorème!de Mazur|)}

%% ══════════════════════════════════════════════════════════════════
\section{Application : minimiseurs en calcul des variations}
\label{sec:calcul-variations}
%% ══════════════════════════════════════════════════════════════════

\index{calcul des variations|(}
\index{méthode directe|(}

Les topologies faibles trouvent une application majeure dans l'existence de minimiseurs
pour des problèmes variationnels. La \emph{méthode directe}\index{méthode directe!du calcul des variations}
du calcul des variations, due à Tonelli\index{Tonelli, Leonida}, repose sur le schéma suivant.

\begin{theoreme}[Méthode directe du calcul des variations]
\label{thm:methode-directe}
\index{méthode directe!théorème}
Soit $E$ un espace de Banach réflexif et $\Phi : E \to \R \cup \{+\infty\}$ une
fonctionnelle vérifiant :
\begin{enumerate}[label=(\roman*)]
  \item $\Phi$ est \emph{coercive}\index{coercivité} :
    $\Phi(x) \to +\infty$ quand $\norm{x} \to \infty$.
  \item $\Phi$ est \emph{séquentiellement faiblement semi-continue inférieurement}\index{semi-continuité inférieure!faible}
    (s.c.i.\ faible) :
    si $x_n \weakto x$, alors $\Phi(x) \le \liminf_{n \to \infty} \Phi(x_n)$.
  \item $\Phi \not\equiv +\infty$ (le domaine effectif est non vide).
\end{enumerate}
Alors $\Phi$ atteint son minimum : il existe $x_0 \in E$ tel que
$\Phi(x_0) = \inf_{x \in E} \Phi(x)$.
\end{theoreme}

\begin{proof}
Posons $m = \inf_{x \in E} \Phi(x) \in [-\infty, +\infty[$. Comme $\Phi \not\equiv +\infty$,
$m < +\infty$.

\medskip
\noindent\textbf{Étape 1 : Construction d'une suite minimisante.}
Il existe une suite $(x_n)$ dans $E$ telle que $\Phi(x_n) \to m$.

\medskip
\noindent\textbf{Étape 2 : Bornitude de la suite.}
Par coercivité, si $\norm{x_n} \to \infty$ (le long d'une sous-suite), alors
$\Phi(x_n) \to +\infty$, ce qui contredit $\Phi(x_n) \to m < +\infty$.
Donc $\sup_n \norm{x_n} < \infty$.

\medskip
\noindent\textbf{Étape 3 : Extraction d'une sous-suite faiblement convergente.}
Comme $E$ est réflexif, la boule fermée $\overline{B}(0, R)$ (avec
$R = \sup_n \norm{x_n}$) est faiblement compacte (Kakutani, théorème~\ref{thm:kakutani}).
Par le théorème d'Eberlein--\v{S}mulian (théorème~\ref{thm:eberlein-smulian}),
la compacité faible séquentielle permet d'extraire une sous-suite
$(x_{n_k})$ telle que $x_{n_k} \weakto x_0$ pour un $x_0 \in E$.

\medskip
\noindent\textbf{Étape 4 : Passage à la limite.}
Par semi-continuité inférieure faible,
\[
  \Phi(x_0) \le \liminf_{k \to \infty} \Phi(x_{n_k}) = m.
\]
Comme $\Phi(x_0) \ge m$ par définition de l'infimum, on conclut $\Phi(x_0) = m$.
En particulier, $m > -\infty$.
\end{proof}

\begin{exemple}[Minimisation d'une fonctionnelle intégrale]
\label{ex:minimisation-integrale}
\index{fonctionnelle intégrale}
Soit $\Omega \subset \R^d$ un ouvert borné et $p > 1$. Considérons l'espace de
Sobolev\index{espace de Sobolev} $W_0^{1,p}(\Omega)$ (qui est réflexif) et la
fonctionnelle
\[
  \Phi(u) = \frac{1}{p} \int_\Omega \abs{\nabla u}^p \, \dd x - \int_\Omega f \, u \, \dd x,
\]
où $f \in L^{p'}(\Omega)$. Alors :
\begin{itemize}
  \item $\Phi$ est coercive par l'inégalité de Poincaré\index{inégalité!de Poincaré}.
  \item La norme $\norm{\nabla \cdot}_p$ est convexe continue, donc faiblement
    s.c.i.\ (par Mazur). Le terme linéaire $u \mapsto \int f u$ est faiblement continu.
    Donc $\Phi$ est faiblement s.c.i.
  \item $\Phi(0) = 0$, donc le domaine est non vide.
\end{itemize}
Par le théorème~\ref{thm:methode-directe}, $\Phi$ admet un minimiseur $u_0 \in W_0^{1,p}(\Omega)$,
solution faible de l'équation du $p$-laplacien\index{p-laplacien@$p$-laplacien}
$-\Delta_p u = f$.
\end{exemple}

\begin{proposition}[Fonctionnelle convexe s.c.i.\ forte $\Rightarrow$ s.c.i.\ faible]
\label{prop:convexe-sci-forte-faible}
\index{semi-continuité inférieure!convexe forte implique faible}
Soit $E$ un espace de Banach et $\Phi : E \to \R \cup \{+\infty\}$ convexe et
semi-continue inférieurement pour la norme. Alors $\Phi$ est séquentiellement
faiblement s.c.i.
\end{proposition}

\begin{proof}
Pour tout $\alpha \in \R$, l'ensemble de sous-niveau
$C_\alpha = \{x \in E : \Phi(x) \le \alpha\}$ est convexe (car $\Phi$ est convexe)
et fermé (car $\Phi$ est s.c.i.\ forte). Par le théorème de Mazur,
$C_\alpha$ est faiblement fermé. Donc $\Phi$ est faiblement s.c.i.\ (car une fonction
est s.c.i.\ si et seulement si tous ses sous-niveaux sont fermés).

Pour la version séquentielle : si $x_n \weakto x$ et $\liminf \Phi(x_n) < \alpha$,
alors pour une sous-suite $x_{n_k} \weakto x$ avec $\Phi(x_{n_k}) \le \alpha$
(pour $k$ assez grand), et $x_{n_k} \in C_\alpha$. Comme $C_\alpha$ est faiblement
fermé, $x \in C_\alpha$, donc $\Phi(x) \le \alpha$. En faisant tendre $\alpha$
vers $\liminf \Phi(x_n)$, on obtient $\Phi(x) \le \liminf \Phi(x_n)$.
\end{proof}

\index{calcul des variations|)}
\index{méthode directe|)}

%% ══════════════════════════════════════════════════════════════════
\section{Compléments sur les topologies faibles}
\label{sec:complements-faibles}
%% ══════════════════════════════════════════════════════════════════

\begin{proposition}[Métrisabilité de la topologie faible sur les bornés]
\label{prop:metrisabilite-faible-borne}
\index{métrisabilité!de la topologie faible}
\index{topologie!faible!métrisabilité}
Si $E'$ est séparable, alors la topologie $\sigma(E,E')$ est métrisable sur les
parties bornées de $E$.
\end{proposition}

\begin{proof}
Soit $(f_n)_{n \ge 1}$ une suite dense dans $\overline{B}_{E'}$. Pour $R > 0$,
sur $R \overline{B}_E$, la métrique
\[
  d(x,y) = \sum_{n=1}^{\infty} 2^{-n} \abs{f_n(x) - f_n(y)}
\]
engendre la topologie $\sigma(E,E')$ restreinte à $R \overline{B}_E$.
En effet, $d(x_k, x) \to 0$ si et seulement si $f_n(x_k) \to f_n(x)$ pour tout $n$,
ce qui équivaut (par densité de $(f_n)$ et bornitude de $(x_k)$) à $f(x_k) \to f(x)$
pour tout $f \in E'$.
\end{proof}

\begin{proposition}[Dual faible]
\label{prop:dual-faible}
\index{topologie!faible!formes linéaires continues}
Les formes linéaires $\sigma(E,E')$-continues sur $E$ sont exactement les éléments
de $E'$. Autrement dit, $(E, \sigma(E,E'))' = E'$.
\end{proposition}

\begin{proof}
Toute $f \in E'$ est $\sigma(E,E')$-continue par définition. Réciproquement, soit
$\varphi : E \to \K$ linéaire et $\sigma(E,E')$-continue. La continuité en $0$
donne : il existe $f_1, \ldots, f_n \in E'$ et $\varepsilon > 0$ tels que
$\abs{f_i(x)} < \varepsilon$ pour $i = 1, \ldots, n$ implique $\abs{\varphi(x)} < 1$.
Par homogénéité, $\abs{\varphi(x)} \le \varepsilon^{-1} \max_i \abs{f_i(x)}$ pour tout $x$.
En particulier, $\bigcap_{i=1}^n \Ker(f_i) \subset \Ker(\varphi)$.

Par un résultat classique d'algèbre linéaire, il existe $\lambda_1, \ldots, \lambda_n \in \K$
tels que $\varphi = \sum_{i=1}^n \lambda_i f_i \in E'$.
\end{proof}

\begin{proposition}[Topologie faible-* et dual]
\label{prop:dual-faible-etoile}
\index{topologie!faible-*!formes linéaires continues}
Les formes linéaires $\sigma(E',E)$-continues sur $E'$ sont exactement les
évaluations $\hat{x}$ pour $x \in E$. Autrement dit, $(E', \sigma(E',E))' = J(E) \cong E$.
\end{proposition}

\begin{proof}
Argument identique au précédent, en remplaçant $E$ par $E'$ et $E'$ par $\{\hat{x} : x \in E\}$.
\end{proof}

%% ══════════════════════════════════════════════════════════════════
\section{Exercices}
\label{sec:exo-ch8}
%% ══════════════════════════════════════════════════════════════════

\begin{exercice}[$\star$]
\label{exo:faible-base}
\index{exercice!base de voisinages faibles}
Montrer que tout voisinage faible de $0$ dans un espace de dimension infinie contient
un sous-espace vectoriel de dimension infinie (donc les voisinages faibles de $0$
ne sont jamais bornés en dimension infinie).
\end{exercice}

\begin{exercice}[$\star$]
\label{exo:faible-suite-bornee}
\index{exercice!suite faiblement convergente bornée}
Soit $E$ un espace de Banach, $(x_n)$ une suite bornée de $E$, et $D \subset E'$
un sous-ensemble dense. Montrer que si $f(x_n) \to f(x)$ pour tout $f \in D$,
alors $x_n \weakto x$.
\end{exercice}

\begin{exercice}[$\star\star$]
\label{exo:faible-fermee}
\index{exercice!boule non faiblement fermée}
Montrer que dans un espace de Banach de dimension infinie, la sphère unité
$S_E = \{x \in E : \norm{x} = 1\}$ n'est pas faiblement fermée.

\textit{Indication :} Montrer que $0$ est dans l'adhérence faible de $S_E$.
\end{exercice}

\begin{exercice}[$\star\star$]
\label{exo:alaoglu-application}
\index{exercice!application de Banach-Alaoglu}
Soit $E$ un espace de Banach séparable et $(f_n)$ une suite bornée dans $E'$.
Montrer qu'il existe une sous-suite $(f_{n_k})$ et $f \in E'$ tels que
$f_{n_k} \weakstarto f$.
\end{exercice}

\begin{exercice}[$\star\star$]
\label{exo:faible-etoile-pas-faible}
\index{exercice!faible-* versus faible}
Soit $E = \ell^1$, de sorte que $E' = \ell^\infty$.
\begin{enumerate}[label=(\alph*)]
  \item Montrer que la suite $(e_n^*)$ des vecteurs canoniques de $\ell^\infty$ converge
    faible-* vers $0$, mais ne converge pas faiblement vers $0$ dans $\ell^\infty$.
  \item En déduire que $\sigma(\ell^\infty, \ell^1) \subsetneq \sigma(\ell^\infty, (\ell^\infty)')$.
\end{enumerate}
\end{exercice}

\begin{exercice}[$\star\star$]
\label{exo:mazur-application}
\index{exercice!application de Mazur}
Soit $(x_n)$ une suite dans un espace de Banach $E$ telle que $x_n \weakto x$.
Montrer qu'il existe une suite $(y_N)$ de combinaisons convexes
$y_N = \sum_{k=1}^{n_N} \lambda_k^{(N)} x_k$ (avec $\lambda_k^{(N)} \ge 0$ et
$\sum_k \lambda_k^{(N)} = 1$) telle que $\norm{y_N - x} \to 0$.
\end{exercice}

\begin{exercice}[$\star\star$]
\label{exo:sci-faible-norme}
\index{exercice!semi-continuité inférieure de la norme}
Soit $E$ un espace de Banach et $x_n \weakto x$.
\begin{enumerate}[label=(\alph*)]
  \item Montrer que $\norm{x} \le \liminf_{n \to \infty} \norm{x_n}$.
  \item Donner un exemple où l'inégalité est stricte.
  \item Montrer que si $\norm{x_n} \to \norm{x}$ et si $E$ est uniformément convexe,
    alors $x_n \to x$ en norme.
\end{enumerate}
\end{exercice}

\begin{exercice}[$\star\star\star$]
\label{exo:goldstine-application}
\index{exercice!application de Goldstine}
Soit $E$ un espace de Banach non réflexif.
\begin{enumerate}[label=(\alph*)]
  \item Montrer qu'il existe $\xi \in E'' \setminus J(E)$ avec $\norm{\xi} = 1$.
  \item En utilisant Goldstine, construire un filet $(x_\alpha)$ dans $\overline{B}_E$
    tel que $J(x_\alpha) \weakstarto \xi$ mais $(x_\alpha)$ n'a pas de sous-filet
    convergent en norme vers un élément de $E$.
  \item En déduire que $\overline{B}_E$ n'est pas faiblement compacte.
\end{enumerate}
\end{exercice}

\begin{exercice}[$\star\star\star$]
\label{exo:minimisation-edp}
\index{exercice!minimisation et EDP}
Soit $\Omega \subset \R^d$ un ouvert borné à bord régulier, et $f \in L^2(\Omega)$.
\begin{enumerate}[label=(\alph*)]
  \item Montrer que la fonctionnelle
    $\Phi(u) = \frac{1}{2}\int_\Omega \abs{\nabla u}^2 \dd x - \int_\Omega f u \, \dd x$
    est coercive et faiblement s.c.i.\ sur $H_0^1(\Omega)$.
  \item En déduire l'existence d'un minimiseur $u_0 \in H_0^1(\Omega)$.
  \item Montrer que $u_0$ est solution faible de $-\Delta u = f$ avec conditions
    de Dirichlet homogènes.
  \item Montrer l'unicité du minimiseur.
\end{enumerate}
\end{exercice}

\begin{exercice}[$\star\star$]
\label{exo:reflexif-suite-minimisante}
\index{exercice!réflexivité et suites minimisantes}
Soit $E$ un espace de Banach réflexif et $C \subset E$ un convexe fermé non vide.
Pour $x \in E$, montrer qu'il existe un unique $y \in C$ tel que
$\norm{x - y} = \inf_{z \in C} \norm{x - z}$ si $E$ est en outre uniformément convexe.

\textit{Indication :} L'existence vient de la méthode directe. Pour l'unicité,
utiliser la convexité uniforme.
\end{exercice}

\begin{exercice}[$\star\star\star$]
\label{exo:schur}
\index{exercice!propriété de Schur}
\index{propriété de Schur}
Démontrer la propriété de Schur\index{Schur, Issai} : dans $\ell^1$, convergence
faible et convergence forte coïncident pour les suites. C'est-à-dire, si
$x_n \weakto x$ dans $\ell^1$, alors $\norm{x_n - x}_1 \to 0$.

\textit{Indication :} Raisonner par l'absurde. Si $\norm{x_n - x}_1 \not\to 0$, extraire
une sous-suite et utiliser un argument de blocs\index{argument de blocs} (lemme de
glissement) pour construire une forme linéaire qui contredit la convergence faible.
\end{exercice}

\begin{exercice}[$\star\star$]
\label{exo:topo-faible-produit}
\index{exercice!topologie faible et produit}
Soient $E$ et $F$ deux espaces de Banach et $T \in \mathcal{L}(E,F)$.
\begin{enumerate}[label=(\alph*)]
  \item Montrer que $T$ est continu de $(E, \sigma(E,E'))$ vers $(F, \sigma(F,F'))$.
  \item En déduire que si $x_n \weakto x$ dans $E$, alors $Tx_n \weakto Tx$ dans $F$.
  \item La réciproque est-elle vraie ? (c'est-à-dire, une application linéaire
    faiblement-faiblement continue est-elle bornée ?)
\end{enumerate}
\end{exercice}

\index{topologie!faible-*|)}
\index{topologie!faible|)}
%%%%%%%%%%%%%%%%%%%%%%%%%%%%%%%%%%%%%%%%%%%%%%%%%%%%%%%%%%%%%%%%%%%%%%%
%  CHAPITRE 9 — Théorie Spectrale des Opérateurs Compacts
%%%%%%%%%%%%%%%%%%%%%%%%%%%%%%%%%%%%%%%%%%%%%%%%%%%%%%%%%%%%%%%%%%%%%%%
\chapter{Théorie Spectrale des Opérateurs Compacts}
\label{chap:spectral-compact}
\minitoc

\section{Opérateurs compacts~: rappels et exemples}
\index{opérateur compact}
\index{opérateur compact!définition}

\begin{definition}[Opérateur compact]
\label{def:op-compact}
Soient $E$ et $F$ deux espaces de Banach. Un opérateur linéaire $T \in \mathcal{L}(E,F)$
est dit \emph{compact}\index{compact|see{opérateur compact}} si l'image de la boule unité
fermée $\overline{B}_E$ est relativement compacte dans $F$, c'est-à-dire si
$\overline{T(\overline{B}_E)}$ est compact dans~$F$.

De manière équivalente, $T$ est compact si et seulement si pour toute suite bornée
$(x_n)$ de $E$, la suite $(Tx_n)$ admet une sous-suite convergente dans $F$.
\end{definition}
\index{opérateur compact!caractérisation séquentielle}

\begin{notation}
On note $\mathcal{K}(E,F)$\index{K(E,F)@$\mathcal{K}(E,F)$} l'ensemble des opérateurs
compacts de $E$ dans $F$, et $\mathcal{K}(E) = \mathcal{K}(E,E)$.
\end{notation}

\begin{exemple}[Opérateurs de rang fini]
\label{ex:rang-fini}
\index{opérateur!de rang fini}
Un opérateur $T \in \mathcal{L}(E,F)$ est de \emph{rang fini} si $\dim \im(T) < \infty$.
Tout opérateur de rang fini est compact~: en effet, $T(\overline{B}_E)$ est un borné d'un
espace de dimension finie, donc relativement compact.

Par exemple, si $\varphi_1, \ldots, \varphi_n \in E'$ et $y_1, \ldots, y_n \in F$,
l'opérateur
\[
  Tx = \sum_{k=1}^n \varphi_k(x) \, y_k
\]
est de rang fini, donc compact.
\end{exemple}

\begin{exemple}[Opérateurs intégraux]
\label{ex:op-integral}
\index{opérateur intégral}
\index{opérateur intégral!noyau de Hilbert--Schmidt}
Soit $\Omega \subset \R^d$ un ouvert borné et $K \in L^2(\Omega \times \Omega)$.
L'opérateur intégral
\[
  (T_K f)(x) = \int_\Omega K(x,y)\, f(y) \, \dd y
\]
définit un opérateur compact $T_K : L^2(\Omega) \to L^2(\Omega)$. On dit que $K$ est le
\emph{noyau} de $T_K$. C'est un opérateur de \emph{Hilbert--Schmidt}.
\index{Hilbert--Schmidt!opérateur de}
\end{exemple}

\begin{exemple}[Injection de Sobolev]
\label{ex:sobolev-compact}
\index{injection de Sobolev}
\index{Rellich--Kondrachov (théorème de)}
Si $\Omega \subset \R^d$ est un ouvert borné régulier, l'injection
$H^1_0(\Omega) \hookrightarrow L^2(\Omega)$ est compacte
(théorème de Rellich--Kondrachov).
\end{exemple}

\begin{exemple}[Opérateur diagonal]
\label{ex:op-diagonal-compact}
\index{opérateur diagonal compact}
Sur $\ell^2(\N)$, considérons l'opérateur diagonal $T(e_n) = \lambda_n e_n$ où
$(\lambda_n)$ est une suite de scalaires. Alors $T$ est compact si et seulement si
$\lambda_n \to 0$.
\end{exemple}

\begin{proposition}[Propriétés élémentaires]
\label{prop:prop-compact}
\index{opérateur compact!propriétés}
Soient $E, F, G$ des espaces de Banach.
\begin{enumerate}[label=(\roman*)]
  \item $\mathcal{K}(E,F)$ est un sous-espace vectoriel de $\mathcal{L}(E,F)$.
  \item Si $T \in \mathcal{K}(E,F)$, $S \in \mathcal{L}(F,G)$ et
        $R \in \mathcal{L}(G,E)$, alors $ST \in \mathcal{K}(E,G)$ et
        $TR \in \mathcal{K}(G,F)$.
  \item $\mathcal{K}(E)$ est un idéal bilatère de $\mathcal{L}(E)$.
\end{enumerate}
\end{proposition}

\begin{proof}
(i) Si $T_1, T_2 \in \mathcal{K}(E,F)$ et $\alpha, \beta \in \K$, soit $(x_n)$ bornée
dans $E$. On extrait une sous-suite telle que $(T_1 x_{n_k})$ converge, puis une
sous-sous-suite telle que $(T_2 x_{n_{k_l}})$ converge. Alors
$((\alpha T_1 + \beta T_2) x_{n_{k_l}})$ converge.

(ii) Si $(x_n)$ est bornée dans $E$, $(Tx_n)$ admet une sous-suite convergente
$(Tx_{n_k}) \to y$. Alors $STx_{n_k} \to Sy$ par continuité de $S$. Donc $ST$ est
compact. Pour $TR$~: si $(z_n)$ est bornée dans $G$, alors $(Rz_n)$ est bornée dans $E$,
donc $(TRz_n)$ admet une sous-suite convergente.

(iii) Conséquence directe de (i) et (ii).
\end{proof}

%-----------------------------------------------------------------------
\section{Fermeture de l'espace des opérateurs compacts}
\index{opérateur compact!fermeture dans L(E)}
%-----------------------------------------------------------------------

\begin{theoreme}[{$\mathcal{K}(E,F)$ est fermé dans $\mathcal{L}(E,F)$}]
\label{thm:compact-ferme}
\index{opérateur compact!espace fermé}
Soient $E, F$ deux espaces de Banach. L'espace $\mathcal{K}(E,F)$ est fermé dans
$\mathcal{L}(E,F)$ pour la topologie de la norme.
\end{theoreme}

\begin{proof}
Soit $(T_n)$ une suite d'opérateurs compacts convergeant vers $T \in \mathcal{L}(E,F)$
en norme~: $\norm{T_n - T} \to 0$. Montrons que $T$ est compact.

Soit $(x_k)$ une suite bornée de $E$, disons $\norm{x_k} \le M$ pour tout $k$. On
procède par extraction diagonale.

\textbf{Étape 1.} Comme $T_1$ est compact, il existe une sous-suite
$(x^{(1)}_k)_{k \ge 1}$ de $(x_k)$ telle que $(T_1 x^{(1)}_k)$ converge. Comme $T_2$
est compact, on extrait de $(x^{(1)}_k)$ une sous-suite $(x^{(2)}_k)$ telle que
$(T_2 x^{(2)}_k)$ converge. Par récurrence, on obtient des sous-suites emboîtées
$(x^{(n)}_k)_k$ telles que $(T_n x^{(n)}_k)_k$ converge pour chaque $n$ fixé.

\textbf{Étape 2.} Posons $y_k = x^{(k)}_k$ (extraction diagonale). Pour tout $n$ fixé,
la suite $(T_n y_k)_{k \ge n}$ est une sous-suite de $(T_n x^{(n)}_k)_k$, donc elle
converge. Ainsi $(T_n y_k)_k$ converge pour tout $n$.

\textbf{Étape 3.} Montrons que $(Ty_k)$ est de Cauchy. Soit $\varepsilon > 0$.
Choisissons $n$ tel que $\norm{T - T_n} < \varepsilon/(3M)$. Comme $(T_n y_k)_k$
converge, il existe $K_0$ tel que pour $j, l \ge K_0$~:
\[
  \norm{T_n y_j - T_n y_l} < \frac{\varepsilon}{3}.
\]
Alors pour $j, l \ge K_0$~:
\begin{align*}
  \norm{Ty_j - Ty_l}
    &\le \norm{Ty_j - T_n y_j} + \norm{T_n y_j - T_n y_l} + \norm{T_n y_l - Ty_l} \\
    &\le \norm{T - T_n}\,\norm{y_j} + \frac{\varepsilon}{3} + \norm{T_n - T}\,\norm{y_l}\\
    &< \frac{\varepsilon}{3M}\cdot M + \frac{\varepsilon}{3} + \frac{\varepsilon}{3M}
       \cdot M = \varepsilon.
\end{align*}
Comme $F$ est complet, $(Ty_k)$ converge. Donc $T$ est compact.
\end{proof}

\begin{corollaire}
\label{cor:limite-rang-fini}
\index{opérateur compact!limite d'opérateurs de rang fini}
\index{propriété d'approximation}
Si $F$ possède la propriété d'approximation (en particulier si $F$ est un espace de
Hilbert), tout opérateur compact $T \in \mathcal{K}(E,F)$ est limite en norme
d'opérateurs de rang fini.
\end{corollaire}

%-----------------------------------------------------------------------
\section{Alternative de Fredholm}
\index{Fredholm!alternative de}
\index{alternative de Fredholm}
%-----------------------------------------------------------------------

L'alternative de Fredholm est un résultat fondamental qui décrit la structure des
solutions de l'équation $x - Tx = y$ lorsque $T$ est compact.

\begin{theoreme}[Alternative de Fredholm]
\label{thm:fredholm}
Soit $E$ un espace de Banach et $T \in \mathcal{K}(E)$. Alors exactement l'une des deux
alternatives suivantes est réalisée~:
\begin{enumerate}[label=\textbf{(\alph*)}]
  \item L'opérateur $\Id - T$ est bijectif~: pour tout $y \in E$, l'équation $x - Tx = y$
        admet une unique solution, et $(\Id - T)^{-1} \in \mathcal{L}(E)$.
  \item L'opérateur $\Id - T$ n'est pas injectif~: $\Ker(\Id - T) \neq \{0\}$, et dans
        ce cas $\dim \Ker(\Id - T) < \infty$ et $\im(\Id - T)$ est un sous-espace fermé
        de codimension finie avec
        $\codim \im(\Id - T) = \dim \Ker(\Id - T)$.
\end{enumerate}
\end{theoreme}

\begin{proof}
La preuve repose sur la théorie de Riesz et le théorème de Schauder.

\textbf{Étape 1~: $\Ker(\Id - T)$ est de dimension finie.}
\index{Riesz!lemme de}
Sur $N = \Ker(\Id - T)$, on a $Tx = x$ pour tout $x \in N$. La boule unité fermée de $N$
est donc contenue dans $T(\overline{B}_E)$, qui est relativement compacte. Ainsi la boule
unité de $N$ est compacte, ce qui implique $\dim N < \infty$ par le théorème de Riesz.

\textbf{Étape 2~: $\im(\Id - T)$ est fermé.}
Comme $\Ker(\Id - T)$ est de dimension finie, il admet un supplémentaire topologique~:
$E = \Ker(\Id - T) \oplus F_0$ avec $F_0$ fermé. La restriction
$(\Id - T)|_{F_0} : F_0 \to E$ est injective.

Montrons qu'il existe $c > 0$ tel que $\norm{x - Tx} \ge c \norm{x}$ pour tout
$x \in F_0$. Supposons par l'absurde que non~: il existe une suite $(x_n) \subset F_0$
avec $\norm{x_n} = 1$ et $\norm{x_n - Tx_n} \to 0$. Comme $T$ est compact, quitte à
extraire, $Tx_{n_k} \to z$. Alors $x_{n_k} \to z$, donc $z \in F_0$ (car $F_0$ est
fermé) et $\norm{z} = 1$. Mais $z - Tz = \lim_{k}(x_{n_k} - Tx_{n_k}) = 0$, donc
$z \in \Ker(\Id - T) \cap F_0 = \{0\}$, contradiction avec $\norm{z} = 1$.

Donc il existe $c > 0$ tel que $\norm{(\Id-T)x} \ge c\norm{x}$ pour tout $x \in F_0$.
Si $(\Id-T)x_n \to y$ avec $x_n \in F_0$, alors $(x_n)$ est de Cauchy, converge vers
$x \in F_0$, et $(\Id-T)x = y$. Donc $\im(\Id-T)$ est fermé.

\textbf{Étape 3~: Noyaux et images itérés.}
\index{noyaux itérés}
Posons $N_k = \Ker(\Id - T)^k$ et $R_k = \im(\Id - T)^k$. On a
$N_0 \subset N_1 \subset N_2 \subset \cdots$ et
$R_0 \supset R_1 \supset R_2 \supset \cdots$, chaque $N_k$ est de dimension finie et
chaque $R_k$ est fermé (par récurrence, en utilisant que $(\Id-T)^k$ est de la forme
$\Id - S_k$ avec $S_k$ compact).

\emph{La suite $(N_k)$ se stabilise.} Sinon, $N_k \subsetneq N_{k+1}$ pour tout $k$.
Par le lemme de Riesz, il existe $x_k \in N_{k+1}$ avec $\norm{x_k} = 1$ et
$\mathrm{dist}(x_k, N_k) \ge 1/2$. Pour $j < k$~:
\[
  Tx_k - Tx_j = x_k - \bigl[(x_k - Tx_k) + x_j - (x_j - Tx_j)\bigr]
              = x_k - \underbrace{\bigl[x_j + (\Id-T)x_k - (\Id-T)x_j\bigr]}_{\in N_k}.
\]
En effet, $(\Id-T)x_k \in N_k$ car $x_k \in N_{k+1}$, et $x_j, (\Id-T)x_j \in N_k$
car $j < k$. Donc $\norm{Tx_k - Tx_j} \ge 1/2$, ce qui contredit la compacité de $T$.

De même, $(R_k)$ se stabilise~: si $R_k \supsetneq R_{k+1}$ pour tout $k$, on obtient
une contradiction similaire.

\textbf{Étape 4~: Indice de Fredholm et codimension.}
\index{Fredholm!indice de}
\index{Schauder!théorème de}
L'opérateur adjoint $T' \in \mathcal{L}(E')$ est aussi compact (théorème de Schauder).
On a la relation
\[
  \im(\Id - T)^\perp = \Ker(\Id - T'),
\]
où $\im(\Id-T)^\perp = \{\varphi \in E' : \varphi|_{\im(\Id-T)} = 0\}$. Comme
$\im(\Id-T)$ est fermé~:
\[
  \codim \im(\Id-T) = \dim \im(\Id-T)^\perp = \dim \Ker(\Id - T').
\]

Il reste à montrer que $\dim \Ker(\Id - T) = \dim \Ker(\Id - T')$. On montre que
l'\emph{indice de Fredholm} $\mathrm{ind}(\Id - T) = \dim \Ker(\Id-T) - \codim \im(\Id-T)$
vaut $0$. Considérons la famille $\Id - tT$ pour $t \in [0,1]$. Pour $t = 0$,
$\mathrm{ind}(\Id) = 0$. L'indice est localement constant sur l'ensemble des opérateurs
de Fredholm (qui est ouvert), et $t \mapsto \Id - tT$ est continue. Donc
$\mathrm{ind}(\Id - T) = \mathrm{ind}(\Id) = 0$.

Ainsi, dans le cas (b)~: $\codim \im(\Id - T) = \dim \Ker(\Id - T) < \infty$.
\end{proof}

\begin{remarque}
\index{Fredholm!alternative de!interprétation duale}
Dans le cas (b), l'équation $x - Tx = y$ a une solution si et seulement si
$\varphi(y) = 0$ pour tout $\varphi \in \Ker(\Id - T')$, c'est-à-dire si $y$ est
orthogonal (au sens de la dualité) à toutes les solutions de l'équation homogène adjointe
$\varphi - T'\varphi = 0$.
\end{remarque}

%-----------------------------------------------------------------------
\section{Théorie de Riesz--Schauder}
\index{Riesz--Schauder (théorie de)}
\index{spectre!opérateur compact}
%-----------------------------------------------------------------------

\begin{theoreme}[Riesz--Schauder]
\label{thm:riesz-schauder}
Soit $E$ un espace de Banach de dimension infinie et $T \in \mathcal{K}(E)$. Alors~:
\begin{enumerate}[label=(\roman*)]
  \item $0 \in \sigma(T)$.
  \item $\sigma(T) \setminus \{0\}$ est constitué uniquement de valeurs propres de $T$.
        \index{valeur propre!opérateur compact}
  \item Chaque valeur propre non nulle $\lambda$ a une multiplicité algébrique finie.
        \index{multiplicité!algébrique finie}
  \item L'ensemble $\sigma(T) \setminus \{0\}$ est au plus dénombrable, et $0$ est le
        seul point d'accumulation possible de $\sigma(T)$.
        \index{spectre!dénombrable}
\end{enumerate}
\end{theoreme}

\begin{proof}
\textbf{(i)} Si $0 \notin \sigma(T)$, alors $T$ est inversible. Donc
$\Id = T^{-1} \circ T$ serait compact (composé d'un opérateur borné et d'un compact).
Or $\Id$ n'est compact que si $\dim E < \infty$. Contradiction.

\textbf{(ii)} Soit $\lambda \in \sigma(T) \setminus \{0\}$. L'opérateur
$\lambda \Id - T = \lambda(\Id - T/\lambda)$. Comme $T/\lambda$ est compact,
l'alternative de Fredholm (théorème~\ref{thm:fredholm}) s'applique à $\Id - T/\lambda$.
Si $\Id - T/\lambda$ était bijectif, alors $\lambda\Id - T$ le serait aussi, et
$\lambda \notin \sigma(T)$~: contradiction. Donc $\Ker(\Id - T/\lambda) \neq \{0\}$,
c'est-à-dire que $\lambda$ est valeur propre de $T$.

\textbf{(iii)} Par l'alternative de Fredholm,
$\dim \Ker(\Id - T/\lambda) < \infty$. Plus généralement,
$\dim \Ker(\Id - T/\lambda)^k < \infty$ pour tout $k$ (car $(\Id - T/\lambda)^k$ est de
la forme $\Id - S$ avec $S$ compact), et la suite des noyaux itérés se stabilise.

\textbf{(iv)} Montrons que pour tout $\varepsilon > 0$, l'ensemble
$\{\lambda \in \sigma(T) : \abs{\lambda} \ge \varepsilon\}$ est fini.

Supposons par l'absurde qu'il existe une suite infinie $(\lambda_n)_{n \ge 1}$ de valeurs
propres distinctes avec $\abs{\lambda_n} \ge \varepsilon > 0$. Soient $e_n$ des vecteurs
propres associés~: $Te_n = \lambda_n e_n$, $\norm{e_n} = 1$. Posons
$F_n = \spn(e_1, \ldots, e_n)$. Comme les $\lambda_n$ sont distincts, les vecteurs
propres sont linéairement indépendants (résultat classique d'algèbre linéaire), donc
$F_{n-1} \subsetneq F_n$ pour tout $n \ge 2$.

Par le lemme de Riesz, il existe $x_n \in F_n$ avec $\norm{x_n} = 1$ et
$\mathrm{dist}(x_n, F_{n-1}) \ge 1/2$. Chaque $F_n$ est invariant par $T$. On écrit
$x_n = \sum_{k=1}^n \alpha_k^{(n)} e_k$, de sorte que
\[
  Tx_n = \sum_{k=1}^n \alpha_k^{(n)} \lambda_k e_k.
\]
Pour $m < n$, calculons~:
\[
  \frac{Tx_n}{\lambda_n} = x_n + \sum_{k=1}^{n-1} \alpha_k^{(n)}
    \Bigl(\frac{\lambda_k}{\lambda_n} - 1\Bigr) e_k
\]
de sorte que $Tx_n / \lambda_n - x_n \in F_{n-1}$. De même,
$Tx_m / \lambda_m \in F_m \subset F_{n-1}$. Donc
\[
  \frac{Tx_n}{\lambda_n} - \frac{Tx_m}{\lambda_m}
    = x_n - \underbrace{\Bigl(x_n - \frac{Tx_n}{\lambda_n}
      + \frac{Tx_m}{\lambda_m}\Bigr)}_{\in F_{n-1}}
\]
ce qui donne $\norm{Tx_n/\lambda_n - Tx_m/\lambda_m} \ge 1/2$. Ainsi
\[
  \norm{Tx_n - Tx_m}
    \ge \frac{1}{2}\min(\abs{\lambda_n}, \abs{\lambda_m})
    \ge \frac{\varepsilon}{2}.
\]
Comme $(x_n)$ est bornée ($\norm{x_n} = 1$), cela contredit la compacité de $T$.

Donc pour tout $\varepsilon > 0$, l'ensemble
$\{\lambda \in \sigma(T) : \abs{\lambda} \ge \varepsilon\}$ est fini. En prenant
$\varepsilon = 1/n$ et en réunissant, on obtient que $\sigma(T) \setminus \{0\}$ est au
plus dénombrable, et $0$ est le seul point d'accumulation possible.
\end{proof}

%-----------------------------------------------------------------------
\section{Opérateurs compacts auto-adjoints~: théorème spectral}
\index{opérateur compact!auto-adjoint}
\index{théorème spectral!opérateurs compacts}
%-----------------------------------------------------------------------

Dans cette section, $H$ désigne un espace de Hilbert séparable de dimension infinie
sur $\K = \R$ ou $\C$.

\begin{lemme}
\label{lem:norme-autoadjoint}
\index{opérateur auto-adjoint!norme}
Soit $T \in \mathcal{L}(H)$ auto-adjoint ($T = T^*$). Alors
\[
  \norm{T} = \sup_{\norm{x}=1} \abs{\inner{Tx}{x}}.
\]
\end{lemme}

\begin{proof}
Posons $M = \sup_{\norm{x}=1} \abs{\inner{Tx}{x}}$. L'inégalité $M \le \norm{T}$ est
immédiate~: $\abs{\inner{Tx}{x}} \le \norm{T}\norm{x}^2 = \norm{T}$.

Pour l'inégalité inverse, utilisons l'identité de polarisation. Comme $T$ est
auto-adjoint, pour tous $u, v \in H$~:
\[
  4\,\Re\inner{Tu}{v} = \inner{T(u+v)}{u+v} - \inner{T(u-v)}{u-v}.
\]
Donc
\[
  4\abs{\Re\inner{Tu}{v}} \le M\bigl(\norm{u+v}^2 + \norm{u-v}^2\bigr)
    = 2M\bigl(\norm{u}^2 + \norm{v}^2\bigr)
\]
par l'identité du parallélogramme\index{identité du parallélogramme}. Pour
$\norm{u} = \norm{v} = 1$, on obtient $\abs{\Re\inner{Tu}{v}} \le M$.

Soit $x \in H$ avec $\norm{x} = 1$ et $Tx \neq 0$. Posons $v = Tx/\norm{Tx}$. Alors
$\norm{v} = 1$ et
\[
  \norm{Tx} = \inner{Tx}{v} = \Re\inner{Tx}{v} \le M.
\]
(On peut toujours choisir la phase pour que $\inner{Tx}{v}$ soit réel positif.)
Donc $\norm{T} = \sup_{\norm{x}=1} \norm{Tx} \le M$.
\end{proof}

\begin{lemme}
\label{lem:vp-extremale}
\index{valeur propre!extrémale}
Soit $T \in \mathcal{K}(H)$ auto-adjoint, $T \neq 0$. Alors $\norm{T}$ ou $-\norm{T}$
est valeur propre de $T$.
\end{lemme}

\begin{proof}
Par le lemme~\ref{lem:norme-autoadjoint}, il existe une suite $(x_n)$ avec
$\norm{x_n} = 1$ et $\inner{Tx_n}{x_n} \to \lambda$ où $\abs{\lambda} = \norm{T}$.
Notons que $\lambda \in \R$ car $\inner{Tx}{x} \in \R$ pour $T$ auto-adjoint.

Calculons~:
\begin{align*}
  \norm{Tx_n - \lambda x_n}^2
    &= \norm{Tx_n}^2 - 2\lambda\inner{Tx_n}{x_n} + \lambda^2\norm{x_n}^2 \\
    &= \norm{Tx_n}^2 - 2\lambda\inner{Tx_n}{x_n} + \lambda^2 \\
    &\le \norm{T}^2 - 2\lambda\inner{Tx_n}{x_n} + \lambda^2.
\end{align*}
Comme $\lambda^2 = \norm{T}^2$ et $\inner{Tx_n}{x_n} \to \lambda$, on obtient
\[
  \norm{Tx_n - \lambda x_n}^2 \le 2\lambda^2 - 2\lambda\inner{Tx_n}{x_n} \to 0.
\]

Comme $T$ est compact et $(x_n)$ bornée, quitte à extraire, $Tx_{n_k} \to y$. Alors
$\lambda x_{n_k} \to y$ (car $\norm{Tx_{n_k} - \lambda x_{n_k}} \to 0$), donc
$x_{n_k} \to e := y/\lambda$ (car $\lambda \neq 0$). On a $\norm{e} = 1$ et
$Te = \lim Tx_{n_k} = y = \lambda e$. Donc $\lambda$ est valeur propre de $T$.
\end{proof}

\begin{theoreme}[Théorème spectral pour les opérateurs compacts auto-adjoints]
\label{thm:spectral-compact}
\index{théorème spectral!décomposition}
\index{décomposition spectrale}
Soit $H$ un espace de Hilbert séparable de dimension infinie et
$T \in \mathcal{K}(H)$ auto-adjoint. Alors il existe une suite (finie ou dénombrable)
$(\lambda_n)_{n \ge 1}$ de valeurs propres réelles non nulles, tendant vers $0$ si la
suite est infinie, et une famille orthonormée $(e_n)_{n \ge 1}$ de vecteurs propres
associés, telle que~:
\[
  Tx = \sum_{n \ge 1} \lambda_n \inner{x}{e_n}\, e_n, \quad \forall x \in H.
\]
De plus, $H = \Ker(T) \oplus \overline{\spn\{e_n : n \ge 1\}}$.
\end{theoreme}

\begin{proof}
\textbf{Étape 1~: Construction itérative.} Si $T = 0$, il n'y a rien à prouver.
Sinon, par le lemme~\ref{lem:vp-extremale}, il existe $\lambda_1$ avec
$\abs{\lambda_1} = \norm{T}$ et un vecteur propre unitaire $e_1$~: $Te_1 = \lambda_1 e_1$.

Posons $H_1 = \{e_1\}^\perp$. Le sous-espace $H_1$ est invariant par $T$~: si
$\inner{x}{e_1} = 0$, alors
$\inner{Tx}{e_1} = \inner{x}{Te_1} = \lambda_1\inner{x}{e_1} = 0$.

La restriction $T_1 = T|_{H_1} : H_1 \to H_1$ est un opérateur compact auto-adjoint sur
l'espace de Hilbert $H_1$. Si $T_1 = 0$, on s'arrête. Sinon, par le
lemme~\ref{lem:vp-extremale} appliqué à $T_1$, il existe $\lambda_2$ avec
$\abs{\lambda_2} = \norm{T_1} \le \norm{T} = \abs{\lambda_1}$ et $e_2 \in H_1$,
$\norm{e_2} = 1$, $Te_2 = \lambda_2 e_2$.

Par récurrence, on construit $\lambda_1, \lambda_2, \ldots$ et $e_1, e_2, \ldots$
orthonormés avec $Te_n = \lambda_n e_n$ et
$\abs{\lambda_1} \ge \abs{\lambda_2} \ge \cdots > 0$.

\textbf{Étape 2~: $\lambda_n \to 0$.} Si la suite est infinie, les $(e_n)$ forment un
système orthonormé, donc $\norm{e_n - e_m} = \sqrt{2}$ pour $n \neq m$. Or
$Te_n = \lambda_n e_n$, d'où
\[
  \norm{Te_n - Te_m}^2 = \lambda_n^2 + \lambda_m^2.
\]
Si $\abs{\lambda_n} \ge \varepsilon > 0$ pour tout $n$, alors
$\norm{Te_n - Te_m} \ge \varepsilon\sqrt{2}$, ce qui contredit la compacité de $T$.
Donc $\lambda_n \to 0$.

\textbf{Étape 3~: Décomposition.} Posons $H_n = \{e_1, \ldots, e_n\}^\perp$. Soit
$x \in H$ et écrivons
\[
  x = \sum_{k=1}^n \inner{x}{e_k}\,e_k + x_n, \qquad x_n \in H_n.
\]
On a $Tx_n = Tx - \sum_{k=1}^n \lambda_k \inner{x}{e_k}\,e_k$ et
$\norm{T|_{H_n}} = \abs{\lambda_{n+1}} \to 0$. Donc
\[
  \norm{Tx - \sum_{k=1}^n \lambda_k \inner{x}{e_k}\,e_k}
    = \norm{Tx_n} \le \abs{\lambda_{n+1}}\,\norm{x_n} \le \abs{\lambda_{n+1}}\,\norm{x}
    \to 0.
\]

Si le processus s'arrête après un nombre fini d'étapes, c'est-à-dire $T|_{H_N} = 0$
pour un certain $N$, la décomposition est une somme finie et $\Ker(T) = H_N$.

Dans tous les cas, si $z \perp e_n$ pour tout $n$, alors $z \in \bigcap_n H_n$ et
$\norm{Tz} \le \abs{\lambda_{n+1}}\norm{z}$ pour tout $n$, donc $Tz = 0$. Ainsi
$\Ker(T) = \overline{\spn\{e_n\}}^\perp$ et
$H = \Ker(T) \oplus \overline{\spn\{e_n : n \ge 1\}}$.
\end{proof}

\begin{remarque}
\label{rem:base-hilbertienne-vp}
\index{base hilbertienne!de vecteurs propres}
Si l'on complète la famille $(e_n)$ par une base hilbertienne de $\Ker(T)$, on obtient
une base hilbertienne de $H$ formée de vecteurs propres de $T$ (avec la valeur propre $0$
pour les vecteurs de $\Ker(T)$).
\end{remarque}

\begin{corollaire}[Min-max de Courant--Fischer]
\label{cor:courant-fischer}
\index{Courant--Fischer (théorème de)}
\index{min-max}
Soit $T \in \mathcal{K}(H)$ auto-adjoint positif ($\inner{Tx}{x} \ge 0$ pour tout $x$).
Ordonnons les valeurs propres~: $\lambda_1 \ge \lambda_2 \ge \cdots \ge 0$. Alors
\[
  \lambda_n = \min_{\substack{V \subset H \\ \codim V = n-1}}
    \sup_{\substack{x \in V \\ \norm{x} = 1}} \inner{Tx}{x}
  = \max_{\substack{W \subset H \\ \dim W = n}}
    \inf_{\substack{x \in W \\ \norm{x} = 1}} \inner{Tx}{x}.
\]
\end{corollaire}

%-----------------------------------------------------------------------
\section{Théorème de Hilbert--Schmidt}
\index{Hilbert--Schmidt!théorème de}
%-----------------------------------------------------------------------

\begin{definition}[Opérateur de Hilbert--Schmidt]
\label{def:hilbert-schmidt}
\index{Hilbert--Schmidt!opérateur de}
\index{norme de Hilbert--Schmidt}
Soit $H$ un espace de Hilbert séparable de base hilbertienne $(e_n)_{n \ge 1}$. Un
opérateur $T \in \mathcal{L}(H)$ est de \emph{Hilbert--Schmidt} si
\[
  \norm{T}_{\mathrm{HS}}^2 := \sum_{n=1}^\infty \norm{Te_n}^2 < \infty.
\]
Cette quantité ne dépend pas du choix de la base hilbertienne~: si $(f_m)$ est une autre
base, alors $\sum_n \norm{Te_n}^2 = \sum_n \sum_m \abs{\inner{Te_n}{f_m}}^2
= \sum_m \sum_n \abs{\inner{e_n}{T^*f_m}}^2 = \sum_m \norm{T^*f_m}^2$, et de même
$\sum_m \norm{Tf_m}^2 = \sum_n \norm{T^*e_n}^2 = \sum_n \norm{Te_n}^2$.
\end{definition}

\begin{theoreme}[Hilbert--Schmidt]
\label{thm:hilbert-schmidt}
Soit $T$ un opérateur de Hilbert--Schmidt sur un espace de Hilbert séparable $H$. Alors~:
\begin{enumerate}[label=(\roman*)]
  \item $T$ est compact.
  \item $\norm{T} \le \norm{T}_{\mathrm{HS}}$.
  \item Si de plus $T$ est auto-adjoint, et $(\lambda_n)$ désigne la suite de ses valeurs
        propres non nulles comptées avec multiplicité, alors
        $\sum_n \lambda_n^2 = \norm{T}_{\mathrm{HS}}^2$.
\end{enumerate}
\end{theoreme}

\begin{proof}
\textbf{(i)} Soit $(e_n)$ une base hilbertienne de $H$. Définissons l'opérateur de rang
fini $T_N$ par~:
\[
  T_N x = T\Bigl(\sum_{n=1}^N \inner{x}{e_n}\,e_n\Bigr)
        = \sum_{n=1}^N \inner{x}{e_n}\, Te_n.
\]
Chaque $T_N$ est de rang au plus $N$, donc compact. Pour $x \in H$ avec $\norm{x} \le 1$~:
\[
  (T - T_N)x = T\Bigl(\sum_{n > N} \inner{x}{e_n}\,e_n\Bigr)
             = \sum_{n > N} \inner{x}{e_n}\, Te_n.
\]
Par l'inégalité de Cauchy--Schwarz~:
\[
  \norm{(T - T_N)x}^2
    \le \Bigl(\sum_{n > N} \abs{\inner{x}{e_n}}^2\Bigr)
        \Bigl(\sum_{n > N} \norm{Te_n}^2\Bigr)
    \le \sum_{n > N} \norm{Te_n}^2 \to 0.
\]
Donc $\norm{T - T_N} \to 0$, et $T$ est compact comme limite d'opérateurs compacts.

\textbf{(ii)} L'inégalité $\norm{T} \le \norm{T}_{\mathrm{HS}}$ découle de~:
$\norm{Tx}^2 = \sum_n \abs{\inner{Tx}{e_n}}^2 = \sum_n \abs{\inner{x}{T^*e_n}}^2
\le \norm{x}^2 \sum_n \norm{T^*e_n}^2 = \norm{x}^2 \norm{T}_{\mathrm{HS}}^2$.

\textbf{(iii)} Par le théorème spectral~\ref{thm:spectral-compact}, choisissons une base
hilbertienne $(f_n)$ de $H$ formée de vecteurs propres de $T$~: $Tf_n = \mu_n f_n$
(avec $\mu_n = 0$ pour les vecteurs de $\Ker(T)$). Alors
\[
  \norm{T}_{\mathrm{HS}}^2 = \sum_n \norm{Tf_n}^2 = \sum_n \mu_n^2 = \sum_n \lambda_n^2
\]
(les $\mu_n$ non nuls étant exactement les $\lambda_n$).
\end{proof}

\begin{exemple}[Noyau de Hilbert--Schmidt]
\label{ex:hs-integral}
\index{Hilbert--Schmidt!noyau}
Soit $\Omega \subset \R^d$ mesurable de mesure finie et $K \in L^2(\Omega \times \Omega)$.
L'opérateur intégral $T_K$ défini par
$(T_K f)(x) = \int_\Omega K(x,y)f(y)\,\dd y$ est de Hilbert--Schmidt sur $L^2(\Omega)$
avec $\norm{T_K}_{\mathrm{HS}} = \norm{K}_{L^2(\Omega \times \Omega)}$.
\end{exemple}

%-----------------------------------------------------------------------
\section{Application~: équations intégrales de Fredholm}
\index{Fredholm!équation intégrale de}
\index{équation intégrale}
%-----------------------------------------------------------------------

Considérons l'équation intégrale de Fredholm de seconde espèce~:
\begin{equation}
\label{eq:fredholm-integrale}
  f(x) - \lambda \int_\Omega K(x,y)\, f(y) \, \dd y = g(x), \quad x \in \Omega,
\end{equation}
où $K \in L^2(\Omega \times \Omega)$, $g \in L^2(\Omega)$ et $\lambda \in \C$.

\begin{corollaire}
\label{cor:fredholm-integrale}
Soit $K \in L^2(\Omega \times \Omega)$ et $\lambda \in \C \setminus \{0\}$.
Alors~:
\begin{enumerate}[label=(\roman*)]
  \item Soit l'équation~\eqref{eq:fredholm-integrale} admet une unique solution
        $f \in L^2(\Omega)$ pour tout $g \in L^2(\Omega)$.
  \item Soit l'équation homogène ($g = 0$) admet un nombre fini $d \ge 1$ de solutions
        linéairement indépendantes $f_1, \ldots, f_d$, et
        l'équation~\eqref{eq:fredholm-integrale} a une solution si et seulement si
        \[
          \int_\Omega g(x)\,\overline{\bar f_j(x)}\,\dd x = 0,
          \qquad j = 1, \ldots, d,
        \]
        où $\bar f_1, \ldots, \bar f_d$ sont les solutions de l'équation homogène
        adjointe~:
        $\bar f(x) - \bar\lambda \int_\Omega \overline{K(y,x)}\,\bar f(y)\,\dd y = 0$.
\end{enumerate}
\end{corollaire}

\begin{proof}
L'équation~\eqref{eq:fredholm-integrale} s'écrit $(\Id - \lambda T_K)f = g$. Comme
$T_K$ est compact, on applique l'alternative de Fredholm à $\lambda T_K$. L'adjoint de
$\lambda T_K$ est $\bar\lambda T_{K^*}$ où $K^*(x,y) = \overline{K(y,x)}$.
\end{proof}

\begin{exemple}[Développement spectral]
\label{ex:fredholm-sym}
\index{Fredholm!noyau symétrique}
Soit $K(x,y) = \overline{K(y,x)}$ avec $K \in L^2([0,1]^2)$. L'opérateur $T_K$ est
auto-adjoint compact sur $L^2([0,1])$. Par le théorème spectral, il existe des valeurs
propres $(\mu_n)_{n \ge 1}$ (réelles, tendant vers $0$) et des fonctions propres
orthonormées $(e_n)$ telles que $T_K f = \sum_n \mu_n \inner{f}{e_n}\,e_n$.

L'équation $f - \lambda T_K f = g$ admet une solution pour tout $g$ si et seulement si
$\lambda \neq 1/\mu_n$ pour tout $n$, et dans ce cas~:
\[
  f = g + \sum_{n=1}^\infty \frac{\lambda \mu_n}{1 - \lambda \mu_n}\,
    \inner{g}{e_n}\, e_n.
\]
\end{exemple}

%-----------------------------------------------------------------------
\section{Exercices}
%-----------------------------------------------------------------------

\begin{exercice}[$\star$]
\index{exercice!opérateur compact}
Montrer que si $E$ est un espace de Banach de dimension infinie, l'identité $\Id_E$ n'est
pas compacte.
\end{exercice}

\begin{exercice}[$\star$]
Soit $T \in \mathcal{K}(H)$ avec $H$ un espace de Hilbert. Montrer que $T^*$ est
compact.
\end{exercice}

\begin{exercice}[$\star\star$]
\index{exercice!opérateur de Volterra}
\index{Volterra!opérateur de}
Montrer que l'opérateur de Volterra $V : L^2([0,1]) \to L^2([0,1])$ défini par
$(Vf)(x) = \int_0^x f(t)\,\dd t$ est compact. Montrer que $\sigma(V) = \{0\}$.

\emph{Indication~: pour le spectre, montrer que $V$ est quasi-nilpotent en estimant
$\norm{V^n}$.}
\end{exercice}

\begin{exercice}[$\star\star$]
\index{exercice!alternative de Fredholm}
Soit $K \in C([0,1]^2)$. Montrer que pour $\abs{\lambda}$ suffisamment petit, l'équation
$f(x) - \lambda \int_0^1 K(x,t)\,f(t)\,\dd t = g(x)$ admet une unique solution continue
pour tout $g \in C([0,1])$. \emph{Indication~: série de Neumann.}
\end{exercice}

\begin{exercice}[$\star\star\star$]
\index{exercice!théorème de Mercer}
\index{Mercer (théorème de)}
(\emph{Théorème de Mercer.}) Soit $K \in C([0,1]^2)$ symétrique et positif. Soient
$(\lambda_n, e_n)$ les couples valeurs propres/vecteurs propres de $T_K$. Montrer que
$K(x,y) = \sum_{n=1}^\infty \lambda_n\,e_n(x)\,\overline{e_n(y)}$ uniformément.
\end{exercice}

\begin{exercice}[$\star\star$]
Soit $T \in \mathcal{K}(H)$ auto-adjoint positif. Montrer que $T$ admet une unique racine
carrée compacte auto-adjointe positive~: il existe un unique $S \in \mathcal{K}(H)$
auto-adjoint positif tel que $S^2 = T$.
\end{exercice}

\begin{exercice}[$\star\star\star$]
\index{exercice!opérateur à trace}
\index{opérateur à trace}
On dit que $T \in \mathcal{L}(H)$ positif est un \emph{opérateur à trace} si
$\sum_n \inner{Te_n}{e_n} < \infty$ pour une base hilbertienne $(e_n)$. Montrer~:
\begin{enumerate}[label=(\alph*)]
  \item Cette définition ne dépend pas du choix de la base.
  \item Tout opérateur à trace est de Hilbert--Schmidt.
  \item Si $T$ est à trace, auto-adjoint positif, de valeurs propres $(\lambda_n)$, alors
        $\Tr(T) := \sum_n \inner{Te_n}{e_n} = \sum_n \lambda_n$.
\end{enumerate}
\end{exercice}


%%%%%%%%%%%%%%%%%%%%%%%%%%%%%%%%%%%%%%%%%%%%%%%%%%%%%%%%%%%%%%%%%%%%%%%
%  CHAPITRE 10 — Théorème Spectral pour les Opérateurs Auto-Adjoints
%%%%%%%%%%%%%%%%%%%%%%%%%%%%%%%%%%%%%%%%%%%%%%%%%%%%%%%%%%%%%%%%%%%%%%%
\chapter{Théorème Spectral pour les Opérateurs Auto-Adjoints}
\label{chap:spectral-autoadjoint}
\minitoc

\section{Calcul fonctionnel continu}
\index{calcul fonctionnel!continu}
\index{opérateur auto-adjoint!calcul fonctionnel}

Soit $H$ un espace de Hilbert complexe et $T \in \mathcal{L}(H)$ un opérateur
auto-adjoint borné. On sait que $\sigma(T) \subset \R$ et plus précisément
$\sigma(T) \subset [m, M]$ où $m = \inf_{\norm{x}=1}\inner{Tx}{x}$ et
$M = \sup_{\norm{x}=1}\inner{Tx}{x}$.

\begin{theoreme}[Calcul fonctionnel continu]
\label{thm:calcul-fonctionnel-continu}
\index{calcul fonctionnel!continu!théorème}
Soit $T \in \mathcal{L}(H)$ auto-adjoint. Il existe un unique morphisme d'algèbres
involutives isométrique
\[
  \Phi : C(\sigma(T)) \to \mathcal{L}(H), \quad f \mapsto f(T),
\]
tel que $\Phi(1) = \Id$ et $\Phi(\mathrm{id}) = T$ (où $\mathrm{id}(\lambda) = \lambda$).
Ce morphisme satisfait~:
\begin{enumerate}[label=(\roman*)]
  \item $(\alpha f + \beta g)(T) = \alpha f(T) + \beta g(T)$ (linéarité).
  \item $(fg)(T) = f(T)\,g(T)$ (multiplicativité).
  \item $\bar f(T) = f(T)^*$ (respect de l'involution).
  \item $\norm{f(T)} = \norm{f}_\infty := \sup_{\lambda \in \sigma(T)} \abs{f(\lambda)}$
        (isométrie).
  \item $\sigma(f(T)) = f(\sigma(T))$ (application spectrale).
\end{enumerate}
\end{theoreme}

\begin{proof}
\textbf{Étape 1~: Polynômes.} Pour un polynôme $p(\lambda) = \sum_{k=0}^n a_k \lambda^k$,
on définit $p(T) = \sum_{k=0}^n a_k T^k$. Les propriétés (i)--(iii) sont immédiates. Pour
(iv), on utilise la formule du rayon spectral\index{rayon spectral} et le fait que pour $T$
auto-adjoint, $\norm{T^n} = \norm{T}^n$~:
\[
  \norm{p(T)} = r(p(T)) = \sup_{\mu \in \sigma(p(T))} \abs{\mu}
    = \sup_{\lambda \in \sigma(T)} \abs{p(\lambda)} = \norm{p}_\infty.
\]
La dernière égalité utilise le théorème d'application spectrale pour les polynômes~:
$\sigma(p(T)) = p(\sigma(T))$.

\textbf{Étape 2~: Extension par densité.} Par le théorème de
Weierstrass\index{Weierstrass!théorème de}, les polynômes sont denses dans
$C(\sigma(T))$. Comme $\Phi$ est une isométrie sur les polynômes, elle se prolonge de
manière unique en une isométrie de $C(\sigma(T))$ dans $\mathcal{L}(H)$.

\textbf{Étape 3~: Propriétés du prolongement.} La linéarité, la multiplicativité et le
respect de l'involution passent à la limite par continuité. Pour l'application spectrale
(v)~: si $f \in C(\sigma(T))$ et $\mu \notin f(\sigma(T))$, alors
$g = (f - \mu)^{-1} \in C(\sigma(T))$ et $g(T)(f(T) - \mu\Id) = \Id$, donc
$\mu \notin \sigma(f(T))$. Réciproquement, si $\mu = f(\lambda_0)$ avec
$\lambda_0 \in \sigma(T)$, on montre que $\mu \in \sigma(f(T))$ par un argument
d'approximation.
\end{proof}

\begin{remarque}
\index{calcul fonctionnel!positivité}
Le calcul fonctionnel continu est \emph{positif}~: si $f \ge 0$ sur $\sigma(T)$, alors
$f(T) \ge 0$ (i.e., $\inner{f(T)x}{x} \ge 0$ pour tout $x$). En effet,
$f = (\sqrt{f})^2$, donc $f(T) = \sqrt{f}(T)^* \sqrt{f}(T) \ge 0$.
\end{remarque}

%-----------------------------------------------------------------------
\section{Mesures spectrales et résolution de l'identité}
\index{mesure spectrale}
\index{résolution de l'identité}
%-----------------------------------------------------------------------

\begin{definition}[Mesure spectrale]
\label{def:mesure-spectrale}
\index{mesure spectrale!définition}
Soit $(\Omega, \mathcal{B})$ un espace mesurable. Une \emph{mesure spectrale} (ou
\emph{mesure à valeurs projections}) sur $(\Omega, \mathcal{B})$ à valeurs dans un espace
de Hilbert $H$ est une application $E : \mathcal{B} \to \mathcal{L}(H)$ telle que~:
\begin{enumerate}[label=(\roman*)]
  \item $E(B)$ est un projecteur orthogonal pour tout $B \in \mathcal{B}$.
  \item $E(\emptyset) = 0$ et $E(\Omega) = \Id$.
  \item $E(B_1 \cap B_2) = E(B_1)\,E(B_2)$ pour tous $B_1, B_2 \in \mathcal{B}$.
  \item Si $(B_n)$ est une suite de boréliens disjoints, alors
        $E\bigl(\bigcup_n B_n\bigr) = \sum_n E(B_n)$ (convergence forte).
\end{enumerate}
\end{definition}

\begin{remarque}
Pour tous $x, y \in H$, l'application $B \mapsto \inner{E(B)x}{y}$ définit une mesure
complexe $\mu_{x,y}$ sur $(\Omega, \mathcal{B})$, et $B \mapsto \inner{E(B)x}{x} = \norm{E(B)x}^2$
est une mesure positive $\mu_x$ de masse totale $\norm{x}^2$.
\index{mesure spectrale!mesure associée}
\end{remarque}

\begin{definition}[Résolution de l'identité]
\label{def:resolution-identite}
\index{résolution de l'identité!définition}
Une \emph{résolution de l'identité} sur $\R$ est une famille $(E_\lambda)_{\lambda \in \R}$
de projecteurs orthogonaux sur $H$ telle que~:
\begin{enumerate}[label=(\roman*)]
  \item $E_\lambda \le E_\mu$ (au sens $E_\lambda H \subset E_\mu H$) pour $\lambda \le \mu$.
  \item $E_\lambda \to 0$ (fortement) quand $\lambda \to -\infty$ et
        $E_\lambda \to \Id$ quand $\lambda \to +\infty$.
  \item $E_\lambda$ est continue à droite~: $E_{\lambda^+} = E_\lambda$.
\end{enumerate}
\end{definition}

%-----------------------------------------------------------------------
\section{Théorème spectral~: version mesure spectrale}
\index{théorème spectral!opérateurs bornés}
\index{théorème spectral!mesure spectrale}
%-----------------------------------------------------------------------

\begin{theoreme}[Théorème spectral --- version mesure spectrale]
\label{thm:spectral-mesure}
Soit $T \in \mathcal{L}(H)$ auto-adjoint. Il existe une unique mesure spectrale $E$ sur
$(\R, \mathcal{B}(\R))$ à valeurs dans $H$, concentrée sur $\sigma(T)$, telle que~:
\[
  T = \int_{\sigma(T)} \lambda \, \dd E(\lambda),
\]
au sens où $\inner{Tx}{y} = \int_{\sigma(T)} \lambda \, \dd\mu_{x,y}(\lambda)$ pour
tous $x, y \in H$. Plus généralement, pour toute fonction borélienne bornée
$f : \sigma(T) \to \C$~:
\[
  f(T) = \int_{\sigma(T)} f(\lambda) \, \dd E(\lambda)
\]
et $\norm{f(T)} = \norm{f}_{L^\infty(\sigma(T), E)}$.
\end{theoreme}

\begin{proof}
\textbf{Étape 1~: Du calcul fonctionnel continu aux mesures.}
Par le théorème de représentation de Riesz\index{Riesz!théorème de représentation},
pour chaque $x \in H$, la forme linéaire positive $f \mapsto \inner{f(T)x}{x}$ sur
$C(\sigma(T))$ est représentée par une unique mesure de Borel positive $\mu_x$ sur
$\sigma(T)$~:
\[
  \inner{f(T)x}{x} = \int_{\sigma(T)} f(\lambda)\, \dd\mu_x(\lambda),
  \quad \forall f \in C(\sigma(T)).
\]
La masse de $\mu_x$ est $\mu_x(\sigma(T)) = \inner{\Id\, x}{x} = \norm{x}^2$.

\textbf{Étape 2~: Mesures complexes.}
Par polarisation, pour $x, y \in H$, on définit la mesure complexe~:
\[
  \mu_{x,y} = \frac{1}{4}\bigl(\mu_{x+y} - \mu_{x-y}
    + i\mu_{x+iy} - i\mu_{x-iy}\bigr),
\]
de sorte que $\inner{f(T)x}{y} = \int f\, \dd\mu_{x,y}$ pour tout $f \in C(\sigma(T))$.

\textbf{Étape 3~: Construction de la mesure spectrale.}
Pour tout borélien $B \subset \sigma(T)$, la forme sesquilinéaire
$(x,y) \mapsto \mu_{x,y}(B)$ est bornée (car $\abs{\mu_{x,y}(B)} \le \norm{x}\norm{y}$).
Par le théorème de représentation de Riesz--Fréchet, il existe un unique opérateur borné
$E(B) \in \mathcal{L}(H)$ tel que $\inner{E(B)x}{y} = \mu_{x,y}(B)$.

\emph{$E(B)$ est un projecteur orthogonal.} Pour $f = \mathbf{1}_B$ (fonction
caractéristique), on a $\mathbf{1}_B^2 = \mathbf{1}_B$ et
$\overline{\mathbf{1}_B} = \mathbf{1}_B$. En étendant le calcul fonctionnel aux fonctions
boréliennes bornées (par passage à la limite monotone et théorèmes de convergence), on
obtient $E(B)^2 = E(B)$ et $E(B)^* = E(B)$.

\emph{$\sigma$-additivité.} Si $B = \bigsqcup_n B_n$ (union disjointe), alors
$\mathbf{1}_B = \sum_n \mathbf{1}_{B_n}$. Pour tout $x$~:
$\norm{E(B)x}^2 = \mu_x(B) = \sum_n \mu_x(B_n) = \sum_n \norm{E(B_n)x}^2$. Comme les
$E(B_n)$ ont des images orthogonales ($E(B_n)E(B_m) = E(B_n \cap B_m) = 0$ pour
$n \neq m$), on obtient $E(B)x = \sum_n E(B_n)x$ (convergence forte).

\textbf{Étape 4~: Calcul fonctionnel borélien.}
Pour toute fonction borélienne bornée $f$, on définit $f(T) = \int f\,\dd E$. Par
construction~: $\inner{f(T)x}{y} = \int f\,\dd\mu_{x,y}$. Les propriétés de morphisme
d'algèbre et l'isométrie $\norm{f(T)} = \norm{f}_{L^\infty(E)}$ se vérifient par
approximation.

\textbf{Unicité.} Si $E'$ est une autre mesure spectrale vérifiant l'énoncé, alors
$\int f\,\dd E = \int f\,\dd E'$ pour tout $f \in C(\sigma(T))$. Par le théorème
d'unicité de Riesz, $\mu_x^E = \mu_x^{E'}$ pour tout $x$, d'où $E = E'$.
\end{proof}

%-----------------------------------------------------------------------
\section{Calcul fonctionnel borélien}
\index{calcul fonctionnel!borélien}
%-----------------------------------------------------------------------

\begin{definition}[Calcul fonctionnel borélien]
\label{def:calcul-borelien}
\index{calcul fonctionnel!borélien!définition}
Soit $T \in \mathcal{L}(H)$ auto-adjoint et $E$ sa mesure spectrale. Pour toute fonction
borélienne bornée $f : \sigma(T) \to \C$, on définit~:
\[
  f(T) = \int_{\sigma(T)} f(\lambda)\, \dd E(\lambda) \in \mathcal{L}(H).
\]
\end{definition}

\begin{proposition}
\label{prop:calcul-borelien-prop}
Le calcul fonctionnel borélien $f \mapsto f(T)$ est un morphisme d'algèbres involutives~:
\begin{enumerate}[label=(\roman*)]
  \item $(fg)(T) = f(T)\,g(T)$.
  \item $\bar f(T) = f(T)^*$.
  \item $\norm{f(T)} \le \norm{f}_\infty$.
  \item Si $f_n \to f$ ponctuellement avec $\sup_n \norm{f_n}_\infty < \infty$, alors
        $f_n(T) \to f(T)$ fortement (convergence dominée spectrale).
        \index{convergence dominée spectrale}
\end{enumerate}
\end{proposition}

\begin{exemple}[Projecteur spectral]
\label{ex:proj-spectral}
\index{projecteur spectral}
Pour un borélien $B \subset \sigma(T)$, on a $\mathbf{1}_B(T) = E(B)$. En particulier,
$E(\{\lambda\})$ est le projecteur orthogonal sur le sous-espace propre associé à
$\lambda$ (nul si $\lambda$ n'est pas valeur propre).
\end{exemple}

%-----------------------------------------------------------------------
\section{Opérateurs positifs et racine carrée}
\index{opérateur positif}
\index{racine carrée!d'un opérateur}
%-----------------------------------------------------------------------

\begin{definition}[Opérateur positif]
\label{def:op-positif}
\index{opérateur positif!définition}
Un opérateur $T \in \mathcal{L}(H)$ est dit \emph{positif} (noté $T \ge 0$) si $T$ est
auto-adjoint et $\inner{Tx}{x} \ge 0$ pour tout $x \in H$. On note $T \ge S$ si
$T - S \ge 0$.
\end{definition}

\begin{theoreme}[Racine carrée]
\label{thm:racine-carree}
\index{racine carrée!existence et unicité}
Soit $T \in \mathcal{L}(H)$ positif. Il existe un unique opérateur positif
$S \in \mathcal{L}(H)$ tel que $S^2 = T$. On le note $S = T^{1/2}$ ou $S = \sqrt{T}$.
De plus, $S$ commute avec tout opérateur qui commute avec $T$.
\end{theoreme}

\begin{proof}
\textbf{Existence.} Comme $T \ge 0$, on a $\sigma(T) \subset [0, \norm{T}]$. La fonction
$f(\lambda) = \sqrt{\lambda}$ est continue sur $\sigma(T)$. On pose $S = f(T) = \sqrt{\cdot}(T)$
par le calcul fonctionnel continu. Alors $S^2 = (\sqrt{\cdot})^2(T) = \mathrm{id}(T) = T$
et $S^* = \overline{\sqrt{\cdot}}(T) = \sqrt{\cdot}(T) = S$ (car $\sqrt{\cdot}$ est
réelle). De plus, $\inner{Sx}{x} = \int \sqrt{\lambda}\,\dd\mu_x(\lambda) \ge 0$.

\textbf{Unicité.} Soit $R \ge 0$ avec $R^2 = T$. Montrons $R = S$. D'abord, $R$ commute
avec $T = R^2$, et donc avec tout polynôme en $T$, et par densité avec $S = f(T)$. Ainsi
$RS = SR$.

Soit $x \in H$ et $y = (S - R)x$. Alors
\[
  \inner{Sy}{y} + \inner{Ry}{y}
    = \inner{(S+R)y}{y}
    = \inner{(S+R)(S-R)x}{y}
    = \inner{(S^2 - R^2)x}{y}
    = \inner{(T - T)x}{y} = 0.
\]
(On a utilisé $SR = RS$ pour que $(S+R)(S-R) = S^2 - R^2$.)
Comme $S, R \ge 0$, on a $\inner{Sy}{y} \ge 0$ et $\inner{Ry}{y} \ge 0$, donc
$\inner{Sy}{y} = \inner{Ry}{y} = 0$. En particulier,
$\inner{Sy}{y} = \norm{S^{1/2}y}^2 = 0$ (en utilisant la racine carrée de $S$), d'où
$S^{1/2}y = 0$, puis $Sy = S^{1/2}(S^{1/2}y) = 0$.

De même, $Ry = 0$. Donc $(S-R)y = 0$, i.e., $(S-R)^2 x = 0$ pour tout $x$. Comme
$S - R$ est auto-adjoint, $\norm{(S-R)x}^2 = \inner{(S-R)^2 x}{x} = 0$, d'où
$S = R$.

\textbf{Commutation.} Si $A$ commute avec $T$, alors $A$ commute avec tout polynôme en $T$,
et par densité (calcul fonctionnel continu), $A$ commute avec $S = \sqrt{\cdot}(T)$.
\end{proof}

\begin{definition}[Module d'un opérateur]
\label{def:module-op}
\index{module d'un opérateur}
\index{valeur absolue!d'un opérateur|see{module d'un opérateur}}
Pour $T \in \mathcal{L}(H)$, on définit le \emph{module} de $T$ par
$\abs{T} = (T^*T)^{1/2}$.
\end{definition}

%-----------------------------------------------------------------------
\section{Décomposition polaire}
\index{décomposition polaire}
%-----------------------------------------------------------------------

\begin{theoreme}[Décomposition polaire]
\label{thm:decomposition-polaire}
\index{décomposition polaire!théorème}
\index{isométrie partielle}
Soit $T \in \mathcal{L}(H)$. Il existe une unique isométrie partielle $U$ telle que
$T = U\abs{T}$ et $\Ker(U) = \Ker(T)$. De plus~:
\begin{enumerate}[label=(\roman*)]
  \item $U^*U$ est le projecteur orthogonal sur $\overline{\im(\abs{T})} = \overline{\im(T^*)} = (\Ker T)^\perp$.
  \item $UU^*$ est le projecteur orthogonal sur $\overline{\im(T)}$.
  \item $\abs{T} = U^*T$.
\end{enumerate}
\end{theoreme}

\begin{proof}
\textbf{Existence.} Pour tout $x \in H$~:
\[
  \norm{\abs{T}x}^2 = \inner{\abs{T}^2 x}{x} = \inner{T^*Tx}{x} = \norm{Tx}^2.
\]
Donc l'application $\abs{T}x \mapsto Tx$ est bien définie et isométrique de
$\im(\abs{T})$ dans $\im(T)$. Elle se prolonge en une isométrie
$U_0 : \overline{\im(\abs{T})} \to \overline{\im(T)}$. On prolonge $U_0$ en
$U \in \mathcal{L}(H)$ en posant $U = 0$ sur $\overline{\im(\abs{T})}^\perp = \Ker(\abs{T}) = \Ker(T)$.

Par construction, $U\abs{T}x = Tx$ pour tout $x$, $\Ker(U) = \Ker(\abs{T}) = \Ker(T)$,
et $U$ est une isométrie partielle d'espace initial
$\overline{\im(\abs{T})} = (\Ker T)^\perp$ et d'espace final $\overline{\im(T)}$.

\textbf{Propriétés.}
(i) $U^*U$ est le projecteur sur l'espace initial de $U$, soit $(\Ker T)^\perp$.
(ii) $UU^*$ est le projecteur sur l'espace final, soit $\overline{\im(T)}$.
(iii) $U^*T = U^*U\abs{T} = P_{(\Ker T)^\perp}\abs{T} = \abs{T}$ (car
$\im(\abs{T}) \subset (\Ker T)^\perp$).

\textbf{Unicité.} Si $T = V\abs{T}$ avec $V$ isométrie partielle et
$\Ker(V) = \Ker(T)$, alors $V$ et $U$ coïncident sur $\im(\abs{T})$ (car
$V\abs{T}x = Tx = U\abs{T}x$), donc sur $\overline{\im(\abs{T})}$ par continuité, et
toutes deux sont nulles sur $\Ker(T) = \overline{\im(\abs{T})}^\perp$. Donc $V = U$.
\end{proof}

%-----------------------------------------------------------------------
\section{Opérateurs non bornés}
\index{opérateur non borné}
\index{opérateur non borné!définition}
%-----------------------------------------------------------------------

\begin{definition}[Opérateur non borné]
\label{def:op-non-borne}
Un \emph{opérateur non borné} sur un espace de Hilbert $H$ est un opérateur linéaire
$T : \dom(T) \to H$ défini sur un sous-espace $\dom(T) \subset H$ appelé le
\emph{domaine} de $T$. On suppose en général que $\dom(T)$ est dense dans $H$.
\index{domaine d'un opérateur}
\end{definition}

\begin{definition}[Graphe d'un opérateur]
\label{def:graphe}
\index{graphe d'un opérateur}
Le \emph{graphe} de $T$ est le sous-espace
$\Gamma(T) = \{(x, Tx) : x \in \dom(T)\} \subset H \times H$.
\end{definition}

\begin{definition}[Opérateur fermé, opérateur fermable]
\label{def:op-ferme}
\index{opérateur fermé}
\index{opérateur fermable}
Un opérateur $T$ est \emph{fermé} si son graphe $\Gamma(T)$ est fermé dans $H \times H$.
Un opérateur $T$ est \emph{fermable} si la fermeture $\overline{\Gamma(T)}$ est encore le
graphe d'un opérateur (noté $\overline{T}$, la \emph{fermeture} de $T$).
\end{definition}

\begin{proposition}
\label{prop:ferme-equiv}
$T$ est fermé si et seulement si~: pour toute suite $(x_n) \subset \dom(T)$ avec
$x_n \to x$ et $Tx_n \to y$, on a $x \in \dom(T)$ et $Tx = y$.
\end{proposition}

\begin{exemple}[Opérateur de dérivation]
\label{ex:derivation}
\index{opérateur de dérivation}
Sur $H = L^2([0,1])$, l'opérateur $Tf = f'$ avec $\dom(T) = H^1([0,1])$ est un opérateur
fermé non borné, densément défini.
\end{exemple}

%-----------------------------------------------------------------------
\section{Opérateurs auto-adjoints non bornés}
\index{opérateur auto-adjoint!non borné}
%-----------------------------------------------------------------------

\begin{definition}[Adjoint d'un opérateur non borné]
\label{def:adjoint-non-borne}
\index{opérateur adjoint!non borné}
Soit $T$ densément défini. L'adjoint $T^*$ est l'opérateur de domaine
\[
  \dom(T^*) = \{y \in H : x \mapsto \inner{Tx}{y} \text{ est continu sur } \dom(T)\}
\]
et défini par $\inner{Tx}{y} = \inner{x}{T^*y}$ pour tout $x \in \dom(T)$.
\end{definition}

\begin{definition}[Opérateur symétrique, auto-adjoint]
\label{def:symetrique-autoadjoint}
\index{opérateur symétrique}
\index{opérateur auto-adjoint!non borné!définition}
Un opérateur densément défini $T$ est~:
\begin{itemize}
  \item \emph{symétrique} si $T \subset T^*$, i.e., $\dom(T) \subset \dom(T^*)$ et
        $T^*x = Tx$ pour $x \in \dom(T)$.
  \item \emph{auto-adjoint} si $T = T^*$, i.e., $T$ est symétrique et
        $\dom(T) = \dom(T^*)$.
\end{itemize}
\end{definition}

\begin{remarque}
\index{opérateur symétrique!vs auto-adjoint}
La distinction entre symétrique et auto-adjoint est cruciale pour la théorie spectrale.
Un opérateur symétrique peut ne pas être auto-adjoint si $\dom(T) \subsetneq \dom(T^*)$.
\end{remarque}

\begin{theoreme}[Critère d'auto-adjonction]
\label{thm:critere-autoadjoint}
\index{opérateur auto-adjoint!critère}
Soit $T$ un opérateur symétrique densément défini. Les conditions suivantes sont
équivalentes~:
\begin{enumerate}[label=(\roman*)]
  \item $T$ est auto-adjoint.
  \item $T$ est fermé et $\Ker(T^* \pm i\Id) = \{0\}$.
  \item $\im(T \pm i\Id) = H$.
\end{enumerate}
\end{theoreme}

\begin{proof}
\textbf{(i) $\Rightarrow$ (ii).} Si $T = T^*$, alors $T$ est fermé car $T^*$ est
toujours fermé. Si $T^*y = iy$, alors $\inner{Ty}{y} = \inner{y}{T^*y} = \inner{y}{iy} = -i\norm{y}^2$ et aussi $\inner{Ty}{y} = \inner{T^*y}{y} = i\norm{y}^2$ (car $T = T^*$). Donc
$\norm{y}^2 = 0$, d'où $y = 0$. Même argument pour $-i$.

\textbf{(ii) $\Rightarrow$ (iii).} Montrons que $\im(T + i\Id)$ est fermé. Si
$(T + i\Id)x_n \to y$ avec $x_n \in \dom(T)$, alors pour tout $n$~:
\[
  \norm{(T + i\Id)x_n}^2 = \norm{Tx_n}^2 + \norm{x_n}^2
\]
(par développement, utilisant $\inner{Tx_n}{ix_n} + \inner{ix_n}{Tx_n} = 0$ car $T$ est
symétrique). Donc $\norm{x_n}^2 \le \norm{(T+i)x_n}^2$, la suite $(x_n)$ est de Cauchy
(car $((T+i)x_n)$ converge), $x_n \to x$, et $Tx_n \to y - ix$. Comme $T$ est fermé,
$x \in \dom(T)$ et $(T+i)x = y$. Donc $\im(T+i)$ est fermé.

Si $\im(T+i) \neq H$, il existe $z \neq 0$ avec $z \perp \im(T+i)$. Pour tout
$x \in \dom(T)$~: $\inner{(T+i)x}{z} = 0$, donc $\inner{Tx}{z} = -i\inner{x}{z}$.
Cela signifie $z \in \dom(T^*)$ et $T^*z = iz$, d'où $z \in \Ker(T^* - i) = \{0\}$~:
contradiction. Donc $\im(T+i) = H$. Même argument pour $T - i$.

\textbf{(iii) $\Rightarrow$ (i).} Soit $y \in \dom(T^*)$. Comme $\im(T+i) = H$, il
existe $x_0 \in \dom(T)$ tel que $(T+i)x_0 = (T^*+i)y$. Comme $T \subset T^*$~:
$(T^*+i)x_0 = (T^*+i)y$, d'où $T^*(y - x_0) = -i(y - x_0)$. Or
$\Ker(T^* + i) = \im(T - i)^\perp = H^\perp = \{0\}$ (car $\im(T-i) = H$). Donc
$y = x_0 \in \dom(T)$. Ainsi $\dom(T^*) \subset \dom(T)$, d'où $T = T^*$.
\end{proof}

\begin{exemple}[Le laplacien $-\Delta$]
\label{ex:laplacien}
\index{laplacien}
\index{opérateur!laplacien}
Sur $\Omega \subset \R^d$ ouvert borné régulier, l'opérateur $-\Delta$ avec domaine
$\dom(-\Delta) = H^2(\Omega) \cap H^1_0(\Omega)$ est auto-adjoint positif sur
$L^2(\Omega)$.
\end{exemple}

\begin{exemple}[{L'opérateur $i\frac{\dd}{\dd x}$}]
\label{ex:derivee-autoadjoint}
\index{opérateur!dérivée}
Sur $L^2(\R)$, l'opérateur $P = -i\frac{\dd}{\dd x}$ (opérateur d'impulsion en
mécanique quantique) avec domaine $\dom(P) = H^1(\R)$ est auto-adjoint.

Sur $L^2([0,1])$, la situation est plus subtile~: l'opérateur $-i\frac{\dd}{\dd x}$ avec
domaine $C_c^\infty((0,1))$ est symétrique mais pas auto-adjoint. Ses extensions
auto-adjointes sont paramétrées par les conditions aux bords
$f(1) = e^{i\theta}f(0)$, $\theta \in [0, 2\pi)$.
\index{extensions auto-adjointes}
\end{exemple}

%-----------------------------------------------------------------------
\section{Théorème spectral pour les opérateurs auto-adjoints non bornés}
\index{théorème spectral!opérateurs non bornés}
%-----------------------------------------------------------------------

\begin{theoreme}[Théorème spectral --- opérateurs non bornés]
\label{thm:spectral-non-borne}
Soit $T$ un opérateur auto-adjoint (non nécessairement borné) sur $H$. Il existe une
unique mesure spectrale $E$ sur $(\R, \mathcal{B}(\R))$, concentrée sur $\sigma(T)$,
telle que~:
\[
  T = \int_\R \lambda \, \dd E(\lambda)
\]
avec $\dom(T) = \{x \in H : \int_\R \lambda^2 \, \dd\mu_x(\lambda) < \infty\}$. Plus
généralement, pour toute fonction borélienne $f : \R \to \C$~:
\[
  f(T) = \int_\R f(\lambda) \, \dd E(\lambda), \quad
  \dom(f(T)) = \Bigl\{x \in H : \int_\R \abs{f(\lambda)}^2\,\dd\mu_x(\lambda) < \infty\Bigr\}.
\]
\end{theoreme}

Nous omettons la preuve complète, qui utilise la transformée de Cayley
\index{Cayley!transformée de}
$U = (T - i)(T + i)^{-1}$, qui est un opérateur unitaire. On applique le théorème
spectral pour les opérateurs unitaires à $U$, puis on revient à $T$ par transformation
inverse.

%-----------------------------------------------------------------------
\section{Application~: mécanique quantique}
\index{mécanique quantique}
%-----------------------------------------------------------------------

En mécanique quantique, un système physique est décrit par un espace de Hilbert $H$
(l'espace des états), et les observables physiques sont représentées par des opérateurs
auto-adjoints sur $H$.

\begin{remarque}[Postulats de la mécanique quantique]
\index{mécanique quantique!postulats}
\begin{enumerate}[label=(\roman*)]
  \item L'\emph{état} d'un système est un vecteur unitaire $\psi \in H$ (à une phase
        près).
  \item Une \emph{observable} est un opérateur auto-adjoint $A$ sur $H$.
  \item La \emph{valeur moyenne} de $A$ dans l'état $\psi$ est
        $\inner{A\psi}{\psi}$.
  \item Les \emph{valeurs possibles} d'une mesure de $A$ sont les éléments de $\sigma(A)$.
  \item La \emph{probabilité} d'obtenir un résultat dans l'ensemble borélien $B$ est
        $\norm{E(B)\psi}^2 = \mu_\psi(B)$.
  \item Après une mesure donnant un résultat dans $B$, l'état devient
        $E(B)\psi / \norm{E(B)\psi}$ (\emph{réduction du paquet d'ondes}).
        \index{réduction du paquet d'ondes}
\end{enumerate}
\end{remarque}

\begin{exemple}[Position et impulsion]
\index{opérateur!position}
\index{opérateur!impulsion}
Sur $H = L^2(\R)$~:
\begin{itemize}
  \item L'\emph{opérateur position}~: $(Qf)(x) = xf(x)$, $\dom(Q) = \{f \in L^2 : xf \in L^2\}$.
  \item L'\emph{opérateur impulsion}~: $(Pf)(x) = -i\hbar f'(x)$, $\dom(P) = H^1(\R)$.
\end{itemize}
Ces opérateurs satisfont la \emph{relation de commutation canonique}\index{relation de commutation canonique}~:
$[Q, P] = QP - PQ = i\hbar\,\Id$ (sur un domaine dense approprié).
Le spectre de $Q$ et de $P$ est $\R$ tout entier (spectre purement continu).
\end{exemple}

\begin{exemple}[Hamiltonien de l'oscillateur harmonique]
\index{oscillateur harmonique}
\index{Hamiltonien}
L'opérateur $H_{\mathrm{osc}} = -\frac{\dd^2}{\dd x^2} + x^2$ (avec $\hbar = m = \omega = 1$) est
auto-adjoint sur $L^2(\R)$. Son spectre est purement ponctuel~:
$\sigma(H_{\mathrm{osc}}) = \{2n+1 : n \in \N\}$, avec des fonctions propres données par
les fonctions de Hermite\index{Hermite!fonctions de}.
\end{exemple}

%-----------------------------------------------------------------------
\section{Exercices}
%-----------------------------------------------------------------------

\begin{exercice}[$\star$]
\index{exercice!spectre auto-adjoint}
Soit $T \in \mathcal{L}(H)$ auto-adjoint. Montrer que $\sigma(T) \subset [m, M]$ où
$m = \inf_{\norm{x}=1} \inner{Tx}{x}$ et $M = \sup_{\norm{x}=1} \inner{Tx}{x}$, et
que $m, M \in \sigma(T)$.
\end{exercice}

\begin{exercice}[$\star$]
\index{exercice!projecteur spectral}
Soit $T$ auto-adjoint borné et $E$ sa mesure spectrale. Montrer que
$\lambda \in \sigma(T)$ si et seulement si $E((a,b)) \neq 0$ pour tout intervalle ouvert
$(a,b)$ contenant $\lambda$.
\end{exercice}

\begin{exercice}[$\star\star$]
\index{exercice!calcul fonctionnel}
Soit $T$ auto-adjoint borné avec $\sigma(T) = \{a, b\}$ (deux points). Montrer que $T$
est de la forme $T = aP + b(\Id - P)$ où $P$ est un projecteur orthogonal. Calculer la
mesure spectrale.
\end{exercice}

\begin{exercice}[$\star\star$]
\index{exercice!racine carrée}
Soit $T \in \mathcal{L}(H)$ positif. Montrer que $\Ker(T^{1/2}) = \Ker(T)$ et
$\overline{\im(T^{1/2})} = \overline{\im(T)}$.
\end{exercice}

\begin{exercice}[$\star\star$]
\index{exercice!décomposition polaire}
Soit $T \in \mathcal{L}(H)$ et $T = U\abs{T}$ sa décomposition polaire. Montrer que
$T^* = U^*\abs{T^*}$ n'est pas en général la décomposition polaire de $T^*$, mais que
$T^* = V\abs{T^*}$ avec $V = U^*$.
\end{exercice}

\begin{exercice}[$\star\star\star$]
\index{exercice!critère auto-adjoint}
\index{Nelson (critère de)}
(\emph{Critère de Nelson.}) Soit $T$ symétrique sur $H$. On dit que $x \in \dom^\infty(T) = \bigcap_{n \ge 1} \dom(T^n)$ est un \emph{vecteur analytique} si
$\sum_{n=0}^\infty \frac{\norm{T^n x}}{n!}\,t^n < \infty$ pour un $t > 0$. Soit
$\mathcal{D} \subset \dom^\infty(T)$ un sous-espace dense de vecteurs analytiques.
Montrer que $T$ est essentiellement auto-adjoint (i.e., $\overline{T}$ est auto-adjoint).
\end{exercice}

\begin{exercice}[$\star\star$]
\index{exercice!opérateur de multiplication}
\index{opérateur de multiplication}
Soit $(\Omega, \mu)$ un espace mesuré $\sigma$-fini et $\varphi : \Omega \to \R$ mesurable.
L'opérateur de multiplication $M_\varphi f = \varphi f$ sur $L^2(\Omega, \mu)$ avec domaine
$\dom(M_\varphi) = \{f \in L^2 : \varphi f \in L^2\}$ est auto-adjoint. Déterminer sa
mesure spectrale et son spectre.
\end{exercice}

\begin{exercice}[$\star\star\star$]
\index{exercice!Stone (théorème de)}
\index{Stone!théorème de}
(\emph{Théorème de Stone.}) Soit $(U_t)_{t \in \R}$ un groupe fortement continu
d'opérateurs unitaires sur $H$. Montrer qu'il existe un unique opérateur auto-adjoint $A$
tel que $U_t = e^{itA}$ pour tout $t \in \R$.
\end{exercice}


%%%%%%%%%%%%%%%%%%%%%%%%%%%%%%%%%%%%%%%%%%%%%%%%%%%%%%%%%%%%%%%%%%%%%%%
%  CHAPITRE 11 — Espaces Lp et Dualité
%%%%%%%%%%%%%%%%%%%%%%%%%%%%%%%%%%%%%%%%%%%%%%%%%%%%%%%%%%%%%%%%%%%%%%%
\chapter{Espaces $L^p$ et Dualité}
\label{chap:Lp}
\minitoc

\section{Inégalités fondamentales}
\index{espace Lp@espace $L^p$}

Soit $(\Omega, \mathcal{A}, \mu)$ un espace mesuré. Pour $1 \le p \le \infty$, on note
$L^p(\Omega) = L^p(\Omega, \mathcal{A}, \mu)$ l'espace des (classes de) fonctions
mesurables $f : \Omega \to \K$ telles que $\norm{f}_p < \infty$, où
\[
  \norm{f}_p = \begin{cases}
    \displaystyle \left(\int_\Omega \abs{f}^p\,\dd\mu\right)^{1/p} & \text{si } 1 \le p < \infty, \\[6pt]
    \displaystyle \inf\{C \ge 0 : \abs{f} \le C \text{ $\mu$-p.p.}\} & \text{si } p = \infty.
  \end{cases}
\]
\index{norme Lp@norme $L^p$}

On rappelle que les \emph{exposants conjugués}\index{exposants conjugués} $p$ et $q$
satisfont $\frac{1}{p} + \frac{1}{q} = 1$, avec la convention $q = \infty$ si $p = 1$.

\begin{theoreme}[Inégalité de Young]
\label{thm:young-ineq}
\index{Young!inégalité de}
Pour tous $a, b \ge 0$ et $1 < p < \infty$ avec $\frac{1}{p} + \frac{1}{q} = 1$~:
\[
  ab \le \frac{a^p}{p} + \frac{b^q}{q}.
\]
L'égalité a lieu si et seulement si $a^p = b^q$.
\end{theoreme}

\begin{proof}
La fonction $\varphi(t) = \log t$ est concave. Par la convexité de l'exponentielle, pour
$s, t > 0$~:
\[
  st = e^{\log s + \log t} = e^{\frac{1}{p}(p\log s) + \frac{1}{q}(q\log t)}
    \le \frac{1}{p}e^{p\log s} + \frac{1}{q}e^{q\log t} = \frac{s^p}{p} + \frac{t^q}{q}
\]
par l'inégalité de Jensen\index{Jensen!inégalité de} (ou directement par convexité de
$e^x$). On applique avec $s = a$, $t = b$.
\end{proof}

\begin{theoreme}[Inégalité de Hölder]
\label{thm:holder}
\index{Hölder!inégalité de}
Soient $1 \le p \le \infty$ et $q$ son exposant conjugué. Pour tous
$f \in L^p(\Omega)$ et $g \in L^q(\Omega)$~:
\[
  \norm{fg}_1 \le \norm{f}_p \, \norm{g}_q.
\]
\end{theoreme}

\begin{proof}
Les cas $p = 1$ ($q = \infty$) et $p = \infty$ ($q = 1$) sont immédiats.

Supposons $1 < p < \infty$. Si $\norm{f}_p = 0$ ou $\norm{g}_q = 0$, le résultat est
trivial. Sinon, posons $\tilde f = f/\norm{f}_p$ et $\tilde g = g/\norm{g}_q$. Par
l'inégalité de Young (théorème~\ref{thm:young-ineq})~:
\[
  \abs{\tilde f(x)} \cdot \abs{\tilde g(x)}
    \le \frac{\abs{\tilde f(x)}^p}{p} + \frac{\abs{\tilde g(x)}^q}{q}
\]
pour $\mu$-presque tout $x$. En intégrant~:
\[
  \int_\Omega \abs{\tilde f}\,\abs{\tilde g}\,\dd\mu
    \le \frac{1}{p}\int_\Omega \abs{\tilde f}^p\,\dd\mu
      + \frac{1}{q}\int_\Omega \abs{\tilde g}^q\,\dd\mu
    = \frac{1}{p} + \frac{1}{q} = 1.
\]
Donc $\int \abs{fg}\,\dd\mu \le \norm{f}_p\,\norm{g}_q$.
\end{proof}

\begin{theoreme}[Inégalité de Minkowski]
\label{thm:minkowski}
\index{Minkowski!inégalité de}
Pour $1 \le p \le \infty$ et $f, g \in L^p(\Omega)$~:
\[
  \norm{f + g}_p \le \norm{f}_p + \norm{g}_p.
\]
\end{theoreme}

\begin{proof}
Les cas $p = 1$ (inégalité triangulaire pour $\abs{\cdot}$) et $p = \infty$ sont
immédiats. Supposons $1 < p < \infty$. On a~:
\[
  \abs{f+g}^p \le \abs{f+g}^{p-1}\abs{f} + \abs{f+g}^{p-1}\abs{g}.
\]
Comme $(p-1)q = p$, la fonction $\abs{f+g}^{p-1} \in L^q(\Omega)$ avec
$\norm{\abs{f+g}^{p-1}}_q = \norm{f+g}_p^{p/q}$. En appliquant Hölder à chaque terme~:
\begin{align*}
  \norm{f+g}_p^p &\le \norm{f}_p\,\norm{f+g}_p^{p/q} + \norm{g}_p\,\norm{f+g}_p^{p/q}\\
  &= (\norm{f}_p + \norm{g}_p)\,\norm{f+g}_p^{p/q}.
\end{align*}
Si $\norm{f+g}_p > 0$, on divise par $\norm{f+g}_p^{p/q}$ et on utilise
$p - p/q = 1$.
\end{proof}

%-----------------------------------------------------------------------
\section{Complétude de $L^p$~: théorème de Riesz--Fischer}
\index{Riesz--Fischer (théorème de)}
\index{espace Lp@espace $L^p$!complétude}
%-----------------------------------------------------------------------

\begin{theoreme}[Riesz--Fischer]
\label{thm:riesz-fischer}
Pour $1 \le p \le \infty$, l'espace $L^p(\Omega, \mu)$ est un espace de Banach.
\end{theoreme}

\begin{proof}
Seule la complétude reste à montrer (la vérification que $\norm{\cdot}_p$ est une norme
est classique via Minkowski). Traitons le cas $1 \le p < \infty$.

\textbf{Méthode par sous-suites.} Soit $(f_n)$ une suite de Cauchy dans $L^p$. Il suffit
de montrer qu'elle admet une sous-suite convergente (car une suite de Cauchy ayant une
sous-suite convergente converge elle-même).

Extrayons une sous-suite $(f_{n_k})$ telle que
$\norm{f_{n_{k+1}} - f_{n_k}}_p \le 2^{-k}$. Posons
\[
  g_N(x) = \sum_{k=1}^{N} \abs{f_{n_{k+1}}(x) - f_{n_k}(x)}, \quad
  g(x) = \sum_{k=1}^{\infty} \abs{f_{n_{k+1}}(x) - f_{n_k}(x)}.
\]
Par Minkowski~:
\[
  \norm{g_N}_p \le \sum_{k=1}^N \norm{f_{n_{k+1}} - f_{n_k}}_p \le \sum_{k=1}^N 2^{-k} \le 1.
\]
Par le lemme de Fatou\index{Fatou!lemme de}~:
$\int g^p\,\dd\mu \le \liminf_N \int g_N^p\,\dd\mu \le 1$. Donc $g < \infty$
$\mu$-presque partout, ce qui signifie que la série télescopique
\[
  f_{n_1}(x) + \sum_{k=1}^\infty \bigl(f_{n_{k+1}}(x) - f_{n_k}(x)\bigr)
\]
converge absolument pour $\mu$-presque tout $x$. Notons $f(x)$ sa limite~: $f_{n_k}(x) \to f(x)$
$\mu$-p.p.

Il reste à montrer que $f \in L^p$ et $\norm{f_{n_k} - f}_p \to 0$.
On a $\abs{f_{n_k}} \le \abs{f_{n_1}} + g$ $\mu$-p.p., donc
$\abs{f} \le \abs{f_{n_1}} + g \in L^p$, d'où $f \in L^p$.
De plus, $\abs{f_{n_k} - f}^p \le (2(\abs{f_{n_1}} + g))^p \in L^1$. Par convergence
dominée\index{convergence dominée}, $\norm{f_{n_k} - f}_p \to 0$.

Le cas $p = \infty$ est plus simple~: si $(f_n)$ est de Cauchy dans $L^\infty$, il
existe un ensemble de mesure nulle $N$ tel que $(f_n(x))$ soit de Cauchy dans $\K$ pour
tout $x \notin N$. La limite ponctuelle $f$ est dans $L^\infty$ et
$\norm{f_n - f}_\infty \to 0$.
\end{proof}

%-----------------------------------------------------------------------
\section{Densité de $C_c^\infty$ dans $L^p$}
\index{densité!de Cc dans Lp@de $C_c^\infty$ dans $L^p$}
\index{fonctions test}
%-----------------------------------------------------------------------

\begin{theoreme}
\label{thm:densite-Cc-Lp}
Soit $\Omega \subset \R^d$ un ouvert et $1 \le p < \infty$. Alors $C_c^\infty(\Omega)$
est dense dans $L^p(\Omega)$ (pour la mesure de Lebesgue).
\end{theoreme}

\begin{proof}
On procède en trois étapes.

\textbf{Étape 1~: Les fonctions simples à support compact sont denses.}
Par définition de l'intégrale de Lebesgue, toute $f \in L^p(\Omega)$ non négative est
limite croissante de fonctions simples $s_n$. Par convergence dominée,
$\norm{f - s_n}_p \to 0$. En tronquant ($s_n \cdot \mathbf{1}_{B(0,n)}$), on obtient
des fonctions simples à support compact.

\textbf{Étape 2~: Les fonctions continues à support compact sont denses.}
\index{Lusin (théorème de)}
Par le théorème de Lusin, pour tout $\varepsilon > 0$, toute fonction mesurable $f$ à
support compact coïncide avec une fonction continue sur un ensemble dont le complémentaire
a mesure $< \varepsilon$. Cela permet d'approcher les fonctions simples à support compact
par des fonctions de $C_c(\Omega)$.

Alternativement, on approche $\mathbf{1}_K$ (pour $K$ compact $\subset \Omega$) par des
fonctions de $C_c(\Omega)$ via le lemme d'Urysohn\index{Urysohn!lemme d'}.

\textbf{Étape 3~: Régularisation par convolution.}
\index{convolution!régularisation}
\index{approximation de l'identité}
Soit $\rho \in C_c^\infty(\R^d)$ avec $\rho \ge 0$, $\supp \rho \subset B(0,1)$ et
$\int \rho = 1$. Posons $\rho_\varepsilon(x) = \varepsilon^{-d}\rho(x/\varepsilon)$.
Pour $f \in C_c(\Omega)$, la fonction $f * \rho_\varepsilon \in C^\infty(\R^d)$ et
$\supp(f * \rho_\varepsilon) \subset \supp(f) + B(0,\varepsilon)$. Pour $\varepsilon$
assez petit, $f * \rho_\varepsilon \in C_c^\infty(\Omega)$ et
$\norm{f * \rho_\varepsilon - f}_p \to 0$ par un argument standard de convergence dominée.
\end{proof}

\begin{remarque}
Pour $p = \infty$, $C_c^\infty(\Omega)$ n'est pas dense dans $L^\infty(\Omega)$~: en
effet, l'adhérence de $C_c(\Omega)$ dans $L^\infty$ est $C_0(\Omega)$ (fonctions
continues tendant vers $0$ à l'infini), qui est strictement inclus dans $L^\infty$.
\end{remarque}

%-----------------------------------------------------------------------
\section{Convolution et approximation de l'identité}
\index{convolution}
\index{approximation de l'identité}
%-----------------------------------------------------------------------

\begin{definition}[Convolution]
\label{def:convolution}
\index{convolution!définition}
Soient $f, g : \R^d \to \K$ mesurables. Le \emph{produit de convolution} est défini (lorsqu'il existe) par~:
\[
  (f * g)(x) = \int_{\R^d} f(x - y)\, g(y)\, \dd y.
\]
\end{definition}

\begin{theoreme}[Inégalité de Young pour la convolution]
\label{thm:young-convolution}
\index{Young!inégalité de convolution}
Soient $1 \le p, q, r \le \infty$ avec $\frac{1}{p} + \frac{1}{q} = 1 + \frac{1}{r}$.
Si $f \in L^p(\R^d)$ et $g \in L^q(\R^d)$, alors $f * g \in L^r(\R^d)$ et
\[
  \norm{f * g}_r \le \norm{f}_p \, \norm{g}_q.
\]
\end{theoreme}

\begin{definition}[Approximation de l'identité]
\label{def:approx-identite}
\index{approximation de l'identité!définition}
Une \emph{approximation de l'identité} est une famille $(\rho_\varepsilon)_{\varepsilon > 0}$
de fonctions dans $L^1(\R^d)$ telle que~:
\begin{enumerate}[label=(\roman*)]
  \item $\int_{\R^d} \rho_\varepsilon\,\dd x = 1$ pour tout $\varepsilon > 0$.
  \item $\sup_\varepsilon \norm{\rho_\varepsilon}_1 < \infty$.
  \item Pour tout $\delta > 0$~:
        $\int_{\abs{x} > \delta} \abs{\rho_\varepsilon(x)}\,\dd x \to 0$ quand
        $\varepsilon \to 0$.
\end{enumerate}
\end{definition}

\begin{proposition}
\label{prop:approx-identite-conv}
Soit $(\rho_\varepsilon)$ une approximation de l'identité et $1 \le p < \infty$. Pour
tout $f \in L^p(\R^d)$~:
\[
  \norm{f * \rho_\varepsilon - f}_p \to 0 \quad (\varepsilon \to 0).
\]
\end{proposition}

\begin{proof}
On a $(f * \rho_\varepsilon)(x) - f(x) = \int [f(x-y) - f(x)]\,\rho_\varepsilon(y)\,\dd y$.
Par Minkowski intégral\index{Minkowski!inégalité intégrale}~:
\[
  \norm{f * \rho_\varepsilon - f}_p
    \le \int_{\R^d} \norm{\tau_y f - f}_p\, \abs{\rho_\varepsilon(y)}\,\dd y
\]
où $\tau_y f(x) = f(x-y)$. On sait que $y \mapsto \norm{\tau_y f - f}_p$ est continue et
vaut $0$ en $y = 0$. Soit $\delta > 0$ tel que $\norm{\tau_y f - f}_p < \varepsilon'$ pour
$\abs{y} < \delta$. Alors~:
\begin{align*}
  \norm{f*\rho_\varepsilon - f}_p
    &\le \varepsilon' \int_{\abs{y} < \delta} \abs{\rho_\varepsilon}\,\dd y
      + 2\norm{f}_p \int_{\abs{y} \ge \delta} \abs{\rho_\varepsilon}\,\dd y \\
    &\le \varepsilon' \norm{\rho_\varepsilon}_1
      + 2\norm{f}_p \int_{\abs{y} \ge \delta} \abs{\rho_\varepsilon}\,\dd y
    \to \varepsilon' \cdot C
\end{align*}
quand $\varepsilon \to 0$. Comme $\varepsilon'$ est arbitraire, on conclut.
\end{proof}

%-----------------------------------------------------------------------
\section{Dual de $L^p$}
\index{dual!de Lp@de $L^p$}
\index{espace Lp@espace $L^p$!dualité}
%-----------------------------------------------------------------------

\begin{theoreme}[Représentation de Riesz pour $L^p$]
\label{thm:dual-Lp}
\index{Riesz!théorème de représentation pour $L^p$}
Soit $(\Omega, \mathcal{A}, \mu)$ un espace mesuré $\sigma$-fini et $1 \le p < \infty$.
Alors l'application
\[
  \Phi : L^q(\Omega) \to (L^p(\Omega))', \quad
  \Phi(g)(f) = \int_\Omega f\,\bar g\,\dd\mu
\]
(où $\frac{1}{p} + \frac{1}{q} = 1$) est une isométrie isomorphe~:
$(L^p(\Omega))' \cong L^q(\Omega)$.
\end{theoreme}

\begin{proof}
\textbf{$\Phi$ est une isométrie.} Par Hölder, $\abs{\Phi(g)(f)} \le \norm{f}_p\norm{g}_q$,
donc $\norm{\Phi(g)} \le \norm{g}_q$. Pour l'égalité, si $g \neq 0$, posons
\[
  f_0 = \frac{\bar g}{\abs{g}} \cdot \abs{g}^{q/p}
    = \frac{\bar g}{\abs{g}} \cdot \abs{g}^{q-1}
\]
(avec $f_0 = 0$ là où $g = 0$). Alors $\abs{f_0}^p = \abs{g}^{(q-1)p} = \abs{g}^q$,
donc $\norm{f_0}_p = \norm{g}_q^{q/p}$. Et
$\Phi(g)(f_0) = \int \abs{g}^q\,\dd\mu = \norm{g}_q^q$. Donc
$\norm{\Phi(g)} \ge \norm{g}_q^q / \norm{g}_q^{q/p} = \norm{g}_q^{q - q/p} = \norm{g}_q$.

\textbf{$\Phi$ est surjective.} Soit $\ell \in (L^p(\Omega))'$. On construit $g$ tel que
$\ell = \Phi(g)$.

\textbf{Cas $p = 2$.} C'est le théorème de représentation de Riesz--Fréchet dans l'espace
de Hilbert $L^2$.

\textbf{Cas $1 \le p < \infty$ général, $\mu$ $\sigma$-finie.}
Définissons une mesure signée $\nu$ sur $(\Omega, \mathcal{A})$ par
$\nu(A) = \ell(\mathbf{1}_A)$ pour tout $A \in \mathcal{A}$ de mesure finie. Vérifions
que $\nu$ est bien définie et $\sigma$-additive~:
\begin{itemize}
  \item $\mathbf{1}_A \in L^p$ car $\mu(A) < \infty$, donc $\nu(A) = \ell(\mathbf{1}_A)$ est défini.
  \item $\abs{\nu(A)} = \abs{\ell(\mathbf{1}_A)} \le \norm{\ell}\,\norm{\mathbf{1}_A}_p = \norm{\ell}\,\mu(A)^{1/p}$. En particulier, $\nu \ll \mu$.
  \item La $\sigma$-additivité de $\nu$ résulte de la continuité de $\ell$ et de la
        convergence dans $L^p$ des sommes partielles $\sum_{k=1}^N \mathbf{1}_{A_k}$.
\end{itemize}

Par le théorème de Radon--Nikodym\index{Radon--Nikodym (théorème de)} (valable car $\mu$ est
$\sigma$-finie), il existe $g \in L^1_{\mathrm{loc}}(\Omega)$ telle que
$\nu(A) = \int_A g\,\dd\mu$ pour tout ensemble mesurable $A$ de mesure finie. Donc~:
\[
  \ell(s) = \int_\Omega s\, g\,\dd\mu
\]
pour toute fonction simple $s$ à support de mesure finie. Par densité des fonctions
simples dans $L^p$ et continuité de $\ell$~:
\[
  \ell(f) = \int_\Omega f\,g\,\dd\mu, \quad \forall f \in L^p(\Omega).
\]

Il reste à montrer que $g \in L^q(\Omega)$. Pour $1 < p < \infty$~: prenons
$f_n = \bar g / \abs{g} \cdot \min(\abs{g}^{q-1}, n) \cdot \mathbf{1}_{\Omega_n}$ où
$\Omega_n \uparrow \Omega$ avec $\mu(\Omega_n) < \infty$. Alors $f_n \in L^p$ et
\[
  \int_{\Omega_n} \min(\abs{g}^q, n\abs{g})\,\dd\mu = \ell(f_n) \le \norm{\ell}\,\norm{f_n}_p.
\]
Un calcul montre $\norm{f_n}_p^p \le \int \min(\abs{g}^q, n\abs{g})\,\dd\mu$, d'où
$\norm{f_n}_p^{p} \le \norm{\ell} \cdot \norm{f_n}_p$, c'est-à-dire
$\norm{f_n}_p^{p-1} \le \norm{\ell}$. En passant à la limite ($n \to \infty$) par
convergence monotone, $\norm{g}_q^{q/p} \le \norm{\ell}$, donc $g \in L^q$.

Pour $p = 1$ ($q = \infty$)~: pour tout ensemble mesurable $A$ avec $\mu(A) > 0$~:
\[
  \frac{1}{\mu(A)}\abs[\bigg]{\int_A g\,\dd\mu}
    = \frac{\abs{\ell(\mathbf{1}_A)}}{\mu(A)}
    \le \frac{\norm{\ell}\,\mu(A)}{\mu(A)} = \norm{\ell}.
\]
Par le théorème de Lebesgue sur la différentiation\index{différentiation de Lebesgue},
$\abs{g(x)} \le \norm{\ell}$ pour $\mu$-presque tout $x$. Donc $g \in L^\infty$ et
$\norm{g}_\infty \le \norm{\ell}$.
\end{proof}

\begin{remarque}
\index{dual!de L infini@de $L^\infty$}
\index{espace Linfini@espace $L^\infty$!dual}
Le dual de $L^\infty(\Omega)$ est strictement plus grand que $L^1(\Omega)$~: c'est
l'espace des mesures finiment additives absolument continues $\mathrm{ba}(\Omega, \mu)$.
En particulier, $L^\infty$ n'est pas réflexif\index{réflexif!L infini@$L^\infty$ non}.
\end{remarque}

%-----------------------------------------------------------------------
\section{$L^2$ comme espace de Hilbert}
\index{espace L2@espace $L^2$!espace de Hilbert}
%-----------------------------------------------------------------------

\begin{remarque}
L'espace $L^2(\Omega, \mu)$ est un espace de Hilbert pour le produit scalaire
\[
  \inner{f}{g} = \int_\Omega f\,\bar g\,\dd\mu.
\]
C'est le seul espace $L^p$ qui soit un espace de Hilbert ($p = q = 2$). La dualité
$(L^2)' \cong L^2$ coïncide avec le théorème de représentation de Riesz--Fréchet pour les
espaces de Hilbert.
\end{remarque}

\begin{exemple}[Bases hilbertiennes classiques]
\label{ex:bases-L2}
\index{base hilbertienne!exemples dans $L^2$}
\begin{enumerate}[label=(\roman*)]
  \item Sur $L^2([0, 2\pi])$~: la famille $\{e^{inx}/\sqrt{2\pi}\}_{n \in \Z}$ est une
        base hilbertienne (système trigonométrique)\index{système trigonométrique}.
  \item Sur $L^2(\R, e^{-x^2}\,\dd x)$~: les polynômes de Hermite normalisés forment une
        base hilbertienne\index{Hermite!polynômes de}.
  \item Sur $L^2([-1,1])$~: les polynômes de Legendre normalisés forment une base
        hilbertienne\index{Legendre!polynômes de}.
\end{enumerate}
\end{exemple}

%-----------------------------------------------------------------------
\section{Théorème de Radon--Nikodym}
\index{Radon--Nikodym (théorème de)}
%-----------------------------------------------------------------------

Le théorème de Radon--Nikodym a été utilisé dans la preuve du dual de $L^p$. Énonçons-le
pour référence.

\begin{theoreme}[Radon--Nikodym]
\label{thm:radon-nikodym}
\index{Radon--Nikodym (théorème de)!énoncé}
\index{dérivée de Radon--Nikodym}
Soit $(\Omega, \mathcal{A}, \mu)$ un espace mesuré $\sigma$-fini et $\nu$ une mesure
signée $\sigma$-finie sur $(\Omega, \mathcal{A})$ avec $\nu \ll \mu$ (absolument
continue). Alors il existe une unique (à modification $\mu$-p.p. près) fonction mesurable
$g : \Omega \to \R$ (ou $\C$) telle que~:
\[
  \nu(A) = \int_A g\,\dd\mu, \quad \forall A \in \mathcal{A}.
\]
La fonction $g$ s'appelle la \emph{dérivée de Radon--Nikodym} de $\nu$ par rapport à
$\mu$ et se note $g = \frac{\dd\nu}{\dd\mu}$.
\end{theoreme}

Nous admettons ce résultat (voir par exemple Rudin \cite{rudin-real-complex} pour une
preuve complète). La preuve classique utilise le théorème de représentation de Riesz dans
$L^2$ ou le lemme de Von Neumann\index{Von Neumann!lemme de}.

%-----------------------------------------------------------------------
\section{Interpolation~: théorème de Riesz--Thorin}
\index{Riesz--Thorin (théorème de)}
\index{interpolation}
%-----------------------------------------------------------------------

\begin{theoreme}[Riesz--Thorin]
\label{thm:riesz-thorin}
\index{Riesz--Thorin (théorème de)!énoncé}
\index{interpolation complexe}
Soient $(\Omega_1, \mu_1)$ et $(\Omega_2, \mu_2)$ deux espaces mesurés $\sigma$-finis.
Soient $1 \le p_0, p_1, q_0, q_1 \le \infty$ et $T$ un opérateur linéaire défini sur
$L^{p_0}(\Omega_1) + L^{p_1}(\Omega_1)$ tel que~:
\[
  \norm{Tf}_{q_0} \le M_0 \norm{f}_{p_0}, \quad
  \norm{Tf}_{q_1} \le M_1 \norm{f}_{p_1}.
\]
Alors pour tout $0 < \theta < 1$, en posant
\[
  \frac{1}{p} = \frac{1-\theta}{p_0} + \frac{\theta}{p_1}, \quad
  \frac{1}{q} = \frac{1-\theta}{q_0} + \frac{\theta}{q_1},
\]
on a $T : L^p(\Omega_1) \to L^q(\Omega_2)$ avec
\[
  \norm{Tf}_q \le M_0^{1-\theta}\, M_1^\theta \, \norm{f}_p.
\]
\end{theoreme}

\begin{proof}
La preuve utilise le \emph{théorème des trois droites}\index{trois droites!théorème des}
de Hadamard. On se ramène au cas $M_0 = M_1 = 1$ par normalisation.

Soit $f = \sum_{j} a_j \mathbf{1}_{A_j}$ une fonction simple sur $\Omega_1$ et
$g = \sum_{k} b_k \mathbf{1}_{B_k}$ une fonction simple sur $\Omega_2$ avec
$\norm{f}_p = \norm{g}_{q'} = 1$ (où $q'$ est le conjugué de $q$). Il suffit de montrer
que $\abs{\int (Tf)\bar g\,\dd\mu_2} \le 1$.

Pour $z \in S = \{z \in \C : 0 \le \Re z \le 1\}$ (la bande), définissons~:
\[
  f_z(x) = \sum_j \abs{a_j}^{p(\frac{1-z}{p_0} + \frac{z}{p_1})}
    \frac{a_j}{\abs{a_j}} \mathbf{1}_{A_j}(x), \quad
  g_z(x) = \sum_k \abs{b_k}^{q'(\frac{1-z}{q_0'} + \frac{z}{q_1'})}
    \frac{b_k}{\abs{b_k}} \mathbf{1}_{B_k}(x).
\]
Posons $F(z) = \int_{\Omega_2} (Tf_z)\,\overline{g_z}\,\dd\mu_2$. La fonction $F$ est
holomorphe sur l'intérieur de $S$, continue et bornée sur $S$.

Sur la droite $\Re z = 0$~: $\norm{f_{it}}_{p_0} = 1$ et $\norm{g_{it}}_{q_0'} = 1$,
donc $\abs{F(it)} \le \norm{Tf_{it}}_{q_0} \le M_0 = 1$.

Sur la droite $\Re z = 1$~: $\norm{f_{1+it}}_{p_1} = 1$ et $\norm{g_{1+it}}_{q_1'} = 1$,
donc $\abs{F(1+it)} \le 1$.

Par le théorème des trois droites~: $\abs{F(z)} \le 1$ pour tout $z \in S$. En
particulier, $\abs{F(\theta)} = \abs{\int (Tf)\bar g\,\dd\mu_2} \le 1$.

Par densité des fonctions simples et dualité $(L^q)' = L^{q'}$~:
$\norm{Tf}_q \le \norm{f}_p$.
\end{proof}

%-----------------------------------------------------------------------
\section{Application~: transformée de Fourier}
\index{transformée de Fourier}
%-----------------------------------------------------------------------

\begin{definition}[Transformée de Fourier sur $L^1$]
\label{def:fourier-L1}
\index{transformée de Fourier!sur $L^1$}
Pour $f \in L^1(\R^d)$, la \emph{transformée de Fourier} de $f$ est~:
\[
  \hat f(\xi) = \mathcal{F}f(\xi) = \int_{\R^d} f(x)\, e^{-2\pi i x \cdot \xi}\,\dd x,
  \quad \xi \in \R^d.
\]
\end{definition}

\begin{proposition}
\label{prop:fourier-L1}
\index{Riemann--Lebesgue (lemme de)}
Pour $f \in L^1(\R^d)$~: $\hat f \in C_0(\R^d)$ (continue, tendant vers $0$ à l'infini)
et $\norm{\hat f}_\infty \le \norm{f}_1$ (lemme de Riemann--Lebesgue).
\end{proposition}

\begin{theoreme}[Plancherel]
\label{thm:plancherel}
\index{Plancherel (théorème de)}
\index{transformée de Fourier!sur $L^2$}
La transformée de Fourier se prolonge en une isométrie surjective
$\mathcal{F} : L^2(\R^d) \to L^2(\R^d)$~:
\[
  \norm{\hat f}_2 = \norm{f}_2, \quad \forall f \in L^2(\R^d).
\]
\end{theoreme}

\begin{proof}
Sur $\mathcal{S}(\R^d)$ (espace de Schwartz\index{Schwartz!espace de}), on montre par
un calcul direct que $\norm{\hat f}_2 = \norm{f}_2$~: par Fubini,
\begin{align*}
  \int \abs{\hat f(\xi)}^2\,\dd\xi
    &= \int \hat f(\xi)\,\overline{\hat f(\xi)}\,\dd\xi
    = \int \hat f(\xi)\,\widehat{\bar f}(-\xi)\,\dd\xi \\
    &= \int f(x)\,\overline{f(x)}\,\dd x = \norm{f}_2^2
\end{align*}
(en utilisant la formule de Parseval\index{Parseval!formule de} pour les fonctions de
Schwartz). Comme $\mathcal{S}(\R^d)$ est dense dans $L^2(\R^d)$, l'isométrie se prolonge
par continuité. La surjectivité résulte de la formule d'inversion de Fourier sur
$\mathcal{S}$.
\end{proof}

\begin{corollaire}[Hausdorff--Young]
\label{cor:hausdorff-young}
\index{Hausdorff--Young (inégalité de)}
Pour $1 \le p \le 2$ et $f \in L^p(\R^d)$~:
$\norm{\hat f}_{p'} \le \norm{f}_p$ (où $\frac{1}{p} + \frac{1}{p'} = 1$).
\end{corollaire}

\begin{proof}
Par Riesz--Thorin~: $\mathcal{F} : L^1 \to L^\infty$ avec norme $\le 1$ et
$\mathcal{F} : L^2 \to L^2$ avec norme $= 1$. Par interpolation avec
$\theta = 2/p' = 2 - 2/p$~: $\mathcal{F} : L^p \to L^{p'}$ avec norme $\le 1$.
\end{proof}

%-----------------------------------------------------------------------
\section{Exercices}
%-----------------------------------------------------------------------

\begin{exercice}[$\star$]
\index{exercice!Hölder}
Montrer le cas d'égalité dans l'inégalité de Hölder~: $\norm{fg}_1 = \norm{f}_p\norm{g}_q$
si et seulement si $\abs{f}^p$ et $\abs{g}^q$ sont proportionnelles $\mu$-p.p.
\end{exercice}

\begin{exercice}[$\star$]
\index{exercice!inclusion Lp}
Soit $\mu(\Omega) < \infty$ et $1 \le p \le r \le \infty$. Montrer que
$L^r(\Omega) \subset L^p(\Omega)$ avec injection continue~:
$\norm{f}_p \le \mu(\Omega)^{1/p - 1/r}\norm{f}_r$.
\end{exercice}

\begin{exercice}[$\star\star$]
\index{exercice!Lp non réflexif}
Montrer que $L^1(\R)$ n'est pas réflexif.
\emph{Indication~: $(L^1)' = L^\infty$ et $(L^\infty)' \supsetneq L^1$.}
\end{exercice}

\begin{exercice}[$\star\star$]
\index{exercice!convergence dans Lp}
Soit $(f_n)$ une suite dans $L^p(\Omega)$ ($1 \le p < \infty$) convergeant vers $f$ dans
$L^p$. Montrer qu'il existe une sous-suite $(f_{n_k})$ et une fonction $g \in L^p$ telles
que $f_{n_k} \to f$ $\mu$-p.p.\ et $\abs{f_{n_k}} \le g$ $\mu$-p.p.\ pour tout $k$.
\end{exercice}

\begin{exercice}[$\star\star$]
\index{exercice!convolution}
Montrer que si $f \in L^1(\R^d)$ et $g \in L^p(\R^d)$ ($1 \le p \le \infty$), alors
$f * g \in L^p(\R^d)$ et $\norm{f*g}_p \le \norm{f}_1\norm{g}_p$.
\end{exercice}

\begin{exercice}[$\star\star\star$]
\index{exercice!Riesz--Thorin application}
Utiliser le théorème de Riesz--Thorin pour montrer l'inégalité de Hausdorff--Young
généralisée~: si $G$ est un groupe abélien localement compact et $1 \le p \le 2$, alors
la transformée de Fourier envoie $L^p(G)$ dans $L^{p'}(\hat G)$ continûment.
\end{exercice}

\begin{exercice}[$\star\star$]
\index{exercice!séparabilité Lp}
Montrer que $L^p(\R^d)$ est séparable pour $1 \le p < \infty$, mais que $L^\infty(\R^d)$
ne l'est pas.
\end{exercice}

\begin{exercice}[$\star\star$]
\index{exercice!convergence faible Lp}
\index{convergence faible!dans $L^p$}
Soit $1 < p < \infty$ et $(f_n)$ bornée dans $L^p(\Omega)$ avec $\mu(\Omega) < \infty$.
Montrer que si $f_n \to f$ $\mu$-p.p., alors $f_n \weakto f$ dans $L^p(\Omega)$.
\emph{Indication~: utiliser la réflexivité de $L^p$ et la convergence dominée.}
\end{exercice}


%%%%%%%%%%%%%%%%%%%%%%%%%%%%%%%%%%%%%%%%%%%%%%%%%%%%%%%%%%%%%%%%%%%%%%%
%  CHAPITRE 12 — Algèbres de Banach et C*-Algèbres — Introduction
%%%%%%%%%%%%%%%%%%%%%%%%%%%%%%%%%%%%%%%%%%%%%%%%%%%%%%%%%%%%%%%%%%%%%%%
\chapter{Algèbres de Banach et $C^*$-Algèbres~: Introduction}
\label{chap:cstar}
\minitoc

\section{Algèbres de Banach~: définition et exemples}
\index{algèbre de Banach}
\index{algèbre de Banach!définition}

\begin{definition}[Algèbre de Banach]
\label{def:alg-banach}
Une \emph{algèbre de Banach} est un espace de Banach $(A, \norm{\cdot})$ muni d'une
multiplication $A \times A \to A$, $(a,b) \mapsto ab$, qui est bilinéaire, associative et
satisfait $\norm{ab} \le \norm{a}\,\norm{b}$ pour tous $a, b \in A$.

Si $A$ possède un élément unité $e$ avec $\norm{e} = 1$, on dit que $A$ est une algèbre
de Banach \emph{unitaire} (ou \emph{unital}).
\index{algèbre de Banach!unitaire}
\end{definition}

\begin{exemple}[Exemples fondamentaux]
\label{ex:alg-banach}
\begin{enumerate}[label=(\roman*)]
  \item $\mathcal{L}(E)$, l'algèbre des opérateurs bornés sur un espace de Banach $E$,
        avec la composition et la norme d'opérateur. L'unité est $\Id$.
        \index{algèbre!des opérateurs bornés}
  \item $C(K)$, l'algèbre des fonctions continues sur un espace compact $K$, avec la
        multiplication ponctuelle et la norme $\norm{f}_\infty = \sup_{x \in K}\abs{f(x)}$.
        \index{algèbre!C(K)@$C(K)$}
  \item $\ell^1(\Z)$ avec le produit de convolution~:
        $(a * b)(n) = \sum_{k \in \Z} a(k)\,b(n-k)$. C'est une algèbre de Banach
        commutative unitaire (l'unité est $\delta_0$).
        \index{algèbre!l1(Z)@$\ell^1(\Z)$}
        \index{convolution!sur $\ell^1(\Z)$}
  \item $L^1(G)$ pour un groupe localement compact $G$, avec le produit de convolution.
        \index{algèbre!L1(G)@$L^1(G)$}
  \item Le \emph{disque algébrique} $A(\mathbb{D})$~: fonctions continues sur
        $\overline{\mathbb{D}}$ et holomorphes sur $\mathbb{D}$, avec la norme
        $\norm{f}_\infty$.
        \index{disque algébrique}
\end{enumerate}
\end{exemple}

%-----------------------------------------------------------------------
\section{Spectre et rayon spectral}
\index{spectre!dans une algèbre de Banach}
\index{rayon spectral!formule du}
%-----------------------------------------------------------------------

\begin{definition}[Spectre dans une algèbre de Banach]
\label{def:spectre-algebre}
\index{spectre!définition algébrique}
\index{résolvante}
Soit $A$ une algèbre de Banach unitaire et $a \in A$. Le \emph{spectre} de $a$ est
\[
  \Spec(a) = \{\lambda \in \C : a - \lambda e \text{ n'est pas inversible dans } A\}.
\]
L'ensemble $\rho(a) = \C \setminus \Spec(a)$ est la \emph{résolvante} de $a$, et
l'application $\lambda \mapsto (a - \lambda e)^{-1}$ est la \emph{fonction résolvante}.
\end{definition}

\begin{proposition}
\label{prop:spectre-compact}
\index{spectre!compact non vide}
Pour tout $a \in A$ ($A$ unitaire), $\Spec(a)$ est un compact non vide de $\C$ contenu
dans le disque $\{\abs{\lambda} \le \norm{a}\}$.
\end{proposition}

\begin{proof}
Si $\abs{\lambda} > \norm{a}$, alors $a - \lambda e = -\lambda(e - a/\lambda)$ et
$\norm{a/\lambda} < 1$, donc $e - a/\lambda$ est inversible par la série de Neumann\index{Neumann!série de}~:
$(e - a/\lambda)^{-1} = \sum_{n=0}^\infty (a/\lambda)^n$. Donc
$\Spec(a) \subset \overline{B}(0, \norm{a})$.

Le spectre est fermé car l'ensemble des éléments inversibles est ouvert (par perturbation
de l'identité). Le spectre est non vide~: si $\Spec(a) = \emptyset$, alors la fonction
résolvante $\lambda \mapsto (a - \lambda e)^{-1}$ est entière (holomorphe sur $\C$
tout entier, à valeurs dans $A$). De plus, pour $\abs{\lambda} \to \infty$,
$\norm{(a - \lambda e)^{-1}} = \frac{1}{\abs{\lambda}}\norm{(e - a/\lambda)^{-1}} \to 0$.
Par le théorème de Liouville\index{Liouville!théorème de} vectoriel, la résolvante est
identiquement nulle, ce qui est absurde.
\end{proof}

\begin{theoreme}[Formule du rayon spectral]
\label{thm:rayon-spectral}
\index{rayon spectral!formule}
Soit $A$ une algèbre de Banach unitaire et $a \in A$. Le \emph{rayon spectral}
$r(a) = \sup\{\abs{\lambda} : \lambda \in \Spec(a)\}$ satisfait~:
\[
  r(a) = \lim_{n \to \infty} \norm{a^n}^{1/n} = \inf_{n \ge 1} \norm{a^n}^{1/n}.
\]
\end{theoreme}

\begin{proof}
Posons $R = \lim_{n} \norm{a^n}^{1/n}$ (la limite existe car la suite
$(\log\norm{a^n})$ est sous-additive, par le lemme de Fekete\index{Fekete!lemme de}).

\textbf{$r(a) \le R$.} Si $\lambda \in \Spec(a)$, alors $\lambda^n \in \Spec(a^n)$
(car si $a^n - \lambda^n e$ est inversible, un argument de factorisation montre que
$a - \lambda e$ est inversible). Donc $\abs{\lambda}^n \le \norm{a^n}$, i.e.,
$\abs{\lambda} \le \norm{a^n}^{1/n}$. En passant à la limite~: $\abs{\lambda} \le R$.

\textbf{$R \le r(a)$.} La série de Laurent de la résolvante au voisinage de l'infini
est~:
\[
  (a - \lambda e)^{-1} = -\sum_{n=0}^\infty \frac{a^n}{\lambda^{n+1}},
  \quad \abs{\lambda} > r(a).
\]
Cette série converge dans $A$ pour $\abs{\lambda} > r(a)$. Pour tout $\varphi \in A'$
(le dual de $A$), la fonction $\lambda \mapsto \varphi((a - \lambda e)^{-1})$ est
holomorphe sur $\{\abs{\lambda} > r(a)\}$ avec le développement
$-\sum_n \varphi(a^n)/\lambda^{n+1}$. Par la théorie des séries de Laurent, cette série
converge pour $\abs{\lambda} > r(a)$, donc
$\abs{\varphi(a^n)}/\abs{\lambda}^{n+1} \to 0$ pour tout $\abs{\lambda} > r(a)$.
Ainsi la suite $(\varphi(a^n/\lambda^n))$ est bornée pour tout $\varphi \in A'$. Par
Banach--Steinhaus\index{Banach--Steinhaus}, $\sup_n \norm{a^n/\lambda^n} < \infty$, d'où
$\norm{a^n}^{1/n} \le C^{1/n}\abs{\lambda}$ pour tout $\abs{\lambda} > r(a)$.
En passant à la limite~: $R \le \abs{\lambda}$ pour tout $\abs{\lambda} > r(a)$, donc
$R \le r(a)$.
\end{proof}

%-----------------------------------------------------------------------
\section{Algèbres de Banach commutatives~: théorie de Gelfand}
\index{Gelfand!théorie de}
\index{algèbre de Banach!commutative}
%-----------------------------------------------------------------------

\begin{definition}[Caractère]
\label{def:caractere}
\index{caractère!d'une algèbre de Banach}
Soit $A$ une algèbre de Banach commutative unitaire. Un \emph{caractère} de $A$ est un
morphisme d'algèbres non nul $\chi : A \to \C$, c'est-à-dire une forme linéaire
multiplicative~: $\chi(ab) = \chi(a)\chi(b)$ et $\chi(e) = 1$.
\end{definition}

\begin{proposition}
\label{prop:caractere-continu}
\index{caractère!automatiquement continu}
Tout caractère $\chi$ est automatiquement continu avec $\norm{\chi} = 1$.
\end{proposition}

\begin{proof}
Si $\abs{\chi(a)} > \norm{a}$, posons $\lambda = \chi(a)$. Alors
$a - \lambda e$ est inversible (car $\abs{\lambda} > \norm{a}$), et
$\chi(a - \lambda e) = \chi(a) - \lambda = 0$, donc $\chi$ annulerait un élément
inversible~: $0 = \chi((a-\lambda e)(a-\lambda e)^{-1}) = \chi(a - \lambda e)\chi((a-\lambda e)^{-1}) = 0$,
ce qui est cohérent mais impossible car
$\chi(e) = \chi((a-\lambda e)(a-\lambda e)^{-1}) = 0 \cdot \chi(\ldots) = 0 \neq 1$.
Contradiction. Donc $\abs{\chi(a)} \le \norm{a}$, i.e., $\norm{\chi} \le 1$. Comme
$\chi(e) = 1$ et $\norm{e} = 1$, on a $\norm{\chi} = 1$.
\end{proof}

\begin{definition}[Espace de Gelfand]
\label{def:espace-gelfand}
\index{Gelfand!espace de}
\index{spectre de Gelfand}
L'\emph{espace de Gelfand} (ou \emph{spectre}) de $A$ est l'ensemble $\hat A$ de tous les
caractères de $A$, muni de la topologie faible-$*$ (restreinte de $\sigma(A', A)$).
\end{definition}

\begin{proposition}
\label{prop:gelfand-compact}
\index{Gelfand!espace de!compact}
$\hat A$ est un espace compact non vide.
\end{proposition}

\begin{proof}
Comme $\norm{\chi} \le 1$ pour tout $\chi \in \hat A$, l'ensemble $\hat A$ est contenu
dans la boule unité fermée de $A'$, qui est compacte pour la topologie faible-$*$ par le
théorème de Banach--Alaoglu\index{Banach--Alaoglu (théorème de)}. Montrons que $\hat A$
est fermé (pour $\sigma(A', A)$)~: si $\chi_\alpha \to \varphi$ dans $\sigma(A', A)$,
alors $\varphi(ab) = \lim \chi_\alpha(ab) = \lim \chi_\alpha(a)\chi_\alpha(b) = \varphi(a)\varphi(b)$
et $\varphi(e) = 1$. Donc $\varphi \in \hat A$. Fermé dans un compact, $\hat A$ est
compact.

Non-vacuité~: tout idéal maximal $M$ de $A$ donne un caractère $A \to A/M \cong \C$ (car
$A/M$ est un corps, et un corps qui est une algèbre de Banach est isomorphe à $\C$ par
le théorème de Gelfand--Mazur\index{Gelfand--Mazur (théorème de)}).
\end{proof}

\begin{definition}[Transformée de Gelfand]
\label{def:transformee-gelfand}
\index{Gelfand!transformée de}
La \emph{transformée de Gelfand} est l'application
\[
  \Gamma : A \to C(\hat A), \quad \Gamma(a)(\chi) = \hat a(\chi) = \chi(a).
\]
\end{definition}

\begin{theoreme}[Propriétés de la transformée de Gelfand]
\label{thm:gelfand-morphisme}
\index{Gelfand!transformée de!propriétés}
L'application $\Gamma : A \to C(\hat A)$ est un morphisme d'algèbres continu avec~:
\begin{enumerate}[label=(\roman*)]
  \item $\widehat{ab} = \hat a \cdot \hat b$ et $\widehat{\alpha a + \beta b} = \alpha\hat a + \beta \hat b$.
  \item $\hat e = 1$.
  \item $\norm{\hat a}_\infty = r(a)$ (rayon spectral de $a$).
  \item $\Spec(a) = \hat a(\hat A) = \{\chi(a) : \chi \in \hat A\}$.
\end{enumerate}
\end{theoreme}

\begin{proof}
(i) et (ii) sont immédiats par définition.

(iii) On a $\lambda \in \Spec(a) \iff a - \lambda e$ n'est pas inversible
$\iff a - \lambda e$ appartient à un idéal maximal $M$
$\iff$ il existe $\chi \in \hat A$ avec $\chi(a - \lambda e) = 0$
$\iff \chi(a) = \lambda$ pour un $\chi \in \hat A$.
Donc $\Spec(a) = \hat a(\hat A)$ et
$\norm{\hat a}_\infty = \sup_\chi \abs{\chi(a)} = \sup\{\abs{\lambda} : \lambda \in \Spec(a)\} = r(a)$.
\end{proof}

%-----------------------------------------------------------------------
\section{$C^*$-algèbres}
\index{C*-algèbre@$C^*$-algèbre}
\index{C*-algèbre@$C^*$-algèbre!définition}
%-----------------------------------------------------------------------

\begin{definition}[$C^*$-algèbre]
\label{def:cstar}
Une \emph{$C^*$-algèbre} est une algèbre de Banach $A$ munie d'une
\emph{involution}\index{involution} $* : A \to A$ (antilinéaire, $(ab)^* = b^*a^*$,
$(a^*)^* = a$) satisfaisant l'\emph{axiome $C^*$}\index{axiome C*@axiome $C^*$}~:
\[
  \norm{a^*a} = \norm{a}^2, \quad \forall a \in A.
\]
\end{definition}

\begin{exemple}[Exemples de $C^*$-algèbres]
\label{ex:cstar}
\index{C*-algèbre@$C^*$-algèbre!exemples}
\begin{enumerate}[label=(\roman*)]
  \item $\mathcal{L}(H)$ avec $T^* =$ adjoint hilbertien. L'axiome $C^*$ est~:
        $\norm{T^*T} = \norm{T}^2$.
  \item $C(K)$ avec $f^*(x) = \overline{f(x)}$. L'axiome $C^*$ est~:
        $\norm{\bar f f}_\infty = \norm{f}_\infty^2$ (évident).
  \item Toute sous-algèbre fermée et stable par $*$ de $\mathcal{L}(H)$ est une
        $C^*$-algèbre.
\end{enumerate}
\end{exemple}

\begin{lemme}
\label{lem:cstar-autoadjoint}
\index{C*-algèbre@$C^*$-algèbre!éléments auto-adjoints}
Dans une $C^*$-algèbre, si $a = a^*$ (auto-adjoint), alors $\norm{a} = r(a)$.
\end{lemme}

\begin{proof}
On a $\norm{a^2} = \norm{a^*a} = \norm{a}^2$. Par récurrence~:
$\norm{a^{2^n}} = \norm{a}^{2^n}$. Donc
$r(a) = \lim_n \norm{a^n}^{1/n} = \lim_k \norm{a^{2^k}}^{1/2^k} = \norm{a}$.
\end{proof}

%-----------------------------------------------------------------------
\section{Théorème de Gelfand--Naimark commutatif}
\index{Gelfand--Naimark!théorème commutatif}
%-----------------------------------------------------------------------

\begin{theoreme}[Gelfand--Naimark commutatif]
\label{thm:gelfand-naimark}
Soit $A$ une $C^*$-algèbre commutative unitaire. Alors la transformée de Gelfand
\[
  \Gamma : A \to C(\hat A), \quad a \mapsto \hat a,
\]
est un isomorphisme isométrique de $C^*$-algèbres. Autrement dit, $A \cong C(\hat A)$.
\end{theoreme}

\begin{proof}
Par le théorème~\ref{thm:gelfand-morphisme}, $\Gamma$ est un morphisme d'algèbres avec
$\norm{\hat a}_\infty = r(a)$.

\textbf{Respect de l'involution.} Pour tout $\chi \in \hat A$ et $a \in A$~: si $a = a^*$,
alors $\Spec(a) \subset \R$ (dans une $C^*$-algèbre, les éléments auto-adjoints ont un
spectre réel\footnote{En effet, si $a = a^*$ et $\lambda = \alpha + i\beta$ avec
$\beta \neq 0$, alors $\norm{a + it e}^2 = \norm{(a+ite)^*(a+ite)} = \norm{a^2 + t^2 e} \le \norm{a}^2 + t^2$, mais aussi $r(a+ite) \ge \abs{\alpha + i(\beta + t)}$. Pour $t$ grand, cela donne une contradiction.}),
donc $\chi(a) \in \R = \overline{\chi(a)}$. Tout $a$ s'écrit
$a = \frac{a+a^*}{2} + i\frac{a - a^*}{2i}$ (somme de deux auto-adjoints), d'où
$\chi(a^*) = \overline{\chi(a)}$, i.e., $\widehat{a^*} = \overline{\hat a}$.

\textbf{Isométrie.} Pour tout $a \in A$~: $a^*a$ est auto-adjoint, donc
$\norm{a}^2 = \norm{a^*a} = r(a^*a) = \norm{\widehat{a^*a}}_\infty = \norm{\overline{\hat a}\cdot\hat a}_\infty = \norm{\hat a}_\infty^2$.
Donc $\norm{\hat a}_\infty = \norm{a}$~: $\Gamma$ est une isométrie.

\textbf{Injectivité.} Conséquence immédiate de l'isométrie.

\textbf{Surjectivité.} L'image $\Gamma(A)$ est une sous-algèbre fermée de $C(\hat A)$
(car $A$ est complet et $\Gamma$ isométrique) qui sépare les points de $\hat A$ (si
$\chi_1 \neq \chi_2$, il existe $a$ avec $\chi_1(a) \neq \chi_2(a)$), contient les
constantes ($\hat e = 1$), et est stable par conjugaison ($\widehat{a^*} = \overline{\hat a}$).
Par le théorème de Stone--Weierstrass\index{Stone--Weierstrass!théorème de},
$\Gamma(A) = C(\hat A)$.
\end{proof}

\begin{remarque}
Ce théorème dit que \emph{toute $C^*$-algèbre commutative unitaire est (isométriquement
isomorphe à) l'algèbre des fonctions continues sur un espace compact}. C'est un résultat
remarquable~: la structure algébrique et l'involution déterminent complètement l'espace
topologique sous-jacent.
\index{Gelfand--Naimark!interprétation}
\end{remarque}

%-----------------------------------------------------------------------
\section{Calcul fonctionnel continu dans les $C^*$-algèbres}
\index{calcul fonctionnel!dans les C*-algèbres@dans les $C^*$-algèbres}
%-----------------------------------------------------------------------

\begin{corollaire}[Calcul fonctionnel continu]
\label{cor:calcul-fonctionnel-cstar}
Soit $A$ une $C^*$-algèbre unitaire et $a \in A$ un élément normal ($a^*a = aa^*$). Alors
il existe un unique $*$-morphisme isométrique
$\Phi : C(\Spec(a)) \to A$ tel que $\Phi(\mathrm{id}) = a$ et $\Phi(1) = e$.
\end{corollaire}

\begin{proof}
La sous-$C^*$-algèbre $C^*(a)$ engendrée par $a$ et $e$ est commutative (car $a$ est
normal). Par le théorème de Gelfand--Naimark, $C^*(a) \cong C(\widehat{C^*(a)})$. On
montre que $\widehat{C^*(a)} \cong \Spec(a)$ via $\chi \mapsto \chi(a)$, ce qui donne
l'identification $C^*(a) \cong C(\Spec(a))$.
\end{proof}

%-----------------------------------------------------------------------
\section{États et formes positives}
\index{état!sur une C*-algèbre@sur une $C^*$-algèbre}
\index{forme positive}
%-----------------------------------------------------------------------

\begin{definition}[Forme positive, état]
\label{def:etat}
\index{forme positive!définition}
\index{état!définition}
Soit $A$ une $C^*$-algèbre unitaire.
\begin{itemize}
  \item Une forme linéaire $\varphi : A \to \C$ est \emph{positive} si
        $\varphi(a^*a) \ge 0$ pour tout $a \in A$.
  \item Un \emph{état} est une forme positive $\varphi$ avec $\varphi(e) = 1$.
\end{itemize}
L'ensemble des états est noté $\mathcal{S}(A)$.
\end{definition}

\begin{proposition}
\label{prop:etat-proprietes}
\index{état!propriétés}
\index{Cauchy--Schwarz!pour les états}
Soit $\varphi$ un état sur $A$. Alors~:
\begin{enumerate}[label=(\roman*)]
  \item $\varphi(a^*) = \overline{\varphi(a)}$ pour tout $a$.
  \item $\abs{\varphi(a)}^2 \le \varphi(a^*a)$ (Cauchy--Schwarz).
  \item $\varphi$ est continu avec $\norm{\varphi} = 1$.
\end{enumerate}
\end{proposition}

\begin{proof}
(i) On écrit $a = h + ik$ avec $h, k$ auto-adjoints. Comme $\varphi(h) \in \R$ pour $h$
auto-adjoint (car $h = h^*$ et $\varphi(h) = \varphi(h^*) = \overline{\varphi(h)}$, montré en écrivant $h$ comme différence de deux éléments positifs), on obtient
$\varphi(a^*) = \varphi(h - ik) = \varphi(h) - i\varphi(k) = \overline{\varphi(a)}$.

(ii) La forme sesquilinéaire $(a,b) \mapsto \varphi(b^*a)$ est positive
($\varphi(a^*a) \ge 0$). L'inégalité de Cauchy--Schwarz donne~:
$\abs{\varphi(b^*a)}^2 \le \varphi(a^*a)\varphi(b^*b)$. Avec $b = e$~:
$\abs{\varphi(a)}^2 \le \varphi(a^*a) \cdot 1$.

(iii) De (ii)~: $\abs{\varphi(a)}^2 \le \varphi(a^*a) \le \norm{a^*a} = \norm{a}^2$
(car $\varphi$ est positive et $\norm{a^*a}\,e - a^*a \ge 0$). Donc
$\norm{\varphi} \le 1$, et $\varphi(e) = 1$ donne $\norm{\varphi} = 1$.
\end{proof}

%-----------------------------------------------------------------------
\section{Construction GNS}
\index{GNS (construction)}
\index{Gelfand--Naimark--Segal}
%-----------------------------------------------------------------------

\begin{theoreme}[Construction GNS]
\label{thm:gns}
\index{GNS (construction)!théorème}
Soit $A$ une $C^*$-algèbre unitaire et $\varphi$ un état sur $A$. Il existe un espace de
Hilbert $H_\varphi$, un $*$-morphisme $\pi_\varphi : A \to \mathcal{L}(H_\varphi)$
(une \emph{représentation}) et un vecteur $\xi_\varphi \in H_\varphi$ tels que~:
\begin{enumerate}[label=(\roman*)]
  \item $\varphi(a) = \inner{\pi_\varphi(a)\xi_\varphi}{\xi_\varphi}$ pour tout $a \in A$.
  \item $\xi_\varphi$ est \emph{cyclique}~:
        $\overline{\pi_\varphi(A)\xi_\varphi} = H_\varphi$.
\end{enumerate}
Ce triplet $(H_\varphi, \pi_\varphi, \xi_\varphi)$ est unique à isomorphisme unitaire
près.
\end{theoreme}

\begin{proof}
\textbf{Étape 1~: Pré-espace de Hilbert.}
\index{GNS (construction)!preuve}
Posons $N_\varphi = \{a \in A : \varphi(a^*a) = 0\}$. Par Cauchy--Schwarz,
$N_\varphi$ est un idéal à gauche fermé de $A$. Sur le quotient $A / N_\varphi$,
la forme sesquilinéaire
\[
  \inner{[a]}{[b]} = \varphi(b^*a)
\]
est bien définie et est un produit scalaire (la positivité résulte de la définition de
$N_\varphi$~: $\inner{[a]}{[a]} = \varphi(a^*a) = 0 \Rightarrow [a] = 0$).

\textbf{Étape 2~: Complétion.}
Soit $H_\varphi$ le complété de $(A/N_\varphi, \inner{\cdot}{\cdot})$.

\textbf{Étape 3~: Représentation.}
Pour $a \in A$, l'application $[b] \mapsto [ab]$ est bien définie sur $A/N_\varphi$
(car $N_\varphi$ est un idéal à gauche). Elle est bornée~:
\[
  \norm{[ab]}^2 = \varphi(b^*a^*ab) \le \norm{a^*a}\,\varphi(b^*b) = \norm{a}^2\norm{[b]}^2
\]
(en utilisant que $\norm{a}^2 e - a^*a \ge 0$, donc
$\varphi(b^*(\norm{a}^2 e - a^*a)b) \ge 0$). Donc elle se prolonge en un opérateur borné
$\pi_\varphi(a) \in \mathcal{L}(H_\varphi)$ avec $\norm{\pi_\varphi(a)} \le \norm{a}$.

On vérifie que $\pi_\varphi$ est un $*$-morphisme~:
\begin{itemize}
  \item $\pi_\varphi(ab)[c] = [abc] = \pi_\varphi(a)[bc] = \pi_\varphi(a)\pi_\varphi(b)[c]$.
  \item $\inner{\pi_\varphi(a^*)[b]}{[c]} = \inner{[a^*b]}{[c]} = \varphi(c^*a^*b) = \inner{[b]}{[ac]} = \inner{[b]}{\pi_\varphi(a)[c]}$.
\end{itemize}
Donc $\pi_\varphi(a^*) = \pi_\varphi(a)^*$.

\textbf{Étape 4~: Vecteur cyclique.}
Posons $\xi_\varphi = [e] \in H_\varphi$. Alors~:
\[
  \inner{\pi_\varphi(a)\xi_\varphi}{\xi_\varphi}
    = \inner{[a]}{[e]} = \varphi(e^*a) = \varphi(a).
\]
Et $\pi_\varphi(A)\xi_\varphi = \{[a] : a \in A\} = A/N_\varphi$, qui est dense dans
$H_\varphi$ par construction.

\textbf{Unicité.} Si $(H', \pi', \xi')$ est un autre triplet satisfaisant (i) et (ii),
l'application $U : \pi_\varphi(a)\xi_\varphi \mapsto \pi'(a)\xi'$ est bien définie et
isométrique~:
$\norm{\pi'(a)\xi'}^2 = \inner{\pi'(a^*a)\xi'}{\xi'} = \varphi(a^*a) = \norm{\pi_\varphi(a)\xi_\varphi}^2$.
Elle se prolonge en un unitaire $U : H_\varphi \to H'$ avec $U\pi_\varphi(a)U^* = \pi'(a)$
et $U\xi_\varphi = \xi'$.
\end{proof}

\begin{corollaire}[Gelfand--Naimark abstrait]
\label{cor:gelfand-naimark-abstrait}
\index{Gelfand--Naimark!théorème abstrait}
Toute $C^*$-algèbre admet une représentation fidèle (injective) sur un espace de Hilbert.
Autrement dit, toute $C^*$-algèbre est isométriquement $*$-isomorphe à une sous-$C^*$-algèbre
de $\mathcal{L}(H)$ pour un certain espace de Hilbert $H$.
\end{corollaire}

\begin{proof}
Pour chaque état $\varphi \in \mathcal{S}(A)$, construisons le triplet GNS
$(H_\varphi, \pi_\varphi, \xi_\varphi)$. Posons
$H = \bigoplus_{\varphi \in \mathcal{S}(A)} H_\varphi$ et
$\pi = \bigoplus_\varphi \pi_\varphi$. Alors $\pi : A \to \mathcal{L}(H)$ est un
$*$-morphisme.

Injectivité~: si $\pi(a) = 0$, alors $\pi_\varphi(a) = 0$ pour tout état $\varphi$,
donc $\varphi(a^*a) = \inner{\pi_\varphi(a^*a)\xi_\varphi}{\xi_\varphi} = \norm{\pi_\varphi(a)\xi_\varphi}^2 = 0$.
Comme les états séparent les éléments positifs (pour tout $b \ge 0$, $b \neq 0$, il
existe un état $\varphi$ avec $\varphi(b) > 0$), on obtient $a^*a = 0$, d'où
$\norm{a}^2 = \norm{a^*a} = 0$ et $a = 0$.
\end{proof}

%-----------------------------------------------------------------------
\section{Application~: mécanique quantique et théorie des représentations}
\index{mécanique quantique!C*-algèbres@et $C^*$-algèbres}
%-----------------------------------------------------------------------

\begin{remarque}[Observables en mécanique quantique]
\index{observable}
La formulation algébrique de la mécanique quantique (due à Segal, Haag, Kastler)
considère les observables comme les éléments auto-adjoints d'une $C^*$-algèbre $A$. Les
états physiques sont les états au sens de la définition~\ref{def:etat}. La construction
GNS fournit alors l'espace de Hilbert et la représentation concrète du système.

Cette approche est essentielle en théorie quantique des champs, où l'espace de Hilbert
n'est pas fixé a priori mais dépend de l'état (représentations inéquivalentes).
\index{théorie quantique des champs}
\end{remarque}

%-----------------------------------------------------------------------
\section{Exercices}
%-----------------------------------------------------------------------

\begin{exercice}[$\star$]
\index{exercice!spectre}
Calculer le spectre de l'élément $f(x) = x$ dans $C([0,1])$. En déduire que
$r(f) = \norm{f}_\infty$.
\end{exercice}

\begin{exercice}[$\star$]
\index{exercice!rayon spectral}
Soit $A$ une algèbre de Banach unitaire et $a, b \in A$. Montrer que
$\Spec(ab) \cup \{0\} = \Spec(ba) \cup \{0\}$.
\end{exercice}

\begin{exercice}[$\star\star$]
\index{exercice!Gelfand--Mazur}
\index{Gelfand--Mazur (théorème de)}
(\emph{Théorème de Gelfand--Mazur.}) Montrer que si $A$ est une algèbre de Banach unitaire
dans laquelle tout élément non nul est inversible, alors $A \cong \C$.
\end{exercice}

\begin{exercice}[$\star\star$]
\index{exercice!transformée de Gelfand}
Calculer l'espace de Gelfand de $\ell^1(\Z)$ et identifier la transformée de Gelfand.
\emph{Indication~: les caractères sont les morphismes $a \mapsto \sum_n a(n) z^n$ pour
$z \in \mathbb{T}$ (le cercle unité).}
\end{exercice}

\begin{exercice}[$\star\star$]
\index{exercice!C*-algèbre}
Soit $A$ une $C^*$-algèbre et $a \in A$ normal. Montrer que
$\norm{a} = r(a)$.
\end{exercice}

\begin{exercice}[$\star\star\star$]
\index{exercice!GNS}
Soit $A = M_n(\C)$ (les matrices $n \times n$) et $\varphi(a) = \frac{1}{n}\Tr(a)$.
Construire explicitement le triplet GNS $(H_\varphi, \pi_\varphi, \xi_\varphi)$.
\end{exercice}

\begin{exercice}[$\star\star$]
\index{exercice!état pur}
\index{état!pur}
Un état $\varphi$ est \emph{pur} s'il est un point extrême de $\mathcal{S}(A)$. Montrer
que $\varphi$ est pur si et seulement si la représentation GNS associée est irréductible.
\end{exercice}

\begin{exercice}[$\star\star\star$]
\index{exercice!Wiener (théorème de)}
\index{Wiener!théorème de}
(\emph{Théorème de Wiener.}) Soit $f \in \ell^1(\Z)$ et
$\tilde f(z) = \sum_{n \in \Z} f(n)z^n$ sa série de Fourier. Montrer~: si $\tilde f(z) \neq 0$
pour tout $z \in \mathbb{T}$, alors $1/\tilde f$ est aussi une série de Fourier absolument
convergente, c'est-à-dire $1/\tilde f = \tilde g$ pour un $g \in \ell^1(\Z)$.

\emph{Indication~: utiliser la théorie de Gelfand.}
\end{exercice}


%%%%%%%%%%%%%%%%%%%%%%%%%%%%%%%%%%%%%%%%%%%%%%%%%%%%%%%%%%%%%%%%%%%%%%%
%  APPENDICE — Rappels d'Analyse et de Topologie
%%%%%%%%%%%%%%%%%%%%%%%%%%%%%%%%%%%%%%%%%%%%%%%%%%%%%%%%%%%%%%%%%%%%%%%
\appendix
\chapter{Rappels d'Analyse et de Topologie}
\label{app:rappels}
\minitoc

Ce chapitre rassemble sans démonstration les résultats d'analyse, de topologie et de
théorie de la mesure utilisés dans le cours.

\section{Théorème de Baire}
\index{Baire (théorème de)}

\begin{theoreme}[Baire]
\label{thm:baire-rappel}
Soit $X$ un espace métrique complet (ou un espace localement compact). Toute intersection
dénombrable d'ouverts denses est dense dans $X$. De manière équivalente, $X$ ne peut pas
être réunion dénombrable de fermés d'intérieur vide.
\end{theoreme}

\begin{remarque}
Ce théorème est le fondement des trois grands théorèmes de l'analyse fonctionnelle
linéaire~: Banach--Steinhaus, théorème de l'application ouverte et théorème du graphe
fermé.
\end{remarque}

\section{Lemme de Zorn}
\index{Zorn (lemme de)}

\begin{theoreme}[Lemme de Zorn]
\label{thm:zorn}
Soit $(X, \le)$ un ensemble ordonné inductif (i.e., toute chaîne admet un majorant).
Alors $X$ possède au moins un élément maximal.
\end{theoreme}

Ce résultat, équivalent à l'axiome du choix\index{axiome du choix}, est utilisé pour~:
\begin{itemize}
  \item L'existence de bases algébriques (bases de Hamel)\index{Hamel!base de}.
  \item Le théorème de Hahn--Banach\index{Hahn--Banach (théorème de)}.
  \item L'existence d'idéaux maximaux dans les algèbres.
  \item Le théorème de Tychonoff.
\end{itemize}

\section{Théorème de Tychonoff}
\index{Tychonoff (théorème de)}

\begin{theoreme}[Tychonoff]
\label{thm:tychonoff}
Tout produit d'espaces topologiques compacts est compact (pour la topologie produit).
\end{theoreme}

Ce théorème est équivalent à l'axiome du choix. Il est utilisé dans la preuve du théorème
de Banach--Alaoglu\index{Banach--Alaoglu (théorème de)} (compacité de la boule unité
du dual pour la topologie faible-$*$).

\section{Mesure de Lebesgue~: rappels}
\index{Lebesgue!mesure de}
\index{mesure de Lebesgue}

On rappelle les propriétés fondamentales de la mesure de Lebesgue $\lambda$ sur $\R^d$~:
\begin{enumerate}[label=(\roman*)]
  \item $\lambda$ est une mesure de Borel régulière, $\sigma$-finie, invariante par
        translation.
  \item $\lambda$ est complète (les sous-ensembles de mesure nulle sont mesurables).
  \item $\lambda$ est la complétion de sa restriction aux boréliens.
  \item Pour tout borélien $B$~: $\lambda(B) = \inf\{\sum_n \abs{I_n} : B \subset \bigcup_n I_n,
        \text{ $I_n$ intervalles}\}$.
\end{enumerate}

\section{Convergence dominée}
\index{convergence dominée}
\index{Lebesgue!théorème de convergence dominée}

\begin{theoreme}[Convergence dominée de Lebesgue]
\label{thm:conv-dominee}
Soit $(\Omega, \mathcal{A}, \mu)$ un espace mesuré et $(f_n)$ une suite de fonctions
mesurables telle que~:
\begin{enumerate}[label=(\roman*)]
  \item $f_n \to f$ $\mu$-presque partout.
  \item Il existe $g \in L^1(\Omega, \mu)$ avec $\abs{f_n} \le g$ $\mu$-p.p.\ pour tout $n$.
\end{enumerate}
Alors $f \in L^1(\Omega, \mu)$ et $\int f_n\,\dd\mu \to \int f\,\dd\mu$. De plus,
$\norm{f_n - f}_1 \to 0$.
\end{theoreme}

\section{Théorème de Fubini}
\index{Fubini (théorème de)}
\index{Tonelli (théorème de)}

\begin{theoreme}[Fubini--Tonelli]
\label{thm:fubini}
Soient $(\Omega_1, \mu_1)$ et $(\Omega_2, \mu_2)$ des espaces mesurés $\sigma$-finis.
\begin{enumerate}[label=(\roman*)]
  \item \textbf{Tonelli.} Si $f : \Omega_1 \times \Omega_2 \to [0, +\infty]$ est
        mesurable, alors
        \[
          \int_{\Omega_1 \times \Omega_2} f\,\dd(\mu_1 \otimes \mu_2)
            = \int_{\Omega_1}\left(\int_{\Omega_2} f(x,y)\,\dd\mu_2(y)\right)\dd\mu_1(x)
            = \int_{\Omega_2}\left(\int_{\Omega_1} f(x,y)\,\dd\mu_1(x)\right)\dd\mu_2(y).
        \]
  \item \textbf{Fubini.} Si $f \in L^1(\Omega_1 \times \Omega_2, \mu_1 \otimes \mu_2)$,
        alors les mêmes égalités sont valides et les intégrales intérieures définissent des
        fonctions intégrables pour $\mu_1$-p.p.\ $x$ (resp.\ $\mu_2$-p.p.\ $y$).
\end{enumerate}
\end{theoreme}

\section{Autres rappels utiles}

\begin{theoreme}[Convergence monotone]
\index{convergence monotone}
\label{thm:conv-monotone}
Si $(f_n)$ est une suite croissante de fonctions mesurables positives, alors
$\int \lim_n f_n\,\dd\mu = \lim_n \int f_n\,\dd\mu$.
\end{theoreme}

\begin{theoreme}[Lemme de Fatou]
\label{thm:fatou}
\index{Fatou!lemme de}
Si $(f_n)$ est une suite de fonctions mesurables positives, alors
$\int \liminf_n f_n\,\dd\mu \le \liminf_n \int f_n\,\dd\mu$.
\end{theoreme}

\begin{theoreme}[Stone--Weierstrass]
\label{thm:stone-weierstrass-rappel}
\index{Stone--Weierstrass!théorème de}
Soit $K$ un espace compact et $\mathcal{A} \subset C(K, \R)$ une sous-algèbre qui sépare
les points et contient les constantes. Alors $\mathcal{A}$ est dense dans $C(K, \R)$.

Version complexe~: si de plus $\mathcal{A}$ est stable par conjugaison, alors
$\mathcal{A}$ est dense dans $C(K, \C)$.
\end{theoreme}

\begin{theoreme}[Ascoli--Arzelà]
\label{thm:ascoli}
\index{Ascoli--Arzelà (théorème d')}
Soit $K$ un espace métrique compact. Un sous-ensemble $\mathcal{F} \subset C(K)$ est
relativement compact si et seulement si $\mathcal{F}$ est borné et équicontinu.
\end{theoreme}


%%%%%%%%%%%%%%%%%%%%%%%%%%%%%%%%%%%%%%%%%%%%%%%%%%%%%%%%%%%%%%%%%%%%%%%
%  BIBLIOGRAPHIE
%%%%%%%%%%%%%%%%%%%%%%%%%%%%%%%%%%%%%%%%%%%%%%%%%%%%%%%%%%%%%%%%%%%%%%%
\begin{thebibliography}{99}

\bibitem{Megginson1998}
R.\,E. Megginson,
\emph{An Introduction to Banach Space Theory},
Graduate Texts in Mathematics~183, Springer-Verlag, 1998.

\bibitem{rudin-real-complex}
W. Rudin,
\emph{Real and Complex Analysis},
3e édition, McGraw-Hill, 1987.

\end{thebibliography}

%%%%%%%%%%%%%%%%%%%%%%%%%%%%%%%%%%%%%%%%%%%%%%%%%%%%%%%%%%%%%%%%%%%%%%%
%  INDEX ET FIN DU DOCUMENT
%%%%%%%%%%%%%%%%%%%%%%%%%%%%%%%%%%%%%%%%%%%%%%%%%%%%%%%%%%%%%%%%%%%%%%%
\printindex
\end{document}
